%% ============================================================
%%  Série Avaliação na Prática — FGV CLEAR
%%  Guia de Fontes de Dados para Avaliação de Políticas Públicas
%% ============================================================

% !TEX program = pdflatex
% !TEX encoding = UTF-8 Unicode

\documentclass[12pt, a4paper, openany]{book}

%% ── Codificação e Idioma ────────────────────────────────────
\usepackage[T1]{fontenc}
\usepackage[utf8]{inputenc}
\usepackage[brazil]{babel}

%% ── Fonte ───────────────────────────────────────────────────
\usepackage{tgpagella}

%% ── Matemática ──────────────────────────────────────────────
\usepackage{amsmath, amssymb, amsfonts, dsfont, amsthm}

%% ── Figuras e Tabelas ───────────────────────────────────────
\usepackage{graphicx, float}
\usepackage[labelfont=bf, font=small, justification=justified]{caption}
\usepackage{subcaption}
\usepackage{booktabs, multirow, multicol, array}
\usepackage{longtable}
\usepackage{tabularx}

%% ── Layout ──────────────────────────────────────────────────
\usepackage[top=2.5cm, bottom=2.5cm, left=2.5cm, right=2.5cm, headheight=15pt]{geometry}
\usepackage{setspace}
\onehalfspacing
\usepackage{fancyhdr}
\usepackage{emptypage}
\usepackage{enumitem}
\setlist[itemize]{leftmargin=*, topsep=4pt, itemsep=2pt}
\setlist[enumerate]{leftmargin=*, topsep=4pt, itemsep=2pt}

%% ── Cores (paleta CLEAR) ────────────────────────────────────
\usepackage{xcolor}
\definecolor{clearblue}{HTML}{1B3A6E}
\definecolor{clearlightblue}{HTML}{5CA4A9}
\definecolor{cleardarkblue}{HTML}{003F72}
\definecolor{cleargray}{HTML}{6E6E6E}
% Variações específicas usadas nas caixas deste guia (preservadas)
\definecolor{clearsalmon}{RGB}{207, 90, 54}
\definecolor{clearlightsalmon}{RGB}{250, 228, 218}
\definecolor{clearlightbluelight}{RGB}{215, 232, 248}
\definecolor{cleargreen}{RGB}{45, 120, 65}
\definecolor{clearlightgreen}{RGB}{218, 240, 222}
\definecolor{clearred}{RGB}{170, 30, 30}
\definecolor{clearlightred}{RGB}{248, 218, 215}
% Paleta padronizada do template (mantida para compatibilidade)
\definecolor{salmonbg}{RGB}{252,235,230}
\definecolor{salmonframe}{RGB}{230,160,140}
\definecolor{bluebg}{RGB}{235,245,251}
\definecolor{blueframe}{RGB}{58,125,181}
\definecolor{orangebg}{RGB}{254,245,231}
\definecolor{orangeframe}{RGB}{230,126,34}
\definecolor{codebg}{RGB}{248,248,248}
\definecolor{codeframe}{RGB}{200,200,200}

%% ── Hyperlinks e Bibliografia ───────────────────────────────
\usepackage{hyperref}
\usepackage{url}
\hypersetup{
  colorlinks=true,
  linkcolor=clearlightblue,
  citecolor=clearlightblue,
  urlcolor=clearlightblue,
  pdfauthor={FGV CLEAR},
  pdftitle={Guia de Fontes de Dados para Avalia\c{c}\~ao de Pol\'iticas P\'ublicas}
}
\usepackage[round, authoryear]{natbib}
\bibliographystyle{plainnat}

%% URLs longos: permitir quebra em qualquer separador comum
\makeatletter
\g@addto@macro{\UrlBreaks}{\do\/\do\-\do\.\do\_\do\?\do\=\do\&\do\#\do\%\do\:}
\makeatother
\emergencystretch=3em
\tolerance=800
\hyphenpenalty=50

%% Hifenização para termos técnicos recorrentes no guia
\hyphenation{
  CEPAL-STAT BA-DE-HOG PLA-TA-FORMA-NILO-PECANHA POR-TAL-DA-TRANS-PARENCIA
  Mi-cro-da-dos ava-li-a-ti-va ava-li-a-ti-vas desa-gre-ga-ção desa-gre-ga-ções
  mi-cro-da-dos multi-di-men-si-o-nais subna-ci-o-nais re-gis-tros es-ta-du-al
  mu-ni-ci-pal es-ta-du-ais me-to-do-ló-gi-cas me-to-do-ló-gi-co comple-men-ta-res
  PNAD-C PNAD FINBRA FINBRA-Siconfi PRODES MapBiomas Datasus
  Eco-nô-mi-co De-sen-vol-vi-men-to in-ter-na-ci-o-nais in-ter-na-ci-o-nal
}

%% ── TColorbox ───────────────────────────────────────────────
\usepackage[most]{tcolorbox}
\tcbuselibrary{breakable, skins}

%% ── Código R ────────────────────────────────────────────────
\usepackage{listings}
\lstset{
  language=R, basicstyle=\small\ttfamily,
  keywordstyle=\color{clearblue}\bfseries,
  commentstyle=\color{cleargray}\itshape,
  stringstyle=\color{cleargreen},
  backgroundcolor=\color{codebg},
  frame=single, rulecolor=\color{codeframe},
  breaklines=true, showstringspaces=false,
  numbers=left, numberstyle=\tiny\color{cleargray}, numbersep=5pt, xleftmargin=12pt
}

%% ── Cabeçalhos e Rodapés ────────────────────────────────────
\pagestyle{fancy}
\fancyhf{}
\fancyhead[R]{\small\color{clearblue}\textit{Guia CLEAR $|$ Fontes de Dados}}
\fancyfoot[C]{\thepage}
\renewcommand{\headrulewidth}{0.4pt}
\renewcommand{\headrule}{\hbox to\headwidth{\color{clearblue}\leaders\hrule height \headrulewidth\hfill}}
\fancypagestyle{plain}{\fancyhf{}\fancyfoot[C]{\thepage}\renewcommand{\headrulewidth}{0pt}}

%% ── Títulos de Seção ────────────────────────────────────────
\usepackage{titlesec}
\titleformat{\chapter}[display]
  {\normalfont\huge\bfseries\color{clearblue}}
  {\normalsize\color{clearlightblue}\chaptertitlename~\thechapter}
  {8pt}{\Huge\color{clearblue}}
\titlespacing*{\chapter}{0pt}{-10pt}{24pt}
\titleformat{\section}
  {\normalfont\Large\bfseries\color{clearblue}}{\thesection}{1em}{}
\titleformat{\subsection}
  {\normalfont\large\bfseries\color{cleardarkblue}}{\thesubsection}{1em}{}
\titleformat{\subsubsection}
  {\normalfont\normalsize\bfseries\color{cleardarkblue}}{\thesubsubsection}{1em}{}

%% ── Wrapper de Título para Caixas com Títulos Longos ────────
\newcommand{\tituloboxlimitado}[1]{%
  \parbox[t]{\dimexpr\linewidth-16mm\relax}{\raggedright\strut #1\strut}%
}

%% ── Caixas Temáticas (semântica preservada deste guia) ──────
% boxdestaque (salmão) — usada para conceitos-chave
\newtcolorbox{boxdestaque}[1]{enhanced,breakable,colback=clearlightsalmon,colframe=clearsalmon,
  coltitle=white,fonttitle=\bfseries\small,
  attach boxed title to top left={yshift=-2mm,xshift=6mm},
  boxed title style={colback=clearsalmon,colframe=clearsalmon,sharp corners,rounded corners=northeast,
    top=2pt,bottom=2pt,left=6pt,right=6pt},
  title={\tituloboxlimitado{#1}},sharp corners,rounded corners=southeast}

% boxatencao (vermelha) — alertas e cuidados
\newtcolorbox{boxatencao}[1]{enhanced,breakable,colback=clearlightred,colframe=clearred,
  coltitle=white,fonttitle=\bfseries\small,
  attach boxed title to top left={yshift=-2mm,xshift=6mm},
  boxed title style={colback=clearred,colframe=clearred,sharp corners,rounded corners=northeast,
    top=2pt,bottom=2pt,left=6pt,right=6pt},
  title={\tituloboxlimitado{#1}},sharp corners,rounded corners=southeast}

% boxdica (azul) — dicas práticas
\newtcolorbox{boxdica}[1]{enhanced,breakable,colback=clearlightbluelight,colframe=clearblue,
  coltitle=white,fonttitle=\bfseries\small,
  attach boxed title to top left={yshift=-2mm,xshift=6mm},
  boxed title style={colback=clearblue,colframe=clearblue,sharp corners,rounded corners=northeast,
    top=2pt,bottom=2pt,left=6pt,right=6pt},
  title={\tituloboxlimitado{#1}},sharp corners,rounded corners=southeast}

% boxexemplo (verde) — exemplos aplicados
\newtcolorbox{boxexemplo}[1]{enhanced,breakable,colback=clearlightgreen,colframe=cleargreen,
  coltitle=white,fonttitle=\bfseries\small,
  attach boxed title to top left={yshift=-2mm,xshift=6mm},
  boxed title style={colback=cleargreen,colframe=cleargreen,sharp corners,rounded corners=northeast,
    top=2pt,bottom=2pt,left=6pt,right=6pt},
  title={\tituloboxlimitado{#1}},sharp corners,rounded corners=southeast}

%% Caixas padronizadas adicionais do template (disponíveis para uso futuro)
\newtcolorbox{destaquebox}[1][Destaque]{
  enhanced, breakable,
  colback=salmonbg, colframe=salmonframe, coltitle=black, fonttitle=\bfseries,
  title=#1, boxrule=0.5pt, arc=2pt,
  left=6pt, right=6pt, top=4pt, bottom=4pt
}
\newtcolorbox{notabox}[1][Nota]{
  enhanced, breakable,
  colback=bluebg, colframe=blueframe, coltitle=white, fonttitle=\bfseries,
  title=#1, boxrule=0.5pt, arc=2pt,
  left=6pt, right=6pt, top=4pt, bottom=4pt
}
\newtcolorbox{conexaobox}[1][Conex\~ao com a s\'erie]{
  enhanced, breakable,
  colback=orangebg, colframe=orangeframe, coltitle=black, fonttitle=\bfseries,
  title=#1, boxrule=0.5pt, arc=2pt,
  left=6pt, right=6pt, top=4pt, bottom=4pt
}
\newtcolorbox{exemplobox}[1][Exemplo]{
  enhanced, breakable,
  colback=salmonbg, colframe=salmonframe, coltitle=black, fonttitle=\bfseries,
  title=#1, boxrule=0.5pt, arc=2pt,
  left=6pt, right=6pt, top=4pt, bottom=4pt
}
\newtcolorbox{hipotesebox}[1][Hip\'otese]{
  enhanced, breakable,
  colback=salmonbg, colframe=salmonframe, coltitle=black, fonttitle=\bfseries,
  title=#1, boxrule=0.5pt, arc=2pt,
  left=6pt, right=6pt, top=4pt, bottom=4pt
}

%% ── Ficha rápida (comando ESSENCIAL deste guia) ─────────────
% URLs longos na linha "Acesso" precisam quebrar em qualquer caractere
\newcommand{\ficharapida}[6]{%
  \begin{tcolorbox}[enhanced,breakable,colback=clearlightbluelight!60,colframe=clearblue,
    coltitle=white,fonttitle=\bfseries\small,title={\tituloboxlimitado{Ficha rápida: #1}},
    attach boxed title to top left={yshift=-2mm,xshift=6mm},
    boxed title style={colback=clearblue,colframe=clearblue,sharp corners,rounded corners=northeast,
      top=2pt,bottom=2pt,left=6pt,right=6pt},
    sharp corners,rounded corners=southeast]
  \small
  \begingroup
    \Urlmuskip=0mu plus 1mu\relax
    \def\UrlBreaks{\do\/\do\-\do\.\do\_\do\?\do\=\do\&\do\#\do\%\do\:\do\a\do\b\do\c\do\d\do\e\do\f\do\g\do\h\do\i\do\j\do\k\do\l\do\m\do\n\do\o\do\p\do\q\do\r\do\s\do\t\do\u\do\v\do\w\do\x\do\y\do\z\do\A\do\B\do\C\do\D\do\E\do\F\do\G\do\H\do\I\do\J\do\K\do\L\do\M\do\N\do\O\do\P\do\Q\do\R\do\S\do\T\do\U\do\V\do\W\do\X\do\Y\do\Z\do\0\do\1\do\2\do\3\do\4\do\5\do\6\do\7\do\8\do\9}%
    \sloppy\hbadness=10000\relax
  \begin{tabular}{@{}p{3.2cm}p{\dimexpr\linewidth-3.5cm}@{}}
    \textbf{Produtor:}      & #2 \\[3pt]
    \textbf{Acesso:}        & {\footnotesize\url{#3}} \\[3pt]
    \textbf{Cobertura:}     & #4 \\[3pt]
    \textbf{Periodicidade:} & #5 \\[3pt]
    \textbf{Unidade:}       & #6 \\
  \end{tabular}
  \endgroup
  \end{tcolorbox}%
}

%% ── Comando \fonte: equivalente local simplificado do \figmeta ──
\newcommand{\fonte}[1]{\vspace{3pt}\par\small\textit{Fonte: #1}}

%% ── Metadados de Figuras/Quadros (padrão da série) ──────────
\newcommand{\figmeta}[3]{%
  \par\vspace{4pt}%
  \begin{minipage}{\linewidth}%
  \footnotesize%
  \textbf{Fonte:} #1%
  \ifx&#2&\else\\\textbf{Elabora\c{c}\~ao:} #2\fi%
  \ifx&#3&\else\\\textbf{Nota:} #3\fi%
  \end{minipage}%
  \vspace{6pt}%
}

%% ── Teoremas ────────────────────────────────────────────────
\newtheorem{definicao}{Defini\c{c}\~ao}[chapter]
\newtheorem{proposicao}{Proposi\c{c}\~ao}[chapter]

%% ── Profundidade de numera\c{c}\~ao ─────────────────────────
\setcounter{secnumdepth}{2}
\setcounter{tocdepth}{2}

% ============================================================
\begin{document}

\frontmatter
% ============================================================
%  cap00_capa.tex — Capa e Ficha Técnica
%  Guia de Fontes de Dados para Avaliação de Políticas Públicas
% ============================================================

\begin{titlepage}
\newgeometry{top=3cm, bottom=2.5cm, left=3cm, right=2.5cm}
\thispagestyle{empty}
\color{clearblue}

% Série
\vspace*{1cm}
{\centering
{\large\bfseries\color{clearlightblue}
\MakeUppercase{Série: Avaliação na Prática}}\par}

\vspace{1.8cm}

% Título principal
{\centering
{\fontsize{36}{42}\selectfont\bfseries\color{clearblue}
Guia de Fontes de Dados}\par}

\vspace{0.4cm}

{\centering
{\fontsize{24}{30}\selectfont\bfseries\color{clearblue}
para Avaliação de}\par}

\vspace{0.2cm}

{\centering
{\fontsize{24}{30}\selectfont\bfseries\color{clearblue}
Políticas Públicas}\par}

\vspace{2cm}

% Instituições
{\centering
{\large\color{clearblue}
FGV CLEAR \,|\, FGV EESP \,|\, Global Evaluation Initiative}\par}

\vspace{0.5cm}

% Data
{\centering{\normalsize\color{cleargray}[Mês/Ano]}\par}

\vspace{2cm}

% Linha separadora
{\color{clearlightblue}\rule{\linewidth}{0.8pt}}

\vspace{0.6cm}

% Resumo/sinopse
\begin{tcolorbox}[
  enhanced, colback=clearlightbluelight!50, colframe=clearblue,
  sharp corners, boxrule=0.5pt, left=8pt, right=8pt, top=6pt, bottom=6pt
]
\small\color{black}
Este guia documenta o ecossistema de fontes de dados disponíveis para a avaliação de
políticas públicas no Brasil. Cobre pesquisas amostrais, registros administrativos,
censos, bases internacionais e ferramentas emergentes (APIs, \textit{web scraping} e
inteligência artificial), organizados por área temática: mercado de trabalho, saúde,
educação, assistência social, segurança pública, infraestrutura, finanças públicas e
fontes internacionais. Para cada fonte são apresentados produtor, cobertura, periodicidade,
unidade de análise, pontos fortes, limitações e indicações de uso em avaliação de impacto.
\end{tcolorbox}

\restoregeometry
\end{titlepage}

% ── Ficha Técnica ──────────────────────────────────────────
\thispagestyle{empty}
\vspace*{2cm}

{\large\bfseries\color{clearblue} Ficha Técnica}

\bigskip

\noindent\textbf{Organizadores:} [Nome 1], [Nome 2]

\smallskip
\noindent\textbf{Equipe técnica:} [Listar membros]

\bigskip

\noindent\textbf{Apresentação.} Este guia integra a série de publicações
\textit{Avaliação na Prática}, desenvolvida pelo FGV CLEAR com o objetivo de
ampliar o acesso a conhecimentos sobre monitoramento e avaliação com foco em
políticas públicas.

\medskip

\noindent Fundado em 2015, o FGV CLEAR tem se dedicado ao fortalecimento da cultura
de gestão orientada por evidências no Brasil e em países lusófonos. Com sede na
Escola de Economia de São Paulo da Fundação Getulio Vargas (FGV EESP), o FGV CLEAR
atua como centro regional da Iniciativa CLEAR
(\textit{Centers for Learning on Evaluation and Results}).

\medskip

\noindent Saiba mais sobre o FGV CLEAR e acesse outras publicações em:
\textbf{www.fgvclear.org}

\vfill

\noindent{\small\color{cleargray}%
Os textos desta publicação são de responsabilidade exclusiva dos autores.
As opiniões expressas não refletem necessariamente a posição das instituições parceiras.}

\clearpage

% ============================================================
%  cap00_apresentacao.tex
%  Apresentação (seção não numerada — frontmatter)
% ============================================================

\chapter*{Apresentação}
\addcontentsline{toc}{chapter}{Apresentação}

Este guia integra a série \textit{Avaliação na Prática}, publicada pelo FGV CLEAR --- Centro de Aprendizado e Pesquisa em Avaliação e Resultados, vinculado à Escola de Economia de São Paulo da Fundação Getulio Vargas (FGV EESP). A série tem como objetivo oferecer materiais de referência práticos e acessíveis para gestores, pesquisadores e técnicos envolvidos na avaliação de políticas públicas no Brasil.

Os volumes anteriores da série trataram de métodos de avaliação --- diferenças em diferenças, regressão descontínua, \textit{matching}, poder estatístico e análise de custo-benefício. Este guia ocupa um lugar diferente: em vez de discutir \textit{como} avaliar, ocupa-se de discutir \textit{com quais dados} avaliar. É um atlas do ecossistema de fontes de dados disponíveis para a avaliação de políticas públicas no Brasil.

\medskip

O Brasil é, sob muitos aspectos, um caso privilegiado para a avaliação quantitativa de políticas públicas. O país dispõe de um sistema estatístico nacional amadurecido, liderado pelo IBGE, de registros administrativos abrangentes consolidados ao longo de décadas, e de uma cultura crescente de abertura de dados que colocou as informações governamentais ao alcance de qualquer pesquisador com acesso à internet. O Censo Demográfico, a PNAD Contínua, o Censo Escolar, o DATASUS, a RAIS, o CadÚnico --- cada uma dessas fontes representou, em seu tempo, um avanço significativo na capacidade de compreender e avaliar as políticas públicas brasileiras.

Ao mesmo tempo, a abundância de fontes traz seus próprios desafios. Cada base foi criada com uma finalidade específica, obedece a sua própria lógica de produção e tem limitações que só se revelam a quem a conhece bem. Usar o Censo Escolar para uma pergunta que exige a PNAD, ou confundir os dados do SIH com uma medida da saúde da população em geral, são erros que comprometem análises tecnicamente sofisticadas. Este guia pretende ser o mapa que evita esses desvios.

\medskip

O guia está organizado em treze capítulos. O primeiro discute por que a escolha das fontes de dados é uma decisão metodológica com consequências diretas sobre a validade das avaliações e apresenta os critérios para selecionar fontes adequadas. Os capítulos seguintes percorrem as principais áreas temáticas --- tipologia de fontes, pesquisas amostrais e censos, mercado de trabalho, saúde, educação, assistência social, segurança pública, infraestrutura e território, finanças públicas e fontes internacionais --- antes de dedicar um capítulo às ferramentas tecnológicas emergentes (APIs, \textit{web scraping} e inteligência artificial) e encerrar com orientações práticas sobre como escolher, acessar, linkar e documentar fontes de dados.

Os Anexos A e B oferecem ferramentas de referência rápida: uma tabela-catálogo consolidada de todas as fontes apresentadas e um glossário dos principais termos técnicos.

\medskip

Este guia foi elaborado pela equipe do FGV CLEAR. Qualquer erro ou omissão é de nossa responsabilidade.

\bigskip

\noindent\textit{São Paulo, 2024}

\noindent\textbf{FGV CLEAR}

% ─── SOBRE ESTE GUIA ─────────────────────────────────────────────────────────
\thispagestyle{empty}
\section*{Sobre este guia}
Toda avaliação começa por uma pergunta. Mas a pergunta só vira evidência
quando há dados que a respondam. E os dados disponíveis num país nunca
são neutros: foram produzidos por alguém, com algum objetivo, com alguma
decisão metodológica embutida que afeta o que pode e o que não pode ser
feito com eles depois.

O Brasil tem uma infraestrutura ampla de dados públicos, construída por
uma sucessão de instituições ao longo de décadas. O IBGE produz o Censo
Demográfico decenal e a PNAD trimestral, base para indicadores sociais e
de mercado de trabalho. O Datasus consolida o sistema de informação em
saúde: SIM para mortalidade, SIH para internações hospitalares, SIA para
atenção ambulatorial. O Ministério do Trabalho mantém o CAGED, registro
de admissões e desligamentos no emprego formal. O Inep cuida da educação
básica e superior, do Censo Escolar ao Saeb ao Enem. O INPE produz dados
ambientais a partir de satélite, do desmatamento amazônico ao uso do solo
agrícola. O Siconfi consolida as contas dos entes federados. E há
dezenas de outras: CGU, IPEA, ANVISA, Suframa, Banco Central, todos
publicando dados, cada um com sua lógica de coleta, frequência, cobertura
e granularidade.

Saber qual dessas fontes serve para qual pergunta é boa parte da batalha
em qualquer avaliação. Querer estimar desigualdade salarial entre regiões:
a PNAD captura, mas a RAIS (limitada ao emprego formal) pode subestimar
a desigualdade real porque deixa de fora o setor informal, onde se
concentra parte importante da heterogeneidade. Querer medir a prevalência
de uma doença crônica não letal: o SIM não serve (registra óbitos), o SIH
pode (registra internações), com a ressalva de que captura apenas os
casos graves o suficiente para chegar ao hospital. Querer estudar política
agrícola em regiões específicas: a Pesquisa Agrícola Municipal do IBGE
existe, mas tem cobertura desigual entre culturas, e para algumas
culturas é mais confiável usar dados de comercialização da Conab. Cada
fonte cobre uma fatia da realidade, e usar a fonte errada para a pergunta
certa produz uma evidência que parece sólida mas mede a coisa errada.

A cada ano também surgem fontes novas, e algumas antigas mudam de
natureza ou de qualidade. APIs governamentais agora permitem extrair
dados em volume sem fazer download manual. Microdados antes restritos
passam a ser abertos. Algumas fontes históricas sofrem rupturas
metodológicas (a PNAD virou PNAD Contínua em 2012, com mudança de painel
e amostra) que tornam séries longas comparáveis só com cuidado.
Ferramentas de IA e web scraping ampliam o que pode ser considerado
``fonte'' hoje, pegando informação de portais de transparência, notícias
e documentos que antes ficavam fora do alcance da avaliação.

Este guia funciona como um catálogo orientado. Para cada eixo temático
(mercado de trabalho e renda, saúde, educação, assistência social,
segurança pública, infraestrutura e território, finanças públicas, fontes
internacionais), descreve as principais fontes brasileiras: quem produz,
o que cobrem, com que frequência são atualizadas, em que unidade
geográfica são disponibilizadas, e que limitações o avaliador precisa
conhecer antes de usá-las. Cobre também ferramentas de acesso (APIs, web
scraping, uso de IA na extração de dados não estruturados) e fontes
internacionais para comparações entre países. Não é um manual de método.
É um guia de fontes: o que existe, onde encontrar, e como decidir qual
delas responde à pergunta em mãos.
\newpage

% ─── COMO LER ESTE GUIA ──────────────────────────────────────────────────────
\thispagestyle{empty}
\section*{Como ler este guia}
Três tipos de caixa estruturam a leitura, identificados por cor:
\vspace{0.3cm}
\begin{boxdica}{Caixa azul --- pontos-chave}
Aparece ao fim de cada capítulo, com sínteses e recomendações para quem
precisa decidir.
\end{boxdica}
\begin{boxexemplo}{Caixa verde --- exemplo aplicado}
Mostra como as ideias do capítulo se traduzem em situações de política
pública. Os casos são fictícios, mas calibrados por padrões recorrentes
na literatura e na prática de gestão.
\end{boxexemplo}
\begin{boxatencao}{Caixa vermelha --- cuidados essenciais}
Sinaliza armadilhas comuns --- pontos em que uma escolha mal fundamentada
compromete toda a análise.
\end{boxatencao}
\vspace{0.5cm}
\noindent O guia faz referência recorrente ao \textit{Guia CLEAR de
Monitoramento e Avaliação de Políticas Públicas: do diagnóstico à decisão}
\citep{Lima2025} --- doravante \textbf{Guia CLEAR de M\&A} ---, que trata
em profundidade as etapas do ciclo da política pública adotado na série.
Quando o texto menciona Teoria da Mudança, plano de monitoramento ou
tipologia de avaliações, é nesse guia que o leitor encontra o
desenvolvimento completo.
\newpage

\tableofcontents


\mainmatter
% ============================================================
%  cap01_introducao.tex
%  Capítulo 1: Por que as fontes de dados importam para
%              avaliação de políticas públicas?
% ============================================================

% ════════════════════════════════════════════════════════════
\chapter{Por que as fontes de dados importam para avaliação de políticas públicas?}
% ════════════════════════════════════════════════════════════

\begin{boxdestaque}{Neste capítulo}
A escolha da fonte de dados não é uma decisão técnica menor --- é um passo metodológico com consequências diretas sobre a validade das conclusões de qualquer avaliação. Este capítulo apresenta os princípios que orientam essa escolha, organizados em torno de três perguntas: (i) \textit{o que distingue dados bons de dados ruins para uma determinada pergunta avaliativa?}; (ii) \textit{como pensar a complementaridade entre diferentes tipos de fontes?}; e (iii) \textit{o que o avaliador precisa verificar antes de depositar confiança em qualquer base?} O capítulo também oferece uma visão panorâmica do ecossistema de dados disponíveis para avaliação de políticas públicas no Brasil.
\end{boxdestaque}

% ────────────────────────────────────────────────────────────
\section{A fonte de dados como escolha metodológica}
% ────────────────────────────────────────────────────────────

Toda avaliação de políticas públicas começa com uma pergunta. Mas entre a pergunta e a resposta existe uma etapa que frequentemente recebe menos atenção do que merece: a escolha dos dados. Essa escolha não é neutra. A fonte de dados selecionada determina o que pode ser medido, com que precisão, para quem e em que período. Quando a fonte é inadequada para a pergunta, nem a mais sofisticada estratégia de identificação causal produz resultados válidos.

Considere um exemplo simples. Um avaliador quer estimar o impacto de um programa de capacitação profissional sobre a empregabilidade dos participantes. Se usar a PNAD Contínua como fonte de dados, estará trabalhando com uma amostra representativa da população geral --- mas provavelmente insuficiente para identificar os próprios beneficiários do programa, cujo número pode ser pequeno. Se usar a RAIS, terá dados administrativos sobre o emprego formal de praticamente toda a força de trabalho formal brasileira --- mas perderá de vista os trabalhadores informais, que frequentemente são o público prioritário desses programas. Se conseguir acesso aos registros administrativos do programa, terá os beneficiários identificados individualmente --- mas precisará de uma fonte externa para construir o grupo de comparação.

Cada uma dessas fontes captura uma dimensão diferente da realidade, e a pergunta avaliativa determina qual delas --- ou qual combinação --- é mais adequada. Essa lógica de adequação entre pergunta e fonte é o fio condutor deste guia.

\begin{boxatencao}{Atenção --- Não existe fonte perfeita}
Toda fonte de dados tem limitações. O Censo Demográfico tem cobertura universal, mas é realizado apenas a cada dez anos. A PNAD Contínua é anual e tempestiva, mas tem tamanho de amostra insuficiente para análises no nível de municípios pequenos. Os registros administrativos têm cobertura ampla dos eventos que registram, mas cobrem apenas quem interage com o sistema --- excluindo populações que precisam de serviços mas não os acessam. Reconhecer as limitações das fontes que se usa não é fraqueza metodológica: é o primeiro passo para interpretar resultados com honestidade.
\end{boxatencao}

% ────────────────────────────────────────────────────────────
\section{Tipos de fontes de dados: uma primeira distinção}
% ────────────────────────────────────────────────────────────

Antes de discutir os critérios de seleção de fontes, é útil entender as grandes categorias em que as fontes disponíveis se organizam. O Capítulo~2 apresenta uma tipologia completa; aqui basta introduzir a distinção fundamental.

\subsection{Registros administrativos}

Registros administrativos são dados coletados como subproduto da operação de sistemas governamentais ou privados: cada empregado registrado na RAIS\footnote{Relação Anual de Informações Sociais, do Ministério do Trabalho e Emprego.}, cada internação no SIH\footnote{Sistema de Informações Hospitalares do SUS, gerido pelo DATASUS.}, cada aluno matriculado no Censo Escolar\footnote{Levantamento censitário anual conduzido pelo Inep junto às escolas de educação básica.}, cada beneficiário do Bolsa Família no CadÚnico\footnote{Cadastro Único para Programas Sociais do Governo Federal, instrumento de identificação e caracterização das famílias de baixa renda, gerido pelo MDS.}. Esses dados não foram coletados para fins de pesquisa, mas para gerir programas e sistemas.

As vantagens são claras: cobertura potencialmente universal da população atendida pelo sistema, baixo custo de coleta para o pesquisador (os dados já existem) e, frequentemente, possibilidade de acompanhamento longitudinal por meio de identificadores individuais (CPF, NIS, código de escola). A limitação central é que o dado registrado é aquele que o sistema precisa registrar --- não necessariamente o que o avaliador quer medir.

\subsection{Pesquisas amostrais}

Pesquisas amostrais --- como a PNAD Contínua, a POF e as pesquisas de saúde e vitimização --- coletam dados diretamente dos informantes (domicílios, pessoas, estabelecimentos) por meio de questionários estruturados aplicados a amostras representativas da população. Permitem medir variáveis que os registros administrativos não capturam: percepções, comportamentos, condições de saúde, consumo, experiências de vitimização.

A limitação central é o tamanho da amostra: pesquisas amostrais são projetadas para produzir estimativas representativas para determinados domínios geográficos e populacionais, e análises que extrapolam esses domínios produzem estimativas com erros padrão proibitivamente grandes.

\subsection{Censos}

Censos são levantamentos que cobrem toda a população de interesse, não uma amostra. O Censo Demográfico do IBGE, o Censo Escolar do Inep e o Censo Agropecuário são exemplos. Combinam a cobertura universal dos registros administrativos com a riqueza temática das pesquisas amostrais --- mas a custos de coleta muito elevados, o que os torna periodicamente irregulares.

% ────────────────────────────────────────────────────────────
\section{Critérios para avaliar a adequação de uma fonte}
% ────────────────────────────────────────────────────────────

A escolha da fonte procede em duas etapas. Primeiro, identifica-se o \textbf{tipo} de fonte potencialmente adequado à pergunta avaliativa --- registro administrativo, pesquisa amostral, censo, sistema setorial etc., conforme a tipologia do Capítulo~2. Em seguida, dentro do tipo escolhido, verifica-se se uma \textbf{fonte específica} é de fato adequada. Essa segunda etapa é organizada em torno de oito critérios.

\subsection{Cobertura geográfica}

A fonte cobre o território relevante para a pergunta? Uma pesquisa representativa para o Brasil como um todo pode não produzir estimativas confiáveis para municípios individuais ou para regiões específicas. A PNAD Contínua, por exemplo, produz estimativas trimestrais para o Brasil, grandes regiões, UFs e regiões metropolitanas --- mas não para municípios. Se a pergunta avaliativa requer análise municipal, é preciso buscar fontes com cobertura universal no nível do município (como o Censo Escolar ou a RAIS) ou construir estratégias de análise compatíveis com o que a fonte oferece.

\subsection{Periodicidade}

A periodicidade da fonte permite capturar o período relevante --- antes e depois da intervenção? Uma política implementada em 2018 que precisa ser avaliada com base em dados de 2016 e 2020 requer fontes que cubram esses dois momentos com metodologia comparável. Quebras metodológicas em séries históricas --- quando a fonte muda o questionário, a amostra ou os procedimentos de coleta entre as edições --- podem impossibilitar comparações ao longo do tempo.

\subsection{Unidade de análise}

A unidade de análise da fonte é compatível com a da pergunta? Se a pergunta é sobre alunos individuais, é preciso microdados de alunos (como o Censo Escolar), não dados agregados por escola ou município. Se a pergunta é sobre \textit{desempenho médio por escola}, tanto o SAEB quanto o Censo Escolar oferecem agregações adequadas --- mas agregar microdados de aluno do Censo Escolar preserva mais flexibilidade analítica do que trabalhar com médias pré-calculadas. Trabalhar com unidades de análise incorretas leva a falácias ecológicas e erros de inferência.

\subsection{Representatividade}

Para pesquisas amostrais: a amostra é representativa para o domínio de interesse? Verificar os domínios de representatividade declarados pela pesquisa antes de qualquer análise. Para registros administrativos: a cobertura do sistema é suficiente para os fins da análise, ou há exclusão sistemática de subgrupos relevantes?

\subsection{Granularidade temática}

A fonte contém as variáveis de que o analista precisa --- ou variáveis a partir das quais elas podem ser construídas? Um estudo sobre o impacto da escolaridade materna sobre saúde infantil precisa de uma fonte que contenha simultaneamente informações sobre a mãe e sobre a criança. Verificar os dicionários de variáveis antes de comprometer o desenho de avaliação com uma fonte específica.

\subsection{Identificadores para integração}

A fonte tem identificadores que permitem cruzar com outras bases de que o analista precisa? CPF, NIS, CNPJ, código de escola do Inep e código de estabelecimento de saúde (CNES) são os principais identificadores que viabilizam integrações no ecossistema de dados brasileiro. A disponibilidade desses identificadores --- e as restrições de acesso que os acompanham --- é uma dimensão crítica na fase de planejamento. Um exemplo concreto ilustra o valor desses identificadores: o código CNES identifica univocamente cada estabelecimento de saúde e está presente tanto no cadastro do CNES propriamente (com informações de infraestrutura e recursos humanos) quanto no SIH (com informações de internações e mortalidade hospitalar), permitindo, por exemplo, relacionar a oferta de leitos de UTI em cada hospital à mortalidade hospitalar observada --- uma análise impossível sem o identificador compartilhado.

\subsection{Qualidade do preenchimento}

Qual é a taxa de não-resposta para as variáveis de interesse? Há subregistro sistemático --- concentrado em determinadas regiões, grupos populacionais ou períodos? O SIM, por exemplo, tem subregistro de óbitos mais elevado no Norte e Nordeste; o SINAN tem subnotificação estrutural que varia muito por tipo de agravo. Ignorar esses padrões de subregistro diferencial leva a conclusões enganosas sobre desigualdades regionais e populacionais.

\subsection{Documentação e historicidade}

A fonte tem documentação metodológica que permita entender o que cada variável mede de fato? Há quebras metodológicas documentadas que afetam comparações ao longo do tempo? Fontes com documentação adequada reduzem o risco de erros de interpretação; fontes com histórico longo permitem análises de tendências e construção de contrafactuais mais robustos.

% CORRIGIDO: título encurtado para caber na caixa
\begin{boxdica}{Dica — Consulte sempre a documentação metodológica}
Todo grande produtor de dados brasileiro publica notas metodológicas, questionários e dicionários de variáveis. O IBGE publica documentação detalhada para a PNAD Contínua, o Censo Demográfico e todas as suas pesquisas. O Inep publica os cadernos de metodologia do Censo Escolar, do SAEB e do ENEM. O DATASUS disponibiliza os manuais de instrução de preenchimento do SIM, do SINASC e dos demais sistemas. Ler essa documentação antes de iniciar a análise é um investimento que poupa erros custosos.
\end{boxdica}

% ────────────────────────────────────────────────────────────
\section{O ecossistema de dados brasileiro: uma visão panorâmica}
% ────────────────────────────────────────────────────────────

O Brasil dispõe de um ecossistema de dados para avaliação de políticas públicas notavelmente rico, organizado em torno de três grandes pilares institucionais.

\textbf{O IBGE} é o produtor central do sistema estatístico nacional. Conduz o Censo Demográfico (decenal), a PNAD Contínua (trimestral/anual), a POF, o Censo Agropecuário e diversas outras pesquisas temáticas. Produz também as malhas territoriais e setores censitários que servem de base para análises espaciais. Para qualquer avaliação que precise de dados representativos da população brasileira em suas dimensões socioeconômicas, demográficas e de condições de vida, o ponto de partida é invariavelmente o IBGE.

\textbf{Os ministérios e sistemas setoriais} produzem os registros administrativos que documentam a operação das políticas públicas nas diferentes áreas. O Ministério da Saúde, via DATASUS, mantém o SIM, o SINASC, o SIH, o SINAN, o CNES e outros sistemas. O Ministério da Educação, via Inep, conduz o Censo Escolar, o SAEB e o ENEM. O Ministério do Trabalho gerencia a RAIS e o CAGED. O Ministério do Desenvolvimento e Assistência Social opera o CadÚnico e os dados do Bolsa Família. A Secretaria do Tesouro Nacional mantém o Siconfi/FINBRA. Cada um desses sistemas é descrito em detalhe nos capítulos temáticos deste guia.

\textbf{As agências reguladoras e outros produtores de dados especializados} cobrem setores específicos com dados que não estão disponíveis no sistema estatístico central: a ANEEL para energia elétrica, a ANATEL para telecomunicações, o INPE para dados de desmatamento, o TSE para dados eleitorais e de financiamento de campanha.

\begin{table}[H]
\centering
\caption{Visão panorâmica do ecossistema de dados para avaliação de políticas públicas}
\label{tab:ecossistema_geral}
\small
\begin{tabularx}{\textwidth}{@{} p{3.5cm} p{3cm} X @{}}
\toprule
\textbf{Pilar} & \textbf{Principais fontes} & \textbf{Uso central em avaliação} \\
\midrule
IBGE           & Censo Demográfico, PNAD-C, POF, Censo Agropecuário & Perfil populacional, condições de vida, representatividade nacional \\
Saúde (MS/DATASUS) & SIM, SINASC, SIH, SINAN, CNES & Mortalidade, nascimentos, internações, infraestrutura de saúde \\
Educação (Inep) & Censo Escolar, SAEB, ENEM & Matrículas, desempenho cognitivo, acesso ao ensino \\
Trabalho (MTE)  & RAIS, CAGED & Emprego formal, salários, rotatividade \\
Proteção social (MDS) & CadÚnico, Bolsa Família, BPC & Famílias vulneráveis, transferências de renda \\
Finanças públicas (STN) & Siconfi/FINBRA, Portal da Transparência & Gasto público, execução orçamentária \\
Agências reguladoras & ANEEL, ANATEL, ANS, ANTT & Infraestrutura setorial, cobertura, qualidade \\
Outros produtores de dados & TSE, INPE, SNIS, Ipea & Eleitoral, ambiental, saneamento, indicadores derivados \\
\bottomrule
\end{tabularx}
\fonte{FGV CLEAR, elaboração própria.}
\end{table}

% ────────────────────────────────────────────────────────────
\section{Complementaridade e integração entre fontes}
% ────────────────────────────────────────────────────────────

Uma das características mais importantes --- e mais subutilizadas --- do ecossistema de dados brasileiro é a possibilidade de \textbf{cruzar bases de dados} por meio de identificadores individuais. O CPF vincula pessoas físicas em qualquer base que o contenha. O NIS conecta beneficiários de programas sociais ao CadÚnico. O código de escola do Inep está presente no Censo Escolar e no SAEB. O CNES identifica estabelecimentos de saúde no SIH, no SIA e no SINAN.

Essas integrações permitem construir perfis multidimensionais que nenhuma fonte isolada poderia fornecer. Um exemplo clássico: cruzar o CadÚnico com o Censo Escolar pelo número de identificação do aluno permite analisar simultaneamente as condições de vida da família (renda, composição familiar, acesso a serviços) e o desempenho e a trajetória escolar das crianças. Cruzar a RAIS com dados de mortalidade do SIM pelo CPF permite estimar o impacto da formalização do emprego sobre expectativa de vida.

A capacidade de fazer essas integrações é, em muitos casos, o que torna avaliações causais rigorosas possíveis no Brasil. O Capítulo~13 discute as boas práticas e as obrigações legais que acompanham o uso de integrações de bases.

% ────────────────────────────────────────────────────────────
\section{Síntese do capítulo}
% ────────────────────────────────────────────────────────────

Este capítulo estabeleceu os fundamentos que norteiam todo o guia. Os pontos centrais são:

\begin{itemize}
  \item A escolha da fonte de dados é uma \textbf{decisão metodológica} com consequências diretas sobre a validade das conclusões: não existe fonte perfeita, mas existem fontes mais ou menos adequadas para cada pergunta avaliativa;

  \item As fontes se organizam em três grandes categorias --- \textbf{registros administrativos}, \textbf{pesquisas amostrais} e \textbf{censos} --- com perfis de vantagens e limitações complementares que o avaliador precisa conhecer;

  \item Oito critérios estruturam a avaliação de adequação de uma fonte: cobertura geográfica, periodicidade, unidade de análise, representatividade, granularidade temática, identificadores para integração, qualidade do preenchimento e documentação metodológica;

  \item O ecossistema brasileiro é organizado em torno do \textbf{IBGE}, dos \textbf{sistemas setoriais} dos ministérios e das \textbf{agências reguladoras e outros produtores de dados especializados} --- cada pilar com fontes distintas e complementares;

  \item A \textbf{integração entre bases} por identificadores individuais é um dos ativos mais poderosos do ecossistema brasileiro, viabilizando análises multidimensionais que nenhuma fonte isolada permitiria.
\end{itemize}

\medskip

O capítulo seguinte aprofunda a tipologia das fontes --- pesquisas amostrais, censos, registros administrativos, sistemas setoriais, dados fiscais, dados territoriais, fontes internacionais e APIs --- antes que os capítulos temáticos cubram fonte por fonte cada área de política pública.
% ============================================================
%  cap02_tipologia.tex
%  Capítulo 2: Tipologia das fontes de dados
% ============================================================

% ════════════════════════════════════════════════════════════
\chapter{Tipologia das Fontes de Dados}
% ════════════════════════════════════════════════════════════

\begin{boxdestaque}{Neste capítulo}
Antes de mergulhar nos catálogos temáticos, este capítulo organiza o ecossistema de fontes de dados disponíveis para avaliação de políticas públicas no Brasil. Apresentamos uma tipologia em sete categorias — pesquisas amostrais, censos, registros administrativos, sistemas de informação setoriais, dados fiscais, dados territoriais e fontes internacionais — e discutimos as características, vantagens e limitações de cada tipo. O objetivo é oferecer uma bússola conceitual que oriente a leitura dos capítulos seguintes.
\end{boxdestaque}

% ────────────────────────────────────────────────────────────
\section{Por que tipologizar?}
% ────────────────────────────────────────────────────────────

O universo de fontes de dados disponíveis para quem trabalha com políticas públicas no Brasil é vasto e heterogêneo. Há pesquisas domiciliares representativas para todo o país, registros administrativos produzidos rotineiramente pela gestão pública, sistemas de informação setoriais com décadas de histórico, bases fiscais municipais e estaduais, dados georreferenciados e um conjunto crescente de fontes internacionais comparadas. A cada ano, novas bases são disponibilizadas, APIs são abertas e ferramentas digitais ampliam o acesso a dados antes restritos.

Navegar nesse ecossistema sem uma estrutura conceitual é ineficiente e arriscado. A tipologia apresentada neste capítulo tem três funções práticas: primeiro, ajuda o avaliador a saber \textit{onde procurar} quando tem uma pergunta em mente; segundo, permite antecipar as características e limitações de uma fonte antes mesmo de consultá-la; terceiro, facilita a escolha de combinações complementares de dados para triangulação metodológica.

A tipologia não é exaustiva nem mutuamente exclusiva — algumas fontes se encaixam em mais de uma categoria. Ela é, antes de tudo, uma ferramenta de orientação.

% ────────────────────────────────────────────────────────────
\section{Pesquisas amostrais domiciliares}
% ────────────────────────────────────────────────────────────

As pesquisas amostrais domiciliares são levantamentos periódicos conduzidos junto a amostras probabilísticas de domicílios e seus moradores. No Brasil, o IBGE é o principal produtor dessas pesquisas, embora ministérios e institutos de pesquisa também conduzam levantamentos específicos.

\textbf{O que as caracteriza.} Essas pesquisas são projetadas para produzir estimativas representativas de uma população-alvo — em geral, a população residente em domicílios particulares permanentes do Brasil ou de recortes geográficos específicos. A representatividade é garantida pelo desenho amostral probabilístico, que atribui a cada domicílio uma probabilidade conhecida de seleção. Os resultados devem ser interpretados com os pesos amostrais fornecidos pelo produtor.

\textbf{Vantagens.} Por cobrirem amostras grandes e demograficamente diversas, essas pesquisas permitem estimar indicadores para múltiplos grupos populacionais e regiões. Têm questionários padronizados e documentação metodológica detalhada, o que facilita comparações ao longo do tempo. Muitas disponibilizam microdados públicos — arquivos com informações individuais anonimizadas — que permitem análises customizadas.

\textbf{Limitações.} O tamanho da amostra impõe limites à desagregação: estimativas para municípios pequenos, grupos étnico-raciais específicos ou subgrupos populacionais raramente têm precisão estatística adequada. Além disso, a periodicidade — trimestral, anual ou quinquenal — impede o acompanhamento de fenômenos de curto prazo. Há também o risco de não-resposta diferencial, que pode introduzir vieses se certos grupos de domicílios recusarem participar mais do que outros.

\textbf{Principais exemplos no Brasil.} PNAD Contínua, Pesquisa de Orçamentos Familiares (POF), Pesquisa Nacional de Saúde (PNS), Pesquisa Mensal de Emprego (descontinuada em 2016).

\begin{boxdica}{Dica — Sempre use os pesos amostrais}
Toda pesquisa amostral domiciliar do IBGE fornece uma variável de peso amostral (\texttt{V1028} na PNAD Contínua, por exemplo). Ignorar esses pesos ao calcular médias, proporções ou totais produz estimativas incorretas. Softwares como R (\texttt{survey}), Stata (\texttt{svyset}) e Python (\texttt{statsmodels}) têm rotinas específicas para análise de dados complexos com pesos.
\end{boxdica}

% ────────────────────────────────────────────────────────────
\section{Censos}
% ────────────────────────────────────────────────────────────

Censos são levantamentos que buscam cobrir \textit{toda} a população ou universo de interesse — não uma amostra, mas o conjunto completo de unidades. No Brasil, os principais censos relevantes para avaliação de políticas públicas são o Censo Demográfico (IBGE), o Censo Escolar (Inep), o Censo Agropecuário (IBGE), a Pesquisa de Assistência Médico-Sanitária (AMS, IBGE) --- de caráter censitário para estabelecimentos de saúde --- e o Censo SUAS (MDS), que cobre os equipamentos da assistência social (CRAS, CREAS e unidades de acolhimento).

\textbf{O que os caracteriza.} Por visarem cobertura universal, os censos permitem desagregações impossíveis em pesquisas amostrais: estimativas para municípios pequenos, bairros, zonas rurais específicas ou grupos populacionais minoritários. O Censo Escolar, por exemplo, registra individualmente cada aluno, cada professor e cada escola do Brasil, permitindo análises extremamente granulares.

\textbf{Vantagens.} A cobertura praticamente total elimina erros de amostragem. Para análises territoriais finas — como identificar onde concentrar uma intervenção — censos são frequentemente insubstituíveis. Permitem também a construção de painéis longitudinais quando os identificadores individuais estão disponíveis (como no Censo Escolar, onde cada aluno tem um código único).

\textbf{Limitações.} A principal limitação dos censos demográficos é a periodicidade longa: o Censo Demográfico é decenal, o que significa que pode haver uma defasagem de até dez anos entre a coleta e o momento da avaliação. Mesmo o Censo Escolar, anual, tem uma defasagem de alguns meses entre o período de referência (geralmente maio) e a publicação dos microdados. Além disso, censos estão sujeitos a subregistro: domicílios em áreas de difícil acesso, populações em situação de rua e comunidades isoladas tendem a ser sub-representados.

\textbf{Principais exemplos no Brasil.} Censo Demográfico (decenal), Censo Escolar (anual), Censo da Educação Superior (anual), Censo Agropecuário (periodicidade irregular), AMS (IBGE) e Censo SUAS (MDS).

% ────────────────────────────────────────────────────────────
\section{Registros administrativos}
% ────────────────────────────────────────────────────────────

Registros administrativos são dados produzidos como subproduto da operação rotineira de serviços e programas públicos. Não foram coletados com fins analíticos, mas contêm informações valiosas sobre a população atendida, os serviços prestados e os recursos utilizados.

\textbf{O que os caracteriza.} Ao contrário das pesquisas, os registros administrativos não têm um questionário padronizado aplicado a uma amostra — eles registram transações, eventos ou cadastros à medida que ocorrem. Um nascimento registrado no SINASC, uma admissão no CAGED, uma matrícula no Censo Escolar: cada um desses é um registro administrativo.

\textbf{Vantagens.} Frequência e cobertura são as principais. O CAGED, por exemplo, é publicado mensalmente e cobre todos os movimentos de emprego formal no país. O Censo Escolar captura cada matrícula individualmente. Essa granularidade e frequência permitem análises que seriam impossíveis com pesquisas amostrais. Além disso, registros administrativos costumam ter baixo custo de acesso — em geral estão disponíveis gratuitamente em portais do governo.

\textbf{Limitações.} A limitação estrutural dos registros administrativos é que eles cobrem apenas o que passa pelo sistema formal. Trabalhadores informais não estão no CAGED. Óbitos sem atendimento médico em áreas remotas podem não chegar ao SIM. Crianças fora da escola não constam no Censo Escolar. Essa é uma forma de viés de cobertura que pode ser especialmente problemática em avaliações de políticas voltadas a populações vulneráveis, que tendem a ter menor presença nos sistemas formais.

Outra limitação importante é a variabilidade na qualidade de preenchimento: como os dados são inseridos por diferentes operadores em milhares de municípios, erros de digitação, campos em branco e inconsistências são frequentes e precisam ser tratados antes da análise.

\textbf{Principais exemplos no Brasil.} CAGED, RAIS, Cadastro Único, Censo Escolar, SINASC, SIM, SIH/SIA, SINAN.

% CORRIGIDO: título encurtado para caber na caixa
\begin{boxatencao}{Atenção — Viés nos registros administrativos}
A qualidade de um registro administrativo depende dos incentivos institucionais de quem o preenche. Secretarias municipais de saúde podem sub-registrar óbitos por causas sensíveis politicamente. Escolas podem registrar alunos ausentes como presentes para evitar penalidades. Esse viés de notificação é estrutural e deve ser considerado na interpretação dos dados. Por outro lado, o arranjo federativo brasileiro cria contrapesos importantes: boa parte das transferências federais para municípios --- no SUS, na educação básica e na assistência social --- está condicionada ao preenchimento regular dos sistemas de informação correspondentes. Esse desenho institucional aumenta os incentivos para registro e tende a mitigar os vieses descritos acima, sobretudo nos sistemas mais maduros. Ainda assim, sempre que possível, compare a fonte administrativa com fontes independentes para checar consistência.
\end{boxatencao}

% ────────────────────────────────────────────────────────────
\section{Sistemas de informação setoriais}
% ────────────────────────────────────────────────────────────

Sistemas de informação setoriais são plataformas institucionais que integram múltiplos registros administrativos de uma área de política pública específica. Diferem dos registros administrativos isolados por terem uma arquitetura de gestão da informação mais estruturada e por produzirem estatísticas consolidadas regularmente.

\textbf{O que os caracteriza.} No Brasil, os principais setores de política pública têm sistemas próprios: o DATASUS na saúde, o Inep na educação, o MDS na assistência social. Esses sistemas agregam dados de produção de serviços, infraestrutura, recursos humanos e resultados em plataformas de consulta pública.

\textbf{Vantagens.} Oferecem acesso centralizado a múltiplas dimensões de um setor. O DATASUS, por exemplo, reúne em um só portal sistemas de mortalidade, nascimentos, internações, ambulatório, vigilância epidemiológica e cadastro de estabelecimentos. Isso facilita a construção de análises multidimensionais sem necessidade de acessar fontes dispersas.

\textbf{Limitações.} A integração nem sempre é plena: sistemas diferentes dentro de um mesmo setor frequentemente não compartilham identificadores comuns, dificultando a integração entre bases. Além disso, a qualidade dos dados varia entre municípios e estados, refletindo diferenças na capacidade administrativa local --- ainda que, em sistemas consolidados como os do DATASUS e do Censo Escolar, essa heterogeneidade venha diminuindo ao longo do tempo.

\textbf{Principais exemplos no Brasil.} DATASUS (saúde), Inep/sistemas educacionais, SUAS/MDS (assistência social), SNIS (saneamento).

% ────────────────────────────────────────────────────────────
\section{Dados fiscais e orçamentários}
% ────────────────────────────────────────────────────────────

Dados fiscais registram receitas, despesas, transferências e execução orçamentária de entes públicos. São fundamentais para análises de eficiência, custo-efetividade e alocação de recursos em políticas públicas.

\textbf{O que os caracteriza.} No Brasil, esses dados são produzidos pela obrigação constitucional de transparência fiscal. Municípios, estados e o governo federal devem publicar suas contas em portais de transparência e reportar dados ao Tesouro Nacional. O FINBRA (Finanças do Brasil) consolida os dados municipais; o SIOPE e o SIOPS fazem o mesmo para gastos em educação e saúde, respectivamente.

\textbf{Vantagens.} Permitem verificar quanto cada ente gasta em cada função de governo, comparar alocações entre municípios e calcular métricas de eficiência como custo por beneficiário ou gasto por aluno. São essenciais em análises de custo-benefício e custo-efetividade.

\textbf{Limitações.} A classificação orçamentária nem sempre reflete com precisão o que foi efetivamente gasto em cada atividade. Um município pode classificar despesas de gestão sob a função educação, inflando artificialmente o gasto per capita. Além disso, a qualidade do preenchimento varia muito entre municípios, e dados de municípios pequenos ou com baixa capacidade administrativa tendem a ter mais inconsistências.

\textbf{Principais exemplos no Brasil.} FINBRA/Siconfi, SIOPE, SIOPS, Portais de Transparência (federal, estaduais e municipais), dados do TSE sobre financiamento eleitoral.

% ────────────────────────────────────────────────────────────
\section{Dados georreferenciados e territoriais}
% ────────────────────────────────────────────────────────────

Dados georreferenciados associam informações a coordenadas espaciais, permitindo análises que levam em conta a dimensão territorial das políticas públicas.

\textbf{O que os caracteriza.} Incluem malhas territoriais (limites de municípios, estados, setores censitários), dados populacionais por setor censitário, mapas de cobertura de serviços públicos, imagens de satélite e índices territoriais. No Brasil, o IBGE é o principal produtor de dados georreferenciados para uso em políticas públicas.

\textbf{Vantagens.} Permitem identificar padrões espaciais de cobertura e exclusão, mapear onde os problemas se concentram, calcular distâncias entre populações e serviços e construir análises de spillover entre territórios vizinhos. Em avaliações de impacto, a dimensão geográfica pode ser usada como fonte de variação exógena (como em regressões descontínuas geográficas).

\textbf{Limitações.} A análise espacial exige competências técnicas específicas (SIG, R espacial, Python com GeoPandas) que nem todas as equipes de avaliação dominam. Além disso, a granularidade dos dados pode não coincidir com a unidade de implementação da política — dados disponíveis por município podem ser insuficientes para avaliar uma política implementada em nível de bairro.

\textbf{Principais exemplos no Brasil.} Malhas territoriais do IBGE, setores censitários, PRODES/INPE (desmatamento), dados georreferenciados do Cadastro Ambiental Rural (CAR).

% ────────────────────────────────────────────────────────────
\section{Fontes internacionais e comparadas}
% ────────────────────────────────────────────────────────────

Fontes internacionais produzem dados comparáveis entre países, permitindo contextualizar indicadores brasileiros em perspectiva global ou regional.

\textbf{O que as caracteriza.} São produzidas por organismos multilaterais (Banco Mundial, ONU, OCDE, OIT, OMS) ou institutos de pesquisa internacionais. Em geral, harmonizam dados de diferentes países para permitir comparações, o que implica algum grau de padronização que pode não capturar especificidades nacionais com total fidelidade.

\textbf{Vantagens.} Permitem situar o desempenho do Brasil em perspectiva comparada, identificar países com trajetórias similares e aprender com experiências internacionais. Também são úteis quando dados nacionais estão indisponíveis ou têm qualidade insuficiente para determinados indicadores.

\textbf{Limitações.} A harmonização metodológica necessária para comparações internacionais frequentemente implica perdas de detalhe. Além disso, os dados costumam ter maior defasagem do que as fontes nacionais originais, pois dependem de processos de coleta e padronização mais longos.

\textbf{Principais exemplos.} World Development Indicators (Banco Mundial), bases da OCDE (saúde, educação, bem-estar), PISA (Inep/OCDE), bases da CEPAL, OIT ILOSTAT, OMS Global Health Observatory.

% ────────────────────────────────────────────────────────────
\section{Dados produzidos com apoio de tecnologia: APIs,
         \textit{web scraping} e inteligência artificial}
% ────────────────────────────────────────────────────────────

Uma categoria crescente e ainda pouco sistematizada é a de dados obtidos ou estruturados com apoio de tecnologias digitais — interfaces de programação (APIs), coleta automatizada de páginas web (\textit{web scraping}) e ferramentas de inteligência artificial.

\textbf{O que os caracteriza.} Não são fontes no sentido tradicional — um órgão produtor com um questionário e uma metodologia documentada — mas formas de acessar ou produzir dados a partir de informações já existentes em formato digital. Portais de transparência, diários oficiais, sistemas de informação governamental e redes sociais são exemplos de fontes de dados não estruturados que podem ser acessados por essas vias.

\textbf{Vantagens.} Permitem obter dados em tempo real ou com alta frequência, acessar informações que não estão disponíveis em formato tabulado e automatizar coletas que seriam inviáveis manualmente. APIs governamentais abertas — como as do IBGE, do Banco Central e do Portal de Dados Abertos — democratizaram o acesso a dados que antes exigiam download manual de dezenas de arquivos.

\textbf{Limitações.} Exigem competências técnicas de programação. Dados obtidos por \textit{scraping} podem ter qualidade variável e estão sujeitos a mudanças na estrutura dos sites-fonte. O uso de IA para estruturar ou sintetizar dados introduz riscos de erros não detectados (``alucinações'') que exigem verificação cuidadosa.

\textbf{Principais exemplos.} API do IBGE (Sidra), API do Banco Central (SGS), Portal de Dados Abertos do Governo Federal (\url{dados.gov.br}), diários oficiais municipais e estaduais.

O Capítulo~12 trata essa categoria com maior profundidade.

% ────────────────────────────────────────────────────────────
\section{A fronteira entre primário e secundário revisitada}
% ────────────────────────────────────────────────────────────

Vale retomar, à luz dessa tipologia, uma distinção que parece simples mas tem nuances importantes. O Censo Escolar, por exemplo, é um registro administrativo produzido pelo Inep para fins de gestão e alocação de recursos — mas é amplamente utilizado como dado de pesquisa. A RAIS foi criada como instrumento de gestão do mercado de trabalho formal, mas tornou-se uma das bases mais usadas em pesquisas econômicas no Brasil. Esses dados são ``primários'' para o Inep e o Ministério do Trabalho, mas ``secundários'' para o avaliador externo que os usa.

O que importa, do ponto de vista do avaliador, não é a classificação abstrata, mas o conhecimento das condições de produção de cada fonte: quem coletou, para que fim, com que método e com que incentivos. Essa consciência é o que permite usar dados administrativos com o rigor que uma avaliação de qualidade exige.

\begin{table}[H]
\centering
\caption{Síntese da tipologia de fontes de dados}
\label{tab:tipologia}
\small
\begin{tabularx}{\textwidth}{@{} p{3.2cm} X p{2.8cm} p{3cm} @{}}
\toprule
\textbf{Tipo} & \textbf{Características principais} & \textbf{Periodicidade típica} & \textbf{Exemplos} \\
\midrule
Pesquisas amostrais & Representativas, microdados públicos, pesos amostrais & Trimestral a quinquenal & PNAD-C, POF, PNS \\
\addlinespace
Censos & Cobertura universal, alta desagregação & Anual a decenal & Censo Demográfico, Censo Escolar \\
\addlinespace
Registros administrativos & Alta frequência, cobertura do sistema formal & Mensal a anual & CAGED, RAIS, CadÚnico \\
\addlinespace
Sistemas setoriais & Integração temática, consulta pública & Contínua a anual & DATASUS, Inep, SNIS \\
\addlinespace
Dados fiscais & Receitas e despesas por ente federativo & Anual & FINBRA, SIOPE, Transparência \\
\addlinespace
Dados territoriais & Coordenadas espaciais, malhas, satélite & Variável & IBGE/Setores, PRODES \\
\addlinespace
Fontes internacionais & Comparabilidade entre países & Anual & WDI, OCDE, PISA \\
\addlinespace
APIs e tecnologia & Automação, tempo real, não estruturados & Contínua & IBGE API, dados.gov.br \\
\bottomrule
\end{tabularx}
\fonte{FGV CLEAR, elaboração própria.}
\end{table}

% ────────────────────────────────────────────────────────────
\section{Síntese do capítulo}
% ────────────────────────────────────────────────────────────

Este capítulo organizou o ecossistema de fontes de dados em oito tipos, cada um com características, vantagens e limitações próprias. A Tabela~\ref{tab:tradeoffs_tipos} sintetiza esses trade-offs.

\begin{table}[H]
\centering
\caption{Trade-offs entre os principais tipos de fontes}
\label{tab:tradeoffs_tipos}
\small
\begin{tabularx}{\textwidth}{@{} >{\raggedright\arraybackslash}p{3.2cm} X X @{}}
\toprule
\textbf{Tipo} & \textbf{Vantagem central} & \textbf{Limitação central} \\
\midrule
Pesquisas amostrais     & Medem variáveis que nenhum registro administrativo captura (percepções, comportamentos, condições de vida) & Tamanho de amostra limita desagregação territorial e por subgrupos \\
Censos                  & Cobertura universal; permitem análises territoriais finas                                                  & Periodicidade longa e defasagem entre coleta e disponibilização \\
Registros administrativos & Alta frequência, granularidade individual, baixo custo                                                    & Cobrem apenas quem passa pelo sistema formal; qualidade heterogênea \\
Sistemas setoriais      & Integração temática, consulta pública, padronização entre unidades federativas                            & Nem sempre integram identificadores entre bases do mesmo sistema \\
Dados fiscais           & Permitem análises de eficiência, custo-efetividade e custo-benefício                                       & Classificação orçamentária pode não refletir o gasto efetivo por atividade \\
Dados territoriais      & Resolução espacial fina; múltiplas fontes complementares                                                    & Exigem competências de geoprocessamento \\
Fontes internacionais   & Comparabilidade entre países                                                                                & Harmonização reduz detalhe nacional \\
APIs e tecnologia       & Tempo real, automação, acesso a dados não estruturados                                                      & Exigem competências técnicas; qualidade variável \\
\bottomrule
\end{tabularx}
\fonte{FGV CLEAR, elaboração própria.}
\end{table}

\noindent Em síntese:

\begin{itemize}
  \item Pesquisas amostrais oferecem representatividade estatística, mas têm limites de desagregação; censos oferecem cobertura universal, mas periodicidade longa.
  \item Registros administrativos têm alta frequência e granularidade, mas cobrem apenas o que passa pelo sistema formal e dependem da qualidade do preenchimento.
  \item Sistemas setoriais integram múltiplas dimensões de uma área de política, mas a qualidade varia entre municípios.
  \item Dados fiscais são essenciais para análises de eficiência, mas a classificação orçamentária pode não refletir com fidelidade os gastos reais por atividade.
  \item Dados georreferenciados abrem dimensões analíticas importantes, mas exigem competências técnicas específicas.
  \item Fontes internacionais permitem contextualização comparada, mas implicam perdas de detalhe pela necessidade de harmonização.
  \item APIs e tecnologias digitais ampliam o acesso a dados, mas exigem cuidado metodológico adicional.
\end{itemize}

Os próximos capítulos percorrem esse território fonte por fonte, começando pelas categorias mais consolidadas e avançando para as mais especializadas. A leitura pode ser feita em ordem ou por saltos temáticos.

\medskip

O capítulo seguinte abre o catálogo pelas pesquisas amostrais e censos do IBGE: PNAD Contínua, POF, PNS, Censo Demográfico --- a base do diagnóstico socioeconômico brasileiro.
% ============================================================
%  cap03_pesquisas_amostrais.tex  [REVISADO]
%  Capítulo 3: Pesquisas Amostrais e Censos Demográficos
% ============================================================

% ════════════════════════════════════════════════════════════
\chapter{Pesquisas Amostrais e Censos Demográficos}
% ════════════════════════════════════════════════════════════

\begin{boxdestaque}{Neste capítulo}
Este capítulo apresenta as principais pesquisas amostrais domiciliares e os censos demográficos disponíveis para o Brasil. Para cada fonte, descrevemos seu histórico, estrutura, cobertura, principais variáveis, forma de acesso aos microdados e aplicações típicas em avaliação de políticas públicas. O capítulo cobre o Censo Demográfico, a PNAD Contínua, as edições anteriores da PNAD, a Pesquisa Mensal de Emprego, a Pesquisa de Orçamentos Familiares e o Censo Agropecuário.
\end{boxdestaque}

% ────────────────────────────────────────────────────────────
\section{Censo Demográfico}
% ────────────────────────────────────────────────────────────

\ficharapida
  {Censo Demográfico}
  {Instituto Brasileiro de Geografia e Estatística (IBGE)}
  {https://www.ibge.gov.br/estatisticas/sociais/populacao/22827-censo-demografico-2022.html}
  {Brasil, Grandes Regiões, Unidades da Federação, municípios, distritos, subdistritos e setores censitários}
  {Decenal: 1872, 1890, 1900, 1920, 1940, 1950, 1960, 1970, 1980, 1991, 2000, 2010, 2022}
  {Domicílio e pessoa}

\subsection{O que é e por que importa}

O Censo Demográfico é o mais abrangente levantamento de informações sobre a população brasileira. Realizado pelo IBGE desde 1940 com periodicidade decenal --- com exceção do intervalo de onze anos entre 2010 e 2022, provocado pela pandemia de COVID-19 ---, o Censo visita \textit{todos} os domicílios do território nacional, tornando-o a única fonte capaz de produzir informações confiáveis para municípios pequenos, áreas rurais remotas, comunidades indígenas e quilombolas, e qualquer outro recorte territorial ou demográfico para o qual pesquisas amostrais não têm representatividade.

Para a avaliação de políticas públicas, o Censo tem um papel insubstituível como fonte de \textbf{linha de base} --- especialmente quando a política foi implementada entre dois censos --- e como base para a \textbf{construção de contrafactuais} em métodos como o pareamento por escore de propensão. É também amplamente usado como fonte de \textbf{variáveis de controle} em desenhos quase-experimentais, já que oferece informações socioeconômicas granulares e temporalmente estáveis que podem ser incorporadas como covariáveis em estratégias como pareamento, diferença-em-diferenças e regressão descontínua. Dados censitários também são amplamente utilizados para calcular taxas populacionais (mortalidade infantil, analfabetismo, cobertura de saneamento) que exigem denominadores populacionais precisos.

\subsection{Estrutura da coleta: questionário básico e amostral}

Desde 1960, o Censo utiliza dois instrumentos de coleta com abrangências distintas:

\begin{itemize}
  \item \textbf{Questionário básico}: aplicado a todos os domicílios do país. Coleta informações elementares sobre os moradores --- nome, sexo, idade, relação com o responsável pelo domicílio e rendimento nominal mensal.

  \item \textbf{Questionário amostral}: aplicado a uma fração dos domicílios (aproximadamente 10\% a 20\%, dependendo do porte do município). Com a aplicação dos pesos amostrais fornecidos pelo IBGE, é representativo no nível municipal, o que amplia consideravelmente o leque de análises locais viáveis a partir do Censo Demográfico. Contém um conjunto muito mais amplo de variáveis, incluindo:
  \begin{itemize}
    \item Características sociodemográficas: escolaridade, raça/cor, religião, migração, fecundidade, nupcialidade;
    \item Características do domicílio: material das paredes, tipo de piso, existência de água encanada, esgotamento sanitário, coleta de lixo, energia elétrica;
    \item Posse de bens: geladeira, automóvel, computador, acesso à internet;
    \item Trabalho e rendimento: ocupação, setor de atividade, renda de todas as fontes.
  \end{itemize}
\end{itemize}

\begin{boxatencao}{Atenção --- Censo 2022: o que mudou}
O Censo 2022 trouxe mudanças metodológicas relevantes em relação a 2010. A coleta passou a ser realizada em dispositivos móveis (tablets), com questionários digitais. Foram incluídas novas perguntas sobre identidade de gênero e orientação sexual, pessoas com deficiência (seguindo o modelo do Washington Group) e uso de internet. A variável de renda passou por revisão conceitual. Ao comparar dados de 2022 com edições anteriores, verifique as notas metodológicas do IBGE para cada variável de interesse.
\end{boxatencao}

\subsection{Microdados: estrutura e acesso}

Os microdados do Censo são disponibilizados publicamente pelo IBGE em duas bases separadas:

\begin{itemize}
  \item \textbf{Base de domicílios}: cada linha é um domicílio; contém variáveis sobre características físicas e de infraestrutura.
  \item \textbf{Base de pessoas}: cada linha é um morador; contém variáveis sociodemográficas e econômicas individuais.
\end{itemize}

As duas bases são linkáveis pelo código do domicílio, permitindo análises que combinam características individuais e domiciliares. Os microdados do questionário amostral estão disponíveis separados por Unidade da Federação, o que exige concatenação para análises nacionais.

O acesso é gratuito e não requer cadastro. Os arquivos estão em formato fixo (FWD) ou CSV nas edições mais recentes, acompanhados de dicionário de variáveis e documentação metodológica.

\begin{boxdica}{Dica --- Pacotes para trabalhar com microdados do Censo}
\begin{itemize}
  \item \textbf{R}: o pacote \texttt{censobr} (disponível no CRAN) facilita o download e a leitura dos microdados do Censo 2022 diretamente no ambiente R.
  \item \textbf{Python}: a biblioteca \texttt{censobrapi} e o pacote \texttt{pandas} permitem manipular os arquivos CSV.
  \item \textbf{Stata}: os microdados podem ser importados com \texttt{infix} (formato fixo) ou \texttt{import delimited} (CSV), usando o dicionário de variáveis como referência.
\end{itemize}
\end{boxdica}

\subsection{Aplicações típicas em avaliação}

\begin{itemize}
  \item Construção de \textbf{indicadores de linha de base} para políticas implementadas no período intercensitário;
  \item \textbf{Caracterização do público-alvo} de políticas focalizadas;
  \item \textbf{Denominadores populacionais} para calcular taxas e coberturas;
  \item \textbf{Variáveis de controle} em modelos de avaliação de impacto;
  \item \textbf{Construção de grupos de comparação} em análises de pareamento;
  \item \textbf{Análises territoriais} com desagregação por setor censitário.
\end{itemize}

\subsection{Limitações e cuidados}

A principal limitação do Censo para avaliação de políticas é sua \textbf{periodicidade decenal}: entre dois censos, não há atualização das informações coletadas no questionário amostral.

Outra limitação é o \textbf{subregistro em populações de difícil acesso}: áreas de conflito, comunidades indígenas isoladas, pessoas em situação de rua e moradores de favelas densas tendem a ser sub-representados.

Por fim, \textbf{mudanças metodológicas entre edições} podem comprometer comparações diretas de certas variáveis ao longo do tempo. A documentação metodológica do IBGE deve ser consultada antes de construir qualquer série histórica com dados censitários.

% ────────────────────────────────────────────────────────────
\section{Pesquisa Nacional por Amostra de Domicílios Contínua
         (PNAD Contínua)}
% ────────────────────────────────────────────────────────────

\ficharapida
  {PNAD Contínua (PNADC)}
  {Instituto Brasileiro de Geografia e Estatística (IBGE)}
  {https://www.ibge.gov.br/estatisticas/sociais/trabalho/17270-pnad-continua.html}
  {Brasil, Grandes Regiões, Unidades da Federação, Regiões Metropolitanas com capitais e capitais}
  {Trimestral (desde 2012); anual (desde 2016)}
  {Domicílio e morador}

\subsection{O que é e por que importa}

A Pesquisa Nacional por Amostra de Domicílios Contínua (PNAD Contínua ou PNADC) é, atualmente, a principal pesquisa domiciliar de acompanhamento contínuo da população brasileira. Iniciada em 2012, ela substituiu a PNAD anual e a Pesquisa Mensal de Emprego (PME) a partir de 2016, reunindo em uma única plataforma amostral informações sobre mercado de trabalho, educação, habitação e acesso a serviços.

O diferencial da PNAD Contínua em relação às pesquisas que substituiu é a \textbf{continuidade e frequência}: enquanto a PNAD era realizada uma vez por ano e a PME em seis regiões metropolitanas, a PNADC produz estimativas trimestrais para o Brasil inteiro e para todas as unidades da federação, além de estimativas mensais para indicadores selecionados do mercado de trabalho.

\subsection{Desenho amostral: painel rotativo e amostra complexa}

O desenho amostral da PNAD Contínua combina duas características que o analista precisa compreender antes de usar os dados.

A primeira é o \textbf{painel rotativo}: cada domicílio selecionado para a pesquisa é entrevistado em cinco ocasiões ao longo de cinco trimestres consecutivos. Após a quinta entrevista, o domicílio é substituído por outro. Esse formato permite análises longitudinais de trajetória no mercado de trabalho que seriam impossíveis em pesquisas puramente transversais.

A segunda --- e frequentemente subestimada --- é a \textbf{amostra complexa}. A PNADC utiliza um desenho de amostragem em dois estágios com estratificação e conglomeração: no primeiro estágio, são selecionados municípios (unidades primárias de amostragem, ou PSUs); no segundo, setores censitários dentro dos municípios selecionados. Isso tem consequências diretas e não triviais para a análise:

\begin{itemize}
  \item \textbf{Os pesos amostrais devem ser sempre aplicados} para produzir estimativas representativas da população. A variável de peso para o trimestre de referência é \texttt{V1028}. Análises sem pesos produzem estimativas incorretas;

  \item \textbf{O desenho amostral completo --- estratos e PSUs --- deve ser declarado} para o cálculo correto dos erros padrão. Ignorar a estrutura de conglomeração produz erros padrão sistematicamente \textit{subestimados}, o que infla artificialmente a significância estatística das estimativas. Em R, use o pacote \texttt{survey} com \texttt{svydesign()}; em Stata, use \texttt{svyset} antes de qualquer estimação; em Python, use \texttt{statsmodels.stats.survey};

  \item \textbf{A PNAD Contínua não deve ser usada para inferências no nível municipal}. O desenho amostral foi planejado para produzir estimativas representativas para o Brasil, grandes regiões, unidades da federação, regiões metropolitanas com capitais e capitais. Estimativas para municípios individuais --- exceto os que compõem as regiões metropolitanas com capitais, e mesmo assim com cautela --- não têm respaldo no desenho amostral da pesquisa e produzem resultados sem precisão estatística adequada;

  \item \textbf{Recortes populacionais muito finos exigem cautela}. Estimativas para subgrupos pequenos (como mulheres negras de 60 a 69 anos em uma UF específica) podem ter coeficientes de variação elevados que tornam as estimativas pouco confiáveis. Sempre reporte os erros padrão ou intervalos de confiança, e verifique o coeficiente de variação (CV) das estimativas antes de interpretá-las. O IBGE recomenda cautela para CVs acima de 15\% e não divulga estimativas com CV acima de 50\%;

  \item \textbf{A agregação temporal pode melhorar a precisão} para recortes finos. Para subgrupos com amostras pequenas, a combinação de múltiplos trimestres pode reduzir o CV das estimativas --- mas exige atenção ao painel rotativo para evitar dupla contagem de domicílios.
\end{itemize}

A documentação metodológica completa da PNADC está \href{ftp://ftp.ibge.gov.br/Trabalho_e_Rendimento/Pesquisa_Nacional_por_Amostra_de_Domicilios_continua/Notas_metodologicas/notas_metodologicas.pdf}{disponível neste link} e deve ser consultada antes de qualquer análise que envolva domínios de estimação não convencionais.

\begin{boxatencao}{Atenção --- PNAD Contínua não é base de inferência municipal}
Um erro frequente em avaliações de políticas públicas é usar a PNAD Contínua para calcular indicadores de mercado de trabalho ou renda no nível municipal, muitas vezes por ausência de alternativa mais conveniente. Esse uso não é metodologicamente válido: a amostra da PNADC não foi desenhada para representar municípios individuais, e as estimativas produzidas para esse nível têm precisão estatística inadequada. Para análises que exijam indicadores no nível municipal, as alternativas são o Censo Demográfico (para variáveis cobertas pelo questionário amostral) ou registros administrativos (para indicadores de mercado de trabalho formal, como CAGED e RAIS).
\end{boxatencao}

\subsection{Estrutura de divulgação: mensal, trimestral e anual}

A PNAD Contínua produz resultados em três periodicidades distintas, cada uma com um escopo temático diferente:

\begin{itemize}
  \item \textbf{Mensal}: conjunto restrito de indicadores do mercado de trabalho apenas para o Brasil como um todo. Serve principalmente para acompanhamento conjuntural;

  \item \textbf{Trimestral}: conjunto ampliado de indicadores do mercado de trabalho para todos os níveis de divulgação da pesquisa. É a divulgação de referência para análises do mercado de trabalho;

  \item \textbf{Anual}: inclui os temas permanentes adicionais --- educação (2º trimestre), acesso à internet e posse de telefone celular (4º trimestre) --- e outros temas variáveis. É a divulgação mais relevante para análises além do mercado de trabalho.
\end{itemize}

\subsection{Principais temas cobertos}

\begin{itemize}
  \item \textbf{Mercado de trabalho}: condição de ocupação, tipo de vínculo empregatício (formal/informal), setor de atividade, ocupação (ISCO/CBO), jornada de trabalho, rendimento do trabalho;
  \item \textbf{Educação}: nível de instrução, frequência escolar, alfabetização (2º trimestre);
  \item \textbf{Habitação}: características do domicílio, acesso a serviços básicos;
  \item \textbf{Tecnologia}: acesso à internet, posse de telefone celular e computador (4º trimestre);
  \item \textbf{Rendimento}: rendimento de todas as fontes, incluindo trabalho, aposentadoria e transferências.
\end{itemize}

\begin{boxatencao}{Atenção --- Quebra metodológica nos dados de rendimento}
A partir do 4º trimestre de 2019, o IBGE implementou mudanças na metodologia de captação do rendimento na PNAD Contínua, alterando a forma de coleta de rendimentos de algumas fontes. Isso introduziu uma \textbf{quebra na série histórica} que impede comparações diretas entre o rendimento médio antes e depois dessa data sem os devidos ajustes. O IBGE disponibiliza notas técnicas explicando a mudança; qualquer análise de tendência do rendimento que cruze essa quebra deve ser feita com cautela e devidamente documentada.
\end{boxatencao}

\subsection{Microdados: acesso e estrutura}

Os microdados trimestrais e anuais da PNAD Contínua são disponibilizados gratuitamente no site do IBGE. As variáveis de peso amostral (\texttt{V1028}) e as variáveis de estratificação e PSU devem ser sempre utilizadas nas análises.

\begin{boxdica}{Dica --- Pacotes para trabalhar com a PNAD Contínua}
\begin{itemize}
  \item \textbf{R}: o pacote \texttt{PNADcIBGE} automatiza o download e a preparação dos microdados já com o objeto de desenho amostral (\texttt{survey design}) corretamente especificado --- com estratos, PSUs e pesos. É a forma mais segura de garantir que o desenho amostral complexo está sendo tratado corretamente;
  \item \textbf{Python}: a biblioteca \texttt{basedosdados} oferece acesso aos microdados da PNADC via BigQuery, com tratamento prévio das variáveis;
  \item \textbf{Stata}: use \texttt{svyset [psu] [strata] [pweight]} para declarar o desenho amostral completo antes de qualquer análise estatística. Não use apenas o peso; declare também os estratos e as PSUs.
\end{itemize}
\end{boxdica}

\subsection{Aplicações típicas em avaliação}

\begin{itemize}
  \item Acompanhamento de indicadores do \textbf{mercado de trabalho} ao longo do tempo para o Brasil, regiões e UFs;
  \item Análises de \textbf{educação}: evolução da escolaridade, acesso à escola por faixa etária;
  \item Estimação de \textbf{indicadores de pobreza e desigualdade};
  \item \textbf{Avaliação de resultados} de políticas de emprego, educação e habitação;
  \item Construção de \textbf{variáveis de contexto} para estados e regiões em estudos de impacto.
\end{itemize}

% ────────────────────────────────────────────────────────────
\section{PNAD --- Pesquisa Nacional por Amostra de Domicílios
         (edições anteriores)}
% ────────────────────────────────────────────────────────────

\ficharapida
  {PNAD (edições 1967--2015)}
  {Instituto Brasileiro de Geografia e Estatística (IBGE)}
  {https://www.ibge.gov.br/estatisticas/sociais/trabalho/9127-pesquisa-nacional-por-amostra-de-domicilios.html}
  {Brasil, Grandes Regiões, Unidades da Federação e nove Regiões Metropolitanas}
  {Anual (1967--2015); não realizada em anos de Censo Demográfico}
  {Domicílio e morador}

\subsection{O que é e por que ainda importa}

A PNAD em sua versão original foi realizada entre 1967 e 2015, tornando-se a principal fonte de dados sobre condições de vida da população brasileira por quase cinco décadas. Embora tenha sido substituída pela PNAD Contínua, as edições anteriores continuam sendo amplamente utilizadas em pesquisas que precisam de séries históricas longas --- especialmente para o período anterior a 2012.

A PNAD cobria anualmente um conjunto fixo de temas (características gerais da população, educação, trabalho, rendimento e habitação) e, em determinados anos, incluía \textbf{suplementos temáticos} com questionários adicionais sobre assuntos específicos.

\textbf{O alerta sobre amostra complexa se aplica igualmente à PNAD anterior}: o desenho amostral também era estratificado e por conglomerados, com pesos e PSUs que devem ser declarados para produzir estimativas válidas e erros padrão corretos. O mesmo limite de representatividade estadual --- sem capacidade de inferência municipal na maior parte dos casos --- se aplicava à PNAD original.

\subsection{Suplementos temáticos relevantes para avaliação de políticas}

Os suplementos temáticos da PNAD são de grande valor para avaliadores, pois oferecem informações sobre dimensões raramente cobertas por outras fontes nacionais. Os principais suplementos incluem:

\begin{table}[H]
\centering
\caption{Principais suplementos temáticos da PNAD por ano}
\label{tab:pnad_suplementos}
\begin{tabularx}{\textwidth}{@{} p{1.5cm} X @{}}
\toprule
\textbf{Ano} & \textbf{Temas do suplemento} \\
\midrule
2015 & Acesso à internet e televisão; relações de trabalho; cuidados de crianças menores de 4 anos; práticas de esporte e atividade física \\
2014 & Acesso ao Cadastro Único e programas de inclusão produtiva; educação e qualificação profissional; mobilidade socioocupacional \\
2013 & Segurança alimentar; acesso à internet e televisão \\
2011 & Acesso à internet e posse de telefone celular \\
2009 & Segurança alimentar; vitimização e acesso à justiça \\
2008 & Aspectos complementares de educação e trabalho infantil \\
2005 & Tabagismo \\
2004 & Aspectos complementares de educação e trabalho infantil \\
1998 & Saúde e acesso a serviços de saúde \\
1995 & Trabalho infantil e contrato de trabalho \\
\bottomrule
\end{tabularx}
\fonte{IBGE, organização FGV CLEAR.}
\end{table}

\begin{boxexemplo}{Exemplo prático --- Usando o suplemento de vitimização da PNAD 2009}
O suplemento ``Características da vitimização e do acesso à Justiça no Brasil'' da PNAD 2009 é uma das poucas fontes nacionais com dados sobre vitimização criminal autoreferida. Pesquisadores utilizaram esses dados para comparar a vitimização declarada com os registros policiais estaduais, evidenciando que apenas uma fração dos crimes chega ao conhecimento da polícia. Para avaliações de políticas de segurança pública, esse suplemento é uma referência metodológica importante sobre os limites dos dados administrativos.

\textit{Fonte: FGV CLEAR, elaboração própria.}
\end{boxexemplo}

\subsection{Cuidados na comparação com a PNAD Contínua}

A transição da PNAD para a PNAD Contínua não foi metodologicamente neutra. As principais diferenças que afetam a comparabilidade das séries incluem:

\begin{itemize}
  \item \textbf{Definição de desocupação}: a PNAD Contínua adota a definição da OIT mais recente;
  \item \textbf{Período de referência}: a PNAD era coletada em setembro/outubro de cada ano; a PNADC é contínua;
  \item \textbf{Cobertura geográfica}: a PNADC cobre a área rural da região Norte, que a PNAD antiga não cobria adequadamente;
  \item \textbf{Metodologia de rendimento}: há diferenças na captação de rendimentos não laborais.
\end{itemize}

Para análises de longo prazo que precisem cruzar as duas pesquisas, recomenda-se consultar os estudos de harmonização produzidos pelo IBGE e as notas técnicas sobre a transição metodológica.

% ────────────────────────────────────────────────────────────
\section{Pesquisa Mensal de Emprego (PME)}
% ────────────────────────────────────────────────────────────

\ficharapida
  {Pesquisa Mensal de Emprego (PME)}
  {Instituto Brasileiro de Geografia e Estatística (IBGE)}
  {https://www.ibge.gov.br/estatisticas/sociais/trabalho/9180-pesquisa-mensal-de-emprego.html}
  {Seis Regiões Metropolitanas: Recife, Salvador, Belo Horizonte, Rio de Janeiro, São Paulo e Porto Alegre}
  {Mensal (1980--março de 2016)}
  {Domicílio e morador}

\subsection{O que é e aplicações}

A PME foi a principal fonte de dados mensais sobre o mercado de trabalho nas seis maiores regiões metropolitanas do Brasil entre 1980 e 2016. Sua relevância atual é principalmente \textbf{histórica}: para análises do mercado de trabalho metropolitano nas décadas de 1980, 1990 e 2000, continua sendo insubstituível.

O formato de painel rotativo da PME --- cada domicílio entrevistado por quatro meses consecutivos, ausente por oito meses e retornando por mais quatro --- permitia análises de transição no mercado de trabalho que influenciaram a metodologia da PNAD Contínua.

\subsection{Cuidados}

A PME cobre apenas seis regiões metropolitanas, o que limita qualquer generalização para o Brasil como um todo ou para cidades médias e pequenas. Como nas demais pesquisas amostrais do IBGE, o uso correto exige declaração do desenho amostral complexo --- não apenas a aplicação dos pesos.

% ────────────────────────────────────────────────────────────
\section{Pesquisa de Orçamentos Familiares (POF)}
% ────────────────────────────────────────────────────────────

\ficharapida
  {Pesquisa de Orçamentos Familiares (POF)}
  {Instituto Brasileiro de Geografia e Estatística (IBGE)}
  {https://www.ibge.gov.br/estatisticas/sociais/saude/24786-pesquisa-de-orcamentos-familiares-2.html}
  {Brasil, Grandes Regiões, Unidades da Federação e situação do domicílio (urbano/rural)}
  {Irregular: 1987/88, 1995/96, 2002/03, 2008/09, 2017/18}
  {Domicílio e morador}

\subsection{O que é e por que importa}

A Pesquisa de Orçamentos Familiares (POF) é o levantamento mais detalhado sobre consumo, despesas, rendimentos e condições de vida das famílias brasileiras. Diferentemente da PNAD Contínua, que captura o rendimento de forma agregada, a POF registra \textit{todas} as despesas e aquisições do domicílio durante um período de referência, permitindo análises de padrões de consumo, estrutura orçamentária e bem-estar que nenhuma outra fonte viabiliza.

Para a avaliação de políticas públicas, a POF é especialmente relevante para:

\begin{itemize}
  \item \textbf{Definição e atualização de linhas de pobreza}: a estrutura de consumo captada pela POF é a base metodológica para o cálculo de linhas de pobreza absolutas no Brasil;
  \item \textbf{Análise de políticas de transferência de renda}: ao capturar o orçamento completo das famílias, permite avaliar como transferências afetam o padrão de consumo;
  \item \textbf{Políticas de alimentação e nutrição}: a edição 2017/18 incluiu módulo detalhado sobre consumo alimentar individual.
\end{itemize}

\subsection{Amostra complexa na POF}

Assim como a PNAD Contínua, a POF utiliza desenho amostral estratificado por conglomerados. O mesmo cuidado se aplica: pesos amostrais e estrutura de estratos e PSUs devem ser declarados antes de qualquer análise estatística. A POF tem domínios de representatividade definidos --- Brasil, grandes regiões, UFs e situação do domicílio (urbano/rural) --- e estimativas para recortes mais finos têm precisão estatística reduzida.

\subsection{Principais módulos}

A POF coleta dados por meio de múltiplos cadernos de registro preenchidos pelos próprios moradores ao longo de vários dias. Os principais módulos incluem:

\begin{itemize}
  \item Despesas coletivas do domicílio (aluguel, contas de serviços, alimentação no domicílio);
  \item Despesas individuais de cada morador (vestuário, transporte, lazer, saúde);
  \item Rendimentos de todas as fontes;
  \item Avaliação das condições de vida (percepção subjetiva de bem-estar);
  \item Consumo alimentar individual (edição 2017/18).
\end{itemize}

\subsection{Limitações}

A principal limitação da POF para avaliação é sua \textbf{periodicidade irregular}: as edições têm ocorrido a cada seis ou sete anos. Além disso, o custo e a complexidade da coleta tornam a POF uma das pesquisas mais caras do sistema estatístico nacional, com amostras menores do que a PNAD Contínua e, portanto, menor capacidade de desagregação.

% ────────────────────────────────────────────────────────────
\section{Censo Agropecuário}
% ────────────────────────────────────────────────────────────

\ficharapida
  {Censo Agropecuário}
  {Instituto Brasileiro de Geografia e Estatística (IBGE)}
  {https://www.ibge.gov.br/estatisticas/economicas/agricultura-e-pecuaria/21814-2017-censo-agropecuario.html}
  {Brasil, Grandes Regiões, Unidades da Federação e municípios}
  {Irregular: 1920, 1940, 1950, 1960, 1970, 1975, 1980, 1985, 1995/96, 2006, 2017}
  {Estabelecimento agropecuário}

\subsection{O que é e aplicações}

O Censo Agropecuário é o levantamento mais completo sobre a estrutura produtiva do campo brasileiro. Cobre todos os estabelecimentos agropecuários do país --- propriedades rurais que desenvolvem atividades de lavoura, pecuária, horticultura, aquicultura ou agroindústria --- independentemente de seu tamanho ou regime de posse da terra.

Para a avaliação de políticas públicas, é a principal referência para:

\begin{itemize}
  \item \textbf{Políticas de reforma agrária e regularização fundiária}: dados sobre estrutura da posse da terra, tamanho dos estabelecimentos e condição do produtor;
  \item \textbf{Políticas de agricultura familiar}: identificação dos estabelecimentos por tipo (familiar vs. não familiar), acesso a crédito, assistência técnica e mercados;
  \item \textbf{Políticas ambientais}: uso do solo, áreas de preservação, práticas de manejo;
  \item \textbf{Segurança alimentar}: produção de alimentos por tipo e região.
\end{itemize}

\subsection{Limitações}

A periodicidade irregular é a principal limitação. Para o acompanhamento anual da produção agrícola, as alternativas são a Pesquisa Agrícola Municipal (PAM) e a Pesquisa Pecuária Municipal (PPM), também do IBGE, que produzem dados anuais mas com menor riqueza de variáveis.

% ────────────────────────────────────────────────────────────
\section{Síntese do capítulo}
% ────────────────────────────────────────────────────────────

Este capítulo apresentou as principais pesquisas amostrais e censos demográficos disponíveis para o Brasil. A Tabela~\ref{tab:sintese_cap3} resume as características essenciais de cada fonte.

\begin{table}[H]
\centering
\caption{Síntese das pesquisas amostrais e censos demográficos}
\label{tab:sintese_cap3}
\small
\begin{tabularx}{\textwidth}{@{} p{2.8cm} p{2.2cm} >{\raggedright\arraybackslash}p{2.5cm} X @{}}
\toprule
\textbf{Fonte} & \textbf{Período} & \textbf{Periodicidade} & \textbf{Principal uso em avaliação} \\
\midrule
Censo Demográfico    & 1940--2022  & Decenal      & Linha de base, denominadores populacionais, contrafactuais \\
PNAD Contínua        & 2012--atual & Trimestral\slash \allowbreak anual & Mercado de trabalho, educação, habitação, pobreza (Brasil, regiões, UFs) \\
PNAD (anterior)      & 1967--2015  & Anual        & Séries históricas longas, suplementos temáticos \\
PME                  & 1980--2016  & Mensal       & Mercado de trabalho metropolitano histórico \\
POF                  & 1987--2018  & Irregular    & Consumo, linhas de pobreza, transferências de renda \\
Censo Agropecuário   & 1920--2017  & Irregular    & Políticas rurais, fundiárias e ambientais \\
\bottomrule
\end{tabularx}
\fonte{FGV CLEAR, elaboração própria.}
\end{table}

Os pontos centrais do capítulo são:

\begin{itemize}
  \item O \textbf{Censo Demográfico} é insubstituível para desagregações territoriais finas e para construção de linha de base, mas sua periodicidade decenal limita seu uso em análises recentes;

  \item A \textbf{PNAD Contínua} é a principal fonte nacional de acompanhamento do mercado de trabalho, educação e condições de vida para o Brasil, regiões e UFs. Seu uso correto exige a declaração do \textbf{desenho amostral complexo} (estratos, PSUs e pesos) --- não apenas a aplicação dos pesos. A pesquisa \textbf{não deve ser usada para inferências no nível municipal};

  \item Todas as pesquisas amostrais do IBGE discutidas neste capítulo --- PNAD Contínua, PNAD anterior, PME e POF --- utilizam desenhos amostrais complexos com estratificação e conglomeração. Recortes muito finos devem ser acompanhados de verificação dos coeficientes de variação e, quando necessário, de agregação temporal;

  \item As edições anteriores da PNAD mantêm relevância para séries históricas e pelos suplementos temáticos que cobrem dimensões raramente disponíveis em outras fontes;

  \item A POF é a referência para análises de consumo e definição de linhas de pobreza, mas tem periodicidade irregular;

  \item A comparação de séries entre pesquisas exige atenção às quebras metodológicas documentadas pelo IBGE.
\end{itemize}

O próximo capítulo aborda as fontes de dados sobre mercado de trabalho e renda, com foco nos registros administrativos --- CAGED, RAIS, eSocial e Cadastro Único.

% ============================================================
%  cap04_mercado_trabalho.tex
%  Capítulo 4: Mercado de Trabalho e Renda
% ============================================================

% ════════════════════════════════════════════════════════════
\chapter{Mercado de Trabalho e Renda}
% ════════════════════════════════════════════════════════════

\begin{boxdestaque}{Neste capítulo}
Este capítulo apresenta as principais fontes de dados administrativos sobre mercado de trabalho e renda no Brasil. Cobrimos o CAGED, a RAIS, o eSocial, o Cadastro Único e o Novo CAGED, detalhando para cada fonte seu histórico, estrutura, variáveis disponíveis, forma de acesso e aplicações típicas em avaliação de políticas públicas. O capítulo também discute a relação entre essas fontes e como combiná-las de forma estratégica.
\end{boxdestaque}

% ────────────────────────────────────────────────────────────
\section{Introdução: o ecossistema de dados do mercado de trabalho}
% ────────────────────────────────────────────────────────────

O mercado de trabalho brasileiro é monitorado por um conjunto relativamente rico de fontes administrativas, produzidas principalmente pelo Ministério do Trabalho e Emprego (MTE) e pelo Ministério do Desenvolvimento e Assistência Social (MDS). Cada fonte cobre uma dimensão específica: o CAGED registra o \textit{fluxo} de admissões e demissões no setor formal; a RAIS registra o \textit{estoque} de vínculos formais ao final de cada ano; o Cadastro Único cobre as famílias de baixa renda, incluindo trabalhadores informais e desocupados que estão fora do alcance do CAGED e da RAIS.

Entender a relação entre essas fontes — o que cada uma cobre, o que exclui e como se complementam — é fundamental para qualquer avaliação de políticas de emprego, renda e proteção social.

\begin{figure}[H]
\centering
\begin{boxdestaque}{O ecossistema de dados do mercado de trabalho formal}
\begin{center}
\begin{tabular}{@{} p{4cm} p{3.5cm} p{5.5cm} @{}}
\toprule
\textbf{Fonte} & \textbf{O que mede} & \textbf{O que exclui} \\
\midrule
CAGED & Fluxo de admissões e demissões & Estatutários, militares, informais \\
RAIS  & Estoque de vínculos formais & Estatutários (parcial), militares, informais \\
eSocial & Obrigações trabalhistas e previdenciárias & Informais, autônomos sem registro \\
Cadastro Único & Famílias de baixa renda & Famílias acima do limite de renda \\
PNAD Contínua & Mercado de trabalho total & Não há: cobre formais e informais \\
\bottomrule
\end{tabular}
\end{center}
\end{boxdestaque}
\end{figure}

% ────────────────────────────────────────────────────────────
\section{CAGED — Cadastro Geral de Empregados e Desempregados}
% ────────────────────────────────────────────────────────────

\ficharapida
  {CAGED}
  {Ministério do Trabalho e Emprego (MTE)}
  {https://www.gov.br/trabalho-e-emprego/pt-br/assuntos/estatisticas-trabalho/caged}
  {Brasil, Grandes Regiões, Unidades da Federação e municípios}
  {Mensal (1992--2019); Novo CAGED a partir de janeiro de 2020}
  {Vínculo empregatício (admissão ou demissão)}

\subsection{O que é e por que importa}

O Cadastro Geral de Empregados e Desempregados (CAGED) é o principal registro administrativo de \textbf{fluxo} do mercado de trabalho formal brasileiro. Instituído pela Lei n.º 4.923/1965 e operacionalizado pelo Ministério do Trabalho e Emprego, o CAGED obriga todos os estabelecimentos que contratam ou dispensam trabalhadores regidos pela Consolidação das Leis do Trabalho (CLT) a comunicar esses movimentos mensalmente ao governo.

O resultado é uma base com cobertura praticamente universal do emprego com carteira assinada no setor privado, atualizada mensalmente e desagregável por município, setor de atividade econômica, ocupação, sexo, faixa etária e escolaridade. Nenhuma outra fonte brasileira oferece essa combinação de frequência, cobertura e desagregação para o emprego formal.

\subsection{Histórico e o Novo CAGED}

O CAGED em sua forma original foi operado entre 1992 e dezembro de 2019. A partir de janeiro de 2020, foi substituído pelo \textbf{Novo CAGED}, que passou a ser alimentado pelos dados do eSocial e do e-Empregador — sistemas de declaração eletrônica unificada de obrigações trabalhistas. Essa transição trouxe mudanças substanciais:

\begin{itemize}
  \item \textbf{Universo de declarantes ampliado}: o Novo CAGED incorporou categorias de trabalhadores antes não cobertas, como trabalhadores avulsos e alguns vínculos de servidores públicos celetistas;
  \item \textbf{Metodologia de competência vs. caixa}: o CAGED antigo registrava movimentos pela data de declaração; o Novo CAGED registra pela data de competência (quando o evento ocorreu de fato), o que melhora a precisão mas pode gerar revisões retroativas;
  \item \textbf{Quebra de série}: a mudança introduziu uma \textbf{descontinuidade metodológica} que impede a comparação direta de dados do Novo CAGED com séries históricas do CAGED antigo sem ajustes.
\end{itemize}

% CORRIGIDO: título encurtado para caber na caixa
\begin{boxatencao}{Atenção — Quebra de série: CAGED e Novo CAGED}
A transição para o Novo CAGED em janeiro de 2020 não é apenas uma atualização tecnológica — é uma mudança metodológica que altera o universo coberto, o conceito de competência e a forma de registro. Análises que precisem de séries históricas longas (anteriores a 2020) devem tratar os dois períodos separadamente ou utilizar fatores de ajuste. O próprio MTE publicou notas técnicas sobre a transição que devem ser consultadas antes de qualquer comparação temporal.
\end{boxatencao}

\subsection{Principais variáveis disponíveis}

O CAGED (e o Novo CAGED) disponibiliza, para cada movimentação registrada:

\begin{itemize}
  \item Tipo de movimentação (admissão ou desligamento);
  \item Motivo do desligamento (demissão sem justa causa, com justa causa, pedido de demissão, aposentadoria, etc.);
  \item Salário de admissão ou demissão;
  \item Sexo, faixa etária e grau de instrução do trabalhador;
  \item Município e UF do estabelecimento;
  \item Setor de atividade econômica (CNAE);
  \item Ocupação (CBO — Classificação Brasileira de Ocupações);
  \item Porte do estabelecimento.
\end{itemize}

\subsection{Acesso aos dados}

Os dados agregados do CAGED são publicados mensalmente pelo MTE e disponibilizados no portal do governo federal. Os microdados — com informações por vínculo individual — são disponibilizados com defasagem de alguns meses e podem ser acessados no portal do MTE ou via plataforma \textit{Base dos Dados} (\url{basedosdados.org}), que oferece os dados já tratados e compatibilizados em BigQuery.

\subsection{Aplicações típicas em avaliação}

\begin{itemize}
  \item \textbf{Avaliação de políticas de geração de emprego}: monitoramento do saldo líquido de empregos formais criados ou destruídos em resposta a uma intervenção;
  \item \textbf{Análise setorial e regional}: identificação de quais setores e regiões responderam a políticas de incentivo fiscal ou crédito;
  \item \textbf{Avaliação de programas de qualificação profissional}: comparação de admissões antes e depois de programas de treinamento, controlando por setor e ocupação;
  \item \textbf{Estudos de impacto de choques econômicos}: o CAGED é amplamente usado para avaliar o impacto de crises, choques setoriais e mudanças regulatórias sobre o emprego formal.
\end{itemize}

% CORRIGIDO: título encurtado para caber na caixa
\begin{boxexemplo}{Exemplo — Polo industrial e emprego formal}
Um estado implementou um programa de atração de investimentos industriais que resultou na instalação de um polo automotivo em um município de médio porte. Para avaliar o impacto sobre o emprego formal local, pesquisadores utilizaram os microdados do CAGED para construir uma série mensal de admissões e desligamentos no município tratado e em municípios de controle com características similares (usando pareamento por escore de propensão). O CAGED permitiu não apenas medir o saldo líquido de empregos, mas também verificar se os postos criados eram de melhor qualidade (maior salário, maior escolaridade exigida) do que os que existiam anteriormente.

\textit{Fonte: FGV CLEAR, elaboração própria.}
\end{boxexemplo}

% ────────────────────────────────────────────────────────────
\section{RAIS — Relação Anual de Informações Sociais}
% ────────────────────────────────────────────────────────────

\ficharapida
  {RAIS}
  {Ministério do Trabalho e Emprego (MTE)}
  {https://www.gov.br/trabalho-e-emprego/pt-br/assuntos/estatisticas-trabalho/rais}
  {Brasil, Grandes Regiões, Unidades da Federação e municípios}
  {Anual (desde 1975)}
  {Vínculo empregatício e estabelecimento; inclui remuneração mensal e total anual, jornada, CBO, CNAE, tempo no emprego e motivo de desligamento}

\subsection{O que é e por que importa}

Enquanto o CAGED registra o \textit{fluxo} de entradas e saídas do emprego formal, a Relação Anual de Informações Sociais (RAIS) registra o \textit{estoque}: todos os vínculos empregatícios ativos em cada estabelecimento ao longo do ano. Instituída pelo Decreto n.º 76.900/1975, a RAIS é uma declaração anual obrigatória para todos os estabelecimentos — públicos e privados — que possuam ou tenham possuído empregados com vínculo formal durante o ano-base.

A RAIS é amplamente considerada a \textbf{fonte mais completa para análise do mercado de trabalho formal brasileiro} em perspectiva anual. Sua cobertura é maior do que a do CAGED porque inclui, além dos trabalhadores celetistas, os \textbf{servidores públicos estatutários} — categoria excluída do CAGED — e informações sobre estabelecimentos que não tiveram movimentação durante o ano (declaração negativa).

\subsection{Duas bases complementares}

A RAIS é disponibilizada em duas bases distintas, com unidades de análise diferentes:

\begin{itemize}
  \item \textbf{Base de vínculos}: cada linha é um vínculo empregatício. Contém informações individuais sobre o trabalhador (identificado pelo CPF, mas anonimizado nos microdados públicos) e sobre o vínculo (salário, ocupação, tempo de emprego, motivo de desligamento, etc.);
  \item \textbf{Base de estabelecimentos}: cada linha é um estabelecimento (identificado pelo CNPJ). Contém informações sobre localização geográfica, setor de atividade, porte, número de empregados e composição da força de trabalho por escolaridade, sexo e faixa etária.
\end{itemize}

As duas bases podem ser cruzadas pelo CNPJ do estabelecimento, permitindo análises que combinam características dos trabalhadores e das firmas.

\subsection{Principais variáveis da base de vínculos}

\begin{itemize}
  \item Identificador do trabalhador (CPF anonimizado);
  \item Município e UF do estabelecimento;
  \item Setor de atividade (CNAE) e ocupação (CBO);
  \item Remuneração média e em dezembro;
  \item Grau de instrução, sexo, faixa etária, raça/cor e nacionalidade;
  \item Tempo de emprego no vínculo;
  \item Mês e ano de admissão e, se houver, de desligamento e motivo;
  \item Tipo de vínculo (CLT, estatutário, temporário, etc.);
  \item Indicador de portador de deficiência.
\end{itemize}

\subsection{Acesso: dados públicos vs. dados restritos}

O acesso à RAIS tem dois níveis:

\begin{itemize}
  \item \textbf{Dados agregados e tabulações}: disponíveis gratuitamente no portal do MTE, sem necessidade de cadastro;
  \item \textbf{Microdados anonimizados}: disponíveis para pesquisadores mediante cadastro no portal do MTE. Os microdados públicos têm o CPF substituído por um identificador anonimizado, o que impede integração com outras bases por CPF, mas permite análises individuais;
  \item \textbf{Microdados com CPF}: disponíveis apenas mediante convênio formal com o MTE, para pesquisadores vinculados a instituições acadêmicas ou governamentais. Com o CPF, é possível cruzar a RAIS com outras bases (como o CadÚnico, o SINASC ou dados de programas específicos), o que abre possibilidades analíticas muito mais ricas.
\end{itemize}

\begin{boxdica}{Dica — RAIS identificada via convênio}
Para pesquisadores que precisam da RAIS com CPF para integração de bases, o processo de convênio com o MTE exige uma proposta de pesquisa, termo de sigilo e aprovação por comitê. O prazo costuma ser de dois a quatro meses. A plataforma \textit{Base dos Dados} oferece a RAIS anonimizada já tratada em BigQuery, o que facilita muito o acesso para análises que não exigem identificação individual.
\end{boxdica}

\subsection{Aplicações típicas em avaliação}

\begin{itemize}
  \item \textbf{Análises salariais e de desigualdade}: a RAIS é a principal fonte para estudos de diferencial salarial por sexo, raça, escolaridade e setor;
  \item \textbf{Avaliação de políticas de qualificação profissional}: integração entre participantes de programas de treinamento (identificados pelo CPF) e seus vínculos na RAIS antes e depois do programa;
  \item \textbf{Estudos de mobilidade ocupacional e setorial}: trajetórias de trabalhadores ao longo do tempo, usando o painel implícito construído pelo CPF;
  \item \textbf{Análise do setor público}: por incluir estatutários, a RAIS é a única fonte administrativa que permite comparar remuneração e características do emprego público e privado;
  \item \textbf{Estudos de firmas}: a base de estabelecimentos permite análises sobre produtividade, crescimento e mortalidade de empresas;
  \item \textbf{Análise da transição entre formal e informal}: combinada à PNAD Contínua (que capta emprego informal) e ao CAGED/Novo CAGED, a RAIS permite estimar efeitos de políticas de formalização sobre a composição do mercado de trabalho --- identificando, por exemplo, se uma política gerou novos vínculos formais ou apenas realocou trabalhadores já formalizados.
\end{itemize}

% ────────────────────────────────────────────────────────────
\section{eSocial}
% ────────────────────────────────────────────────────────────

\ficharapida
  {eSocial}
  {Governo Federal (MTE, Receita Federal, Previdência Social)}
  {https://www.gov.br/esocial}
  {Brasil}
  {Contínuo (implementação gradual a partir de 2018)}
  {Evento trabalhista (admissão, folha, férias, desligamento, acidente, etc.)}

\subsection{O que é}

O eSocial é um sistema de escrituração digital unificada das obrigações fiscais, previdenciárias e trabalhistas do empregador. Não é uma pesquisa nem um registro administrativo no sentido tradicional — é uma plataforma de declaração eletrônica que integra informações que antes eram prestadas em sistemas separados (CAGED, RAIS, GFIP, SEFIP, entre outros).

A implementação foi gradual: grandes empresas começaram a declarar em 2018; médias e pequenas empresas em 2019 e 2020; empregadores domésticos e pessoas físicas a partir de 2020.

\subsection{Relevância para pesquisa e avaliação}

O eSocial ainda está em fase de consolidação como fonte de dados para pesquisa. Seu principal impacto já visível é ter alimentado o Novo CAGED a partir de 2020. No médio prazo, espera-se que o eSocial permita:

\begin{itemize}
  \item \textbf{Acompanhamento em tempo quase real} de vínculos trabalhistas, incluindo eventos como acidentes de trabalho e afastamentos;
  \item \textbf{Cobertura ampliada} de categorias antes sub-representadas nos registros administrativos, como trabalhadores domésticos formalizados;
  \item \textbf{Integração com dados previdenciários}, permitindo análises que conectam trajetórias de emprego com benefícios recebidos.
\end{itemize}

Por ora, o acesso aos microdados do eSocial para fins de pesquisa ainda é restrito e os protocolos de acesso estão em desenvolvimento. Pesquisadores interessados devem acompanhar as publicações do MTE e da Receita Federal sobre abertura dos dados.

% ────────────────────────────────────────────────────────────
\section{Cadastro Único para Programas Sociais (CadÚnico)}
% ────────────────────────────────────────────────────────────

\ficharapida
  {Cadastro Único (CadÚnico)}
  {Ministério do Desenvolvimento e Assistência Social (MDS)}
  {https://www.gov.br/mds/pt-br/acoes-e-programas/cadastro-unico}
  {Brasil, Unidades da Federação e municípios}
  {Contínuo (cadastramento e atualização permanentes; desde 2001)}
  {Família e pessoa}

\subsection{O que é e por que importa}

O Cadastro Único para Programas Sociais (CadÚnico) é o principal instrumento de identificação e caracterização socioeconômica das famílias de baixa renda no Brasil. Criado em 2001 e reformulado em 2011 para unificar informações antes dispersas entre diferentes programas, o CadÚnico serve como porta de entrada para a maioria dos programas sociais federais — Bolsa Família/Auxílio Brasil/Bolsa Família (novo), Tarifa Social de Energia Elétrica, Minha Casa Minha Vida (faixa 1), entre outros.

\textbf{Famílias elegíveis} para cadastramento são aquelas com renda per capita de até meio salário mínimo ou renda familiar total de até três salários mínimos. O cadastramento é voluntário e realizado pelos municípios, que são responsáveis pela inserção e atualização dos dados.

Para a avaliação de políticas públicas, o CadÚnico é especialmente valioso porque cobre exatamente a população que tende a ser o público-alvo de políticas sociais — e que, por isso mesmo, está sub-representada em pesquisas amostrais (que podem não ter amostras suficientemente grandes nos estratos de menor renda) e ausente nos registros do mercado formal (CAGED e RAIS).

\subsection{Estrutura e principais variáveis}

O CadÚnico é organizado em duas unidades de análise vinculadas:

\begin{itemize}
  \item \textbf{Família}: composição familiar, endereço, características do domicílio (tipo de construção, acesso a água, esgoto, energia), renda familiar per capita;
  \item \textbf{Pessoa}: para cada membro da família — nome, CPF, NIS (Número de Identificação Social), data de nascimento, sexo, raça/cor, escolaridade, situação no mercado de trabalho, renda individual de todas as fontes.
\end{itemize}

O NIS é o identificador único de cada pessoa no CadÚnico e pode ser usado para cruzar o cadastro com bases de benefícios (Bolsa Família, BPC) e, mediante convênio, com outras bases administrativas.

\subsection{Atualização e qualidade dos dados}

As famílias devem atualizar seus dados a cada dois anos ou sempre que houver alteração relevante nas condições socioeconômicas (mudança de endereço, nascimento, morte, alteração de renda). Na prática, a qualidade e a regularidade da atualização variam muito entre municípios, o que introduz um grau de defasagem e inconsistência que o analista deve considerar.

O MDS publica mensalmente estatísticas agregadas do CadÚnico por município, UF e Brasil, incluindo número de famílias cadastradas, perfil de renda, cobertura de programas sociais e características domiciliares. Os microdados do CadÚnico têm acesso restrito e são disponibilizados mediante convênio com o MDS para pesquisadores e gestores públicos.

\begin{boxatencao}{Atenção — CadÚnico não é censo da pobreza}
O CadÚnico cobre apenas as famílias que se cadastraram voluntariamente. Famílias elegíveis que nunca buscaram os programas sociais — por desconhecimento, estigma, barreiras de acesso ou outros motivos — não constam na base. Estimativas sobre a cobertura sugerem que uma parcela significativa das famílias que seriam elegíveis não está cadastrada, especialmente em áreas rurais remotas e entre populações nômades ou em situação de rua. Portanto, o CadÚnico não deve ser interpretado como uma listagem completa da população em situação de pobreza.
\end{boxatencao}

\subsection{Acesso aos dados}

\begin{itemize}
  \item \textbf{Dados agregados públicos}: disponíveis no portal CECAD (\url{cecad.cidadania.gov.br}) e no VISDATA (\url{aplicacoes.mds.gov.br/sagi/vis/data3}), com tabelas e mapas por município;
  \item \textbf{Microdados anonimizados}: disponíveis para pesquisadores mediante convênio com o MDS. O processo exige proposta de pesquisa, termo de confidencialidade e aprovação institucional;
  \item \textbf{Integração com Bolsa Família}: o MDS disponibiliza dados de beneficiários do Bolsa Família (atual e histórico) no Portal da Transparência, com NIS como identificador.
\end{itemize}

\subsection{Aplicações típicas em avaliação}

\begin{itemize}
  \item \textbf{Avaliação de programas de transferência de renda}: o CadÚnico é a base de elegibilidade do Bolsa Família, tornando-o essencial para identificar o grupo de tratamento em qualquer avaliação desse programa;
  \item \textbf{Focalização de políticas sociais}: análise de cobertura — quantas famílias elegíveis estão sendo atendidas e quem está sendo excluído;
  \item \textbf{Caracterização multidimensional da pobreza}: a riqueza de variáveis sobre domicílio, educação e trabalho permite análises de pobreza multidimensional;
  \item \textbf{Integração com registros de saúde e educação}: mediante convênio, o CadÚnico pode ser cruzado com o SINASC, o SIM ou o Censo Escolar pelo CPF ou NIS, permitindo avaliar como condições socioeconômicas se associam a resultados de saúde e educação.
\end{itemize}

% CORRIGIDO: título encurtado para caber na caixa
\begin{boxexemplo}{Exemplo — Integração CadÚnico + Censo Escolar: Bolsa Família}
Uma das avaliações mais citadas do Programa Bolsa Família utilizou a integração entre o CadÚnico (que identifica as famílias beneficiárias pelo NIS) e o Censo Escolar (que registra a frequência de cada aluno pelo CPF). Ao combinar as duas bases, os pesquisadores puderam comparar a frequência escolar de crianças de famílias beneficiárias com a de crianças de famílias elegíveis mas não beneficiárias, controlando por características socioeconômicas observáveis do CadÚnico. Essa integração tornou possível uma avaliação robusta sem necessidade de coleta primária de dados.

\textit{Fonte: FGV CLEAR, elaboração própria.}
\end{boxexemplo}

% ────────────────────────────────────────────────────────────
\section{Combinando as fontes: uma matriz de decisão}
% ────────────────────────────────────────────────────────────

As fontes apresentadas neste capítulo cobrem dimensões distintas e complementares do mercado de trabalho e da renda. A escolha entre elas — ou a combinação estratégica — depende da pergunta avaliativa e do público de interesse.

\begin{table}[H]
\centering
\caption{Matriz de decisão: fontes de dados para mercado de trabalho e renda}
\label{tab:matriz_trabalho}
\small
\begin{tabularx}{\textwidth}{@{} p{4.5cm} X @{}}
\toprule
\textbf{Pergunta avaliativa} & \textbf{Fonte(s) recomendada(s)} \\
\midrule
Quantos empregos formais foram criados no setor X após a política? & CAGED / Novo CAGED \\
\addlinespace
Quais são as características dos trabalhadores formais em determinado setor ou região? & RAIS (base de vínculos) \\
\addlinespace
Como as firmas responderam à política (tamanho, setor, localização)? & RAIS (base de estabelecimentos) \\
\addlinespace
Qual é o perfil socioeconômico do público-alvo de uma política social? & Cadastro Único \\
\addlinespace
A política aumentou a taxa de formalização entre trabalhadores de baixa renda? & CadÚnico + CAGED/RAIS (integração por CPF) \\
\addlinespace
Como evoluiu o rendimento médio dos trabalhadores ao longo do tempo? & PNAD Contínua (séries longas) \\
\addlinespace
Quais trabalhadores transitaram do desemprego para o emprego após o programa? & CAGED (admissões) + CadÚnico (perfil pré-programa) \\
\addlinespace
Quanto representa o emprego público no total do mercado formal? & RAIS (inclui estatutários) \\
\bottomrule
\end{tabularx}
\fonte{FGV CLEAR, elaboração própria.}
\end{table}

% ────────────────────────────────────────────────────────────
\section{Síntese do capítulo}
% ────────────────────────────────────────────────────────────

Este capítulo apresentou as principais fontes administrativas de dados sobre mercado de trabalho e renda no Brasil. Os pontos centrais são:

\begin{itemize}
  \item O \textbf{CAGED} registra o fluxo mensal de admissões e demissões no emprego formal celetista; o \textbf{Novo CAGED} (a partir de 2020) ampliou a cobertura mas introduziu uma quebra de série que impede comparação direta com o período anterior sem ajustes;
  \item A \textbf{RAIS} registra o estoque anual de vínculos formais, incluindo estatutários, e é a principal fonte para análises de diferencial salarial, mobilidade ocupacional e características de firmas; o acesso com CPF, via convênio, abre possibilidades de integração com outras bases;
  \item O \textbf{eSocial} unifica obrigações trabalhistas e deve se tornar progressivamente uma fonte central de dados administrativos, mas seu uso para pesquisa ainda está em consolidação;
  \item O \textbf{Cadastro Único} é a principal fonte de dados sobre famílias de baixa renda e trabalhadores informais, cobrindo exatamente o público que o CAGED e a RAIS não alcançam; sua limitação é que cobre apenas quem se cadastrou voluntariamente;
  \item A \textbf{combinação estratégica} dessas fontes — especialmente via integração por CPF ou NIS — permite análises muito mais ricas do que qualquer fonte isolada.
\end{itemize}

O próximo capítulo apresenta as fontes de dados sobre saúde: os sistemas do DATASUS (SIM, SINASC, SIH, SIA, SINAN, CNES e outros) e a Pesquisa Nacional de Saúde, que cobrem nascimentos, óbitos, internações, vigilância epidemiológica e a experiência de saúde da população brasileira como um todo.
% ============================================================
%  cap05_saude.tex — Capítulo 5: Saúde
% ============================================================

% ════════════════════════════════════════════════════════════
\chapter{Saúde}
% ════════════════════════════════════════════════════════════

\begin{boxdestaque}{Neste capítulo}
Este capítulo apresenta as principais fontes de dados sobre saúde disponíveis para o Brasil. Cobrimos o DATASUS e seus sistemas componentes --- SIM, SINASC, SIH, SIA, SINAN, CNES, SIOPS, SISVAN e SISAB ---, a Pesquisa Nacional de Saúde (PNS) e a Pesquisa de Assistência Médico-Sanitária (AMS). Para cada fonte, detalhamos histórico, estrutura, variáveis disponíveis, forma de acesso e aplicações típicas em avaliação de políticas de saúde pública.
\end{boxdestaque}

% ────────────────────────────────────────────────────────────
\section{Introdução: o ecossistema de dados de saúde no Brasil}
% ────────────────────────────────────────────────────────────

O Brasil dispõe de um dos sistemas de informação em saúde mais abrangentes do mundo em desenvolvimento, construído ao longo de décadas a partir da estrutura do Sistema Único de Saúde (SUS). O DATASUS --- Departamento de Informática do SUS, criado em 1991 --- é o principal repositório desses dados, reunindo em uma única plataforma sistemas de mortalidade, nascimentos, internações, ambulatório, vigilância epidemiológica, cadastro de estabelecimentos e orçamento público em saúde.

Essa riqueza de fontes é um ativo extraordinário para a avaliação de políticas de saúde. Ao mesmo tempo, o ecossistema é complexo: cada sistema tem sua própria lógica de produção, periodicidade, cobertura e qualidade. Entender o que cada sistema mede --- e, igualmente importante, o que não mede --- é o ponto de partida para qualquer análise séria.

A Figura~\ref{fig:datasus_sistemas} apresenta uma visão esquemática dos principais sistemas do DATASUS e sua relação com os diferentes pontos de contato do cidadão com o sistema de saúde.

\begin{table}[H]
\centering
\caption{Principais sistemas do DATASUS e suas coberturas}
\label{fig:datasus_sistemas}
\small
\begin{tabularx}{\textwidth}{@{} p{1.8cm} p{3.5cm} >{\raggedright\arraybackslash}p{2.5cm} X @{}}
\toprule
\textbf{Sistema} & \textbf{O que registra} & \textbf{Desde} & \textbf{Principal uso em avaliação} \\
\midrule
SIM    & Óbitos                        & 1979 & Mortalidade, causas de morte, violência \\
SINASC & Nascimentos vivos             & 1990 & Saúde materna, prematuridade, desigualdades \\
SIH    & Internações hospitalares SUS  & 1992 & Uso hospitalar, procedimentos, custos \\
SIA    & Atendimentos ambulatoriais SUS & 1994 & Produção ambulatorial, cobertura \\
SINAN  & Agravos de notificação        & 1993 & Doenças infecciosas, vigilância epidemiológica \\
CNES   & Estabelecimentos de saúde     & 2005 & Infraestrutura, RH, capacidade instalada \\
SIOPS  & Orçamento e execução em saúde    & 2000 & Gasto público, financiamento do SUS \\
SISVAN & Vigilância alimentar/nutricional & 2008 & Desnutrição, obesidade, segurança alimentar \\
SISAB  & Atenção primária à saúde      & 2013 & Cobertura de AB, ESF, ações de saúde na AB \\
\bottomrule
\end{tabularx}
\fonte{FGV CLEAR, elaboração própria com base em DATASUS.}
\end{table}

% ────────────────────────────────────────────────────────────
\section{SIM --- Sistema de Informações sobre Mortalidade}
% ────────────────────────────────────────────────────────────

\ficharapida
  {SIM --- Sistema de Informações sobre Mortalidade}
  {Ministério da Saúde / DATASUS}
  {https://datasus.saude.gov.br/mortalidade-desde-1996-pela-cid-10}
  {Brasil, Unidades da Federação e municípios}
  {Anual (desde 1979; dados consistentes a partir de 1996 com a CID-10)}
  {Óbito}

\subsection{O que é e por que importa}

O Sistema de Informações sobre Mortalidade (SIM) é o registro nacional de óbitos no Brasil. Criado em 1979, o SIM coleta informações de todas as Declarações de Óbito (DO) emitidas no país, processadas pelas Secretarias Estaduais de Saúde e consolidadas pelo Ministério da Saúde. A partir de 1996, com a adoção da Classificação Internacional de Doenças em sua décima revisão (CID-10), as séries do SIM tornaram-se mais comparáveis internacionalmente e consistentes para análises de longo prazo.

Para a avaliação de políticas públicas, o SIM é uma fonte insubstituível porque a morte é um desfecho objetivo, verificável e registrado de forma relativamente padronizada --- ao contrário de muitos outros indicadores de saúde que dependem de autoavaliação ou de acesso a serviços. Políticas de saúde materno-infantil, de atenção primária, de combate à violência e de doenças crônicas têm no SIM uma das principais fontes de avaliação de resultados e impacto.

\subsection{Estrutura e principais variáveis}

Cada registro do SIM corresponde a uma Declaração de Óbito e contém:

\begin{itemize}
  \item \textbf{Características do falecido}: sexo, idade, raça/cor, escolaridade, estado civil e ocupação --- a qualidade de preenchimento dessas variáveis é heterogênea; em particular, escolaridade, raça/cor e ocupação apresentam taxas expressivas de \textit{missing} e exigem tratamento cuidadoso antes da análise;
  \item \textbf{Causa da morte}: causa básica e causas associadas, codificadas pela CID-10 (adotada no Brasil em 1996, com implementação plena nos registros do SIM a partir de 1998; séries anteriores usam a CID-9 e exigem compatibilização para análises históricas); causa externa (para mortes violentas);
  \item \textbf{Localização}: município de residência e município de ocorrência do óbito;
  \item \textbf{Data e local do óbito}: data, hora, local (hospital, domicílio, via pública, etc.);
  \item \textbf{Para óbitos infantis e fetais}: informações adicionais sobre a mãe (idade, número de filhos, gestação) e sobre o recém-nascido (peso ao nascer, duração da gestação).
\end{itemize}

\subsection{Aplicações típicas em avaliação}

\begin{itemize}
  \item \textbf{Mortalidade infantil e materna}: o SIM, combinado com o SINASC (para os denominadores de nascidos vivos), é a fonte padrão para calcular as taxas de mortalidade infantil e materna por município;
  \item \textbf{Atlas da Violência}: o Instituto de Pesquisa Econômica Aplicada (Ipea) e o Fórum Brasileiro de Segurança Pública produzem anualmente o Atlas da Violência a partir de mortes por causas externas (homicídios, acidentes de trânsito) registradas no SIM;
  \item \textbf{Avaliação do Programa Saúde da Família}: um dos usos mais citados do SIM em avaliação de impacto é a estimação dos efeitos da expansão do PSF sobre mortalidade infantil e materna em municípios brasileiros;
  \item \textbf{Doenças crônicas não transmissíveis}: mortalidade por doenças cardiovasculares, câncer e diabetes para avaliação de programas de rastreamento e tratamento.
\end{itemize}

% CORRIGIDO: título encurtado para caber na caixa
\begin{boxatencao}{Atenção --- Subregistro e causa básica no SIM}
O SIM tem dois problemas de qualidade conhecidos que o analista deve considerar. Primeiro, o \textbf{subregistro de óbitos}: historicamente, uma parcela dos óbitos --- especialmente em municípios pequenos, áreas rurais e regiões Norte e Nordeste --- não era registrada, subestimando a mortalidade real nessas áreas. Esse quadro melhorou substancialmente nas últimas duas décadas, com a expansão da Estratégia Saúde da Família, iniciativas como a Busca Ativa de Óbitos e a qualificação das equipes municipais de vigilância; ainda assim, desigualdades regionais persistem em municípios menores e mais isolados. Segundo, a \textbf{qualidade do preenchimento da causa básica}: uma proporção dos óbitos é classificada como ``causas mal definidas'' (capítulo XVIII da CID-10), especialmente em regiões com menor cobertura de serviços de saúde --- proporção que também vem caindo de forma expressiva ao longo do tempo. Análises de mortalidade por causa específica devem sempre verificar a proporção de causas mal definidas no universo analisado e, se necessário, aplicar redistribuição proporcional.
\end{boxatencao}

\subsection{Acesso aos dados}

Os microdados do SIM são disponibilizados gratuitamente pelo DATASUS em formato DBC (comprimido) ou DBF. O portal TabNet (\url{tabnet.datasus.gov.br}) permite tabulações online sem necessidade de download. Para análises com microdados, a plataforma \textit{Base dos Dados} disponibiliza o SIM já tratado e compatibilizado em formato Parquet/BigQuery.

\begin{boxdica}{Dica --- Pacote \texttt{microdatasus} no R}
O pacote \texttt{microdatasus} (disponível no CRAN) facilita enormemente o download e a leitura dos microdados do SIM e de outros sistemas do DATASUS diretamente no R, convertendo automaticamente os arquivos DBC para formato utilizável e aplicando os dicionários de variáveis.
\end{boxdica}

% ────────────────────────────────────────────────────────────
\section{SINASC --- Sistema de Informações sobre Nascidos Vivos}
% ────────────────────────────────────────────────────────────

\ficharapida
  {SINASC --- Sistema de Informações sobre Nascidos Vivos}
  {Ministério da Saúde / DATASUS}
  {https://datasus.saude.gov.br/nascidos-vivos-desde-1994}
  {Brasil, Unidades da Federação e municípios}
  {Anual (desde 1990; dados mais consistentes a partir de 1994)}
  {Nascimento vivo}

\subsection{O que é e principais variáveis}

O SINASC registra todos os nascimentos vivos ocorridos no Brasil a partir das Declarações de Nascido Vivo (DNV) preenchidas nas maternidades e unidades de saúde. É a principal fonte de dados sobre saúde perinatal e materna no país.

Cada registro contém:

\begin{itemize}
  \item \textbf{Características do recém-nascido}: sexo, peso ao nascer, duração da gestação (semanas), índice de Apgar, presença de anomalia congênita;
  \item \textbf{Características da mãe}: idade, escolaridade, raça/cor, município de residência, número de filhos anteriores (nascidos vivos e mortos), número de consultas de pré-natal, tipo de parto (vaginal ou cesáreo);
  \item \textbf{Local do nascimento}: município de ocorrência, estabelecimento de saúde (quando aplicável) e local do parto (hospitalar, domiciliar, outro estabelecimento de saúde, via pública, outros).
\end{itemize}

\subsection{Aplicações típicas em avaliação}

\begin{itemize}
  \item \textbf{Denominador para taxas de mortalidade infantil}: o SINASC fornece o número de nascidos vivos por município, necessário para calcular a taxa de mortalidade infantil com dados do SIM;
  \item \textbf{Avaliação de políticas de pré-natal}: o número de consultas de pré-natal registrado no SINASC permite avaliar a cobertura e a qualidade do pré-natal por município;
  \item \textbf{Desigualdades no nascimento}: o cruzamento de variáveis como peso ao nascer, prematuridade e características socioeconômicas da mãe permite análises de equidade em saúde materna;
  \item \textbf{Impacto de políticas sobre cesáreas}: o SINASC é a fonte de referência para monitorar a proporção de partos cesáreos --- indicador de uso excessivo de intervenção cirúrgica no parto.
\end{itemize}

% ────────────────────────────────────────────────────────────
\section{SIH e SIA --- Internações e Atendimento Ambulatorial}
% ────────────────────────────────────────────────────────────

\ficharapida
  {SIH/SIA --- Sistemas de Informações Hospitalares e Ambulatoriais}
  {Ministério da Saúde / DATASUS}
  {https://datasus.saude.gov.br/producao-hospitalar-desde-1994}
  {Brasil, Unidades da Federação e municípios}
  {Mensal (SIH desde 1992; SIA desde 1994)}
  {Autorização de Internação Hospitalar (SIH) e procedimento ambulatorial (SIA)}

\subsection{O que são e o que cobrem}

O Sistema de Informações Hospitalares (SIH) e o Sistema de Informações Ambulatoriais (SIA) registram a \textbf{produção de serviços de saúde} financiados pelo SUS --- internações hospitalares e atendimentos ambulatoriais, respectivamente. Não cobrem serviços prestados pela saúde suplementar (planos de saúde privados) nem atendimentos particulares.

O SIH registra cada internação por meio da Autorização de Internação Hospitalar (AIH), que contém:
\begin{itemize}
  \item Diagnóstico principal e secundário (CID-10);
  \item Procedimento realizado;
  \item Tempo de permanência;
  \item Valor pago pelo SUS;
  \item Características do paciente (sexo, idade, município de residência);
  \item Estabelecimento de saúde (CNES).
\end{itemize}

O SIA registra cada procedimento ambulatorial realizado em estabelecimentos credenciados ao SUS, com informações sobre o tipo de procedimento, especialidade, estabelecimento e município.

\subsection{Aplicações típicas e limitações}

Os sistemas SIH e SIA são especialmente úteis para:
\begin{itemize}
  \item Avaliar \textbf{mudanças no padrão de utilização hospitalar} decorrentes de políticas de atenção primária (como a expansão do PSF, que teoricamente deveria reduzir internações evitáveis);
  \item Calcular \textbf{custos de procedimentos} para análises de custo-efetividade;
  \item Monitorar a \textbf{produção de serviços} de saúde por região e estabelecimento.
\end{itemize}

A limitação central é que SIH e SIA cobrem \textbf{apenas o SUS}: aproximadamente 25\% da população brasileira possui plano de saúde privado e seus atendimentos não constam nessas bases. Para análises da população como um todo, é preciso combinar dados do SUS com fontes que cubram a saúde suplementar.

% ────────────────────────────────────────────────────────────
\section{SINAN --- Sistema de Informação de Agravos de Notificação}
% ────────────────────────────────────────────────────────────

\ficharapida
  {SINAN --- Sistema de Informação de Agravos de Notificação}
  {Ministério da Saúde / DATASUS}
  {https://datasus.saude.gov.br/doencas-e-agravos-de-notificacao-de-2007-em-diante-sinan}
  {Brasil, Unidades da Federação e municípios}
  {Contínuo (desde 1993; integração com VIVA a partir de 2009)}
  {Caso notificado de agravo ou doença de notificação compulsória}

\subsection{O que é e principais doenças cobertas}

O SINAN reúne as notificações compulsórias de doenças e agravos à saúde --- ou seja, condições que, por lei, os profissionais e serviços de saúde são obrigados a comunicar às autoridades sanitárias. A lista de doenças de notificação compulsória é definida pelo Ministério da Saúde e inclui atualmente mais de 70 agravos.

Entre os mais relevantes para avaliação de políticas:

\begin{itemize}
  \item \textbf{Doenças infecciosas e parasitárias}: dengue, chikungunya, zika, malária, tuberculose, hanseníase, leishmaniose, hepatites virais;
  \item \textbf{Doenças imunopreveníveis}: sarampo, rubéola, coqueluche, poliomielite;
  \item \textbf{HIV/AIDS}: casos novos, perfil epidemiológico;
  \item \textbf{Violências e acidentes} (via integração com o VIVA): violência doméstica, sexual, autoprovocada, acidentes de trabalho.
\end{itemize}

\subsection{Aplicações típicas em avaliação}

\begin{itemize}
  \item \textbf{Avaliação de programas de vacinação}: a cobertura vacinal pode ser combinada com dados do SINAN sobre incidência das doenças-alvo para avaliar efetividade;
  \item \textbf{Políticas de controle de endemias}: monitoramento de dengue, malária e leishmaniose em resposta a intervenções de saneamento e controle vetorial;
  \item \textbf{Políticas de combate à tuberculose}: o SINAN é a referência para avaliação de programas de tratamento e controle da tuberculose, incluindo taxas de cura, abandono e resistência;
  \item \textbf{Violência doméstica e sexual}: desde a integração com o VIVA, o SINAN passou a ser uma das principais fontes sobre violência atendida em serviços de saúde.
\end{itemize}

% CORRIGIDO: título encurtado para caber na caixa
\begin{boxatencao}{Atenção --- Subnotificação no SINAN}
A notificação compulsória é obrigatória por lei, mas o cumprimento da obrigação varia enormemente entre doenças, regiões e tipos de serviço. Estima-se que para dengue, por exemplo, apenas uma fração dos casos reais é notificada --- especialmente em períodos de alta demanda, quando os serviços ficam sobrecarregados. Para doenças com maior estigma social (como HIV/AIDS e hepatites), a subnotificação pode ser ainda maior. Análises com dados do SINAN devem sempre discutir o potencial de subnotificação e, quando possível, triangular com outras fontes.
\end{boxatencao}

% ────────────────────────────────────────────────────────────
\section{CNES --- Cadastro Nacional de Estabelecimentos de Saúde}
% ────────────────────────────────────────────────────────────

\ficharapida
  {CNES --- Cadastro Nacional de Estabelecimentos de Saúde}
  {Ministério da Saúde / DATASUS}
  {https://cnes.datasus.gov.br}
  {Brasil, Unidades da Federação e municípios}
  {Contínuo (atualização mensal; desde 2005)}
  {Estabelecimento de saúde}

\subsection{O que é e principais variáveis}

O CNES é o cadastro de todos os estabelecimentos de saúde do Brasil --- públicos e privados, com e sem fins lucrativos, vinculados ou não ao SUS. Cada estabelecimento é identificado por um código único (CNES) que serve de chave de integração com os sistemas SIH, SIA e SINAN.

As principais variáveis disponíveis incluem:

\begin{itemize}
  \item \textbf{Tipo e natureza do estabelecimento}: hospital, UBS, UPA, clínica, laboratório, etc.; público, privado, filantrópico;
  \item \textbf{Localização}: município, UF, endereço georreferenciado;
  \item \textbf{Infraestrutura}: número de leitos por tipo (SUS, não-SUS, UTI), equipamentos disponíveis;
  \item \textbf{Recursos humanos}: profissionais de saúde vinculados por categoria (médicos, enfermeiros, técnicos), carga horária e vínculo;
  \item \textbf{Serviços oferecidos}: especialidades, procedimentos habilitados, serviços de apoio diagnóstico.
\end{itemize}

\subsection{Aplicações típicas em avaliação}

\begin{itemize}
  \item \textbf{Avaliação de acesso a serviços de saúde}: o CNES permite calcular indicadores de disponibilidade de serviços (leitos por mil habitantes, médicos por habitante) por município, essenciais para análises de equidade no acesso;
  \item \textbf{Avaliação de políticas de expansão da rede}: monitoramento da abertura ou fechamento de estabelecimentos em resposta a políticas de investimento em infraestrutura;
  \item \textbf{Variável de controle em avaliações de impacto}: características da oferta de saúde (número de leitos, cobertura de ESF) são frequentemente usadas como variáveis de controle em estudos que avaliam outros desfechos de saúde.
\end{itemize}

% ────────────────────────────────────────────────────────────
\section{SIOPS --- Sistema de Informações sobre Orçamentos Públicos em Saúde}
% ────────────────────────────────────────────────────────────

\ficharapida
  {SIOPS}
  {Ministério da Saúde}
  {https://siops.datasus.gov.br}
  {Brasil, Unidades da Federação e municípios}
  {Anual (desde 2000)}
  {Ente federativo (município, estado ou União)}

\subsection{O que é e aplicações}

O SIOPS reúne dados sobre o orçamento público destinado à saúde pelos municípios, estados e União, conforme a obrigação constitucional de aplicação mínima de recursos próprios na área. Cada ente declarante informa suas receitas totais e os gastos realizados em ações e serviços públicos de saúde (ASPS).

Para avaliação de políticas, o SIOPS é fundamental para:
\begin{itemize}
  \item Analisar o \textbf{nível de financiamento público da saúde} por município e UF ao longo do tempo;
  \item Calcular \textbf{gasto per capita em saúde} e comparar entre entes federativos;
  \item Avaliar o cumprimento dos \textbf{pisos constitucionais} de gasto em saúde (municípios devem aplicar mínimo de 15\% da receita própria; estados, 12\%).
\end{itemize}

% ────────────────────────────────────────────────────────────
\section{SISVAN --- Sistema de Vigilância Alimentar e Nutricional}
% ────────────────────────────────────────────────────────────

\ficharapida
  {SISVAN}
  {Ministério da Saúde / DATASUS}
  {https://sisvan.saude.gov.br}
  {Brasil, Unidades da Federação e municípios}
  {Contínuo (desde 2008 com cobertura ampliada)}
  {Indivíduo atendido na atenção básica}

\subsection{O que é e aplicações}

O SISVAN registra o estado nutricional e o consumo alimentar de indivíduos acompanhados pela atenção básica do SUS. Os dados são coletados nas consultas de rotina nas Unidades Básicas de Saúde e cobrem principalmente crianças menores de 5 anos, gestantes, mulheres em idade fértil e idosos.

As principais aplicações em avaliação incluem:
\begin{itemize}
  \item Monitoramento de \textbf{desnutrição e obesidade} por município, especialmente em populações vulneráveis atendidas pelo SUS;
  \item Avaliação de \textbf{políticas de segurança alimentar} e programas de suplementação nutricional;
  \item Acompanhamento nutricional de \textbf{beneficiários do Bolsa Família}, que têm condicionalidades de saúde vinculadas ao SISVAN.
\end{itemize}

A principal limitação do SISVAN é que cobre apenas indivíduos que comparecem às consultas na atenção básica --- portanto, não é representativo da população geral.

% ────────────────────────────────────────────────────────────
\section{SISAB --- Sistema de Informação em Saúde para a Atenção Básica}
% ────────────────────────────────────────────────────────────

\ficharapida
  {SISAB}
  {Ministério da Saúde / SAPS}
  {https://sisab.saude.gov.br}
  {Brasil, Unidades da Federação, municípios e equipes}
  {Contínuo (desde 2013, substituindo o SIAB)}
  {Indivíduo atendido, equipe e estabelecimento de atenção básica}

\subsection{O que é e aplicações}

O SISAB consolida a produção das equipes de atenção básica do SUS --- Estratégia Saúde da Família (ESF), Núcleos Ampliados de Saúde da Família (eMulti, sucessor do NASF-AB) e demais modelos de Unidade Básica de Saúde (UBS). Substituiu o antigo SIAB em 2013 e é alimentado pelo prontuário eletrônico do cidadão (PEC) e por formulários de CDS (Coleta de Dados Simplificada).

Para a avaliação de políticas públicas, o SISAB permite:
\begin{itemize}
  \item Acompanhar a \textbf{cobertura populacional da atenção básica} por município e por equipe, variável-chave em avaliações de impacto da ESF sobre desfechos de saúde;
  \item Construir indicadores de \textbf{produção de ações da AB} (consultas, visitas domiciliares, procedimentos, ações coletivas, visitas de ACS);
  \item Cruzar dados da atenção básica a sistemas hospitalares (SIH) e de mortalidade (SIM) via código CNES, permitindo análises do efeito da AB sobre hospitalizações sensíveis à atenção primária e mortalidade evitável.
\end{itemize}

Como no SISVAN, a limitação central é a cobertura restrita aos usuários efetivos da atenção básica e a heterogeneidade municipal no preenchimento.

% ────────────────────────────────────────────────────────────
\section{Pesquisa Nacional de Saúde (PNS)}
% ────────────────────────────────────────────────────────────

\ficharapida
  {Pesquisa Nacional de Saúde (PNS)}
  {IBGE em parceria com o Ministério da Saúde / Fiocruz}
  {https://www.ibge.gov.br/estatisticas/sociais/saude/9160-pesquisa-nacional-de-saude.html}
  {Brasil, Grandes Regiões, Unidades da Federação e capitais}
  {Quinquenal: 2013 e 2019}
  {Domicílio e morador}

\subsection{O que é e por que importa}

A Pesquisa Nacional de Saúde (PNS) é o principal inquérito domiciliar de saúde do Brasil, realizado pelo IBGE em parceria com o Ministério da Saúde e a Fiocruz. Diferentemente dos sistemas do DATASUS --- que captam apenas quem utiliza serviços de saúde ---, a PNS é uma pesquisa amostral representativa de toda a população residente em domicílios particulares, o que permite estimar a prevalência de condições de saúde na população geral, independentemente do uso de serviços.

\subsection{Principais temas cobertos}

A PNS abrange um conjunto amplo de dimensões da saúde individual e coletiva:

\begin{itemize}
  \item \textbf{Percepção do estado de saúde}: autoavaliação da saúde, limitações funcionais;
  \item \textbf{Doenças crônicas não transmissíveis}: hipertensão, diabetes, doenças cardíacas, AVC, câncer, problemas respiratórios crônicos;
  \item \textbf{Saúde mental}: depressão, ansiedade (edição 2019 com módulo específico);
  \item \textbf{Acesso e uso de serviços de saúde}: consultas médicas, internações, uso de plano de saúde, acesso à atenção básica;
  \item \textbf{Fatores de risco e proteção}: tabagismo, consumo de álcool, atividade física, alimentação;
  \item \textbf{Saúde da mulher}: mamografia, preventivo, pré-natal;
  \item \textbf{Características do domicílio e socioeconômicas}: escolaridade, renda, trabalho.
\end{itemize}

\subsection{Aplicações típicas em avaliação}

\begin{itemize}
  \item \textbf{Estimação de prevalências populacionais} de doenças crônicas e fatores de risco, impossível com os sistemas do DATASUS;
  \item \textbf{Avaliação de acesso a serviços}: comparação entre grupos populacionais (por renda, raça, região) no uso de consultas médicas, exames e internações;
  \item \textbf{Linha de base para políticas de saúde}: as edições de 2013 e 2019 fornecem um retrato do estado de saúde da população antes e depois de mudanças de política relevantes;
  \item \textbf{Desigualdades em saúde}: a PNS permite analisar gradientes socioeconômicos em saúde de forma muito mais rica do que as fontes administrativas.
\end{itemize}

\begin{boxatencao}{Atenção --- Periodicidade quinquenal}
A PNS é realizada a cada cinco anos, o que limita seu uso para avaliações de políticas de curto prazo. Entre edições, os sistemas do DATASUS são as principais fontes para acompanhamento de indicadores de saúde. A edição 2024/25 estava em planejamento no momento da elaboração deste guia; recomenda-se verificar o site do IBGE para informações atualizadas.
\end{boxatencao}

% ────────────────────────────────────────────────────────────
\section{AMS --- Pesquisa de Assistência Médico-Sanitária}
% ────────────────────────────────────────────────────────────

\ficharapida
  {AMS --- Pesquisa de Assistência Médico-Sanitária}
  {Instituto Brasileiro de Geografia e Estatística (IBGE)}
  {https://www.ibge.gov.br/estatisticas/sociais/saude/9067-pesquisa-de-assistencia-medico-sanitaria.html}
  {Brasil, Unidades da Federação e municípios}
  {Cross-sections: 1999, 2002, 2005, 2009 (descontinuada; substituída parcialmente pelo CNES)}
  {Estabelecimento de saúde}

\subsection{O que é e aplicações}

A AMS era um censo de todos os estabelecimentos de saúde do Brasil --- públicos e privados ---, realizado pelo IBGE. Diferentemente do CNES (que é um cadastro administrativo atualizado mensalmente), a AMS era uma pesquisa estruturada com questionários padronizados aplicados a cada estabelecimento, coletando informações detalhadas sobre infraestrutura, equipamentos, recursos humanos e serviços oferecidos.

Com a consolidação do CNES como fonte administrativa de referência, a AMS foi descontinuada após 2009. Para análises históricas da oferta de serviços de saúde no período 1999--2009, a AMS ainda é uma fonte relevante, especialmente para comparações que incluam estabelecimentos privados não vinculados ao SUS --- categoria que o CNES cobre com menor profundidade analítica.

% ────────────────────────────────────────────────────────────
\section{Síntese do capítulo}
% ────────────────────────────────────────────────────────────

Este capítulo percorreu o ecossistema de dados de saúde disponível para o Brasil. A Tabela~\ref{tab:sintese_cap5} resume as características essenciais de cada fonte.

\begin{table}[H]
\centering
\caption{Síntese das fontes de dados de saúde}
\label{tab:sintese_cap5}
\small
\begin{tabularx}{\textwidth}{@{} p{2cm} p{3.5cm} >{\raggedright\arraybackslash}p{2.5cm} X @{}}
\toprule
\textbf{Fonte} & \textbf{O que registra} & \textbf{Cobertura} & \textbf{Principal uso em avaliação} \\
\midrule
SIM    & Óbitos                     & Total (com subregistro) & Mortalidade, causas externas, violência \\
SINASC & Nascimentos vivos          & Total (com subregistro) & Saúde materna, prematuridade \\
SIH    & Internações SUS            & SUS apenas & Uso hospitalar, procedimentos \\
SIA    & Atendimentos ambulatoriais SUS & SUS apenas & Produção ambulatorial \\
SINAN  & Agravos notificáveis       & Notificados (subnotificação) & Vigilância epidemiológica \\
CNES   & Estabelecimentos de saúde  & Todos (públicos e privados) & Infraestrutura, acesso \\
SIOPS  & Orçamento e execução em saúde & Entes federativos & Financiamento do SUS \\
SISVAN & Estado nutricional         & Usuários da atenção básica & Nutrição, segurança alimentar \\
SISAB  & Produção da atenção básica & Usuários da atenção básica & Cobertura da ESF, ações da AB \\
PNS    & Saúde da população geral   & Amostra representativa & Prevalências, acesso, desigualdades \\
AMS    & Estabelecimentos (histórico) & Todos (descontinuada) & Séries históricas 1999--2009 \\
\bottomrule
\end{tabularx}
\fonte{FGV CLEAR, elaboração própria.}
\end{table}

Os pontos centrais do capítulo são:

\begin{itemize}
  \item O \textbf{SIM} e o \textbf{SINASC} são as fontes de referência para mortalidade e nascimentos, com cobertura nacional e séries históricas longas, mas sujeitos a subregistro diferencial por região;
  \item O \textbf{SIH} e o \textbf{SIA} registram a produção de serviços do SUS, mas não cobrem a saúde suplementar privada;
  \item O \textbf{SINAN} é fundamental para vigilância epidemiológica, mas a subnotificação é estrutural e varia muito entre agravos;
  \item O \textbf{CNES} é a referência para infraestrutura e recursos humanos em saúde, com atualização mensal;
  \item A \textbf{PNS} é a única fonte que permite estimar prevalências de condições de saúde na população geral, independentemente do uso de serviços --- mas sua periodicidade quinquenal limita o acompanhamento de curto prazo;
  \item A \textbf{combinação estratégica} das fontes --- por exemplo, SIM + SINASC para mortalidade infantil, ou CNES + SIH para análises de acesso e uso --- é o que permite avaliações mais robustas do setor saúde.
\end{itemize}

\medskip

O próximo capítulo apresenta as fontes sobre educação: o Censo Escolar, o SAEB, o ENEM, o Censo da Educação Superior e o SIOPE --- o ecossistema coordenado pelo Inep, que documenta matrículas, desempenho cognitivo, acesso ao ensino e financiamento da educação básica.
% ============================================================
%  cap06_educacao.tex
%  Capítulo 6: Educação
% ============================================================

% ════════════════════════════════════════════════════════════
\chapter{Educação}
% ════════════════════════════════════════════════════════════

\begin{boxdestaque}{Neste capítulo}
Este capítulo apresenta as principais fontes de dados sobre educação disponíveis para o Brasil. Cobrimos o Censo Escolar, o Censo da Educação Superior, o Sistema de Avaliação da Educação Básica (SAEB), o ENEM, o SIOPE e outras fontes complementares. Para cada fonte, detalhamos histórico, estrutura, variáveis disponíveis, forma de acesso e aplicações típicas em avaliação de políticas educacionais.
\end{boxdestaque}

% ────────────────────────────────────────────────────────────
\section{Introdução: o ecossistema de dados de educação no Brasil}
% ────────────────────────────────────────────────────────────

A educação brasileira é monitorada por um sistema de informações relativamente maduro, coordenado principalmente pelo Instituto Nacional de Estudos e Pesquisas Educacionais Anísio Teixeira (Inep), autarquia vinculada ao Ministério da Educação. O Inep produz e mantém os principais censos, avaliações e indicadores educacionais do país, combinando registros administrativos (como o Censo Escolar) com avaliações de desempenho em larga escala (como o SAEB e o ENEM).

Esse ecossistema oferece ao avaliador uma combinação rara: dados administrativos com cobertura universal (o Censo Escolar registra cada aluno, professor e escola individualmente), avaliações cognitivas padronizadas com séries históricas de mais de duas décadas, e dados de financiamento que permitem análises de eficiência e alocação de recursos.

\begin{table}[H]
\centering
\caption{Principais fontes de dados de educação e suas coberturas}
\label{tab:ecossistema_edu}
\small
\begin{tabularx}{\textwidth}{@{} p{2.5cm} p{3.5cm} >{\raggedright\arraybackslash}p{2.5cm} X @{}}
\toprule
\textbf{Fonte} & \textbf{O que registra} & \textbf{Desde} & \textbf{Principal uso em avaliação} \\
\midrule
Censo Escolar   & Matrículas, escolas, docentes & 1932 & Cobertura, acesso, infraestrutura \\
Censo Ed. Superior & IES, cursos, alunos, docentes & 2001 & Acesso ao ensino superior \\
SAEB            & Desempenho cognitivo (básica) & 1990 & Qualidade do aprendizado \\
ENEM            & Desempenho no ensino médio  & 1998 & Acesso ao ensino superior, desigualdade \\
SIOPE           & Gasto público em educação   & 1999 & Financiamento, custo-aluno \\
IDEB            & Índice sintético de qualidade & 2005 & Monitoramento da educação básica \\
\bottomrule
\end{tabularx}
\fonte{FGV CLEAR, elaboração própria.}
\end{table}

% ────────────────────────────────────────────────────────────
\section{Censo Escolar}
% ────────────────────────────────────────────────────────────

\ficharapida
  {Censo Escolar da Educação Básica}
  {Instituto Nacional de Estudos e Pesquisas Educacionais Anísio Teixeira (Inep)}
  {https://www.gov.br/inep/pt-br/areas-de-atuacao/pesquisas-estatisticas-e-indicadores/censo-escolar}
  {Brasil, Unidades da Federação, municípios e escolas}
  {Anual (desde 1932; microdados públicos a partir de 1995)}
  {Escola, turma, aluno e docente}

\subsection{O que é e por que importa}

O Censo Escolar da Educação Básica é o mais abrangente levantamento estatístico da educação brasileira. Realizado anualmente pelo Inep desde 1932 — tornando-o uma das séries históricas mais longas do sistema estatístico nacional —, o Censo Escolar registra individualmente cada escola, cada turma, cada aluno e cada profissional de educação do Brasil, cobrindo tanto a rede pública (federal, estadual e municipal) quanto a rede privada.

A coleta ocorre anualmente com data de referência na última quarta-feira de maio. As informações são prestadas pelos próprios estabelecimentos de ensino por meio do sistema Educacenso, que substituiu os questionários em papel a partir de 2007. Essa transição para a coleta eletrônica melhorou substancialmente a qualidade e a tempestividade dos dados.

Para a avaliação de políticas educacionais, o Censo Escolar é insubstituível por três razões: cobertura universal (não é amostra — registra todos os alunos e escolas), granularidade (permite análises no nível da escola individual, da turma ou do aluno), e longitudinalidade (o código único de cada aluno permite acompanhar sua trajetória escolar ao longo dos anos).

\subsection{Estrutura e unidades de análise}

O Censo Escolar é organizado em quatro arquivos principais, cada um com uma unidade de análise distinta:

\begin{itemize}
  \item \textbf{Escolas}: cada linha é uma escola. Contém informações sobre localização (município, zona urbana/rural), dependência administrativa (federal, estadual, municipal, privada), infraestrutura (laboratórios, biblioteca, quadra, acesso à internet, saneamento), equipamentos disponíveis e etapas de ensino ofertadas;

  \item \textbf{Turmas}: cada linha é uma turma. Contém informações sobre etapa e modalidade de ensino, turno, tamanho da turma e código da escola;

  \item \textbf{Alunos}: cada linha é uma matrícula. Contém informações sobre o aluno (sexo, data de nascimento, raça/cor, deficiência, município de residência) e sobre a matrícula (etapa, série, turno, situação ao final do ano — aprovado, reprovado, abandonou). O código único do aluno (\texttt{CO\_PESSOA\_FISICA}) permite construir um painel longitudinal seguindo o mesmo aluno ao longo de vários anos;

  \item \textbf{Docentes}: cada linha é um vínculo de docente. Contém informações sobre o professor (sexo, idade, raça/cor, escolaridade, formação) e sobre o vínculo (escola, disciplinas lecionadas, carga horária, tipo de contrato).
\end{itemize}

\subsection{Principais variáveis de interesse}

\textbf{No arquivo de alunos:}
\begin{itemize}
  \item Etapa e modalidade: educação infantil, ensino fundamental (anos iniciais e finais), ensino médio, EJA, educação especial, educação profissional;
  \item Situação de rendimento ao final do ano (aprovado, reprovado, transferido, falecido, abandono);
  \item Indicadores de vulnerabilidade: se recebe Bolsa Família, se é quilombola, indígena, com deficiência, em situação de itinerância;
  \item Transporte escolar público.
\end{itemize}

\textbf{No arquivo de escolas:}
\begin{itemize}
  \item Infraestrutura: existência de água tratada, esgoto, energia elétrica, internet banda larga, laboratório de informática, laboratório de ciências, biblioteca/sala de leitura, quadra esportiva;
  \item Localização: município, UF, zona (urbana/rural), se localizada em área quilombola ou indígena;
  \item Alimentação escolar, atendimento educacional especializado.
\end{itemize}

\subsection{O IDEB — Índice de Desenvolvimento da Educação Básica}

O IDEB é o principal indicador sintético de qualidade da educação básica brasileira, calculado pelo Inep a partir de 2005 com base em dois componentes:

\begin{enumerate}
  \item \textbf{Desempenho}: médias de proficiência em Língua Portuguesa e Matemática nas avaliações do SAEB/Prova Brasil;
  \item \textbf{Fluxo}: taxa de aprovação, calculada a partir dos dados do Censo Escolar.
\end{enumerate}

A fórmula do IDEB penaliza tanto o baixo desempenho quanto a alta reprovação ou abandono, capturando a ideia de que uma escola que aprova todos os alunos sem ensiná-los, ou que ensina bem mas reprova muitos, não tem boa qualidade educacional.

O IDEB é calculado e publicado para os anos iniciais do ensino fundamental (5º ano), anos finais (9º ano) e ensino médio (3ª série), com desagregação por escola, município, UF e Brasil. Está disponível no site do Inep e pode ser acessado diretamente ou por meio de pacotes de R como o \texttt{idebr}.

% CORRIGIDO: título encurtado para caber na caixa
\begin{boxexemplo}{Exemplo — Infraestrutura escolar e rendimento}
Um estado implementou um programa de reforma e ampliação de escolas rurais, priorizando municípios com infraestrutura mais precária. Para avaliar o impacto do programa sobre a frequência e o rendimento escolar, pesquisadores utilizaram o painel longitudinal do Censo Escolar: identificaram as escolas beneficiadas pelo programa (grupo de tratamento) e escolas com características similares que não foram beneficiadas (grupo de controle, via pareamento por escore de propensão com variáveis de infraestrutura do Censo Escolar de anos anteriores). A variável de resultado foi a taxa de abandono e a taxa de aprovação nos dois anos seguintes à reforma. O código único do aluno permitiu controlar pela composição do alunado em cada escola ao longo do tempo.

\textit{Fonte: FGV CLEAR, elaboração própria.}
\end{boxexemplo}

\subsection{Acesso aos microdados}

Os microdados do Censo Escolar são disponibilizados gratuitamente pelo Inep, com os quatro arquivos (escolas, turmas, alunos, docentes) em formato CSV. A defasagem entre o período de referência (maio) e a publicação dos microdados costuma ser de seis a doze meses.

\begin{boxdica}{Dica — Pacotes para trabalhar com o Censo Escolar}
\begin{itemize}
  \item \textbf{R}: o pacote \texttt{educaR} facilita o download e a leitura dos microdados do Censo Escolar, com dicionários de variáveis já integrados;
  \item \textbf{Python}: a plataforma \textit{Base dos Dados} disponibiliza o Censo Escolar compatibilizado e tratado em BigQuery, com tabelas linkáveis por código de escola e de aluno;
  \item \textbf{Stata}: os arquivos CSV podem ser importados com \texttt{import delimited}; o dicionário de variáveis do Inep deve ser consultado para decodificar as variáveis categóricas.
\end{itemize}
\end{boxdica}

\subsection{Limitações e cuidados}

\begin{itemize}
  \item \textbf{Data de referência única}: o Censo Escolar captura um retrato de maio de cada ano. Alunos que se matriculam após essa data — ou que abandonam antes — podem não ser adequadamente captados;
  \item \textbf{Mudanças no sistema de coleta}: a transição para o Educacenso em 2007 introduziu mudanças que limitam comparações com edições anteriores para algumas variáveis;
  \item \textbf{Qualidade do preenchimento}: embora melhorada pela coleta eletrônica, ainda há inconsistências no preenchimento de campos como raça/cor e deficiência, especialmente em redes municipais menores.
\end{itemize}

% ────────────────────────────────────────────────────────────
\section{Censo da Educação Superior}
% ────────────────────────────────────────────────────────────

\ficharapida
  {Censo da Educação Superior}
  {Instituto Nacional de Estudos e Pesquisas Educacionais Anísio Teixeira (Inep)}
  {https://www.gov.br/inep/pt-br/areas-de-atuacao/pesquisas-estatisticas-e-indicadores/censo-da-educacao-superior}
  {Brasil, Unidades da Federação, municípios e instituições}
  {Anual (desde 2001)}
  {Instituição de Educação Superior (IES), curso, aluno e docente}

\subsection{O que é e principais variáveis}

O Censo da Educação Superior registra anualmente todas as Instituições de Educação Superior (IES) do Brasil — públicas e privadas — e seus cursos de graduação e sequenciais. Criado em 2001 pelo Inep, o censo produz estatísticas sobre:

\begin{itemize}
  \item \textbf{Instituições}: categoria administrativa (federal, estadual, municipal, privada com e sem fins lucrativos), organização acadêmica (universidade, centro universitário, faculdade), município sede, número de cursos e vagas ofertadas;
  \item \textbf{Cursos}: área do conhecimento, grau (bacharelado, licenciatura, tecnológico), modalidade (presencial ou a distância), número de vagas, ingressantes, matriculados e concluintes;
  \item \textbf{Alunos}: sexo, idade, raça/cor, situação socioeconômica (para alunos do ProUni e FIES), modalidade de ingresso (ENEM/SISU, vestibular, outros);
  \item \textbf{Docentes}: titulação (graduado, especialista, mestre, doutor), regime de trabalho, sexo e raça/cor.
\end{itemize}

\subsection{Aplicações típicas em avaliação}

\begin{itemize}
  \item \textbf{Avaliação de políticas de expansão do acesso ao ensino superior}: o Censo da Educação Superior é a principal fonte para avaliar o impacto do REUNI, do ProUni e do FIES sobre matrículas e perfil dos estudantes;
  \item \textbf{Análise de desigualdades no acesso}: evolução da proporção de estudantes negros, de baixa renda e de escolas públicas no ensino superior ao longo do tempo;
  \item \textbf{Avaliação de políticas de cotas}: comparação do perfil dos ingressantes antes e depois da Lei de Cotas (Lei n.º 12.711/2012);
  \item \textbf{Monitoramento da educação a distância}: expansão das matrículas EAD e seu perfil demográfico.
\end{itemize}

% ────────────────────────────────────────────────────────────
\section{SAEB — Sistema de Avaliação da Educação Básica}
% ────────────────────────────────────────────────────────────

\ficharapida
  {SAEB — Sistema de Avaliação da Educação Básica}
  {Instituto Nacional de Estudos e Pesquisas Educacionais Anísio Teixeira (Inep)}
  {https://www.gov.br/inep/pt-br/areas-de-atuacao/avaliacao-e-exames-educacionais/saeb}
  {Brasil, Unidades da Federação, municípios e escolas (para a Prova Brasil)}
  {Bienal (desde 1990; Prova Brasil desde 2005)}
  {Aluno (com agregações para escola, município e UF)}

\subsection{O que é e estrutura}

O Sistema de Avaliação da Educação Básica (SAEB) é o principal sistema de avaliação em larga escala da educação básica brasileira. Criado em 1990, passou por diversas reformulações ao longo dos anos. Em sua configuração atual, o SAEB engloba três avaliações principais:

\begin{itemize}
  \item \textbf{ANEB} (Avaliação Nacional da Educação Básica): amostral, aplicada a estudantes do 5º e 9º anos do ensino fundamental e da 3ª série do ensino médio em escolas públicas e privadas. Produz resultados representativos para o Brasil, regiões e UFs;

  \item \textbf{ANRESC/Prova Brasil}: censitária para escolas públicas urbanas com mais de 20 alunos nas séries avaliadas. Produz resultados para cada escola individualmente, o que a torna muito mais útil para avaliações de impacto no nível de escola;

  \item \textbf{ANA} (Avaliação Nacional da Alfabetização): aplicada ao 3º ano do ensino fundamental, avalia alfabetização em Língua Portuguesa e Matemática. Incorporada ao SAEB a partir de 2019.
\end{itemize}

As avaliações medem proficiência em \textbf{Língua Portuguesa} (leitura e interpretação) e \textbf{Matemática}, com escalas padronizadas que permitem comparações ao longo do tempo. Cada aluno avaliado também responde a um \textbf{questionário contextual} sobre características socioeconômicas, hábitos de estudo e percepção da escola — informações fundamentais para análises de equidade.

\subsection{Escala de proficiência e comparabilidade}

As notas do SAEB são expressas em uma escala de proficiência que vai de 0 a 500 pontos (aproximadamente), com âncoras que descrevem o que os alunos em diferentes níveis são capazes de fazer. A escala é construída por Teoria de Resposta ao Item (TRI), o que teoricamente permite comparar o desempenho entre anos diferentes.

\begin{boxatencao}{Atenção — Mudanças metodológicas no SAEB}
O SAEB passou por reformulações metodológicas significativas ao longo de sua história — em particular em 1995, 2005 e 2019 — que afetam a comparabilidade das séries históricas. A edição de 2019 incorporou a ANA e trouxe novos itens e matrizes de referência. Análises de tendência de longo prazo devem verificar cuidadosamente quais edições são metodologicamente comparáveis. O Inep publica documentação técnica detalhada sobre cada mudança.
\end{boxatencao}

\subsection{Aplicações típicas em avaliação}

\begin{itemize}
  \item \textbf{Avaliação de impacto de políticas educacionais}: a Prova Brasil — por produzir resultados no nível de escola — é amplamente usada em avaliações de programas que intervêm em escolas específicas (como programas de formação docente, de gestão escolar ou de infraestrutura);
  \item \textbf{Análise de desigualdades educacionais}: o questionário contextual permite decompor as diferenças de desempenho por características socioeconômicas, raça/cor e localização;
  \item \textbf{Monitoramento do IDEB}: as notas do SAEB são um dos dois componentes do IDEB, tornando-o essencial para entender a evolução do índice;
  \item \textbf{Estudos de eficácia escolar}: identificação de escolas que produzem desempenho acima do esperado dado o perfil socioeconômico de seus alunos.
\end{itemize}

\subsection{Acesso aos microdados}

Os microdados do SAEB são disponibilizados gratuitamente pelo Inep após cada edição da avaliação. Os arquivos contêm as notas individuais dos alunos, as respostas ao questionário contextual e os identificadores de escola e município. A anonimização impede a identificação direta dos alunos, mas os identificadores de escola permitem a integração com o Censo Escolar.

% ────────────────────────────────────────────────────────────
\section{ENEM — Exame Nacional do Ensino Médio}
% ────────────────────────────────────────────────────────────

\ficharapida
  {ENEM — Exame Nacional do Ensino Médio}
  {Instituto Nacional de Estudos e Pesquisas Educacionais Anísio Teixeira (Inep)}
  {https://www.gov.br/inep/pt-br/areas-de-atuacao/avaliacao-e-exames-educacionais/enem}
  {Brasil (participação individual, não representativa da população)}
  {Anual (desde 1998; reestruturado em 2009)}
  {Participante individual}

\subsection{O que é e evolução}

O Exame Nacional do Ensino Médio (ENEM) foi criado em 1998 com o objetivo de avaliar o desempenho dos estudantes ao final da educação básica. Em 2009, passou por uma reestruturação significativa: adotou a TRI como metodologia de cálculo das notas (tornando-as comparáveis entre anos) e passou a ser utilizado como principal critério de acesso ao ensino superior por meio do SISU (Sistema de Seleção Unificada), substituindo progressivamente os vestibulares tradicionais das universidades federais.

O ENEM avalia quatro áreas do conhecimento (Linguagens, Matemática, Ciências da Natureza e Ciências Humanas) e uma redação, com notas em escala TRI de 0 a 1000 para cada área.

\subsection{Questionário socioeconômico}

Um dos ativos mais valiosos do ENEM para pesquisa em políticas públicas é o \textbf{questionário socioeconômico} respondido por todos os inscritos, que coleta informações detalhadas sobre:

\begin{itemize}
  \item Renda familiar mensal (em faixas);
  \item Escolaridade dos pais;
  \item Raça/cor autodeclarada;
  \item Tipo de escola cursada no ensino médio (pública estadual, federal, privada);
  \item Acesso à internet e posse de bens no domicílio;
  \item Se trabalha ou já trabalhou.
\end{itemize}

Esse questionário, combinado com as notas, permite análises riquíssimas sobre desigualdades educacionais por origem social.

\subsection{Limitações importantes}

O ENEM tem duas limitações estruturais que o analista deve ter em mente:

\begin{itemize}
  \item \textbf{Não é representativo}: a participação é voluntária. O perfil dos inscritos é sistematicamente diferente do perfil de todos os concluintes do ensino médio — há viés de seleção em direção a estudantes com maior motivação para o ensino superior e de escolas privadas. Análises que tratem o ENEM como censo do desempenho educacional estão erradas;

  \item \textbf{Mudança de uso em 2009}: a adoção do ENEM como critério de acesso ao ensino superior alterou profundamente o perfil dos participantes e seus incentivos. Comparações de desempenho entre o ENEM anterior e posterior a 2009 são problemáticas.
\end{itemize}

\subsection{Aplicações típicas em avaliação}

\begin{itemize}
  \item \textbf{Avaliação de políticas de cotas e acesso ao ensino superior}: o ENEM/SISU permite analisar como mudanças nas regras de acesso afetam o perfil dos ingressantes nas universidades federais;
  \item \textbf{Análises de desigualdade educacional}: o questionário socioeconômico detalhado permite decompor diferenças de desempenho por renda, raça e tipo de escola;
  \item \textbf{Avaliação de escolas}: para escolas com número suficiente de participantes, é possível construir indicadores de desempenho médio e compará-los com o perfil socioeconômico dos alunos;
  \item \textbf{Impacto de programas de preparação}: avaliações de cursinhos populares, programas de reforço e outras intervenções voltadas ao acesso ao ensino superior.
\end{itemize}

% ────────────────────────────────────────────────────────────
\section{SIOPE — Sistema de Informações sobre Orçamentos Públicos em Educação}
% ────────────────────────────────────────────────────────────

\ficharapida
  {SIOPE}
  {Fundo Nacional de Desenvolvimento da Educação (FNDE) / MEC}
  {https://www.fnde.gov.br/siope}
  {Brasil, Unidades da Federação e municípios}
  {Anual (desde 1999)}
  {Ente federativo (município, estado ou União)}

\subsection{O que é e aplicações}

O SIOPE reúne informações sobre receitas e despesas com educação de todos os entes federativos brasileiros, em cumprimento à obrigação constitucional de aplicação mínima de recursos próprios em educação (municípios: 25\% da receita; estados e União: 18\%). Operacionalizado pelo FNDE, o sistema permite acompanhar o cumprimento dos pisos constitucionais e comparar o esforço fiscal de diferentes entes.

Para avaliação de políticas educacionais, o SIOPE é fundamental para:

\begin{itemize}
  \item Calcular o \textbf{gasto público por aluno} por município e UF, indicador central de análises de custo-efetividade na educação;
  \item Avaliar o impacto de mudanças no \textbf{FUNDEB} (Fundo de Manutenção e Desenvolvimento da Educação Básica) sobre o gasto educacional municipal;
  \item Analisar a \textbf{composição do gasto} por subfunção (ensino fundamental, médio, infantil, EJA) e natureza (pessoal, custeio, investimento).
\end{itemize}

% CORRIGIDO: título encurtado para caber na caixa
\begin{boxexemplo}{Exemplo — SIOPE e Censo Escolar para calcular custo-aluno}
Para avaliar a eficiência de diferentes redes municipais de educação, pesquisadores combinaram dados do SIOPE (gasto total em educação fundamental por município) com dados do Censo Escolar (número de matrículas no ensino fundamental) para calcular o custo-aluno efetivo de cada rede. Esse indicador foi então cruzado com o IDEB dos municípios para identificar redes que obtinham resultados acima da média com gasto abaixo da média — e investigar o que essas redes faziam de diferente.

\textit{Fonte: FGV CLEAR, elaboração própria.}
\end{boxexemplo}

% ────────────────────────────────────────────────────────────
\section{Fontes complementares de dados educacionais}
% ────────────────────────────────────────────────────────────

\subsection{PISA — Programme for International Student Assessment}

O PISA é uma avaliação internacional coordenada pela OCDE que mede o desempenho de estudantes de 15 anos em Leitura, Matemática e Ciências. Realizado a cada três anos (2000, 2003, 2006, 2009, 2012, 2015, 2018, 2022), com um domínio principal em cada edição, o PISA permite comparar o desempenho educacional do Brasil com o de mais de 80 países participantes.

Para a avaliação de políticas, o PISA é útil para contextualizar o desempenho educacional brasileiro em perspectiva internacional e para avaliar o quanto o Brasil avançou — ou regrediu — em relação a países com trajetórias similares. Sua limitação é que a amostra brasileira é representativa do grupo etário (15 anos) a nível nacional, sem desagregação estadual ou municipal.

\subsection{Censo do Professor e Pesquisa do Professor}

O Inep realizou, em anos específicos, levantamentos focados nos docentes da educação básica que complementam as informações do Censo Escolar sobre o perfil e as condições de trabalho dos professores brasileiros. Esses dados são especialmente úteis para avaliações de políticas de valorização docente e de formação inicial e continuada.

\subsection{Plataforma Nilo Peçanha (PNP)}

A Plataforma Nilo Peçanha reúne dados sobre a Rede Federal de Educação Profissional, Científica e Tecnológica (Institutos Federais, CEFETs e outras instituições federais de educação profissional). Disponibiliza anualmente informações sobre matrículas, cursos, docentes e servidores técnico-administrativos de toda a rede federal de educação profissional.

% ────────────────────────────────────────────────────────────
\section{Síntese do capítulo}
% ────────────────────────────────────────────────────────────

Este capítulo apresentou o ecossistema de dados de educação disponível para o Brasil. A Tabela~\ref{tab:sintese_cap6} resume as características essenciais de cada fonte.

\begin{table}[H]
\centering
\caption{Síntese das fontes de dados de educação}
\label{tab:sintese_cap6}
\small
\begin{tabularx}{\textwidth}{@{} p{2.5cm} p{3cm} >{\raggedright\arraybackslash}p{2.5cm} X @{}}
\toprule
\textbf{Fonte} & \textbf{O que registra} & \textbf{Cobertura} & \textbf{Principal uso em avaliação} \\
\midrule
Censo Escolar      & Matrículas, escolas, docentes & Universal & Acesso, fluxo, infraestrutura, painel de alunos \\
Censo Ed. Superior & IES, cursos, alunos           & Universal & Acesso ao ensino superior, cotas, perfil \\
SAEB/Prova Brasil  & Desempenho cognitivo          & Censitária (escolas públicas) & Qualidade do aprendizado, IDEB \\
ENEM               & Desempenho no EM              & Voluntária (não representativa) & Acesso ao superior, desigualdade \\
SIOPE              & Gasto público em educação     & Entes federativos & Financiamento, custo-aluno \\
PISA               & Desempenho (15 anos)          & Amostral nacional & Comparação internacional \\
\bottomrule
\end{tabularx}
\fonte{FGV CLEAR, elaboração própria.}
\end{table}

Os pontos centrais do capítulo são:

\begin{itemize}
  \item O \textbf{Censo Escolar} é a fonte fundamental da educação brasileira: cobertura universal, granularidade no nível do aluno e série histórica longa permitem análises que nenhuma outra fonte viabiliza;
  \item O \textbf{SAEB/Prova Brasil} é a referência para avaliações de desempenho cognitivo, especialmente por produzir resultados no nível de escola — essencial para estudos de impacto de intervenções escolares;
  \item O \textbf{ENEM} tem um questionário socioeconômico muito rico, mas não é representativo: seu uso para inferências sobre a população de concluintes do ensino médio exige cuidados metodológicos;
  \item O \textbf{SIOPE} é indispensável para análises de financiamento e custo-efetividade na educação;
  \item A \textbf{combinação estratégica} dessas fontes — como Censo Escolar + SAEB para o IDEB, ou SIOPE + Censo Escolar para custo-aluno — é o que permite avaliações mais robustas das políticas educacionais.
\end{itemize}

O próximo capítulo apresenta as fontes sobre assistência social e transferência de renda: CadÚnico, Bolsa Família, BPC e Censo SUAS --- as bases que permitem identificar quem são as famílias atendidas pelos programas sociais, mapear sua cobertura e avaliar o impacto dos benefícios.
% ============================================================
%  cap07_assistencia_social.tex
%  Capítulo 7: Assistência Social e Transferência de Renda
% ============================================================

% ════════════════════════════════════════════════════════════
\chapter{Assistência Social e Transferência de Renda}
% ════════════════════════════════════════════════════════════

\begin{boxdestaque}{Neste capítulo}
Este capítulo apresenta as principais fontes de dados sobre assistência social e transferência de renda no Brasil. Cobrimos o Bolsa Família e seu histórico de dados, o Benefício de Prestação Continuada (BPC), o Censo SUAS e os sistemas de informação do Sistema Único de Assistência Social. Para cada fonte, detalhamos histórico, estrutura, variáveis disponíveis, forma de acesso e aplicações típicas em avaliação de políticas sociais.
\end{boxdestaque}

% ────────────────────────────────────────────────────────────
\section{Introdução: o ecossistema de dados da assistência social}
% ────────────────────────────────────────────────────────────

A política de assistência social brasileira é estruturada em torno do Sistema Único de Assistência Social (SUAS), criado em 2005, que organiza serviços, benefícios e programas em uma rede integrada de proteção social. O Ministério do Desenvolvimento e Assistência Social (MDS) é o principal produtor de dados dessa área, operando sistemas que cobrem desde o cadastro das famílias beneficiárias até o monitoramento dos equipamentos da rede socioassistencial.

Para a avaliação de políticas públicas, a assistência social oferece um conjunto de fontes administrativas de qualidade crescente. O Cadastro Único — já tratado no Capítulo~4 como fonte de dados sobre mercado de trabalho e renda — é a espinha dorsal desse sistema: praticamente todos os programas sociais federais usam o CadÚnico como base de elegibilidade. Os dados do Bolsa Família, do BPC e do SUAS complementam esse cadastro com informações sobre benefícios recebidos, condicionalidades cumpridas e serviços acessados.

\begin{table}[H]
\centering
\caption{Principais fontes de dados da assistência social}
\label{tab:ecossistema_suas}
\small
\begin{tabularx}{\textwidth}{@{} p{3cm} p{3.5cm} p{2cm} X @{}}
\toprule
\textbf{Fonte} & \textbf{O que registra} & \textbf{Desde} &
\textbf{Principal uso em avaliação} \\
\midrule
CadÚnico         & Famílias de baixa renda     & 2001 & Elegibilidade, perfil do público-alvo \\
Bolsa Família    & Beneficiários e pagamentos  & 2004 & Avaliação de transferência de renda \\
BPC              & Beneficiários idosos e PcD  & 1996 & Proteção social não contributiva \\
Censo SUAS       & Equipamentos e equipes      & 2007 & Infraestrutura socioassistencial \\
RMA/SUAS         & Atendimentos mensais        & 2012 & Produção de serviços socioassistenciais \\
\bottomrule
\end{tabularx}
\fonte{FGV CLEAR, elaboração própria.}
\end{table}

% ────────────────────────────────────────────────────────────
\section{Bolsa Família — histórico e dados disponíveis}
% ────────────────────────────────────────────────────────────

\ficharapida
  {Bolsa Família / Auxílio Brasil / Bolsa Família (novo)}
  {Ministério do Desenvolvimento e Assistência Social (MDS)}
  {https://www.gov.br/mds/pt-br/acoes-e-programas/transferencia-de-renda/bolsa-familia}
  {Brasil, Unidades da Federação e municípios}
  {Mensal (desde 2004)}
  {Família beneficiária}

\subsection{Histórico do programa e das fontes de dados}

O Programa Bolsa Família foi criado em outubro de 2003 (Lei n.º 10.836/2004), unificando quatro programas de transferência de renda então existentes: Bolsa Escola, Bolsa Alimentação, Auxílio Gás e Cartão Alimentação. Desde sua criação, tornou-se o maior programa de transferência condicionada de renda do mundo em número de beneficiários, chegando a cobrir mais de 14 milhões de famílias em seu auge.

Em 2021, o programa foi renomeado \textbf{Auxílio Brasil} e passou por reformulações nas regras de elegibilidade e nos valores dos benefícios. Em 2023, foi recriado como \textbf{Bolsa Família} com novas regras, incluindo o benefício mínimo de R\$~600,00 por família e o adicional por criança.

Para fins analíticos, é importante distinguir três períodos:

\begin{enumerate}
  \item \textbf{Bolsa Família original (2004--2021)}: série histórica mais longa e metodologicamente consistente;
  \item \textbf{Auxílio Brasil (2021--2022)}: período de transição com mudanças nas regras de elegibilidade e valores, o que limita comparações diretas com o período anterior;
  \item \textbf{Bolsa Família novo (2023--atual)}: nova reformulação com regras distintas de benefício por composição familiar.
\end{enumerate}

\subsection{Fontes de dados disponíveis}

Os dados do Bolsa Família estão disponíveis em múltiplas plataformas com diferentes níveis de desagregação:

\subsubsection{Portal da Transparência do Governo Federal}

O Portal da Transparência (\url{portaldatransparencia.gov.br}) disponibiliza mensalmente os dados de pagamentos do Bolsa Família com as seguintes informações por beneficiário:

\begin{itemize}
  \item NIS do responsável familiar;
  \item Nome do responsável (até 2020; posteriormente anonimizado);
  \item Município e UF de residência;
  \item Valor do benefício recebido no mês;
  \item Quantidade de membros da família.
\end{itemize}

Esses dados permitem construir séries históricas municipais de cobertura (número de famílias beneficiárias) e valor médio dos benefícios, além de análises individuais quando linkados com o CadÚnico pelo NIS.

\subsubsection{CECAD — Consulta, Seleção e Extração de Informações do CadÚnico}

O CECAD (\url{cecad.cidadania.gov.br}) é a plataforma de acesso a dados agregados do CadÚnico e do Bolsa Família por município. Permite consultar, sem necessidade de convênio:

\begin{itemize}
  \item Número de famílias beneficiárias por município e perfil de renda;
  \item Composição familiar dos beneficiários (número de crianças, gestantes, idosos);
  \item Taxas de cobertura estimadas em relação à população elegível.
\end{itemize}

\subsubsection{VISDATA — Visualização de Dados do MDS}

O VISDATA (\url{aplicacoes.mds.gov.br/sagi/vis/data3}) oferece painéis interativos com indicadores do Bolsa Família, do CadÚnico e de outros programas do MDS, com desagregação municipal e séries históricas. É a ferramenta mais acessível para gestores que precisam de indicadores sem trabalhar diretamente com microdados.

\subsection{Condicionalidades: dados de saúde e educação}

Uma característica central do Bolsa Família é o sistema de \textbf{condicionalidades}: para receber o benefício, as famílias devem cumprir compromissos nas áreas de saúde (vacinação, pré-natal, acompanhamento nutricional) e educação (frequência escolar mínima de 85\% para crianças de 6 a 15 anos e 75\% para jovens de 16 e 17 anos).

O acompanhamento das condicionalidades gera dados específicos que são fontes valiosas para avaliação:

\begin{itemize}
  \item \textbf{Condicionalidade de educação}: registros de frequência escolar dos beneficiários, coletados semestralmente pelas escolas e consolidados pelo MEC/Inep. Esses dados podem ser cruzados com o Censo Escolar pelo código do aluno;
  \item \textbf{Condicionalidade de saúde}: registros de acompanhamento do calendário vacinal, consultas de pré-natal e acompanhamento nutricional, coletados nas UBS e registrados no SISVAN e no sistema de acompanhamento de condicionalidades do MDS.
\end{itemize}

% CORRIGIDO: título encurtado para caber na caixa
\begin{boxexemplo}{Exemplo — Bolsa Família e frequência escolar}
Uma das formas mais rigorosas de avaliar o impacto do Bolsa Família sobre a frequência escolar usa a \textbf{regressão descontínua} na renda familiar per capita: famílias imediatamente abaixo do limite de elegibilidade recebem o benefício; famílias imediatamente acima não recebem. Comparando a frequência escolar de crianças dessas duas famílias — assumindo que a única diferença relevante entre elas é o recebimento do benefício — é possível estimar o efeito causal da condicionalidade. Essa estratégia foi implementada usando os dados de condicionalidade de educação do MDS linkados com o Censo Escolar pelo código do aluno.

\textit{Fonte: FGV CLEAR, elaboração própria.}
\end{boxexemplo}

\subsection{Aplicações típicas em avaliação}

\begin{itemize}
  \item \textbf{Avaliação de impacto sobre pobreza e desigualdade}: o Bolsa Família é um dos programas mais avaliados do mundo; suas bases de dados permitem estudos com diferentes estratégias de identificação (RDD na renda, DID na expansão municipal, IV na elegibilidade);
  \item \textbf{Efeitos sobre educação}: frequência escolar, progressão de série e conclusão do ensino fundamental e médio entre beneficiários;
  \item \textbf{Efeitos sobre saúde}: mortalidade infantil, cobertura vacinal, pré-natal e estado nutricional;
  \item \textbf{Efeitos sobre mercado de trabalho}: impacto da transferência sobre oferta de trabalho dos adultos da família — um dos debates mais ricos na literatura sobre o programa;
  \item \textbf{Análise de cobertura e focalização}: proporção das famílias elegíveis que estão sendo atendidas e perfil das famílias excluídas.
\end{itemize}

% ────────────────────────────────────────────────────────────
\section{BPC — Benefício de Prestação Continuada}
% ────────────────────────────────────────────────────────────

\ficharapida
  {BPC — Benefício de Prestação Continuada}
  {Ministério do Desenvolvimento e Assistência Social (MDS) / INSS}
  {https://www.gov.br/mds/pt-br/acoes-e-programas/assistencia-social/beneficios-assistenciais/beneficio-de-prestacao-continuada-bpc}
  {Brasil, Unidades da Federação e municípios}
  {Contínuo (desde 1996)}
  {Beneficiário individual}

\subsection{O que é e quem cobre}

O Benefício de Prestação Continuada (BPC) é o principal benefício de proteção social não contributiva do Brasil. Garantido pela Constituição Federal de 1988 e regulamentado pela Lei Orgânica da Assistência Social (LOAS), o BPC assegura um salário mínimo mensal a dois grupos populacionais:

\begin{itemize}
  \item \textbf{Idosos}: pessoas com 65 anos ou mais cuja renda familiar per capita seja inferior a um quarto do salário mínimo;
  \item \textbf{Pessoas com deficiência (PcD)}: pessoas com deficiência de qualquer idade cuja renda familiar per capita também seja inferior a um quarto do salário mínimo e cuja deficiência impeça a participação plena e efetiva na sociedade em igualdade de condições com as demais pessoas.
\end{itemize}

Diferentemente do Bolsa Família, o BPC não tem condicionalidades e é operado pelo INSS — embora seja um benefício assistencial, não previdenciário.

\subsection{Dados disponíveis e acesso}

Os dados do BPC estão disponíveis em múltiplas fontes:

\begin{itemize}
  \item \textbf{Portal da Transparência}: pagamentos mensais por beneficiário (NIS, município, valor), com histórico desde 2013;
  \item \textbf{Anuário Estatístico da Previdência Social (AEPS)}: publicação anual do MPS/MDS com estatísticas agregadas do BPC por UF, município, faixa etária e tipo (idoso ou PcD);
  \item \textbf{SUIBE — Sistema Único de Informações de Benefícios}: sistema administrativo do INSS com dados detalhados sobre cada benefício concedido, mantido, suspenso ou cancelado. O acesso aos microdados requer convênio com o INSS.
\end{itemize}

\subsection{Aplicações típicas em avaliação}

\begin{itemize}
  \item \textbf{Avaliação de impacto sobre pobreza de idosos}: o BPC é o principal mecanismo de proteção de idosos de baixa renda fora do sistema previdenciário; análises de regressão descontínua no critério de renda são a estratégia mais comum;
  \item \textbf{Impacto sobre saúde dos beneficiários}: linkagem com dados do DATASUS (SIH, SIM) para analisar efeitos sobre hospitalização e mortalidade de idosos e PcD;
  \item \textbf{Análise de cobertura}: proporção de idosos e PcD elegíveis que estão recebendo o benefício por município;
  \item \textbf{Efeitos de transbordamento}: impacto do BPC sobre outros membros do domicílio do beneficiário (como educação de crianças que moram com um idoso beneficiário).
\end{itemize}

% ────────────────────────────────────────────────────────────
\section{Censo SUAS}
% ────────────────────────────────────────────────────────────

\ficharapida
  {Censo SUAS}
  {Ministério do Desenvolvimento e Assistência Social (MDS)}
  {https://www.gov.br/mds/pt-br/acoes-e-programas/assistencia-social/monitoramento-e-avaliacao/censo-suas}
  {Brasil, Unidades da Federação e municípios}
  {Anual (desde 2007)}
  {Equipamento socioassistencial (CRAS, CREAS, Centro POP, etc.)}

\subsection{O que é e estrutura}

O Censo SUAS é o levantamento anual das condições de funcionamento dos equipamentos da rede socioassistencial brasileira. Realizado pelo MDS desde 2007, o Censo SUAS coleta informações de todos os equipamentos cadastrados no sistema, organizados por tipo:

\begin{itemize}
  \item \textbf{CRAS} — Centro de Referência de Assistência Social: unidade pública estatal de proteção social básica, responsável pela oferta do PAIF (Serviço de Proteção e Atendimento Integral à Família) e pelo encaminhamento ao CadÚnico;
  \item \textbf{CREAS} — Centro de Referência Especializado de Assistência Social: oferta serviços de proteção social especial para famílias e indivíduos em situação de risco pessoal e social;
  \item \textbf{Centro POP}: atendimento especializado para pessoas em situação de rua;
  \item \textbf{Gestão municipal e estadual}: informações sobre a estrutura administrativa das secretarias de assistência social.
\end{itemize}

\subsection{Principais variáveis coletadas}

Para cada equipamento, o Censo SUAS coleta:

\begin{itemize}
  \item \textbf{Infraestrutura}: tipo de imóvel (próprio, cedido, alugado), condições físicas, acessibilidade para PcD, equipamentos disponíveis (computadores, internet, veículo);
  \item \textbf{Recursos humanos}: número de profissionais por categoria (assistente social, psicólogo, advogado, auxiliar administrativo), vínculo empregatício (estatutário, CLT, comissionado, terceirizado), carga horária;
  \item \textbf{Serviços ofertados}: tipos de serviços, número de famílias e indivíduos acompanhados, grupos e atividades realizadas no mês;
  \item \textbf{Articulação intersetorial}: parcerias com saúde, educação e outros setores.
\end{itemize}

\subsection{Aplicações típicas em avaliação}

\begin{itemize}
  \item \textbf{Avaliação da expansão da rede SUAS}: monitoramento do crescimento do número de CRAS e CREAS por município ao longo do tempo e seu impacto sobre indicadores de vulnerabilidade social;
  \item \textbf{Análise de capacidade instalada}: identificação de municípios com cobertura insuficiente da rede de proteção básica em relação ao número de famílias vulneráveis cadastradas no CadÚnico;
  \item \textbf{Variável de controle em avaliações de impacto}: a presença e a qualidade dos equipamentos do SUAS são variáveis de contexto importantes em avaliações de políticas sociais mais amplas;
  \item \textbf{Avaliação de gestão}: análise da qualidade da gestão municipal do SUAS a partir de indicadores compostos como o Índice de Gestão Descentralizada (IGD-SUAS).
\end{itemize}

% ────────────────────────────────────────────────────────────
\section{RMA — Registro Mensal de Atendimentos}
% ────────────────────────────────────────────────────────────

\ficharapida
  {RMA — Registro Mensal de Atendimentos}
  {Ministério do Desenvolvimento e Assistência Social (MDS)}
  {https://www.gov.br/mds/pt-br/acoes-e-programas/assistencia-social/monitoramento-e-avaliacao/rede-suas}
  {Brasil, Unidades da Federação e municípios}
  {Mensal (desde 2012)}
  {Atendimento realizado por CRAS ou CREAS}

\subsection{O que é e aplicações}

O Registro Mensal de Atendimentos (RMA) é o sistema de monitoramento da produção de serviços socioassistenciais realizados pelos CRAS e CREAS. Diferentemente do Censo SUAS — que captura uma fotografia anual das condições estruturais dos equipamentos —, o RMA registra mensalmente o volume e o tipo de atendimentos realizados.

As principais informações disponíveis incluem:
\begin{itemize}
  \item Número de famílias atendidas no mês pelo PAIF e pelo PAEFI;
  \item Número de visitas domiciliares realizadas;
  \item Grupos e atividades coletivas oferecidas;
  \item Atendimentos individualizados por tipo de demanda.
\end{itemize}

O RMA é especialmente útil para monitorar a \textbf{regularidade e a intensidade} dos serviços prestados, complementando o Censo SUAS que mede a estrutura disponível. Para avaliações de fidelidade à implementação — verificar se os serviços estão sendo entregues conforme o planejado — o RMA é uma fonte fundamental.

% ────────────────────────────────────────────────────────────
\section{Outros programas: dados específicos}
% ────────────────────────────────────────────────────────────

\subsection{Minha Casa Minha Vida (MCMV)}

O programa Minha Casa Minha Vida — e sua versão anterior, o Programa de Aceleração do Crescimento habitacional — gera dados administrativos sobre contratos, beneficiários e unidades habitacionais entregues. Os dados estão disponíveis parcialmente no Portal da Transparência e, de forma mais detalhada, via convênio com a Caixa Econômica Federal e o Ministério das Cidades.

Para avaliações de impacto do MCMV, as principais fontes de dados são:
\begin{itemize}
  \item \textbf{Dados de contratos}: disponíveis no Portal da Transparência, com localização, faixa de renda e data de assinatura do contrato por beneficiário (NIS);
  \item \textbf{CadÚnico}: para caracterização socioeconômica dos beneficiários da Faixa 1 (população de menor renda);
  \item \textbf{RAIS e CAGED}: para análises de impacto sobre mercado de trabalho local (efeitos da construção sobre o emprego formal).
\end{itemize}

\subsection{Tarifa Social de Energia Elétrica (TSEE)}

A Tarifa Social de Energia Elétrica concede descontos na conta de luz para famílias de baixa renda cadastradas no CadÚnico. Os dados de beneficiários estão disponíveis na ANEEL e podem ser linkados com o CadÚnico pelo NIS para análises de cobertura e focalização.

\subsection{Programas estaduais e municipais}

Além dos programas federais, estados e municípios operam programas próprios de assistência social e transferência de renda. A disponibilidade de dados varia amplamente: alguns estados têm sistemas de informação bem estruturados e dados abertos; outros têm dados dispersos ou indisponíveis. Para avaliações de programas subnacionais, recomenda-se contato direto com as secretarias responsáveis para verificar a existência e o formato dos dados administrativos disponíveis.

% ────────────────────────────────────────────────────────────
\section{Combinando fontes para avaliação de políticas sociais}
% ────────────────────────────────────────────────────────────

A avaliação de políticas de assistência social e transferência de renda se beneficia enormemente da \textbf{linkagem entre bases}. O CadÚnico, por usar o NIS como identificador único de cada pessoa, serve como elo central que conecta:

\begin{itemize}
  \item \textbf{CadÚnico $\leftrightarrow$ Bolsa Família}: identifica quais famílias elegíveis estão recebendo o benefício e quais não estão — essencial para análises de cobertura e para construir grupos de controle em avaliações de impacto;
  \item \textbf{CadÚnico $\leftrightarrow$ Censo Escolar}: permite analisar o desempenho escolar de crianças de famílias beneficiárias em relação a crianças de famílias com perfil similar não beneficiárias;
  \item \textbf{CadÚnico $\leftrightarrow$ DATASUS}: linkagem por CPF ou NIS permite analisar desfechos de saúde (mortalidade, internações, vacinação) em função do recebimento de benefícios;
  \item \textbf{CadÚnico $\leftrightarrow$ RAIS/CAGED}: permite rastrear a inserção no mercado formal de trabalho de beneficiários ao longo do tempo.
\end{itemize}

% CORRIGIDO: título encurtado para caber na caixa
\begin{boxdica}{Dica — Plataforma de Dados Sociais do MDS}
O MDS disponibiliza, para pesquisadores com convênio, uma plataforma de análise de dados que permite trabalhar com os microdados do CadÚnico e do Bolsa Família em ambiente seguro, sem necessidade de transferir os arquivos. Isso facilita enormemente as análises que envolvem dados sensíveis de beneficiários. Informações sobre como solicitar acesso estão disponíveis no portal do MDS.
\end{boxdica}

% ────────────────────────────────────────────────────────────
\section{Síntese do capítulo}
% ────────────────────────────────────────────────────────────

Este capítulo apresentou o ecossistema de dados da assistência social e transferência de renda no Brasil. Os pontos centrais são:

\begin{itemize}
  \item O \textbf{CadÚnico} é a espinha dorsal do sistema de dados de assistência social: serve como porta de entrada para os principais programas e como base para linkagem com outras fontes;
  \item Os dados do \textbf{Bolsa Família} — disponíveis no Portal da Transparência, no CECAD e no VISDATA — permitem análises de cobertura, focalização e impacto com diferentes estratégias de identificação causal;
  \item O \textbf{BPC} cobre idosos e pessoas com deficiência de baixa renda e seus dados, disponíveis via Portal da Transparência e SUIBE, permitem avaliações de impacto com regressão descontínua na renda;
  \item O \textbf{Censo SUAS} e o \textbf{RMA} monitoram a estrutura e a produção dos equipamentos da rede socioassistencial, sendo fundamentais para avaliações de implementação e fidelidade;
  \item A \textbf{linkagem entre bases} — especialmente CadÚnico com Censo Escolar, DATASUS e RAIS — é o que permite as avaliações mais robustas de políticas sociais.
\end{itemize}

O CadÚnico mapeia a vulnerabilidade socioeconômica com uma precisão geográfica notável — e não surpreende que os mapas de cobertura dos programas sociais se sobreponham fortemente aos mapas de violência. Pobreza, exclusão e crime não são a mesma coisa, mas compartilham territórios. O próximo capítulo aborda as fontes de dados sobre segurança pública: o Atlas da Violência, o Anuário FBSP, os registros estaduais e as pesquisas de vitimização, ferramentas essenciais para quem quer entender — e avaliar intervenções sobre — a violência no Brasil.
% ============================================================
%  cap08_seguranca.tex
%  Capítulo 8: Segurança Pública
% ============================================================

% ════════════════════════════════════════════════════════════
\chapter{Segurança Pública}
% ════════════════════════════════════════════════════════════

\begin{boxdestaque}{Neste capítulo}
Este capítulo apresenta as principais fontes de dados sobre segurança pública disponíveis para o Brasil. Cobrimos o Atlas da Violência, o Anuário Brasileiro de Segurança Pública, os registros de ocorrências policiais estaduais, os dados do sistema penitenciário (DEPEN/MJSP) e os suplementos de vitimização da PNAD. Para cada fonte, detalhamos histórico, estrutura, variáveis disponíveis, forma de acesso, aplicações típicas em avaliação e — especialmente importante nessa área — as limitações e os vieses de cobertura.
\end{boxdestaque}

% ────────────────────────────────────────────────────────────
\section{Introdução: os desafios dos dados de segurança pública}
% ────────────────────────────────────────────────────────────

A segurança pública é uma das áreas em que a qualidade e a comparabilidade dos dados entre unidades federativas apresentam os maiores desafios do sistema estatístico brasileiro. Ao contrário da saúde — onde o DATASUS oferece sistemas nacionais padronizados — ou da educação — onde o Inep coordena um sistema unificado —, os dados de segurança pública no Brasil são produzidos principalmente pelas Secretarias de Segurança Pública estaduais, com metodologias, definições e sistemas de informação que variam significativamente entre estados.

Essa fragmentação tem consequências diretas para o avaliador: comparações entre estados devem ser feitas com cautela redobrada, séries históricas podem ser interrompidas por mudanças metodológicas estaduais, e a ausência de um identificador único nacional para ocorrências criminais impede a construção de painéis longitudinais comparáveis.

Há, no entanto, exceções importantes. A mortalidade por causas externas registrada no SIM — especialmente os homicídios — oferece uma série histórica nacional relativamente confiável, porque a morte é um evento objetivo e o registro de óbito é obrigatório. É sobre essa base que o Atlas da Violência é construído, tornando-o a referência mais sólida para análises nacionais de violência letal.

\begin{boxatencao}{Atenção — O problema da cifra oculta}
Em segurança pública, a diferença entre o crime que ocorre e o crime que é registrado é estrutural e se chama \textbf{cifra oculta} (ou \textit{dark figure of crime}). Estimativas internacionais sugerem que apenas uma fração dos crimes é reportada à polícia --- proporção que varia muito por tipo de crime (homicídios têm taxa de registro mais alta; crimes patrimoniais menores e violência doméstica têm taxas muito mais baixas). No Brasil, pesquisas de vitimização sugerem que menos de 40\% dos crimes são registrados em boletim de ocorrência. Qualquer análise baseada em registros policiais está, portanto, analisando o crime \textit{registrado} --- não o crime \textit{ocorrido}. Essa distinção é crítica para a interpretação dos resultados.
\end{boxatencao}

% ────────────────────────────────────────────────────────────
\section{Atlas da Violência}
% ────────────────────────────────────────────────────────────

\ficharapida
  {Atlas da Violência}
  {Instituto de Pesquisa Econômica Aplicada (Ipea) e Fórum Brasileiro de Segurança Pública (FBSP)}
  {https://www.ipea.gov.br/atlasviolencia}
  {Brasil, Unidades da Federação e municípios}
  {Anual (desde 2016; dados retroativos a 1980)}
  {Município (agregação de dados individuais do SIM)}

\subsection{O que é e base metodológica}

O Atlas da Violência é uma publicação anual produzida pelo Ipea em parceria com o Fórum Brasileiro de Segurança Pública (FBSP) desde 2016. Seu diferencial em relação a outras fontes de dados de segurança pública é a base metodológica: os indicadores de violência letal do Atlas são construídos a partir dos dados do \textbf{Sistema de Informações sobre Mortalidade (SIM)} do Ministério da Saúde --- não a partir dos registros policiais estaduais.

Essa escolha metodológica tem implicações importantes. O SIM registra \textit{todos} os óbitos, independentemente de terem sido comunicados à polícia. Para homicídios --- classificados pela CID-10 como ``agressões'' (X85 a Y09) ---, a concordância entre os registros do SIM e os boletins de ocorrência é razoavelmente alta, porque um corpo encontrado gera automaticamente um registro de óbito. Para outros tipos de violência (lesões corporais, ameaças, roubos), o SIM não é útil --- mas para violência letal, é a fonte mais confiável disponível.

\subsection{Principais indicadores disponíveis}

O Atlas da Violência disponibiliza, por município e ano:

\begin{itemize}
  \item \textbf{Taxa de homicídios}: número de óbitos por agressão por 100 mil habitantes, desagregada por sexo, raça/cor e faixa etária;
  \item \textbf{Taxa de homicídios de jovens negros}: indicador central para análises de desigualdade racial na violência;
  \item \textbf{Taxa de feminicídios}: óbitos femininos por agressão, com tendência histórica;
  \item \textbf{Mortes no trânsito}: óbitos por acidentes de transporte terrestre;
  \item \textbf{Mortes por intervenção policial}: registradas no SIM como ``intervenção legal''.
\end{itemize}

Os dados cobrem o período de 1980 até os últimos dois anos disponíveis no SIM (há defasagem de aproximadamente dois anos na publicação dos dados de mortalidade).

\subsection{Acesso aos dados}

O Atlas da Violência disponibiliza seus dados em três formatos:
\begin{itemize}
  \item \textbf{Plataforma interativa}: painéis de visualização com mapas e séries históricas em \url{www.ipea.gov.br/atlasviolencia};
  \item \textbf{Planilhas para download}: tabelas com todos os indicadores por município e ano, em formato Excel e CSV;
  \item \textbf{Publicação anual}: relatório com análises temáticas aprofundadas, metodologia e contextualização dos dados.
\end{itemize}

Para análises com microdados de mortalidade --- necessárias para estudos mais detalhados por características individuais das vítimas ---, o pesquisador deve acessar diretamente os microdados do SIM no DATASUS (ver Capítulo~5).

\subsection{Aplicações típicas em avaliação}

\begin{itemize}
  \item \textbf{Avaliação de políticas de segurança pública estaduais}: monitoramento da taxa de homicídios antes e depois da implementação de programas como o Pacto pela Vida (Pernambuco) ou o Programa Estado Presente (Espírito Santo);
  \item \textbf{Análises de desigualdade racial na violência}: a desagregação por raça/cor permite avaliar se políticas de segurança afetam de forma diferente a mortalidade de negros e brancos;
  \item \textbf{Avaliação de programas de prevenção da violência juvenil}: foco na taxa de homicídios de jovens (15 a 29 anos) por município;
  \item \textbf{Estudos de impacto de intervenções urbanísticas}: relação entre urbanização, iluminação pública, equipamentos culturais e esportivos e taxas de violência letal.
\end{itemize}

\begin{boxexemplo}{Exemplo --- Pacto pela Vida e homicídios em Pernambuco}
O Pacto pela Vida, lançado em Pernambuco em 2007, é frequentemente citado como um dos programas estaduais de segurança pública mais bem avaliados do Brasil. Pesquisadores utilizaram os dados do SIM para construir uma série histórica de taxas de homicídio nos municípios pernambucanos e em municípios de outros estados com tendências pré-tratamento similares (construindo um controle sintético). Os resultados sugeriram uma redução expressiva nos homicídios em Pernambuco após 2007, não observada nos estados de comparação. A base metodológica no SIM --- e não nos registros policiais --- foi fundamental para a credibilidade do estudo, pois eliminou a possibilidade de que a redução observada fosse apenas um artefato de mudanças na forma de registro das ocorrências pela polícia.

\textit{Fonte: FGV CLEAR, elaboração própria.}
\end{boxexemplo}

% ────────────────────────────────────────────────────────────
\section{Anuário Brasileiro de Segurança Pública}
% ────────────────────────────────────────────────────────────

\ficharapida
  {Anuário Brasileiro de Segurança Pública}
  {Fórum Brasileiro de Segurança Pública (FBSP)}
  {https://forumseguranca.org.br/anuario-brasileiro-seguranca-publica}
  {Brasil e Unidades da Federação}
  {Anual (desde 2007)}
  {Unidade da Federação}

\subsection{O que é e metodologia de coleta}

O Anuário Brasileiro de Segurança Pública é a principal publicação nacional que consolida dados de segurança pública provenientes das Secretarias Estaduais de Segurança Pública. Produzido anualmente pelo Fórum Brasileiro de Segurança Pública (FBSP) desde 2007, o Anuário reúne informações que não estão disponíveis em nenhuma base federal centralizada --- justamente porque, no Brasil, as polícias são estaduais e cada estado opera seus próprios sistemas de registro.

A metodologia de coleta do Anuário combina dois mecanismos:
\begin{itemize}
  \item \textbf{Solicitações via Lei de Acesso à Informação (LAI)}: o FBSP encaminha anualmente pedidos de informação padronizados a todas as Secretarias Estaduais de Segurança Pública, solicitando dados sobre ocorrências registradas por tipo de crime;
  \item \textbf{Convênios e acordos de compartilhamento}: em alguns estados, acordos específicos facilitam o acesso a dados mais detalhados.
\end{itemize}

\subsection{Indicadores disponíveis}

O Anuário publica anualmente, para cada Unidade da Federação, 24 indicadores de segurança pública, incluindo:

\begin{itemize}
  \item Homicídios dolosos (registros policiais);
  \item Latrocínios (roubo seguido de morte);
  \item Lesões corporais dolosas;
  \item Estupros e estupros de vulnerável;
  \item Roubos (total, de veículos, a estabelecimentos comerciais, a transeuntes);
  \item Furtos;
  \item Mortes decorrentes de intervenção policial (civis mortos por policiais em serviço e fora de serviço);
  \item Policiais mortos em serviço e fora de serviço;
  \item Número de presos provisórios e definitivos.
\end{itemize}

\subsection{Limitações críticas}

O Anuário é uma fonte valiosa, mas suas limitações são substanciais e devem ser explicitadas em qualquer análise:

\begin{itemize}
  \item \textbf{Heterogeneidade metodológica entre estados}: cada estado define e registra os crimes de forma diferente. O que é classificado como ``homicídio doloso'' em São Paulo pode ser registrado de forma diferente no Pará. Comparações diretas entre estados exigem cautela;
  \item \textbf{Cobertura incompleta}: nem todos os estados respondem a todos os pedidos LAI com a mesma completude. Em alguns anos, certos estados não fornecem dados para determinados indicadores;
  \item \textbf{Cifra oculta}: como discutido na introdução do capítulo, os registros policiais capturam apenas os crimes reportados --- fração que varia por tipo de crime e por estado;
  \item \textbf{Mudanças nos critérios de registro}: estados podem alterar seus critérios de classificação ao longo do tempo, introduzindo quebras nas séries históricas que não refletem mudanças reais na criminalidade.
\end{itemize}

\begin{boxatencao}{Atenção --- Comparações estaduais: verificar definições}
Antes de comparar indicadores do Anuário entre estados diferentes, verifique se as definições operacionais são compatíveis. O FBSP publica notas metodológicas sobre as diferenças entre estados em cada edição do Anuário. Para alguns indicadores --- especialmente homicídios ---, a comparação com os dados do SIM/Atlas da Violência é uma forma de verificar a consistência dos registros policiais estaduais.
\end{boxatencao}

% ────────────────────────────────────────────────────────────
\section{Registros de Ocorrências Policiais Estaduais}
% ────────────────────────────────────────────────────────────

\ficharapida
  {Registros de Ocorrências Policiais}
  {Secretarias Estaduais de Segurança Pública}
  {Varia por estado (ver portais estaduais de dados abertos)}
  {Município e delegacia (varia por estado)}
  {Variável por estado (mensal a anual)}
  {Ocorrência policial ou boletim de ocorrência}

\subsection{O que são e heterogeneidade}

Os registros de ocorrências policiais são os dados primários da segurança pública estadual: cada boletim de ocorrência (BO) registrado em uma delegacia é uma unidade de dado. Agregados por tipo de crime, município e período, esses registros são a principal fonte para análises de criminalidade abaixo do nível estadual --- especialmente para crimes que não resultam em morte e, portanto, não aparecem no SIM.

A disponibilidade e o formato desses dados variam enormemente entre estados:

\begin{itemize}
  \item \textbf{Estados com dados abertos bem estruturados}: São Paulo (SSP-SP), Minas Gerais (CINDS-MG) e Rio de Janeiro (ISP-RJ) disponibilizam microdados de ocorrências com desagregação por delegacia, mês e tipo de crime, em formatos abertos (CSV, API);
  \item \textbf{Estados com dados parcialmente disponíveis}: algumas UFs publicam tabelas agregadas em seus portais de transparência, mas sem microdados;
  \item \textbf{Estados com dados indisponíveis ou de difícil acesso}: em vários estados, os dados só são obtidos via LAI, com resposta variável em qualidade e tempestividade.
\end{itemize}

\subsection{Instituto de Segurança Pública do Rio de Janeiro (ISP-RJ)}

O ISP-RJ merece destaque como referência de boas práticas em transparência de dados de segurança pública. O instituto disponibiliza:

\begin{itemize}
  \item Microdados mensais de ocorrências por tipo de crime e circunscrição policial (AISP), desde 2003;
  \item Dados de letalidade policial com desagregação por raça/cor e circunstâncias;
  \item Dados de apreensão de armas e drogas;
  \item APIs de acesso programático para download automatizado.
\end{itemize}

Esses dados têm sido amplamente utilizados em pesquisas sobre o impacto das Unidades de Polícia Pacificadora (UPPs) e de outras políticas de segurança no Rio de Janeiro.

\subsection{Secretaria de Segurança Pública de São Paulo (SSP-SP)}

A SSP-SP disponibiliza dados mensais de ocorrências para todos os municípios do estado, com desagregação por delegacia, desde 2001. Os dados cobrem homicídios dolosos, tentativas de homicídio, latrocínios, estupros, roubos, furtos e tráfico de drogas, e são disponibilizados em planilhas Excel no portal da SSP-SP e via API no portal de dados abertos do estado.

\subsection{Aplicações típicas em avaliação}

\begin{itemize}
  \item \textbf{Avaliação de políticas de policiamento}: impacto de mudanças no patrulhamento, na distribuição de efetivo ou em operações específicas sobre taxas de crime;
  \item \textbf{Avaliação de intervenções urbanísticas}: relação entre iluminação pública, câmeras de segurança, revitalização de espaços públicos e criminalidade;
  \item \textbf{Análise de deslocamento de crime}: verificar se intervenções reduziram o crime na área tratada ou apenas deslocaram para áreas vizinhas;
  \item \textbf{Estudos de reincidência}: em estados com identificadores individuais nos registros, é possível construir históricos criminais para análises de reincidência.
\end{itemize}

% ────────────────────────────────────────────────────────────
\section{Sistema Penitenciário — DEPEN e MJSP}
% ────────────────────────────────────────────────────────────

\ficharapida
  {Dados do Sistema Penitenciário (DEPEN/Infopen)}
  {Departamento Penitenciário Nacional (DEPEN) / Ministério da Justiça e Segurança Pública (MJSP)}
  {https://www.gov.br/senappen/pt-br/servicos/sisdepen}
  {Brasil, Unidades da Federação e estabelecimentos penais}
  {Semestral (desde 2004; anual em edições anteriores)}
  {Estabelecimento penal e preso}

\subsection{O que é e estrutura}

Os dados do sistema penitenciário brasileiro são coletados pelo Departamento Penitenciário Nacional (DEPEN), vinculado ao Ministério da Justiça e Segurança Pública, por meio do \textbf{Levantamento Nacional de Informações Penitenciárias (Infopen)}. Publicado semestralmente desde 2004, o Infopen consolida informações de todos os estabelecimentos penais do país sobre a população carcerária e as condições de encarceramento.

Em 2021, o Infopen foi substituído pelo \textbf{SISDEPEN} (Sistema de Informações do Departamento Penitenciário Nacional), que opera em tempo quase real com dados inseridos pelos próprios estabelecimentos penais. O SISDEPEN permite consultas mais frequentes e desagregadas do que o antigo Infopen.

\subsection{Principais variáveis disponíveis}

Para cada estabelecimento penal:
\begin{itemize}
  \item Capacidade e lotação (número de vagas vs. número de presos);
  \item Regime de cumprimento (fechado, semiaberto, aberto);
  \item Tipo de estabelecimento (penitenciária, casa de detenção, cadeia pública, etc.);
  \item Natureza (federal, estadual);
  \item Existência de serviços de saúde, educação e trabalho dentro do estabelecimento.
\end{itemize}

Para a população carcerária agregada por estabelecimento:
\begin{itemize}
  \item Número de presos por sexo, faixa etária e raça/cor;
  \item Situação jurídica: presos provisórios (sem condenação definitiva) vs. condenados;
  \item Tipo penal: crimes contra a pessoa, contra o patrimônio, tráfico de drogas, etc.;
  \item Nível de escolaridade da população carcerária.
\end{itemize}

\subsection{Aplicações típicas em avaliação}

\begin{itemize}
  \item \textbf{Avaliação de políticas de desencarceramento}: monitoramento do impacto de medidas alternativas à prisão, como tornozeleiras eletrônicas, sobre o tamanho da população carcerária;
  \item \textbf{Análise de superlotação}: identificação de estabelecimentos com maior nível de superlotação e sua relação com reincidência e violência intramuros;
  \item \textbf{Políticas de educação e trabalho no sistema prisional}: cobertura de programas de educação e qualificação profissional dentro dos estabelecimentos;
  \item \textbf{Análise de desigualdades no encarceramento}: perfil racial e socioeconômico da população carcerária em comparação com a população geral.
\end{itemize}

\begin{boxatencao}{Atenção --- Sistema penitenciário: dados variáveis}
A qualidade dos dados do DEPEN/SISDEPEN varia entre estados, refletindo diferenças na capacidade administrativa dos sistemas penitenciários estaduais. Alguns estados têm registros mais completos e atualizados; outros apresentam inconsistências e lacunas. Além disso, o número de presos provisórios --- pessoas detidas sem condenação definitiva --- é frequentemente subestimado, pois nem sempre os registros são atualizados quando há mudança de situação jurídica. Análises com esses dados devem sempre considerar a possibilidade de subregistro diferencial entre estados.
\end{boxatencao}

% ────────────────────────────────────────────────────────────
\section{Pesquisas de Vitimização}
% ────────────────────────────────────────────────────────────

\ficharapida
  {Suplementos de Vitimização da PNAD}
  {Instituto Brasileiro de Geografia e Estatística (IBGE)}
  {https://www.ibge.gov.br/estatisticas/sociais/justica-e-seguranca/9127-pesquisa-nacional-por-amostra-de-domicilios.html}
  {Brasil, Grandes Regiões e Unidades da Federação}
  {Pontual: 1988 e 2009}
  {Domicílio e morador}

\subsection{O que são e por que importam}

Pesquisas de vitimização são levantamentos domiciliares que perguntam diretamente às pessoas se foram vítimas de crimes em determinado período --- independentemente de terem registrado boletim de ocorrência. São a principal ferramenta para estimar a \textbf{cifra oculta} e para compreender a experiência real da população com a criminalidade, em contraste com os registros policiais.

No Brasil, as duas principais pesquisas de vitimização de abrangência nacional foram os suplementos da PNAD de 1988 e 2009. O suplemento de 2009 --- ``Características da vitimização e do acesso à Justiça no Brasil'' --- é o mais detalhado e metodologicamente mais robusto, cobrindo:

\begin{itemize}
  \item Vitimização por tipo de crime (roubo, furto, agressão física, ameaça, estupro) nos últimos 12 meses;
  \item Características da ocorrência (local, horário, uso de arma, relação com o agressor);
  \item Taxa de registro policial e razões para não registrar;
  \item Acesso à Justiça: experiências com o sistema de justiça criminal (polícia, promotoria, tribunal).
\end{itemize}

\subsection{Outras pesquisas de vitimização}

Além dos suplementos da PNAD, há outras fontes de dados de vitimização relevantes:

\begin{itemize}
  \item \textbf{PNAD Contínua --- Módulo de Segurança Pública (2017 e 2021)}: o IBGE realizou módulos específicos de segurança pública na PNAD Contínua, com metodologia atualizada e maior cobertura temática;
  \item \textbf{Pesquisas estaduais}: alguns estados realizaram pesquisas de vitimização próprias --- São Paulo (2008--2010), Rio de Janeiro (várias edições pelo ISP-RJ), Minas Gerais. Essas pesquisas permitem análises mais detalhadas no nível municipal, mas não são comparáveis entre estados por diferenças metodológicas;
  \item \textbf{Pesquisas internacionais}: o \textit{International Crime Victims Survey} (ICVS) incluiu o Brasil em algumas de suas rodadas, permitindo comparações internacionais de vitimização.
\end{itemize}

\subsection{Aplicações típicas em avaliação}

\begin{itemize}
  \item \textbf{Estimação da cifra oculta}: comparação entre vitimização declarada e registros policiais para estimar a proporção de crimes não reportados por tipo e perfil da vítima;
  \item \textbf{Análise de subnotificação diferencial}: identificar se grupos específicos (mulheres, negros, moradores de favelas) reportam menos crimes à polícia e por quê;
  \item \textbf{Avaliação de confiança nas instituições}: as pesquisas de vitimização geralmente incluem perguntas sobre percepção de segurança e confiança na polícia, úteis para avaliar o impacto de reformas institucionais;
  \item \textbf{Análise de vitimização por violência doméstica}: dado que a subnotificação da violência doméstica é estruturalmente alta, as pesquisas de vitimização são frequentemente a única fonte para estimar sua prevalência real.
\end{itemize}

% ────────────────────────────────────────────────────────────
\section{Sistema de Vigilância de Violências e Acidentes (VIVA)}
% ────────────────────────────────────────────────────────────

\ficharapida
  {VIVA — Sistema de Vigilância de Violências e Acidentes}
  {Ministério da Saúde / DATASUS}
  {https://datasus.saude.gov.br/viva-desde-2006}
  {Brasil, Unidades da Federação e municípios}
  {Contínuo (desde 2006; integrado ao SINAN a partir de 2009)}
  {Atendimento em unidade de saúde por violência ou acidente}

\subsection{O que é e aplicações}

O VIVA registra atendimentos realizados em unidades de saúde sentinela por vítimas de violências e acidentes. Diferentemente do SIM --- que registra apenas óbitos --- o VIVA captura a violência não fatal que demanda atendimento de saúde: lesões corporais, violência sexual, violência doméstica, acidentes de trânsito e acidentes de trabalho atendidos em prontos-socorros e UPAs.

A partir de 2009, o VIVA foi integrado ao SINAN como parte das notificações compulsórias, ampliando sua cobertura para além das unidades sentinela. As principais aplicações em avaliação incluem:

\begin{itemize}
  \item \textbf{Violência doméstica e sexual}: o VIVA é uma das poucas fontes com dados sobre violência não fatal atendida em serviços de saúde, complementando as pesquisas de vitimização;
  \item \textbf{Avaliação de políticas de proteção à mulher}: monitoramento do impacto da Lei Maria da Penha e de outras políticas sobre atendimentos por violência contra a mulher;
  \item \textbf{Acidentes de trânsito}: comparação com os dados de mortalidade do SIM para avaliar políticas de segurança viária.
\end{itemize}

% ────────────────────────────────────────────────────────────
\section{Síntese do capítulo}
% ────────────────────────────────────────────────────────────

Este capítulo apresentou as principais fontes de dados sobre segurança pública disponíveis para o Brasil. A Tabela~\ref{tab:sintese_cap8} resume suas características essenciais.

\begin{table}[H]
\centering
\caption{Síntese das fontes de dados de segurança pública}
\label{tab:sintese_cap8}
\small
\begin{tabularx}{\textwidth}{@{} >{\raggedright\arraybackslash}p{2.8cm} >{\raggedright\arraybackslash}p{3.2cm} >{\raggedright\arraybackslash}p{2.8cm} X @{}}
\toprule
\textbf{Fonte} & \textbf{O que registra} & \textbf{Cobertura} &
\textbf{Principal uso em avaliação} \\
\midrule
Atlas da Violência  & Mortes violentas (SIM)       & Nacional (municipal) & Violência letal, tendências históricas \\
Anuário FBSP        & Ocorrências policiais        & Nacional (estadual)  & Criminalidade registrada, comparações \\
Registros estaduais & Boletins de ocorrência       & Estadual\slash \allowbreak municipal   & Análises locais, policiamento \\
DEPEN\slash \allowbreak SISDEPEN      & População carcerária         & Nacional             & Sistema prisional, encarceramento \\
Vitimização PNAD    & Crimes sofridos (declarado)  & Nacional (amostral)  & Cifra oculta, acesso à Justiça \\
VIVA/SINAN          & Atendimentos por violência   & Serviços de saúde    & Violência doméstica, sexual \\
\bottomrule
\end{tabularx}
\fonte{FGV CLEAR, elaboração própria.}
\end{table}

Os pontos centrais do capítulo são:

\begin{itemize}
  \item A segurança pública é a área com \textbf{maior heterogeneidade} de dados entre estados brasileiros: comparações diretas entre UFs exigem cautela metodológica redobrada;
  \item O \textbf{Atlas da Violência}, baseado no SIM, é a fonte mais confiável para análises de violência letal nacional porque contorna o problema da cifra oculta e da heterogeneidade dos registros policiais;
  \item O \textbf{Anuário FBSP} consolida dados de ocorrências policiais estaduais, mas as limitações de comparabilidade entre estados são estruturais;
  \item Os \textbf{registros estaduais} --- especialmente SP, RJ e MG --- são as fontes mais ricas para avaliações de políticas de segurança no nível local, mas sua disponibilidade e qualidade variam;
  \item As \textbf{pesquisas de vitimização} são indispensáveis para estimar a cifra oculta e para capturar dimensões da violência que os registros policiais sistematicamente subestimam;
  \item A \textbf{triangulação} entre fontes --- registros policiais, mortalidade (SIM) e vitimização --- é a abordagem mais robusta para qualquer análise séria de segurança pública.
\end{itemize}

O próximo capítulo apresenta as fontes que sustentam análises espaciais e territoriais: malhas territoriais do IBGE, dados de saneamento do SNIS, indicadores de agências reguladoras (ANEEL, ANATEL, ANTT) e dados ambientais do INPE (PRODES, DETER) e do MapBiomas.
% ============================================================
%  cap09_infraestrutura_territorio.tex  [REVISADO]
%  Capítulo 9: Infraestrutura, Território e Meio Ambiente
% ============================================================

% ════════════════════════════════════════════════════════════
\chapter{Infraestrutura, Território e Meio Ambiente}
% ════════════════════════════════════════════════════════════

\begin{boxdestaque}{Neste capítulo}
Este capítulo apresenta as principais fontes de dados sobre infraestrutura, território e meio ambiente disponíveis para o Brasil. Cobrimos o SNIS (saneamento), as pesquisas de capacidade municipal do IBGE (Munic e Estadic), os dados das agências reguladoras (ANEEL, ANATEL, ANS), as malhas territoriais e setores censitários do IBGE, as fontes de dados ambientais como PRODES, DETER e MapBiomas, e --- com destaque especial --- as bases municipais sobre habitação e mercado imobiliário, incluindo ITBI, cadastro imobiliário, IPTU, Planta Genérica de Valores e registros de licenciamento urbanístico. Para cada fonte, detalhamos estrutura, cobertura, acesso e aplicações em avaliação de políticas públicas.
\end{boxdestaque}

% ────────────────────────────────────────────────────────────
\section{Introdução: por que infraestrutura e território importam
         para avaliação}
% ────────────────────────────────────────────────────────────

Infraestrutura e território são dimensões transversais em praticamente toda avaliação de política pública. A cobertura de saneamento básico afeta diretamente a saúde da população; o acesso à energia elétrica condiciona o aprendizado escolar e a produtividade econômica; a conectividade de telecomunicações determina o acesso a serviços digitais de governo; a distância entre populações e serviços públicos é um determinante crítico do uso efetivo desses serviços.

Do ponto de vista metodológico, dados de infraestrutura e território cumprem dois papéis distintos na avaliação:

\begin{itemize}
  \item \textbf{Como variáveis de resultado}: quando a política avaliada tem por objetivo expandir a cobertura de um serviço de infraestrutura (saneamento, energia, internet, habitação);
  \item \textbf{Como variáveis de controle ou de contexto}: quando a infraestrutura disponível é um determinante dos resultados de outras políticas (saúde, educação, renda) e precisa ser controlada para isolar o efeito da intervenção de interesse.
\end{itemize}

Além disso, dados georreferenciados --- malhas territoriais, setores censitários, coordenadas de estabelecimentos --- abrem uma dimensão analítica específica: a \textbf{análise espacial}, que permite identificar padrões territoriais de cobertura e exclusão, calcular distâncias entre populações e serviços, e explorar variações geográficas como fonte de identificação causal.

Uma observação importante antes de prosseguir: várias das fontes discutidas neste capítulo só revelam plenamente seu potencial quando integradas a malhas territoriais, setores censitários, bairros, distritos, CEPs ou outros identificadores geográficos. Essa integração espacial é especialmente decisiva em habitação, mercado imobiliário, infraestrutura urbana, saúde, crime e educação. A seção~\ref{sec:malhas} sobre malhas territoriais do IBGE deve ser lida em conjunto com as demais seções do capítulo.

% ────────────────────────────────────────────────────────────
\section{SNIS --- Sistema Nacional de Informações sobre Saneamento}
% ────────────────────────────────────────────────────────────

\ficharapida
  {SNIS --- Sistema Nacional de Informações sobre Saneamento}
  {Ministério das Cidades / Secretaria Nacional de Saneamento}
  {https://www.snis.gov.br}
  {Brasil, Unidades da Federação e municípios}
  {Anual (desde 1995 para água e esgoto; desde 2002 para resíduos sólidos; desde 2015 para drenagem)}
  {Município e prestador de serviço}

\subsection{O que é e por que importa}

O Sistema Nacional de Informações sobre Saneamento (SNIS) é a principal fonte de dados sobre os serviços de saneamento básico no Brasil --- abastecimento de água, esgotamento sanitário, manejo de resíduos sólidos urbanos e drenagem e manejo de águas pluviais. Operado pelo Ministério das Cidades, o SNIS coleta anualmente informações declaradas pelos prestadores de serviços de saneamento sobre cobertura, qualidade, tarifas e finanças.

\subsection{Estrutura e principais variáveis}

\subsubsection{Água e Esgoto}

\begin{itemize}
  \item \textbf{Cobertura}: população atendida com abastecimento de água e com coleta de esgoto; índice de atendimento urbano e total;
  \item \textbf{Qualidade}: índice de conformidade das amostras de água;
  \item \textbf{Operação}: volume de água produzido, consumido e faturado; índice de perdas na distribuição;
  \item \textbf{Finanças}: receita operacional, despesas totais, resultado operacional por prestador;
  \item \textbf{Investimentos}: valor investido em expansão e reposição da rede.
\end{itemize}

\subsubsection{Resíduos Sólidos}

\begin{itemize}
  \item Cobertura da coleta domiciliar urbana;
  \item Tipo de destinação final (aterro sanitário, aterro controlado, lixão);
  \item Existência de coleta seletiva e de triagem;
  \item Quantidade de resíduos coletados e destinados.
\end{itemize}

\subsection{Limitações e cuidados}

A principal limitação do SNIS é a \textbf{cobertura heterogênea entre municípios}: a adesão ao sistema é voluntária para os prestadores, e municípios pequenos têm taxas de resposta mais baixas e qualidade de preenchimento mais variável. Além disso, os dados são \textbf{autodeclarados pelos prestadores}, o que pode gerar superestimação dos indicadores de cobertura.

\begin{boxdica}{Dica --- Combinando SNIS com Censo Demográfico}
O SNIS reporta cobertura a partir das declarações dos prestadores; o Censo Demográfico reporta cobertura a partir de perguntas feitas diretamente aos moradores. As duas fontes frequentemente divergem. Para análises robustas, compare as duas fontes e documente as divergências encontradas.
\end{boxdica}

\subsection{Acesso aos dados}

\begin{itemize}
  \item \textbf{Portal SNIS}: consultas online com tabelas e gráficos por município em \url{www.snis.gov.br};
  \item \textbf{Série Histórica}: planilhas Excel com todos os indicadores por município desde o início da série;
  \item \textbf{API}: o SNIS disponibiliza uma API REST para acesso programático.
\end{itemize}

\subsection{Aplicações típicas em avaliação}

\begin{itemize}
  \item Avaliação de programas de universalização do saneamento;
  \item Impacto do saneamento sobre mortalidade infantil e sobre indicadores sensíveis a saneamento --- diarreias, outras doenças de veiculação hídrica e hospitalizações por gastroenterites, em geral mensurados a partir de SIM/SIH;
  \item Análise de eficiência dos prestadores;
  \item Avaliação do Marco Legal do Saneamento (Lei n.º 14.026/2020).
\end{itemize}

% ────────────────────────────────────────────────────────────
\section{Munic e Estadic --- Pesquisas de Capacidade Institucional
         do IBGE}
% ────────────────────────────────────────────────────────────

\ficharapida
  {Munic --- Pesquisa de Informações Básicas Municipais}
  {Instituto Brasileiro de Geografia e Estatística (IBGE)}
  {https://www.ibge.gov.br/estatisticas/sociais/educacao/10586-pesquisa-de-informacoes-basicas-municipais.html}
  {Brasil e municípios}
  {Irregular (edições principais: 1999--2020)}
  {Município}

\subsection{O que é e por que importa}

A Pesquisa de Informações Básicas Municipais (Munic) é um censo de todos os municípios brasileiros, realizado pelo IBGE com periodicidade irregular desde 1999. O questionário é respondido pelos próprios gestores municipais e cobre a estrutura institucional, a capacidade administrativa e a oferta de serviços públicos nos municípios.

Para a avaliação de políticas públicas, a Munic ocupa um lugar especial: é uma das poucas fontes com informações sistemáticas sobre \textbf{capacidade estatal local} --- dimensão raramente disponível em outras fontes e frequentemente relevante para entender por que políticas funcionam melhor em alguns municípios do que em outros. Sua importância vai além das análises de infraestrutura: a Munic é uma referência para qualquer avaliação que envolva a implementação de políticas pelo nível municipal.

\subsection{Principais temas cobertos}

\begin{itemize}
  \item \textbf{Estrutura administrativa}: existência de secretarias por área (saúde, educação, assistência social, meio ambiente, habitação), número de servidores, plano plurianual;
  \item \textbf{Legislação e planejamento}: plano diretor, lei de uso e ocupação do solo, código de obras, planos setoriais municipais;
  \item \textbf{Serviços públicos}: equipamentos culturais, esportivos, de saúde e assistência social mantidos pelo município;
  \item \textbf{Tecnologia da informação}: uso de sistemas de informação, presença digital, acesso à internet nas secretarias;
  \item \textbf{Meio ambiente}: conselho municipal de meio ambiente, legislação ambiental, coleta seletiva;
  \item \textbf{Participação social}: conselhos municipais por área e regularidade das reuniões;
  \item \textbf{Habitação e urbanismo}: existência de secretaria de habitação, de plano local de habitação de interesse social (PLHIS), de legislação de regularização fundiária, de sistemas de informação habitacional.
\end{itemize}

A \textbf{Estadic} (Pesquisa de Informações Básicas Estaduais), disponível em \url{www.ibge.gov.br/estatisticas/sociais/saude/16770-pesquisa-de-informacoes-basicas-estaduais.html}, é o equivalente da Munic para os estados e o Distrito Federal.

\begin{boxdica}{Dica --- Suplementos temáticos da Munic e da Estadic}
Além do núcleo fixo de perguntas, a Munic e a Estadic incluem, em edições específicas, \textbf{suplementos temáticos} que ampliam consideravelmente o escopo analítico. Entre os mais usados em avaliações: Cultura (2006, 2009, 2014), Esporte (2003, 2015), Juventude (2009, 2013), Assistência Social (2005, 2009, 2013, 2018), Habitação (2001, 2008, 2017), Saúde (2001, 2005, 2009, 2014), Recursos Humanos (diversas edições) e Meio Ambiente (2002, 2017). Para uma análise que dependa de um suplemento específico, é imprescindível verificar no portal do IBGE o ano exato em que cada tema foi investigado antes de desenhar a estratégia empírica.
\end{boxdica}

\subsection{Aplicações típicas em avaliação}

\begin{itemize}
  \item \textbf{Variável de capacidade institucional}: existência de secretaria específica, de plano setorial ou de conselho municipal como variável de controle em avaliações de políticas descentralizadas;
  \item \textbf{Análise de heterogeneidade de impacto}: verificar se políticas funcionam melhor em municípios com maior capacidade institucional;
  \item \textbf{Avaliação de políticas de fortalecimento institucional}: monitoramento da evolução da capacidade municipal ao longo do tempo;
  \item \textbf{Avaliação de políticas habitacionais e urbanísticas}: verificar se municípios com planos diretores, PLHIS ou legislação de regularização fundiária têm resultados distintos em políticas de habitação.
\end{itemize}

% ────────────────────────────────────────────────────────────
\section{Dados de Agências Reguladoras}
% ────────────────────────────────────────────────────────────

\subsection{ANEEL --- Agência Nacional de Energia Elétrica}

\ficharapida
  {Dados da ANEEL}
  {Agência Nacional de Energia Elétrica (ANEEL)}
  {https://dadosabertos.aneel.gov.br}
  {Brasil, Unidades da Federação e municípios}
  {Variável por indicador (mensal a anual)}
  {Unidade consumidora, município e distribuidora}

A ANEEL disponibiliza dados sobre cobertura, qualidade e tarifas do setor elétrico, com relevância direta para avaliação de políticas de acesso à energia. Os principais indicadores incluem número de unidades consumidoras por classe e município, indicadores de continuidade (DEC e FEC) por distribuidora, tarifas por classe de consumo e dados do Programa Luz para Todos.

\subsection{ANATEL --- Agência Nacional de Telecomunicações}

\ficharapida
  {Dados da ANATEL}
  {Agência Nacional de Telecomunicações (ANATEL)}
  {https://dados.anatel.gov.br}
  {Brasil, Unidades da Federação e municípios}
  {Trimestral a anual}
  {Município e prestadora}

A ANATEL disponibiliza dados sobre cobertura de sinal móvel (2G, 3G, 4G, 5G) e banda larga fixa por município, incluindo mapa georreferenciado de cobertura. Para avaliações de políticas de inclusão digital e governo digital, os dados da ANATEL são variáveis de contexto fundamentais.

\subsection{ANS --- Agência Nacional de Saúde Suplementar}

\ficharapida
  {Dados da ANS}
  {Agência Nacional de Saúde Suplementar (ANS)}
  {https://dadosabertos.ans.gov.br}
  {Brasil, Unidades da Federação e municípios}
  {Trimestral}
  {Beneficiário de plano de saúde e operadora}

A ANS disponibiliza dados sobre beneficiários de planos de saúde por município, faixa etária e sexo, operadoras ativas por modalidade e indicadores de qualidade assistencial (Idss). Para avaliações de políticas de saúde que afetam tanto o SUS quanto a saúde suplementar, os dados da ANS são indispensáveis para capturar o universo completo do setor. Um uso recorrente é a análise de efeitos de \textit{crowding-out} entre público e privado: quando a oferta pública de serviços aumenta --- por exemplo, com a expansão da atenção primária ou a abertura de novos hospitais públicos ---, é possível testar se há redução correspondente na adesão a planos privados no mesmo território, combinando microdados da ANS com indicadores de oferta do CNES/SIH.

% ────────────────────────────────────────────────────────────
\section{Malhas Territoriais e Dados Georreferenciados do IBGE}
\label{sec:malhas}
% ────────────────────────────────────────────────────────────

\ficharapida
  {Malhas Territoriais e Setores Censitários}
  {Instituto Brasileiro de Geografia e Estatística (IBGE)}
  {https://www.ibge.gov.br/geociencias/organizacao-do-territorio/malhas-territoriais.html}
  {Brasil (todos os níveis territoriais)}
  {Atualização periódica (versões anuais das malhas municipais)}
  {Setor censitário, município, UF, região}

\subsection{O que são e por que importam}

As malhas territoriais do IBGE são os arquivos cartográficos que definem os limites geográficos das unidades territoriais brasileiras em diferentes níveis de desagregação: países, grandes regiões, unidades da federação, mesorregiões, microrregiões, municípios, distritos, subdistritos e \textbf{setores censitários} --- a menor unidade territorial do sistema estatístico brasileiro.

Para a avaliação de políticas públicas, as malhas territoriais cumprem funções essenciais:

\begin{itemize}
  \item \textbf{Visualização cartográfica}: produção de mapas de cobertura, distribuição de beneficiários e indicadores de resultado por território;
  \item \textbf{Análise espacial}: cálculo de distâncias entre pontos (domicílios e equipamentos públicos), identificação de áreas de influência, análise de autocorrelação espacial e spillovers entre territórios;
  \item \textbf{Integração territorial}: combinação de bases de dados que usam diferentes identificadores geográficos. Muitas bases relevantes para avaliação --- como cadastros imobiliários municipais com endereços, registros de ITBI com CEP, ou dados de licenciamento urbanístico --- não compartilham identificadores administrativos comuns com bases como o Censo Demográfico ou o Censo Escolar. A integração geoespacial, por meio de operações de intersecção ou join espacial, frequentemente é o único caminho para combiná-las.
\end{itemize}

\subsection{Setores censitários}

Os setores censitários são a unidade territorial mais granular do sistema de informações geográficas do IBGE. Cada setor agrupa aproximadamente 200 a 300 domicílios em áreas urbanas e é a unidade de coleta do Censo Demográfico.

Os setores censitários têm dois usos principais em avaliação:

\begin{enumerate}
  \item \textbf{Análise inframunicipal}: para políticas implementadas em nível de bairro ou comunidade, os dados do Censo Demográfico agregados por setor permitem análises com desagregação muito maior do que o nível municipal;
  \item \textbf{Estratégia de identificação geográfica}: em avaliações de impacto, a descontinuidade nos limites de setores ou municípios pode ser explorada como fonte de variação exógena (regressão descontínua geográfica).
\end{enumerate}

\subsection{Principais arquivos disponíveis}

\begin{itemize}
  \item \textbf{Malhas municipais}: polígonos dos 5.570 municípios brasileiros, com atualização anual. Disponíveis em formato Shapefile, GeoJSON e GPKG;
  \item \textbf{Malha de setores censitários}: polígonos de todos os setores censitários do Censo 2022, com código identificador linkável aos microdados do Censo;
  \item \textbf{Grade estatística}: malha regular de células de 1 km² com dados populacionais agregados, útil para análises que não seguem limites administrativos;
  \item \textbf{Localidades}: pontos georreferenciados de sedes municipais, distritos e aglomerados rurais;
  \item \textbf{CNEFE} (Cadastro Nacional de Endereços para Fins Estatísticos): cadastro georreferenciado de todos os endereços do Brasil identificados durante a coleta censitária, valioso para análises de acessibilidade e proximidade a serviços públicos.
\end{itemize}

\begin{boxdica}{Dica --- Pacotes para análise espacial}
\begin{itemize}
  \item \textbf{R}: os pacotes \texttt{sf} (leitura e manipulação de dados espaciais), \texttt{geobr} (download de malhas do IBGE diretamente no R) e \texttt{spdep} (análise de autocorrelação espacial) formam o conjunto mais utilizado para análise espacial em políticas públicas;
  \item \textbf{Python}: as bibliotecas \texttt{geopandas} e \texttt{shapely} oferecem funcionalidades equivalentes; \texttt{pygeobr} facilita o download das malhas do IBGE;
  \item \textbf{QGIS}: software de SIG gratuito e de código aberto, sem necessidade de programação, indicado para visualizações cartográficas e análises espaciais mais simples.
\end{itemize}
\end{boxdica}

% ────────────────────────────────────────────────────────────
\section{Habitação e Mercado Imobiliário: bases municipais}
\label{sec:habitacao}
% ────────────────────────────────────────────────────────────

As avaliações de políticas habitacionais, de regulação urbana e de mercado imobiliário dependem de fontes que raramente estão disponíveis em bases nacionais padronizadas. A informação mais granular e relevante sobre transações imobiliárias, valores de imóveis e características urbanísticas está dispersa em sistemas municipais --- administrados pelas prefeituras no exercício de suas competências tributárias e de planejamento urbano. Esta seção apresenta as principais fontes disponíveis nesse ecossistema, seus usos e suas limitações.

\subsection{Bases de ITBI --- Imposto sobre Transmissão de Bens Imóveis}

O ITBI é um tributo de competência municipal (art. 156, II da Constituição Federal) incidente sobre a transmissão onerosa de bens imóveis entre vivos --- compras e vendas, dações em pagamento, permutas. Cada transação imobiliária gera um registro administrativo no sistema tributário da prefeitura, e esses registros constituem uma das fontes mais ricas para análises de mercado imobiliário disponíveis no Brasil.

A depender do município, os microdados de ITBI podem conter:

\begin{itemize}
  \item Data da transação;
  \item Valor declarado pelo comprador (base de cálculo do imposto);
  \item Valor venal do imóvel (valor atribuído pela prefeitura para fins fiscais);
  \item Tipo de imóvel (residencial, comercial, terreno, galpão);
  \item Área do imóvel (construída e/ou de terreno, quando disponível);
  \item Localização: endereço completo, CEP, bairro --- permitindo georreferenciamento para análise espacial;
  \item Tipo de financiamento (à vista, financiamento bancário, FGTS), em municípios com dados mais detalhados.
\end{itemize}

Essas bases são \textbf{centrais em avaliações urbanas} porque permitem acompanhar transações imobiliárias com granularidade muito superior à das pesquisas nacionais. Para avaliações de impacto de projetos de infraestrutura urbana --- extensões de metrô, corredores de BRT, requalificação de áreas centrais, programas habitacionais como o Minha Casa Minha Vida, obras de saneamento, parques lineares --- sobre os preços dos imóveis; de políticas de regularização fundiária sobre o mercado formal; ou de programas habitacionais sobre a dinâmica do entorno, as bases de ITBI são frequentemente insubstituíveis.

\textbf{Acesso}: os microdados de ITBI não estão disponíveis em nenhum repositório nacional centralizado. O acesso deve ser obtido diretamente junto às prefeituras --- via portais de dados abertos municipais (quando existentes) ou por pedido LAI. Municípios como São Paulo, Belo Horizonte e Porto Alegre disponibilizam dados de ITBI em seus portais de dados abertos com graus variados de detalhe.

\begin{boxatencao}{Atenção --- Valor declarado vs. valor de mercado}
O valor declarado para fins de ITBI é frequentemente \textit{diferente} do preço efetivo de mercado da transação. Em muitos casos, compradores e vendedores têm incentivos para declarar valores menores do que o preço efetivo pago, reduzindo a base de cálculo do imposto. Em municípios onde a alíquota incide sobre o maior entre o valor declarado e o valor venal, esse comportamento é atenuado --- mas não eliminado. Análises que usam o valor declarado de ITBI como proxy de preço de mercado devem discutir esse potencial viés explicitamente.
\end{boxatencao}

\subsection{Cadastro Imobiliário Municipal e IPTU}

O cadastro imobiliário municipal é o registro de todos os imóveis situados no território do município, mantido pela prefeitura para fins de administração tributária --- especialmente do IPTU (Imposto Predial e Territorial Urbano). É uma fonte com potencial analítico enorme para avaliações de habitação e uso do solo.

As principais informações que um cadastro imobiliário bem estruturado pode conter incluem:

\begin{itemize}
  \item \textbf{Identificação do imóvel}: número de inscrição cadastral, endereço, coordenadas geográficas;
  \item \textbf{Características físicas}: área de terreno, área construída, número de pavimentos, tipo de uso (residencial/comercial/misto/industrial), padrão construtivo;
  \item \textbf{Situação jurídica}: proprietário ou possuidor, tipo de domínio (pleno, enfiteuse, aforamento), existência de registro em cartório;
  \item \textbf{Informações tributárias}: valor venal do terreno e da edificação calculado segundo a Planta Genérica de Valores (PGV); alíquota aplicada; histórico de pagamento do IPTU.
\end{itemize}

A \textbf{Planta Genérica de Valores (PGV)} é o instrumento municipal que define os valores unitários de terrenos e edificações por logradouro ou zona, usados para calcular o valor venal dos imóveis. Para avaliações que precisam de um índice de valor imobiliário com cobertura ampla (não apenas imóveis transacionados, como no ITBI), a PGV pode ser uma alternativa útil, embora reflita valores administrativos e não necessariamente de mercado.

\subsection{Registros de Licenciamento Urbanístico}

As prefeituras mantêm registros administrativos de todas as solicitações e aprovações de obras e intervenções no território: alvarás de construção, reformas, demolições, aprovação de projetos e emissão de Habite-se (documento que atesta que a obra foi concluída conforme o projeto aprovado).

Esses registros são relevantes para avaliações que precisam rastrear a dinâmica de produção do espaço urbano:

\begin{itemize}
  \item \textbf{Avaliação de políticas de regularização fundiária}: verificar se a regularização de assentamentos informais induz investimentos em melhorias habitacionais (proxied pelo número de alvarás de reforma na área regularizada);
  \item \textbf{Análise de adensamento e verticalização}: acompanhar a aprovação de projetos de múltiplos pavimentos em áreas sujeitas a mudanças de zoneamento;
  \item \textbf{Avaliação de programas de habitação de interesse social}: monitorar a emissão de Habite-se em unidades do MCMV para verificar a efetiva entrega das unidades.
\end{itemize}

\subsection{Bases de Zoneamento e Uso do Solo}

Municípios com sistemas de informação mais estruturados disponibilizam, em formatos georreferenciados, suas legislações de zoneamento --- definindo as zonas de uso permitido (residencial, comercial, industrial, misto), os índices urbanísticos (coeficiente de aproveitamento, taxa de ocupação, recuo mínimo) e as áreas de proteção ambiental e de especial interesse social (AEIS).

Para avaliações de políticas de regulação urbana --- como mudanças de zoneamento, criação de AEIS ou adoção de instrumentos do Estatuto da Cidade como IPTU progressivo e operações urbanas consorciadas ---, as bases de zoneamento georreferenciadas são insumos essenciais.

\subsection{Alerta: heterogeneidade e comparabilidade entre municípios}

As bases municipais sobre mercado imobiliário e urbanismo são extraordinariamente valiosas, mas apresentam graus de heterogeneidade que não têm equivalente em fontes nacionais padronizadas. Antes de realizar comparações entre municípios ou análises longitudinais dentro de um mesmo município, o pesquisador deve investigar:

\begin{itemize}
  \item \textbf{Layout e estrutura}: cada município usa seu próprio sistema de gestão tributária, com campos, nomenclaturas e formatos distintos. A mesma variável (área construída, tipo de uso) pode ter definições e unidades diferentes entre sistemas;
  \item \textbf{Cobertura}: o cadastro imobiliário municipal pode estar desatualizado ou ter cobertura parcial --- especialmente em áreas de ocupação informal ou em municípios com menor capacidade administrativa;
  \item \textbf{Mudanças de regra tributária}: alterações na PGV, nas alíquotas do IPTU ou nos critérios de avaliação imobiliária ao longo do tempo introduzem quebras de série que precisam ser identificadas antes de análises longitudinais;
  \item \textbf{Mudanças de sistema}: migrações de sistema informatizado (comuns em prefeituras de médio porte) podem gerar inconsistências entre períodos;
  \item \textbf{Padronização de endereços}: endereços coletados em sistemas administrativos municipais frequentemente precisam de tratamento para georreferenciamento --- abreviações, erros de digitação, logradouros sem CEP ou com CEP desatualizado são comuns.
\end{itemize}

Em comparações entre municípios, a harmonização dessas dimensões é um trabalho metodológico substantivo que deve ser documentado com o mesmo rigor que o próprio modelo de avaliação.

% ────────────────────────────────────────────────────────────
\section{Dados Ambientais: PRODES, DETER e MapBiomas}
% ────────────────────────────────────────────────────────────

\subsection{PRODES --- Monitoramento do Desmatamento}

\ficharapida
  {PRODES --- Projeto de Monitoramento do Desmatamento na Amazônia Legal por Satélite}
  {Instituto Nacional de Pesquisas Espaciais (INPE)}
  {https://data.inpe.br/biomasbr/prodes-monitoramento-anual-da-supressao-de-vegetacao-nativa/}
  {Amazônia Legal (dados mais recentes para outros biomas)}
  {Anual (desde 1988)}
  {Polígono de desmatamento (georreferenciado)}

O PRODES é o sistema oficial de monitoramento do desmatamento na Amazônia Legal, operado pelo INPE desde 1988. Utiliza imagens de satélite para identificar e quantificar as áreas de floresta desmatadas a cada ano. Os dados são disponibilizados em formato georreferenciado (Shapefile e GeoTIFF) e incluem polígonos de desmatamento por ano com área calculada, taxa anual de desmatamento por município e estado, e dados históricos desde 1988 para análises de longo prazo.

Para avaliações de políticas ambientais --- como o Plano de Ação para Prevenção e Controle do Desmatamento na Amazônia Legal (PPCDAm), a criação de unidades de conservação ou a implementação do CAR ---, o PRODES é a referência nacional e internacionalmente reconhecida.

\subsection{DETER --- Sistema de Detecção do Desmatamento em Tempo Real}

\ficharapida
  {DETER --- Monitoramento Diário da Supressão e Degradação de Vegetação Nativa}
  {Instituto Nacional de Pesquisas Espaciais (INPE)}
  {https://data.inpe.br/biomasbr/deter-monitoramento-diario-da-supressao-e-degradacao-de-vegetacao-nativa/}
  {Amazônia Legal e Cerrado}
  {Contínuo (alertas em dias a semanas)}
  {Alerta georreferenciado}

O DETER complementa o PRODES com detecção mais frequente de alertas de desmatamento e degradação florestal. É a fonte utilizada pelo Ibama para acionar operações de fiscalização em tempo real. Tem aplicações em avaliações de fiscalização ambiental e de compliance de políticas de proteção florestal.

\subsection{MapBiomas}

\ficharapida
  {MapBiomas --- Mapeamento Anual de Uso e Cobertura da Terra no Brasil}
  {Iniciativa MapBiomas}
  {https://mapbiomas.org}
  {Brasil (todos os biomas)}
  {Anual (desde 1985)}
  {Pixel de 30m (georreferenciado)}

O MapBiomas é uma iniciativa colaborativa entre universidades, ONGs e empresas de tecnologia que produz mapas anuais de uso e cobertura do solo para todo o território brasileiro, desde 1985. Diferentemente do PRODES --- focado no desmatamento ---, o MapBiomas classifica cada pixel do território em categorias de uso do solo (floresta, pastagem, agricultura, área urbana, água, etc.).

Para avaliações de políticas de uso do solo, de agricultura e de meio ambiente, o MapBiomas oferece uma perspectiva de longo prazo sobre as mudanças no território que nenhuma outra fonte disponibiliza com essa cobertura e granularidade. A série desde 1985 permite analisar trajetórias de conversão do uso do solo em resposta a políticas ambientais, agrícolas e de infraestrutura.

% ────────────────────────────────────────────────────────────
\section{Cadastro Ambiental Rural (CAR)}
% ────────────────────────────────────────────────────────────

\ficharapida
  {Cadastro Ambiental Rural (CAR)}
  {Serviço Florestal Brasileiro (SFB) / Ministério do Meio Ambiente}
  {https://www.gov.br/gestao/pt-br/assuntos/cadastro-ambiental-rural-car-1}
  {Brasil}
  {Contínuo (desde 2013; obrigatoriedade a partir de 2016)}
  {Imóvel rural}

O CAR é o registro eletrônico obrigatório para todos os imóveis rurais brasileiros, criado pelo Código Florestal de 2012. Cada imóvel registrado tem seus limites georreferenciados e a identificação das áreas de preservação permanente (APP), reserva legal e uso consolidado.

Para avaliações de políticas ambientais e agrícolas, o CAR é a principal fonte de dados sobre a situação fundiária e ambiental dos imóveis rurais, sendo especialmente relevante para avaliação da adesão ao Código Florestal, análise de políticas de regularização ambiental rural (PRA) e crédito rural condicionado à regularidade ambiental.

% ────────────────────────────────────────────────────────────
\section{Outras fontes de infraestrutura e território}
% ────────────────────────────────────────────────────────────

\subsection{DNIT --- Malha Rodoviária}

O Departamento Nacional de Infraestrutura de Transportes (DNIT) disponibiliza a malha rodoviária federal georreferenciada, com classificação das rodovias por tipo de pista, superfície e condição. Para avaliações de políticas de transporte e de acesso a serviços em áreas remotas, a malha rodoviária é uma variável de contexto relevante.

\subsection{CNT --- Pesquisa Rodoviária}

A Confederação Nacional do Transporte (CNT) realiza anualmente a Pesquisa Rodoviária, que avalia as condições de pavimentação, sinalização e geometria de rodovias federais e estaduais. Os dados são disponibilizados com georreferenciamento e permitem análises de qualidade da infraestrutura de transportes ao longo do tempo.

\subsection{CNEFE --- Cadastro Nacional de Endereços para Fins Estatísticos}

O CNEFE, produzido pelo IBGE a partir dos dados do Censo Demográfico, é um cadastro georreferenciado de todos os endereços do Brasil identificados durante a coleta censitária. É uma fonte valiosa para análises de acessibilidade e proximidade a serviços públicos, especialmente quando combinada com dados georreferenciados de estabelecimentos de saúde (CNES), escolas (Censo Escolar) e equipamentos do SUAS.

% ────────────────────────────────────────────────────────────
\section{Síntese do capítulo}
% ────────────────────────────────────────────────────────────

Este capítulo apresentou as principais fontes de dados sobre infraestrutura, território e meio ambiente disponíveis para o Brasil. A Tabela~\ref{tab:sintese_cap9} resume suas características essenciais.

\begin{table}[H]
\centering
\caption{Síntese das fontes de dados de infraestrutura, território e meio ambiente}
\label{tab:sintese_cap9}
\small
\begin{tabularx}{\textwidth}{@{} >{\raggedright\arraybackslash}p{2.8cm} >{\raggedright\arraybackslash}p{3cm} >{\raggedright\arraybackslash}p{2.5cm} X @{}}
\toprule
\textbf{Fonte} & \textbf{O que registra} & \textbf{Cobertura} &
\textbf{Principal uso em avaliação} \\
\midrule
SNIS         & Saneamento básico           & Municipal (adesão voluntária) & Cobertura de água/esgoto, qualidade \\
Munic/Estadic & Capacidade institucional   & Municipal\slash \allowbreak estadual & Variável de controle institucional, habitação \\
ANEEL        & Energia elétrica            & Municipal          & Cobertura, qualidade, tarifas \\
ANATEL       & Telecomunicações            & Municipal          & Cobertura de internet e telefonia \\
ANS          & Saúde suplementar           & Municipal          & Beneficiários de planos de saúde \\
Malhas IBGE  & Limites territoriais        & Todos os níveis    & Análise espacial, visualização, integração \\
ITBI         & Transações imobiliárias     & Municipal          & Mercado imobiliário, políticas urbanas \\
Cadastro/IPTU & Imóveis e valores venais   & Municipal          & Habitação, regulação urbana, PGV \\
Licenciamento & Aprovação de obras          & Municipal          & Produção habitacional, adensamento \\
PRODES/INPE  & Desmatamento                & Amazônia\slash \allowbreak Brasil    & Políticas ambientais \\
DETER/INPE   & Alertas de desmatamento     & Amazônia\slash \allowbreak Cerrado   & Fiscalização ambiental \\
MapBiomas    & Uso e cobertura do solo     & Brasil (desde 1985) & Mudanças no uso do solo \\
CAR          & Imóveis rurais              & Brasil             & Código Florestal, crédito rural \\
\bottomrule
\end{tabularx}
\fonte{FGV CLEAR, elaboração própria.}
\end{table}

Os pontos centrais do capítulo são:

\begin{itemize}
  \item O \textbf{SNIS} é a referência para dados de saneamento básico, mas sua cobertura é heterogênea entre municípios e os dados são autodeclarados pelos prestadores;

  \item A \textbf{Munic} e a \textbf{Estadic} oferecem informações sobre capacidade institucional municipal e estadual que são raramente disponíveis em outras fontes e muito úteis como variáveis de controle em avaliações de políticas descentralizadas --- incluindo políticas habitacionais e urbanísticas;

  \item As \textbf{bases municipais de ITBI, cadastro imobiliário e IPTU} são fontes centrais para avaliações de habitação e mercado imobiliário, com granularidade muito superior às pesquisas nacionais --- mas apresentam heterogeneidade substancial em layout, cobertura e qualidade entre municípios, o que exige harmonização cuidadosa antes de comparações;

  \item Os \textbf{registros de licenciamento urbanístico} e as \textbf{bases de zoneamento} permitem rastrear a dinâmica da produção do espaço urbano em resposta a políticas de habitação e regulação;

  \item As \textbf{malhas territoriais e setores censitários} do IBGE são a base para qualquer análise espacial e, frequentemente, o único caminho para cruzar fontes que não compartilham identificadores administrativos comuns;

  \item O \textbf{PRODES/INPE}, o \textbf{DETER} e o \textbf{MapBiomas} são as referências para dados ambientais, com séries históricas longas e cobertura nacional ou de biomas específicos;

  \item A \textbf{integração espacial} das fontes --- via malhas territoriais, setores censitários, CEPs ou coordenadas --- é especialmente importante em habitação, mercado imobiliário, infraestrutura urbana, saúde, crime e educação, e deve ser tratada como parte central do desenho metodológico, não como passo técnico secundário.
\end{itemize}

O próximo capítulo aborda as fontes de dados sobre finanças públicas, com foco no FINBRA/Siconfi, nos portais de transparência e nos dados do TSE.

% ============================================================
%  cap10_financas_publicas.tex  [REVISADO]
%  Capítulo 10: Finanças Públicas
% ============================================================

% ════════════════════════════════════════════════════════════
\chapter{Finanças Públicas}
% ════════════════════════════════════════════════════════════

\begin{boxdestaque}{Neste capítulo}
Este capítulo apresenta as principais fontes de dados sobre finanças públicas disponíveis para o Brasil. Cobrimos o FINBRA/Siconfi e seus instrumentos de reporte (RREO, RGF, balanços), os portais de transparência federal e subnacionais, o SIOPE, o SIOPS, os dados do TSE sobre financiamento eleitoral e os dados do Banco Central. Dedicamos atenção especial às finanças públicas subnacionais --- receitas tributárias municipais, capacidade fiscal e estrutura de gasto local ---, dimensão central em avaliações de políticas descentralizadas. O capítulo discute também os principais desafios de comparabilidade e qualidade dos dados fiscais brasileiros.
\end{boxdestaque}

% ────────────────────────────────────────────────────────────
\section{Introdução: por que dados fiscais importam para avaliação}
% ────────────────────────────────────────────────────────────

Dados de finanças públicas são fundamentais para a avaliação de políticas públicas por pelo menos três razões. Primeira, o \textbf{gasto público é frequentemente a política}: em muitos casos, o que se avalia não é uma intervenção específica com regras de elegibilidade bem definidas, mas o impacto de diferentes níveis ou composições do gasto público sobre resultados de interesse. Segunda, dados fiscais permitem calcular \textbf{custos de políticas}, insumo indispensável para análises de custo-efetividade e custo-benefício. Terceira, a \textbf{capacidade fiscal} dos entes federativos é um determinante importante da qualidade e da cobertura dos serviços públicos --- e, portanto, uma variável de contexto relevante em avaliações de políticas descentralizadas.

O Brasil tem um sistema de transparência fiscal relativamente bem desenvolvido para os padrões internacionais, com obrigações de publicação de dados fiscais para todos os entes federativos desde a Lei de Responsabilidade Fiscal (Lei Complementar n.º 101/2000) e a Lei de Transparência (Lei Complementar n.º 131/2009). Na prática, no entanto, a qualidade e a completude dos dados variam consideravelmente entre municípios, estados e o governo federal.

\begin{table}[H]
\centering
\caption{Principais fontes de dados de finanças públicas}
\label{tab:ecossistema_fiscal}
\small
\begin{tabularx}{\textwidth}{@{} p{3cm} p{3.5cm} p{2cm} X @{}}
\toprule
\textbf{Fonte} & \textbf{O que registra} & \textbf{Desde} &
\textbf{Principal uso em avaliação} \\
\midrule
FINBRA/Siconfi   & Receitas e despesas municipais e estaduais & 1986 & Gasto municipal e estadual por função \\
RREO             & Execução orçamentária bimestral  & 2000 & Acompanhamento infraanual do gasto \\
RGF              & Limites e alertas fiscais (LRF)  & 2000 & Capacidade fiscal, endividamento \\
Portal Transparência & Execução orçamentária federal & 2004 & Gasto federal por programa \\
SIOPE            & Gasto em educação               & 1999 & Financiamento, custo-aluno \\
SIOPS            & Gasto em saúde                  & 2000 & Gasto per capita em saúde \\
TSE              & Financiamento eleitoral         & 2002 & Recursos de campanha, doadores \\
BCB/SGS          & Séries macroeconômicas          & Variável & Deflacionamento, contexto macro \\
BCB/Estban       & Crédito por município           & Mensal & Acesso ao crédito, políticas de crédito \\
\bottomrule
\end{tabularx}
\fonte{FGV CLEAR, elaboração própria.}
\end{table}

% ────────────────────────────────────────────────────────────
\section{FINBRA e Siconfi --- Finanças dos Municípios e Estados}
% ────────────────────────────────────────────────────────────

\ficharapida
  {FINBRA / Siconfi}
  {Secretaria do Tesouro Nacional (STN)}
  {https://siconfi.tesouro.gov.br}
  {Brasil, Unidades da Federação e municípios}
  {Anual (FINBRA desde 1986; Siconfi desde 2013)}
  {Ente federativo (município ou estado)}

\subsection{O que é e evolução histórica}

O FINBRA (Finanças do Brasil) foi, por décadas, a principal base de dados sobre receitas e despesas dos municípios brasileiros, produzida pela Secretaria do Tesouro Nacional (STN). Em 2013, o FINBRA foi integrado ao \textbf{Siconfi} (Sistema de Informações Contábeis e Fiscais do Setor Público Brasileiro), plataforma unificada que passou a consolidar as declarações fiscais de municípios, estados e União em um único sistema.

O Siconfi opera por meio de declarações anuais obrigatórias que todos os entes federativos devem prestar à STN. Além das declarações anuais, os entes devem publicar periodicamente dois instrumentos de acompanhamento infraanual que são fontes complementares ao FINBRA/Siconfi:

\begin{itemize}
  \item \textbf{RREO --- Relatório Resumido da Execução Orçamentária}: publicado bimestralmente por municípios, estados e União. Contém dados de execução de receitas e despesas no bimestre e no acumulado do exercício, discriminados por fonte e por função. É a fonte mais atualizada para acompanhar o ritmo de execução do gasto público ao longo do ano, antes da consolidação anual no Siconfi;

  \item \textbf{RGF --- Relatório de Gestão Fiscal}: publicado quadrimestralmente. Contém os indicadores de cumprimento dos limites estabelecidos pela Lei de Responsabilidade Fiscal --- despesa com pessoal (em relação à receita corrente líquida), dívida consolidada líquida, garantias concedidas e operações de crédito. O RGF é a fonte para análises de capacidade fiscal, endividamento e risco fiscal dos entes subnacionais.
\end{itemize}

Ambos os instrumentos estão disponíveis no portal Siconfi e nos portais de transparência de cada ente.

\subsection{Estrutura das despesas: classificações orçamentárias}

Para o uso analítico dos dados do Siconfi/FINBRA, é fundamental entender as classificações orçamentárias brasileiras. As despesas públicas são classificadas de três formas complementares:

\begin{itemize}
  \item \textbf{Por função}: classifica o gasto pela área de atuação governamental --- Saúde (função 10), Educação (função 12), Assistência Social (função 08), Segurança Pública (função 06), Habitação (função 16), Infraestrutura Urbana (função 15), etc. É a classificação mais usada em análises comparativas;

  \item \textbf{Por natureza}: classifica o gasto pelo tipo de despesa --- Pessoal e Encargos Sociais (grupo 1), Juros e Encargos da Dívida (grupo 2), Outras Despesas Correntes (grupo 3), Investimentos (grupo 4), Inversões Financeiras (grupo 5), Amortização da Dívida (grupo 6). Permite analisar a composição do gasto (corrente vs. capital);

  \item \textbf{Por programa}: classifica o gasto pelo programa de governo ao qual está associado. A comparabilidade entre municípios é menor porque cada um define seus próprios programas.
\end{itemize}

\begin{boxatencao}{Atenção --- Classificação por função não é equivalente a gasto efetivo em cada política}
A classificação orçamentária por função é uma aproximação do gasto em cada área de política, não uma medida precisa. Um município pode classificar despesas administrativas gerais sob a função Educação, inflando artificialmente o gasto educacional. Antes de usar dados de gasto por função como variável de resultado ou de controle, verifique se a classificação é consistente entre os entes comparados e ao longo do tempo.
\end{boxatencao}

\subsection{Receitas: composição e estrutura}

Os dados de receitas do Siconfi permitem analisar a composição da arrecadação municipal e estadual:

\begin{itemize}
  \item \textbf{Receitas próprias}: tributos de competência do ente. Para municípios, as principais são:
  \begin{itemize}
    \item \textbf{IPTU} (Imposto Predial e Territorial Urbano): imposto sobre a propriedade imobiliária urbana, de periodicidade anual. A arrecadação agregada de IPTU por município está disponível no Siconfi e é um indicador de capacidade tributária local e de dinâmica do mercado imobiliário. Municípios com maior base imobiliária formal tendem a ter maior arrecadação de IPTU per capita;
    \item \textbf{ISS} (Imposto Sobre Serviços): principal fonte de receita própria municipal em cidades com setor de serviços desenvolvido;
    \item \textbf{ITBI} (Imposto sobre Transmissão de Bens Imóveis): incide sobre transações imobiliárias. A arrecadação agregada de ITBI é um indicador da intensidade do mercado imobiliário local --- municípios com maior número de transações imobiliárias arrecadam mais ITBI. Séries históricas de arrecadação de ITBI por município (disponíveis no Siconfi) podem ser usadas como proxy da atividade do mercado imobiliário quando os microdados de transações não estão disponíveis;
  \end{itemize}
  \item \textbf{Transferências constitucionais}: FPM, cota-parte do ICMS, royalties e outras;
  \item \textbf{Transferências voluntárias}: convênios e contratos de repasse da União;
  \item \textbf{Operações de crédito}: empréstimos e financiamentos.
\end{itemize}

\subsection{Capacidade fiscal e autonomia tributária}

A \textbf{capacidade fiscal} de um município --- sua capacidade de gerar receitas próprias independentemente de transferências --- é uma variável importante em avaliações de políticas descentralizadas. Municípios com maior autonomia tributária têm maior margem para adaptar o gasto às preferências locais e menor vulnerabilidade a cortes nas transferências federais.

O Siconfi permite calcular indicadores de capacidade fiscal, como:

\begin{itemize}
  \item \textbf{Receita própria per capita}: receitas tributárias (IPTU + ISS + ITBI + taxas) divididas pela população estimada;
  \item \textbf{Grau de dependência de transferências}: proporção da receita total proveniente de transferências constitucionais e voluntárias;
  \item \textbf{Esforço fiscal}: arrecadação própria realizada em relação ao potencial fiscal estimado (requer modelagem adicional).
\end{itemize}

Para análises de capacidade fiscal que exijam maior granularidade do que o Siconfi oferece, o \textbf{Ipeadata} (\url{ipeadata.gov.br}) disponibiliza séries históricas de indicadores fiscais municipais e estaduais derivados do FINBRA, com algumas variáveis já deflacionadas e normalizadas por população.

\subsection{Acesso aos dados}

\begin{itemize}
  \item \textbf{Portal Siconfi}: consultas online e download por ente, exercício e classificação orçamentária em \url{siconfi.tesouro.gov.br};
  \item \textbf{FINBRA histórico}: séries 1986--2012 disponíveis para download na STN em formato Excel;
  \item \textbf{Base dos Dados}: a plataforma \textit{Base dos Dados} (\url{basedosdados.org}) disponibiliza o FINBRA e o Siconfi compatibilizados em BigQuery, com tabelas anuais por município desde 1986;
  \item \textbf{Ipeadata}: séries históricas de indicadores fiscais municipais e estaduais com normalização e deflacionamento parcial.
\end{itemize}

\begin{boxdica}{Dica --- Deflacionando séries históricas de gasto}
Comparações de gasto público ao longo do tempo requerem deflacionamento para valores reais. O deflator mais utilizado para séries de gasto público no Brasil é o \textbf{IPCA} (Índice de Preços ao Consumidor Amplo), disponível no IBGE e no Banco Central (API SGS). Para gastos de investimento em infraestrutura, o \textbf{INCC} (Índice Nacional de Custo da Construção) pode ser mais adequado. Sempre especifique o deflator utilizado e o ano-base nas análises.
\end{boxdica}

\subsection{Aplicações típicas em avaliação}

\begin{itemize}
  \item \textbf{Análise do impacto de transferências fiscais}: avaliação de como variações no FPM ou em outras transferências afetam o gasto público municipal e indicadores de bem-estar;
  \item \textbf{Avaliação de políticas de descentralização}: análise de como a autonomia fiscal afeta a qualidade dos serviços públicos locais;
  \item \textbf{Custo-efetividade de políticas setoriais}: divisão do gasto por função pelo número de beneficiários ou unidades de serviço entregues;
  \item \textbf{Análise de ciclo político-orçamentário}: variação do gasto público municipal em anos eleitorais versus não eleitorais;
  \item \textbf{Efeito flypaper}: análise de se transferências intergovernamentais têm efeito multiplicador sobre o gasto local maior do que um aumento equivalente na renda dos moradores;
  \item \textbf{Análise de capacidade fiscal e dependência de transferências}: identificação de municípios mais vulneráveis a cortes nas transferências federais e análise de seus efeitos sobre a provisão de serviços.
\end{itemize}

\begin{boxexemplo}{Exemplo prático --- Avaliando o impacto do FPM sobre gasto municipal e resultados de saúde}
O Fundo de Participação dos Municípios (FPM) distribui recursos federais aos municípios com base em faixas populacionais discretas: municípios que cruzam determinados limiares populacionais recebem mais recursos por habitante do que municípios imediatamente abaixo do limiar. Pesquisadores utilizaram essa descontinuidade como estratégia de regressão descontínua: municípios logo acima e logo abaixo do limiar são comparáveis em tudo, exceto pelo valor das transferências recebidas. Combinando dados do FINBRA (receita do FPM e gasto por função) com dados do DATASUS (mortalidade infantil) e do Censo Escolar (matrículas e IDEB), foi possível estimar o efeito causal de um real adicional de transferência federal sobre o gasto em saúde e educação e sobre os respectivos indicadores de resultado.

\textit{Fonte: FGV CLEAR, elaboração própria.}
\end{boxexemplo}

% ────────────────────────────────────────────────────────────
\section{Portais de Transparência}
% ────────────────────────────────────────────────────────────

\ficharapida
  {Portal da Transparência do Governo Federal}
  {Controladoria-Geral da União (CGU)}
  {https://portaldatransparencia.gov.br}
  {Governo Federal}
  {Diário (desde 2004)}
  {Despesa pública, transferência, beneficiário de programa social}

\subsection{Portal da Transparência Federal}

O Portal da Transparência do Governo Federal, operado pela CGU, disponibiliza em nível de detalhe muito maior do que o Siconfi:

\begin{itemize}
  \item \textbf{Despesas por empenho}: cada empenho individual, com identificação do credor (CNPJ ou CPF), valor, programa, ação e unidade orçamentária;
  \item \textbf{Transferências a estados e municípios}: convênios, contratos de repasse e transferências fundo a fundo, com objeto, valor e situação;
  \item \textbf{Beneficiários de programas sociais}: pagamentos do Bolsa Família, BPC e outros programas por beneficiário (NIS);
  \item \textbf{Servidores públicos federais}: remuneração, lotação e cargo de cada servidor ativo, aposentado e pensionista do Poder Executivo federal;
  \item \textbf{Contratos e licitações}: contratos firmados pela administração federal com fornecedores, com objeto, valor e vigência.
\end{itemize}

\subsection{Portais estaduais e municipais}

A Lei de Transparência (LC n.º 131/2009) obrigou estados e municípios a disponibilizarem suas execuções orçamentárias em tempo real na internet. Na prática, a qualidade e a funcionalidade dos portais variam enormemente:

\begin{itemize}
  \item \textbf{Estados com portais bem estruturados}: São Paulo, Minas Gerais, Rio Grande do Sul e outros estados maiores têm portais de transparência com APIs e dados abertos em formatos padronizados;
  \item \textbf{Municípios}: a qualidade dos portais municipais varia dramaticamente pelo porte do município e pela capacidade técnica da administração. Para análises que precisem de dados de execução orçamentária municipal mais detalhados do que o Siconfi oferece, a alternativa é acessar diretamente o portal de transparência de cada município --- o que pode exigir automação via \textit{web scraping} para estudos com muitos municípios.
\end{itemize}

% ────────────────────────────────────────────────────────────
\section{SIOPE --- Sistema de Informações sobre Orçamentos
         Públicos em Educação}
% ────────────────────────────────────────────────────────────

O SIOPE foi detalhado no Capítulo~6 (Educação), ao qual remetemos o leitor para a descrição completa da fonte. Para fins de completude desta seção sobre finanças públicas, destacamos os principais usos do SIOPE em análises fiscais:

\begin{itemize}
  \item Cálculo do \textbf{gasto por aluno} por município e UF, combinando os dados de despesa do SIOPE com os dados de matrícula do Censo Escolar;
  \item Análise do \textbf{cumprimento do piso constitucional} de gasto em educação (25\% das receitas para municípios; 18\% para estados e União);
  \item Avaliação do impacto do \textbf{FUNDEB} sobre o gasto educacional municipal: a redistribuição de recursos do FUNDEB entre municípios de uma mesma UF cria variação exógena no gasto por aluno que pode ser explorada como estratégia de identificação.
\end{itemize}

% ────────────────────────────────────────────────────────────
\section{SIOPS --- Sistema de Informações sobre Orçamentos
         Públicos em Saúde}
% ────────────────────────────────────────────────────────────

O SIOPS foi detalhado no Capítulo~5 (Saúde). Para o contexto deste capítulo, destacamos:

\begin{itemize}
  \item O SIOPS permite calcular o \textbf{gasto per capita em saúde} por município, indicador central para análises de eficiência e equidade no financiamento da saúde pública;
  \item A \textbf{vinculação constitucional} de 15\% das receitas municipais à saúde cria uma relação mecânica entre receita e gasto em saúde que pode ser explorada como instrumento para identificar o impacto causal do gasto em saúde sobre resultados sanitários;
  \item A combinação de SIOPS + DATASUS (SIM, SINASC, SIH) + Censo Demográfico (denominadores populacionais) é a tríade de fontes para análises de eficiência do gasto público em saúde municipal.
\end{itemize}

% ────────────────────────────────────────────────────────────
\section{Dados do TSE --- Financiamento Eleitoral}
% ────────────────────────────────────────────────────────────

\ficharapida
  {Dados Eleitorais e de Financiamento de Campanha}
  {Tribunal Superior Eleitoral (TSE)}
  {https://dadosabertos.tse.jus.br}
  {Brasil, Unidades da Federação e municípios}
  {Por eleição (municipais e gerais, desde 2002)}
  {Candidato, doador e prestação de contas}

\subsection{O que é e estrutura}

O TSE disponibiliza um conjunto abrangente de dados eleitorais e de financiamento de campanha no portal de dados abertos \url{dadosabertos.tse.jus.br}. Para a avaliação de políticas públicas, os dados mais relevantes são:

\begin{itemize}
  \item \textbf{Prestações de contas de candidatos}: receitas e despesas de campanha por candidato, com identificação dos doadores (pessoas físicas e jurídicas, até 2015; após a reforma eleitoral de 2015, apenas pessoas físicas e fundos partidários) e dos fornecedores de serviços eleitorais;
  \item \textbf{Resultados eleitorais}: votação por candidato, seção eleitoral, município e UF;
  \item \textbf{Perfil dos candidatos}: partido, cargo disputado, sexo, idade, raça/cor autodeclarada, escolaridade, ocupação;
  \item \textbf{Filiações partidárias}: número de filiados por partido e município ao longo do tempo.
\end{itemize}

\subsection{Aplicações típicas em avaliação}

\begin{itemize}
  \item \textbf{Ciclo político-orçamentário}: análise de como a proximidade de eleições afeta o gasto público municipal (combinando dados do TSE com FINBRA);
  \item \textbf{Impacto do financiamento de campanha sobre políticas}: avaliação de se candidatos que receberam doações de determinados setores econômicos favoreceram esses setores quando eleitos;
  \item \textbf{Efeito da alternância política}: análise de como mudanças no partido do prefeito afetam a alocação do gasto municipal;
  \item \textbf{Representatividade e políticas públicas}: análise de se a eleição de candidatas mulheres, negros ou de determinados partidos afeta a composição do gasto público;
  \item \textbf{Avaliação de reformas eleitorais}: impacto das mudanças nas regras de financiamento de campanha sobre o perfil dos candidatos eleitos e sobre as políticas adotadas.
\end{itemize}

\begin{boxexemplo}{Exemplo prático --- Avaliando o impacto da proibição de doações empresariais sobre o perfil dos eleitos}
A decisão do STF em 2015 que proibiu doações de empresas a campanhas eleitorais foi implementada a partir das eleições municipais de 2016. Pesquisadores usaram essa mudança exógena nas regras de financiamento como um experimento natural: compararam o perfil dos vereadores e prefeitos eleitos em 2016 (sem doações empresariais) com os eleitos em 2012 (com doações empresariais), controlando por características dos municípios. Os dados do TSE sobre receitas de campanha por fonte de financiamento foram fundamentais para identificar quais candidatos foram mais afetados pela mudança.

\textit{Fonte: FGV CLEAR, elaboração própria.}
\end{boxexemplo}

% ────────────────────────────────────────────────────────────
\section{Dados do Banco Central}
% ────────────────────────────────────────────────────────────

\ficharapida
  {Sistema Gerenciador de Séries Temporais (SGS) e outras bases do BCB}
  {Banco Central do Brasil (BCB)}
  {https://www.bcb.gov.br/estatisticas}
  {Brasil (séries macroeconômicas) e municípios (crédito)}
  {Variável (mensal a anual)}
  {Série macroeconômica ou município}

\subsection{O que é e principais bases}

O Banco Central do Brasil disponibiliza um conjunto abrangente de dados macroeconômicos e financeiros com relevância para a avaliação de políticas públicas:

\begin{itemize}
  \item \textbf{SGS --- Sistema Gerenciador de Séries Temporais}: base central de séries macroeconômicas do BCB, acessível via API. Inclui IPCA, IGP, taxa Selic, câmbio, PIB, índices de preços setoriais e centenas de outras séries. É a fonte de referência para deflacionamento de séries históricas de gasto público;

  \item \textbf{Estban --- Estatística Bancária por Município}: dados mensais de depósitos e operações de crédito por município e por tipo (crédito pessoal, habitacional, rural, para empresas). Fundamental para análises de acesso ao crédito e de impacto de políticas de crédito rural (como o Pronaf) por município;

  \item \textbf{SCR --- Sistema de Informações de Crédito}: dados individualizados de operações de crédito (acesso restrito via convênio). Permite análises de acesso ao crédito, inadimplência e perfil dos tomadores;

  \item \textbf{PIX e meios de pagamento}: dados de transações financeiras digitais por município, relevantes para análises de inclusão financeira.
\end{itemize}

\subsection{Aplicações típicas em avaliação}

\begin{itemize}
  \item \textbf{Avaliação de políticas de crédito rural}: os dados do Estban (crédito rural por município) combinados com dados do Censo Agropecuário permitem avaliar o impacto do Pronaf sobre produção agrícola e renda rural;
  \item \textbf{Deflacionamento de séries de gasto}: o SGS é a fonte de referência para obter o IPCA e outros índices de preços para deflacionar séries históricas de receitas e despesas municipais;
  \item \textbf{Inclusão financeira}: análise de acesso a serviços bancários por município, combinada com dados socioeconômicos do Censo.
\end{itemize}

% ────────────────────────────────────────────────────────────
\section{Desafios de qualidade e comparabilidade dos dados fiscais}
% ────────────────────────────────────────────────────────────

Trabalhar com dados fiscais brasileiros exige atenção especial a três categorias de problemas que são mais agudos nessa área do que em outras fontes:

\subsection{Heterogeneidade de classificação}

A classificação orçamentária brasileira --- especialmente para municípios --- oferece considerável discricionariedade na alocação de despesas entre funções, subfunções e programas. Isso cria heterogeneidade na classificação que pode ser sistematicamente correlacionada com características dos municípios (porte, capacidade técnica, orientação política), comprometendo comparações.

\subsection{Competência vs. caixa}

Os dados orçamentários podem ser registrados em regime de \textbf{competência} (quando a obrigação é gerada) ou de \textbf{caixa} (quando o pagamento é efetivamente realizado). O Brasil adota formalmente o regime de competência para a despesa, mas o RREO e outros instrumentos reportam também o regime de caixa. Análises que misturam os dois regimes produzem resultados inconsistentes.

\subsection{Consistência entre fontes}

Um mesmo indicador fiscal pode diferir entre o Siconfi, o portal de transparência e os dados da STN, por razões que incluem diferentes momentos de extração, revisões retroativas e diferenças nos critérios de consolidação. Sempre que possível, verifique a consistência entre fontes e documente qual fonte foi utilizada e em que momento foi extraída.

\subsection{Qualidade variável entre municípios}

A capacidade técnica das equipes de finanças municipais varia dramaticamente entre municípios de diferentes portes e regiões. Municípios muito pequenos ou com menor capacidade administrativa tendem a ter dados de menor qualidade, com mais campos em branco, classificações incorretas e inconsistências entre as diferentes peças orçamentárias. Em análises com amostra ampla de municípios, é boa prática verificar a distribuição dos dados e identificar possíveis outliers antes da análise.

% ────────────────────────────────────────────────────────────
\section{Síntese do capítulo}
% ────────────────────────────────────────────────────────────

Este capítulo apresentou as principais fontes de dados de finanças públicas disponíveis para o Brasil. A Tabela~\ref{tab:sintese_cap10} resume suas características essenciais.

\begin{table}[H]
\centering
\caption{Síntese das fontes de dados de finanças públicas}
\label{tab:sintese_cap10}
\small
\begin{tabularx}{\textwidth}{@{} p{3cm} p{3.5cm} p{2cm} X @{}}
\toprule
\textbf{Fonte} & \textbf{O que registra} & \textbf{Cobertura} &
\textbf{Principal uso em avaliação} \\
\midrule
FINBRA/Siconfi   & Receitas e despesas anuais    & Municipal/estadual & Gasto por função, capacidade fiscal \\
RREO             & Execução bimestral            & Municipal/estadual & Acompanhamento infraanual \\
RGF              & Limites LRF                   & Municipal/estadual & Endividamento, risco fiscal \\
Portal Transparência & Execução orçamentária     & Federal            & Programas, transferências, beneficiários \\
SIOPE            & Gasto em educação             & Municipal/estadual & Custo-aluno, FUNDEB \\
SIOPS            & Gasto em saúde                & Municipal/estadual & Gasto per capita em saúde \\
TSE              & Financiamento eleitoral       & Municipal/estadual & Ciclo político, reforma eleitoral \\
BCB/SGS          & Séries macroeconômicas        & Nacional           & Deflacionamento, índices de preços \\
BCB/Estban       & Crédito por município         & Municipal          & Acesso ao crédito, políticas de crédito \\
\bottomrule
\end{tabularx}
\fonte{FGV CLEAR, elaboração própria.}
\end{table}

Os pontos centrais do capítulo são:

\begin{itemize}
  \item O \textbf{FINBRA/Siconfi} é a fonte central para dados de gasto municipal e estadual, com série histórica longa. O \textbf{RREO} e o \textbf{RGF} são instrumentos complementares essenciais: o primeiro para acompanhamento infraanual da execução; o segundo para análise de capacidade fiscal e cumprimento dos limites da LRF;

  \item As \textbf{receitas tributárias municipais} --- em especial IPTU, ISS e ITBI no Siconfi --- permitem analisar capacidade fiscal local e, no caso do ITBI agregado, servem como proxy da atividade do mercado imobiliário quando os microdados de transações não estão disponíveis;

  \item Os \textbf{portais de transparência} oferecem granularidade muito maior do que o Siconfi --- até o nível do empenho individual ---, mas a qualidade e a acessibilidade dos dados variam muito entre entes federativos;

  \item O \textbf{SIOPE} e o \textbf{SIOPS} são as referências setoriais para gasto em educação e saúde, respectivamente, e são fundamentais para análises de custo-efetividade e de cumprimento dos pisos constitucionais;

  \item Os dados do \textbf{TSE} permitem análises da interface entre política eleitoral e políticas públicas, especialmente sobre ciclo político-orçamentário e impacto de reformas eleitorais;

  \item Os dados do \textbf{Banco Central} --- especialmente o SGS e o Estban --- complementam as fontes fiscais com informações sobre crédito, preços e condições macroeconômicas;

  \item A \textbf{heterogeneidade de classificação}, os desafios de comparabilidade entre fontes e a variação na qualidade do dado entre municípios são riscos específicos dos dados fiscais que exigem cuidado metodológico redobrado.
\end{itemize}

O próximo capítulo aborda as fontes internacionais e comparadas, com foco nas bases do Banco Mundial, da OCDE, do PISA e de organismos multilaterais relevantes para políticas públicas no Brasil.

% ============================================================
%  cap11_fontes_internacionais.tex
%  Capítulo 11: Fontes Internacionais e Comparadas
% ============================================================

% ════════════════════════════════════════════════════════════
\chapter{Fontes Internacionais e Comparadas}
% ════════════════════════════════════════════════════════════

\begin{boxdestaque}{Neste capítulo}
Este capítulo apresenta as principais fontes internacionais e comparadas disponíveis para pesquisadores e avaliadores de políticas públicas no Brasil. Cobrimos as bases do Banco Mundial, da OCDE, do PISA, da OIT, da OMS, do PNUD e da CEPAL, além de bases temáticas especializadas. Para cada fonte, detalhamos o que contém, como acessar, quando usar e quais são as limitações de comparabilidade com dados nacionais.
\end{boxdestaque}

% ────────────────────────────────────────────────────────────
\section{Introdução: quando e por que usar fontes internacionais}
% ────────────────────────────────────────────────────────────

Fontes internacionais cumprem um papel específico e insubstituível na avaliação de políticas públicas: permitem contextualizar o desempenho do Brasil em perspectiva comparada, identificar trajetórias de países com características similares, aprender com experiências internacionais e verificar se os resultados observados no Brasil são consistentes com evidências de outros contextos.

Há pelo menos quatro situações em que o avaliador deve considerar seriamente o uso de fontes internacionais:

\begin{enumerate}
  \item \textbf{Benchmarking}: quando a pergunta é ``como o Brasil se compara a outros países em determinado indicador?'' --- taxas de mortalidade infantil, desempenho educacional, gasto público em saúde, índice de desigualdade;
  \item \textbf{Evidências externas}: quando se quer verificar se o efeito estimado de uma política no Brasil é plausível à luz de avaliações similares em outros países;
  \item \textbf{Dados não disponíveis nacionalmente}: quando uma variável de interesse não está disponível em fontes nacionais, mas está em bases internacionais que cobrem o Brasil (como alguns indicadores ambientais ou de governança);
  \item \textbf{Análises comparativas cross-country}: quando a própria pergunta de pesquisa é comparativa --- por exemplo, o impacto de diferentes estruturas de proteção social sobre a desigualdade em países da América Latina.
\end{enumerate}

Ao mesmo tempo, fontes internacionais têm limitações estruturais que precisam ser consideradas. A harmonização metodológica necessária para comparações entre países implica perdas de detalhe e pode não capturar especificidades do contexto brasileiro. Além disso, os dados costumam ter maior defasagem do que as fontes nacionais originais, e a cobertura histórica é frequentemente menor.

\begin{boxatencao}{Atenção --- Fontes internacionais e dados nacionais}
A maior parte das bases internacionais não coleta dados primários no Brasil: ela usa como insumo os dados produzidos pelo IBGE, pelo Ministério da Saúde, pelo Inep e por outras fontes nacionais, aplicando ajustes de harmonização. Isso significa que, quando há divergência entre um indicador nacional e o equivalente em uma base internacional, a base nacional é geralmente mais precisa para análises do Brasil especificamente. Use fontes internacionais quando a comparação com outros países for o objetivo --- não como substituto para as fontes nacionais.
\end{boxatencao}

% ────────────────────────────────────────────────────────────
\section{Banco Mundial}
% ────────────────────────────────────────────────────────────

\subsection{World Development Indicators (WDI)}

\ficharapida
  {World Development Indicators (WDI)}
  {Banco Mundial}
  {https://databank.worldbank.org/source/world-development-indicators}
  {Mais de 200 países e territórios}
  {Anual (séries históricas variáveis por indicador, em geral desde 1960)}
  {País}

O WDI é a principal base de dados comparados do Banco Mundial, reunindo mais de 1.400 indicadores de desenvolvimento para mais de 200 países. Para o Brasil, os principais indicadores disponíveis incluem:

\begin{itemize}
  \item \textbf{Desenvolvimento econômico}: PIB per capita (em dólares correntes, PPC e constantes), crescimento do PIB, inflação, balança comercial, dívida pública;
  \item \textbf{Pobreza e desigualdade}: taxas de pobreza a diferentes linhas internacionais (US\$ 2,15 e US\$ 6,85 por dia em PPC 2017), índice de Gini, participação na renda por décil;
  \item \textbf{Saúde}: mortalidade infantil, mortalidade materna, esperança de vida ao nascer, cobertura de vacinação, prevalência de doenças;
  \item \textbf{Educação}: taxa de matrícula por nível, taxa de conclusão, alfabetização, razão aluno-professor;
  \item \textbf{Infraestrutura}: acesso à água tratada, saneamento, eletricidade, internet;
  \item \textbf{Governança e instituições}: indicadores do projeto Worldwide Governance Indicators (WGI).
\end{itemize}

\textbf{Acesso}: os dados estão disponíveis gratuitamente no portal DataBank do Banco Mundial, com opções de download em CSV e Excel. A API do Banco Mundial (\url{api.worldbank.org}) permite acesso programático, com pacotes disponíveis em R (\texttt{WDI}) e Python (\texttt{wbgapi}).

\subsection{PovcalNet e Poverty and Inequality Platform (PIP)}

\ficharapida
  {Poverty and Inequality Platform (PIP)}
  {Banco Mundial}
  {https://pip.worldbank.org}
  {Mais de 100 países}
  {Por pesquisa domiciliar disponível (irregular)}
  {País e, em alguns casos, área urbana/rural}

O PIP é a plataforma do Banco Mundial para estimativas de pobreza e desigualdade baseadas em microdados de pesquisas domiciliares. Para o Brasil, o PIP usa os microdados da PNAD Contínua e da POF como insumo, produzindo estimativas compatíveis internacionalmente.

É a fonte de referência para comparações de pobreza e desigualdade entre o Brasil e outros países em desenvolvimento, especialmente para análises que usam as linhas de pobreza internacionais do Banco Mundial (US\$ 2,15/dia para pobreza extrema e US\$ 6,85/dia para pobreza moderada em PPC 2017).

\subsection{Banco de Dados de Projetos do Banco Mundial}

Além das bases estatísticas, o Banco Mundial disponibiliza dados sobre seus próprios projetos de financiamento --- incluindo projetos no Brasil --- com informações sobre valores comprometidos e desembolsados, objetivos, setores e resultados reportados. Útil para avaliações de projetos financiados pelo Banco no Brasil.

% ────────────────────────────────────────────────────────────
\section{OCDE}
% ────────────────────────────────────────────────────────────

\ficharapida
  {Bases de Dados da OCDE}
  {Organização para a Cooperação e Desenvolvimento Econômico (OCDE)}
  {https://stats.oecd.org}
  {38 países membros + parceiros (inclui Brasil como parceiro)}
  {Anual (séries variáveis por indicador)}
  {País}

\subsection{O que é e relevância para o Brasil}

A OCDE é uma organização de países desenvolvidos, e o Brasil não é membro pleno --- mas participa como país parceiro e candidato à adesão (processo em curso desde 2022). Essa posição confere ao Brasil uma presença crescente nas bases de dados da OCDE, tornando-a uma fonte relevante para comparações com países de renda alta e para indicadores setoriais nos quais a OCDE tem metodologia superior.

\subsection{Principais bases temáticas}

\subsubsection{OECD Health Statistics}

A base de estatísticas de saúde da OCDE cobre os países membros e alguns parceiros (incluindo o Brasil) com indicadores como:
\begin{itemize}
  \item Gasto em saúde como proporção do PIB e per capita (em PPC);
  \item Número de médicos, enfermeiros e leitos hospitalares por habitante;
  \item Esperança de vida e anos de vida saudável;
  \item Mortalidade por causas específicas comparável entre países.
\end{itemize}

Para o Brasil, é especialmente útil para contextualizar o nível de gasto em saúde e a eficiência do sistema de saúde em perspectiva internacional.

\subsubsection{OECD Education at a Glance}

A publicação anual \textit{Education at a Glance} reúne indicadores educacionais comparados para os países da OCDE e parceiros:
\begin{itemize}
  \item Taxas de matrícula e conclusão por nível de ensino;
  \item Gasto público em educação como proporção do PIB e por aluno;
  \item Salários de professores em relação à renda média;
  \item Retornos à educação no mercado de trabalho.
\end{itemize}

\subsubsection{OECD Employment and Labour Market Statistics}

Dados comparados sobre mercado de trabalho, incluindo taxas de desemprego com definições harmonizadas pela OIT, participação laboral por sexo e faixa etária, e indicadores de qualidade do emprego.

\subsubsection{OECD Revenue Statistics}

Dados comparados sobre carga tributária por país, com decomposição por tipo de tributo (imposto sobre renda, consumo, propriedade, contribuições sociais). Fundamental para análises comparativas sobre o tamanho e a estrutura tributária do Brasil.

% ────────────────────────────────────────────────────────────
\section{PISA --- Programme for International Student Assessment}
% ────────────────────────────────────────────────────────────

\ficharapida
  {PISA --- Programme for International Student Assessment}
  {OCDE / Inep (coordenação nacional)}
  {https://www.oecd.org/pisa}
  {Mais de 80 países e economias participantes}
  {Trienal: 2000, 2003, 2006, 2009, 2012, 2015, 2018, 2022}
  {Estudante de 15 anos (amostra representativa nacional)}

\subsection{O que é e estrutura}

O PISA é uma avaliação internacional coordenada pela OCDE que mede o desempenho de estudantes de 15 anos em três domínios: Leitura, Matemática e Ciências. Cada edição tem um domínio principal --- com avaliação mais aprofundada --- e dois domínios secundários. O Brasil participa do PISA desde a primeira edição, em 2000, sendo um dos países em desenvolvimento com a série histórica mais longa.

O PISA aplica testes padronizados internacionalmente e questiona os estudantes sobre seu contexto socioeconômico, motivação e estratégias de aprendizagem. Também inclui questionários para diretores de escola e, em algumas edições, para professores e pais.

\subsection{Por que o PISA é especialmente relevante para o Brasil}

O Brasil apresenta, nas edições do PISA, um desempenho sistematicamente abaixo da média da OCDE nos três domínios avaliados --- embora com tendência de melhora ao longo das edições. Esse contraste com países de renda similar torna o PISA uma fonte essencial para:

\begin{itemize}
  \item \textbf{Diagnosticar a magnitude do deficit educacional} brasileiro em perspectiva internacional;
  \item \textbf{Identificar fatores associados ao desempenho}: o questionário socioeconômico do PISA permite análises de quais características de alunos, escolas e sistemas educacionais estão associadas a melhores resultados;
  \item \textbf{Acompanhar tendências}: a série desde 2000 permite analisar como o desempenho brasileiro evoluiu em relação a outros países ao longo de mais de duas décadas de reformas educacionais.
\end{itemize}

\subsection{Microdados e acesso}

O Inep disponibiliza os microdados do PISA para o Brasil no seu portal, com os resultados por estudante e os questionários contextuais. Os microdados internacionais completos (todos os países) estão disponíveis no site da OCDE. O R tem o pacote \texttt{PISA} que facilita o acesso e a análise dos microdados com o desenho amostral correto.

\begin{boxatencao}{Atenção --- PISA: representatividade nacional}
A amostra do PISA é desenhada para ser representativa da população de estudantes de 15 anos no nível nacional, não estadual ou municipal. Análises que tentem produzir estimativas estaduais a partir dos microdados do PISA para o Brasil estarão extrapolando além do que o desenho amostral permite. Para análises estaduais de desempenho educacional, o SAEB (especialmente a Prova Brasil) é a fonte adequada.
\end{boxatencao}

% ────────────────────────────────────────────────────────────
\section{Organização Internacional do Trabalho (OIT)}
% ────────────────────────────────────────────────────────────

\ficharapida
  {ILOSTAT --- Base de Dados da OIT}
  {Organização Internacional do Trabalho (OIT)}
  {https://ilostat.ilo.org}
  {Mais de 200 países}
  {Anual (séries variáveis)}
  {País}

\subsection{O que é e principais indicadores}

O ILOSTAT é a principal base de dados de trabalho e emprego da OIT, reunindo indicadores harmonizados para mais de 200 países. Para o Brasil, as principais séries disponíveis incluem:

\begin{itemize}
  \item Taxa de desemprego (definição OIT), participação na força de trabalho e emprego por setor;
  \item Trabalho informal: proporção de trabalhadores em emprego informal, por sexo e setor;
  \item Condições de trabalho: salário mínimo como proporção do salário mediano, horas trabalhadas;
  \item Trabalho infantil: incidência por faixa etária, sexo e tipo de atividade;
  \item Seguridade social: cobertura de proteção social por tipo de benefício.
\end{itemize}

O ILOSTAT é especialmente útil para comparações sobre trabalho informal e trabalho infantil, dimensões em que o Brasil tem indicadores que diferem dos países da OCDE e que são melhor contextualizados em comparação com países de renda média.

% ────────────────────────────────────────────────────────────
\section{Organização Mundial da Saúde (OMS)}
% ────────────────────────────────────────────────────────────

\ficharapida
  {Global Health Observatory (GHO)}
  {Organização Mundial da Saúde (OMS)}
  {https://www.who.int/data/gho}
  {Países membros da OMS (194)}
  {Anual (séries variáveis)}
  {País}

\subsection{O que é e principais indicadores}

O Global Health Observatory (GHO) é a plataforma de dados de saúde da OMS, reunindo indicadores para os países membros com metodologias harmonizadas. Para o Brasil, os indicadores mais relevantes incluem:

\begin{itemize}
  \item \textbf{Mortalidade e morbidade}: taxas de mortalidade por causa (padronizadas por idade), carga de doença em DALYs (anos de vida ajustados por incapacidade), esperança de vida saudável;
  \item \textbf{Doenças específicas}: tuberculose (incidência, prevalência, mortalidade), HIV/AIDS, malária, doenças não transmissíveis;
  \item \textbf{Cobertura de serviços de saúde}: cobertura universal de saúde (UHC index), vacinação, pré-natal;
  \item \textbf{Determinantes da saúde}: tabagismo, obesidade, saneamento, poluição do ar.
\end{itemize}

Uma vantagem específica da OMS sobre as fontes nacionais é que suas estimativas de mortalidade são corrigidas para subregistro e causas mal definidas --- o que as torna mais comparáveis entre países, mas também as afasta dos dados brutos do SIM. Para análises do Brasil especificamente, o SIM é a fonte primária; para comparações internacionais, o GHO é a referência.

% ────────────────────────────────────────────────────────────
\section{PNUD --- Programa das Nações Unidas para o Desenvolvimento}
% ────────────────────────────────────────────────────────────

\ficharapida
  {Human Development Data (HDR)}
  {Programa das Nações Unidas para o Desenvolvimento (PNUD)}
  {https://hdr.undp.org/data-center}
  {Países membros da ONU}
  {Anual}
  {País}

\subsection{IDH e índices derivados}

O PNUD é o produtor do \textbf{Índice de Desenvolvimento Humano (IDH)}, o indicador sintético de desenvolvimento mais conhecido mundialmente, que combina dimensões de saúde (esperança de vida), educação (anos de escolaridade esperados e médios) e renda (RNB per capita em PPC).

Além do IDH, o PNUD publica índices derivados relevantes para avaliação:
\begin{itemize}
  \item \textbf{IDHI --- IDH ajustado pela Desigualdade}: penaliza países com maior desigualdade nas três dimensões;
  \item \textbf{IDG --- Índice de Desigualdade de Gênero}: mede desigualdades entre homens e mulheres em saúde reprodutiva, empoderamento e mercado de trabalho;
  \item \textbf{IPM --- Índice de Pobreza Multidimensional}: mede pobreza além da renda, incorporando privações em saúde, educação e padrão de vida.
\end{itemize}

Para análises subnacionais, o PNUD Brasil publica o \textbf{IDHM} --- IDH Municipal --- calculado para todos os municípios brasileiros com base nos dados do Censo Demográfico. O IDHM é amplamente utilizado como variável de desenvolvimento municipal em avaliações de políticas públicas no Brasil.

% ────────────────────────────────────────────────────────────
\section{CEPAL --- Comissão Econômica para a América Latina e o Caribe}
% ────────────────────────────────────────────────────────────

\ficharapida
  {CEPALSTAT}
  {Comissão Econômica para a América Latina e o Caribe (CEPAL)}
  {https://statistics.cepal.org/portal/cepalstat}
  {América Latina e Caribe (33 países)}
  {Anual (séries variáveis)}
  {País}

\subsection{O que é e relevância}

O CEPALSTAT é o portal de dados estatísticos da CEPAL para a América Latina e o Caribe. Diferentemente das bases globais do Banco Mundial e da ONU, o CEPALSTAT foca especificamente na região, o que o torna especialmente útil para comparações do Brasil com seus vizinhos regionais.

As principais bases disponíveis incluem:

\begin{itemize}
  \item \textbf{BADEHOG --- Base de Datos de Hogares}: microdados harmonizados de pesquisas domiciliares de países latino-americanos, incluindo o Brasil. Permite comparações de pobreza, desigualdade e mercado de trabalho com metodologias padronizadas;
  \item \textbf{Indicadores macroeconômicos}: PIB, inflação, balança de pagamentos, dívida pública por país;
  \item \textbf{Indicadores sociais}: pobreza, desigualdade (Gini), gasto social por função (educação, saúde, proteção social, habitação);
  \item \textbf{Gênero}: indicadores de desigualdade de gênero no mercado de trabalho, violência contra a mulher, participação política.
\end{itemize}

\begin{boxexemplo}{Exemplo --- CEPALSTAT e desigualdade na América Latina}
O Brasil tem consistentemente um dos índices de Gini mais elevados da América Latina, que por sua vez é a região mais desigual do mundo. Para contextualizar essa posição e avaliar se as políticas de transferência de renda implementadas no país produziram convergência em relação à média regional, pesquisadores utilizaram os dados do CEPALSTAT sobre o índice de Gini para todos os países latino-americanos no período 2000--2020. A comparação mostrou que, embora o Gini brasileiro tenha caído substancialmente entre 2003 e 2015, o Brasil ainda permanecia entre os países mais desiguais da região ao final do período --- sugerindo que as melhorias, embora reais, foram insuficientes para alterar a posição relativa do país.

\textit{Fonte: FGV CLEAR, elaboração própria.}
\end{boxexemplo}

% ────────────────────────────────────────────────────────────
\section{Outras fontes internacionais relevantes}
% ────────────────────────────────────────────────────────────

\subsection{FAO --- Organização das Nações Unidas para Alimentação e Agricultura}

\ficharapida
  {FAOSTAT}
  {Organização das Nações Unidas para Alimentação e Agricultura (FAO)}
  {https://www.fao.org/faostat}
  {Mais de 245 países e territórios}
  {Anual (desde 1961 para produção agrícola)}
  {País}

O FAOSTAT é a principal base de dados sobre produção agrícola, pecuária, pesca, florestas e uso do solo em perspectiva comparada. Para o Brasil --- um dos maiores produtores agropecuários do mundo ---, o FAOSTAT é relevante para:
\begin{itemize}
  \item Comparações de produtividade agrícola e pecuária com outros países exportadores;
  \item Análise de segurança alimentar: subalimentação, disponibilidade calórica per capita;
  \item Dados de uso do solo e desmatamento em perspectiva global.
\end{itemize}

\subsection{Penn World Table (PWT)}

\ficharapida
  {Penn World Table}
  {Groningen Growth and Development Centre (University of Groningen) e University of Pennsylvania}
  {https://www.rug.nl/ggdc/productivity/pwt/}
  {Mais de 180 países}
  {Anual (desde 1950; atualização a cada 2--3 anos)}
  {País}

A Penn World Table é a base de referência para contas nacionais comparáveis entre países em termos reais, amplamente utilizada em macroeconomia e economia do desenvolvimento. Traz séries longas de PIB em paridade de poder de compra, capital, produtividade total dos fatores, trabalho e preços relativos, todas construídas com metodologia harmonizada que permite comparações robustas ao longo do tempo e entre países. Para avaliações de políticas brasileiras que exigem contexto macroeconômico comparado --- estudos sobre convergência, produtividade setorial, efeitos de choques externos ou retornos a fatores ---, a PWT é frequentemente a fonte primária.

\subsection{Human Mortality Database (HMD)}

\ficharapida
  {Human Mortality Database}
  {UC Berkeley e Max Planck Institute for Demographic Research}
  {https://www.mortality.org/}
  {41 países com sistemas estatísticos maduros}
  {Anual (séries longas, muitas iniciando no século XIX)}
  {País, sexo e idade (tábuas de vida detalhadas)}

A Human Mortality Database reúne tábuas de vida detalhadas por idade, sexo e ano-calendário para países com estatísticas vitais de alta qualidade. Os dados passam por harmonização rigorosa e são acompanhados de documentação metodológica por país. Para análises de mortalidade comparada, projeções demográficas e avaliação de políticas de longevidade, é a fonte de referência. O Brasil não está atualmente coberto pela HMD devido aos níveis históricos de subregistro, mas a base é útil como contraponto de alta qualidade para análises comparadas com países de renda alta.

\subsection{WHO Mortality Database}

\ficharapida
  {WHO Mortality Database}
  {Organização Mundial da Saúde (OMS)}
  {https://www.who.int/data/data-collection-tools/who-mortality-database}
  {Países membros da OMS que reportam óbitos codificados por CID}
  {Anual}
  {País, sexo, idade e causa (CID-9/CID-10)}

A WHO Mortality Database consolida os registros de óbito reportados à OMS pelos países membros, codificados por CID. Complementa a HMD ao cobrir um conjunto muito maior de países --- incluindo o Brasil, que reporta os dados do SIM --- embora com maior heterogeneidade de qualidade. É a principal fonte para análises comparadas de mortalidade por causa específica (homicídios, suicídios, doenças crônicas, mortalidade materna), permitindo contextualizar avaliações de políticas de saúde brasileiras em perspectiva internacional.

\subsection{Banco Interamericano de Desenvolvimento (BID)}

O BID disponibiliza dados sobre seus projetos de financiamento na América Latina e o Caribe, bem como bases de dados temáticas sobre infraestrutura, educação, saúde e governança na região. O portal \textit{IDB Data} (\url{data.iadb.org}) oferece acesso gratuito a indicadores regionais e dados de projetos.

\subsection{Varieties of Democracy (V-Dem)}

O projeto V-Dem, coordenado pela Universidade de Gotemburgo, produz indicadores de qualidade democrática e governança para mais de 200 países desde 1789. Inclui dimensões como liberdade de expressão, eleições livres e justas, Estado de Direito, participação civil e igualdade política. Para avaliações de políticas que dependem de contexto institucional e democrático, o V-Dem é uma fonte de referência internacional.

\subsection{Freedom House e The Economist Intelligence Unit}

O Freedom House publica anualmente o índice \textit{Freedom in the World}, que classifica países em ``livres'', ``parcialmente livres'' e ``não livres'' com base em indicadores de direitos políticos e liberdades civis. O EIU publica o \textit{Democracy Index}, com classificação em quatro tipos de regime (democracia plena, democracia imperfeita, regime híbrido e regime autoritário). Ambos são usados em estudos comparativos sobre a relação entre qualidade democrática e indicadores de bem-estar.

\subsection{Notre Dame Global Adaptation Initiative (ND-GAIN)}

O ND-GAIN produz um índice de vulnerabilidade e prontidão para adaptação às mudanças climáticas para mais de 180 países. Para avaliações de políticas de adaptação climática e de resiliência a desastres naturais, é uma referência internacional útil para contextualizar a posição do Brasil.

% ────────────────────────────────────────────────────────────
\section{Como usar fontes internacionais com rigor metodológico}
% ────────────────────────────────────────────────────────────

O uso de fontes internacionais em avaliações de políticas públicas exige atenção a um conjunto específico de questões metodológicas:

\subsection{Verificar a fonte original dos dados para o Brasil}

Como mencionado na introdução, a maioria das bases internacionais usa dados nacionais como insumo. Antes de usar um indicador internacional para o Brasil, verifique qual fonte nacional o organismo internacional utilizou. Se a fonte original for a PNAD Contínua ou o Censo Demográfico, pode ser mais eficiente trabalhar diretamente com essas fontes nacionais para análises focadas no Brasil.

\subsection{Atentar para diferenças de definição}

Conceitos aparentemente idênticos podem ter definições operacionais diferentes entre países e organismos internacionais. A ``taxa de desemprego'' da OIT difere da taxa calculada pelo IBGE em alguns aspectos metodológicos. A ``linha de pobreza'' do Banco Mundial (US\$ 2,15/dia em PPC 2017) é diferente das linhas nacionais utilizadas pelo IBGE e pelo governo brasileiro. Sempre que usar um indicador internacional para comparações, verifique se as definições são comparáveis.

\subsection{Cuidado com a defasagem temporal}

Bases internacionais frequentemente publicam dados com um ou dois anos de defasagem em relação às fontes nacionais originais. Para análises que requerem dados recentes, as fontes nacionais são geralmente mais atualizadas.

\subsection{Usar fontes internacionais para o que elas fazem de melhor}

Fontes internacionais são superiores às nacionais em dois cenários específicos: (1) quando a comparação com outros países é o objetivo da análise; (2) quando o organismo internacional aplicou ajustes metodológicos que tornam o dado mais comparável ou mais correto do que o dado bruto nacional (como as estimativas de mortalidade corrigidas para subregistro da OMS). Fora desses cenários, as fontes nacionais são geralmente preferíveis.

% ────────────────────────────────────────────────────────────
\section{Síntese do capítulo}
% ────────────────────────────────────────────────────────────

Este capítulo apresentou as principais fontes internacionais e comparadas disponíveis para avaliadores de políticas públicas. A Tabela~\ref{tab:sintese_cap11} resume suas características essenciais.

\begin{table}[H]
\centering
\caption{Síntese das principais fontes internacionais}
\label{tab:sintese_cap11}
\small
\begin{tabularx}{\textwidth}{@{} >{\raggedright\arraybackslash}p{3cm} >{\raggedright\arraybackslash}p{2.3cm} >{\raggedright\arraybackslash}p{2.3cm} X @{}}
\toprule
\textbf{Fonte} & \textbf{Produtor} & \textbf{Cobertura} &
\textbf{Principal uso} \\
\midrule
WDI               & Banco Mundial        & 200+ países        & Benchmarking geral de desenvolvimento \\
PIP               & Banco Mundial        & 100+ países        & Pobreza e desigualdade comparadas \\
OECD Stats        & OCDE                 & 38+ países         & Saúde, educação, trabalho, tributação \\
PISA              & OCDE/Inep            & 80+ países         & Desempenho educacional comparado \\
ILOSTAT           & OIT                  & 200+ países        & Trabalho, informalidade, trabalho infantil \\
GHO               & OMS                  & 194 países         & Saúde, mortalidade, doenças \\
HDR/IDHM          & PNUD                 & 190+ países        & IDH, pobreza multidimensional \\
CEPALSTAT/\newline BADEHOG & CEPAL                & 33 países AL       & Comparações regionais, microdados \\
FAOSTAT           & FAO                  & 245+ países        & Agropecuária, segurança alimentar \\
V-Dem             & Univ.\ Gotemburgo    & 200+ países        & Qualidade democrática, governança \\
\bottomrule
\end{tabularx}
\fonte{FGV CLEAR, elaboração própria.}
\end{table}

Os pontos centrais do capítulo são:

\begin{itemize}
  \item Fontes internacionais são superiores para \textbf{benchmarking} e \textbf{comparações cross-country}; para análises focadas no Brasil, as fontes nacionais são geralmente mais precisas e atualizadas;
  \item A maior parte das bases internacionais usa \textbf{dados nacionais como insumo} --- verificar qual fonte original foi utilizada é o primeiro passo antes de usar qualquer indicador internacional;
  \item O \textbf{WDI} (Banco Mundial) e o \textbf{CEPALSTAT} (CEPAL) são as referências para comparações macroeconômicas e sociais gerais; o \textbf{PISA} (OCDE) é a referência para comparações educacionais; o \textbf{ILOSTAT} (OIT) para trabalho; o \textbf{GHO} (OMS) para saúde;
  \item O \textbf{IDHM} do PNUD é amplamente utilizado como variável de desenvolvimento municipal em avaliações brasileiras, sendo derivado do Censo Demográfico com metodologia própria;
  \item A \textbf{harmonização metodológica} necessária para comparações internacionais implica perdas de detalhe que precisam ser consideradas; diferenças de definição entre organismos internacionais são comuns e devem ser verificadas antes de qualquer comparação.
\end{itemize}

\medskip

O próximo capítulo apresenta um conjunto emergente de fontes e ferramentas que complementam as bases institucionais: APIs governamentais, \textit{web scraping} e inteligência artificial. Não substituem as fontes tradicionais --- preenchem lacunas, aceleram a coleta e abrem possibilidades analíticas novas que seriam impraticáveis sem elas. O capítulo discute como aproveitá-las com rigor metodológico.
% ============================================================
%  cap12_ia_apis_scraping.tex
%  Capítulo 12: Dados com Apoio de Tecnologia
% ============================================================

% ════════════════════════════════════════════════════════════
\chapter{Dados com Apoio de Tecnologia: APIs, \textit{Web Scraping} e Inteligência Artificial}
% ════════════════════════════════════════════════════════════

\begin{boxdestaque}{Neste capítulo}
Este capítulo apresenta uma categoria crescente e ainda pouco sistematizada de fontes de dados: aquelas obtidas ou estruturadas com apoio de tecnologias digitais --- interfaces de programação (APIs), coleta automatizada de páginas web (\textit{web scraping}) e ferramentas de inteligência artificial (IA). Para cada abordagem, discutimos o que é, quando usar, como usar com rigor metodológico e quais são os riscos e limitações. O capítulo termina com orientações práticas para incorporar essas ferramentas de forma responsável na prática avaliativa.
\end{boxdestaque}

% ────────────────────────────────────────────────────────────
\section{Introdução: por que tecnologia digital importa para dados de políticas públicas}
% ────────────────────────────────────────────────────────────

Por décadas, o acesso a dados para avaliação de políticas públicas dependia de um conjunto relativamente estável de fontes institucionais: pesquisas do IBGE, registros administrativos de ministérios, censos do Inep. A obtenção de dados era um processo manual e frequentemente lento --- download de arquivos, preenchimento de formulários, espera por convênios.

Esse panorama mudou substancialmente na última década. Três desenvolvimentos tecnológicos transformaram o ecossistema de dados disponíveis para avaliadores:

\begin{enumerate}
  \item \textbf{Abertura de APIs governamentais}: portais como \url{dados.gov.br}, as APIs do IBGE e do Banco Central democratizaram o acesso a dados que antes exigiam processos manuais demorados, tornando possível a coleta automatizada e atualizada de dezenas de indicadores;

  \item \textbf{Maturação do \textit{web scraping}}: a consolidação de ferramentas de raspagem de dados permitiu extrair informações de portais governamentais, diários oficiais, sistemas de licitação e outras fontes que disponibilizam dados em formato web mas não oferecem download estruturado;

  \item \textbf{Inteligência artificial}: o avanço dos modelos de linguagem de grande escala (LLMs) e de outras ferramentas de IA abriu possibilidades novas de organizar e estruturar grandes volumes de texto não estruturado --- atas de reuniões, processos administrativos, relatórios de auditoria --- transformando-os em dados analisáveis.
\end{enumerate}

Essas ferramentas não substituem as fontes tradicionais. Elas as complementam, preenchem lacunas e, em alguns casos, criam possibilidades analíticas genuinamente novas. Mas exigem competências técnicas específicas e cuidados metodológicos que este capítulo busca sistematizar.

% ────────────────────────────────────────────────────────────
\section{APIs governamentais abertas}
% ────────────────────────────────────────────────────────────

\subsection{O que é uma API e por que importa}

Uma API (\textit{Application Programming Interface}) é uma interface que permite que programas de computador se comuniquem entre si de forma padronizada. No contexto de dados públicos, uma API governamental permite que o pesquisador solicite dados diretamente de um sistema governamental por meio de código --- sem necessidade de navegar manualmente em portais, clicar em botões de download ou aguardar o envio de arquivos por e-mail.

As vantagens práticas são significativas:

\begin{itemize}
  \item \textbf{Automação}: uma vez escrito o código de coleta, ele pode ser executado repetidamente para atualizar os dados sem intervenção manual;
  \item \textbf{Atualidade}: APIs costumam fornecer dados mais recentes do que os arquivos disponíveis para download, pois refletem o estado atual do banco de dados do provedor;
  \item \textbf{Eficiência}: é possível solicitar apenas os dados necessários (por município, período ou indicador), evitando o download de arquivos grandes para extrair poucas informações;
  \item \textbf{Reprodutibilidade}: o código de coleta é documentável e reproduzível, o que facilita a verificação e a replicação da análise.
\end{itemize}

\subsection{Principais APIs governamentais abertas no Brasil}

\subsubsection{API do IBGE (SIDRA e outras)}

O IBGE disponibiliza diversas APIs para acesso programático a seus dados:

\begin{itemize}
  \item \textbf{SIDRA API}: acesso aos dados do Sistema IBGE de Recuperação Automática, que reúne tabelas de todas as pesquisas do IBGE (Censo, PNAD, POF, Censo Agropecuário, entre outras). A API SIDRA permite solicitar tabelas específicas com parâmetros de período, nível geográfico e variável;

  \item \textbf{API de Malhas}: acesso às malhas territoriais do IBGE em formato GeoJSON, com parâmetros de nível territorial e ano de referência;

  \item \textbf{API de Localidades}: acesso à lista de municípios, estados e regiões com seus códigos e metadados.
\end{itemize}

Em R, o pacote \texttt{sidrar} encapsula a API SIDRA em funções de alto nível, e o pacote \texttt{geobr} faz o mesmo para as malhas territoriais. Em Python, as bibliotecas \texttt{sidrapy} e \texttt{pygeobr} oferecem funcionalidades equivalentes.

\subsubsection{API do Banco Central (SGS)}

O Banco Central disponibiliza o Sistema Gerenciador de Séries Temporais (SGS) via API REST, com acesso a mais de 20.000 séries macroeconômicas --- IPCA, taxa Selic, câmbio, crédito, índices de atividade econômica, entre outras.

\begin{itemize}
  \item \textbf{Endpoint}: \nolinkurl{api.bcb.gov.br/dados/serie/bcdata.sgs.{codigo}/dados}
  \item Em R, o pacote \texttt{rbcb} facilita o acesso; em Python, a biblioteca \texttt{python-bcb} oferece funcionalidade equivalente.
\end{itemize}

\subsubsection{Portal de Dados Abertos do Governo Federal}

O portal \url{dados.gov.br} é o catálogo central de dados abertos do governo federal brasileiro, reunindo datasets de mais de 200 órgãos. Não é uma API unificada, mas um catálogo de fontes --- cada dataset tem seu próprio formato e método de acesso (CSV para download direto, APIs específicas de cada órgão, etc.).

Para o avaliador, o \url{dados.gov.br} é o ponto de partida para descobrir quais dados estão disponíveis em formato aberto. A qualidade e a atualidade dos datasets variam enormemente entre órgãos.

\subsubsection{API do Portal da Transparência (CGU)}

A Controladoria-Geral da União disponibiliza uma API REST para acesso programático aos dados do Portal da Transparência:

\begin{itemize}
  \item Despesas por órgão, programa e natureza;
  \item Transferências a estados e municípios;
  \item Contratos e licitações;
  \item Beneficiários de programas sociais (Bolsa Família, BPC).
\end{itemize}

A documentação está disponível em \url{api.portaldatransparencia.gov.br}. O acesso requer cadastro gratuito para obtenção de chave de API (\textit{token}).

\subsubsection{API do DATASUS}

O DATASUS não disponibiliza uma API REST padrão para seus sistemas (SIM, SINASC, SIH, etc.), mas o portal TabNet (\url{tabnet.datasus.gov.br}) permite consultas estruturadas com parâmetros de período, local e variável. Para acesso programático, o pacote \texttt{microdatasus} em R e a plataforma \textit{Base dos Dados} são as principais alternativas.

\begin{boxdica}{Dica --- Base dos Dados: dados integrados}
A plataforma \textit{Base dos Dados} (\url{basedosdados.org}) é um projeto brasileiro de infraestrutura de dados que disponibiliza dezenas de bases públicas --- PNAD Contínua, Censo Escolar, RAIS, CAGED, SIM, SINASC, FINBRA e muitas outras --- já tratadas, compatibilizadas e disponíveis em BigQuery (Google Cloud). O acesso é gratuito para uso não comercial, com pacotes em R (\texttt{basedosdados}) e Python que permitem consultas SQL diretamente do código analítico. Para o avaliador que trabalha com múltiplas bases, a \textit{Base dos Dados} reduz substancialmente o tempo de preparação dos dados.
\end{boxdica}

\subsection{Boas práticas no uso de APIs}

\begin{itemize}
  \item \textbf{Documente a data de coleta}: dados de APIs são dinâmicos e podem mudar. Sempre registre a data em que os dados foram coletados e, quando possível, salve uma cópia local dos dados brutos obtidos;
  \item \textbf{Respeite os limites de requisição}: APIs governamentais geralmente impõem limites no número de requisições por minuto ou por hora (\textit{rate limits}). Implemente pausas entre requisições para evitar bloqueios;
  \item \textbf{Trate erros}: conexões de rede falham, servidores ficam indisponíveis, respostas chegam malformadas. Implemente tratamento de erros no código de coleta para lidar com essas situações;
  \item \textbf{Verifique a consistência}: compare uma amostra dos dados obtidos via API com os dados disponíveis no portal de origem para verificar se não há discrepâncias.
\end{itemize}

% ────────────────────────────────────────────────────────────
\section{\textit{Web Scraping} para dados não estruturados}
% ────────────────────────────────────────────────────────────

\subsection{O que é e quando usar}

\textit{Web scraping} é a extração automatizada de informações de páginas web. Enquanto APIs fornecem dados em formato estruturado (JSON, XML, CSV), o \textit{web scraping} extrai informações de páginas HTML --- que foram projetadas para leitura humana, não para processamento por máquinas.

O uso de \textit{web scraping} em pesquisa sobre políticas públicas é justificado quando:

\begin{itemize}
  \item Os dados de interesse estão disponíveis em portais web mas não há API ou download estruturado;
  \item Os dados são publicados em formato de tabela HTML ou em documentos PDF embarcados em páginas web;
  \item A coleta manual seria inviável pelo volume de páginas ou pela frequência necessária de atualização.
\end{itemize}

Exemplos de fontes relevantes para avaliação de políticas públicas que frequentemente requerem \textit{scraping}:

\begin{itemize}
  \item \textbf{Diários oficiais municipais e estaduais}: publicações de atos administrativos, contratos, nomeações e exonerações que não estão disponíveis em formato estruturado;
  \item \textbf{Portais de licitações e contratos}: sistemas como o Comprasnet e portais estaduais de compras públicas;
  \item \textbf{Portais de transparência municipal}: muitos municípios publicam execução orçamentária em formato HTML sem opção de download;
  \item \textbf{Atas de câmaras municipais e conselhos}: documentos textuais com informações sobre deliberações e votações.
\end{itemize}

\subsection{Ferramentas principais}

\begin{itemize}
  \item \textbf{Python}: as bibliotecas \texttt{BeautifulSoup} e \texttt{Scrapy} são as mais utilizadas para extração de HTML; \texttt{Selenium} e \texttt{Playwright} permitem automatizar navegadores para lidar com páginas que carregam conteúdo dinamicamente via JavaScript;
  \item \textbf{R}: o pacote \texttt{rvest} oferece funcionalidades de \textit{scraping} integradas ao ecossistema do \texttt{tidyverse};
  \item \textbf{Extração de PDF}: para documentos em PDF, as bibliotecas \texttt{pdfplumber} (Python) e \texttt{pdftools} (R) extraem texto e tabelas com qualidade variável dependendo do tipo de PDF (nativo vs.\ escaneado).
\end{itemize}

\subsection{Limites éticos, legais e metodológicos}

O uso de \textit{web scraping} levanta questões que o pesquisador deve considerar antes de iniciar a coleta:

\begin{itemize}
  \item \textbf{Termos de uso}: alguns sites proíbem explicitamente o \textit{scraping} em seus termos de uso. Para dados governamentais brasileiros, a Lei de Acesso à Informação (LAI) garante o direito ao acesso, mas não o direito à coleta automatizada --- verifique se o órgão tem políticas específicas;

  \item \textbf{Robustez da coleta}: páginas web mudam frequentemente. Uma mudança no HTML do portal pode quebrar o código de \textit{scraping} sem aviso. Scripts de coleta precisam de manutenção regular e verificação periódica de consistência;

  \item \textbf{Qualidade dos dados extraídos}: a extração automatizada de HTML e PDF é imperfeita. Tabelas com células mescladas, PDFs escaneados e layouts complexos produzem erros de extração que precisam ser identificados e corrigidos na etapa de limpeza;

  \item \textbf{Carga nos servidores}: coletar dados em alta velocidade pode sobrecarregar servidores de órgãos públicos com capacidade limitada. Implemente pausas entre requisições (\textit{rate limiting}) e evite coletar os mesmos dados repetidamente.
\end{itemize}

\begin{boxexemplo}{Exemplo --- \textit{Scraping} de diários oficiais}
Pesquisadores interessados em avaliar a qualidade da gestão pública municipal coletaram, via \textit{web scraping}, todas as publicações dos diários oficiais de 200 municípios ao longo de quatro anos. O objetivo era construir um indicador de transparência baseado na regularidade e na completude das publicações obrigatórias (contratos, licitações, prestações de contas). O código em Python usou \texttt{Scrapy} para navegar pelos portais de cada município e \texttt{pdfplumber} para extrair texto dos documentos publicados. O resultado foi um painel longitudinal de indicadores de transparência que permitiu avaliar se políticas de apoio à gestão municipal melhoraram a regularidade das publicações nos municípios beneficiados.

\textit{Fonte: FGV CLEAR, elaboração própria.}
\end{boxexemplo}

% ────────────────────────────────────────────────────────────
\section{Inteligência Artificial como ferramenta de apoio à pesquisa}
% ────────────────────────────────────────────────────────────

\subsection{O novo papel da IA na pesquisa em políticas públicas}

O avanço dos modelos de linguagem de grande escala (LLMs --- \textit{Large Language Models}) como o GPT-4, o Claude e outros abriu possibilidades novas para pesquisadores de políticas públicas. Diferentemente das ferramentas discutidas anteriormente --- que organizam dados já existentes de forma mais eficiente ---, os LLMs têm capacidade de \textit{processar, resumir e estruturar} grandes volumes de texto não estruturado, transformando fontes que antes eram de difícil análise quantitativa em dados analisáveis.

É importante, no entanto, ser preciso sobre o que a IA faz de melhor e onde seus riscos são mais sérios. Esta seção busca oferecer uma visão equilibrada --- nem excessivamente entusiasmada, nem excessivamente cética.

\subsection{Usos bem estabelecidos e de baixo risco}

\subsubsection{Busca e descoberta de fontes}

LLMs podem ser utilizados como ferramentas de busca avançada para identificar fontes de dados relevantes, entender a estrutura de bases desconhecidas e localizar documentação metodológica. Para um avaliador que está mapeando o ecossistema de dados disponíveis para uma nova área temática, uma conversa com um LLM pode acelerar substancialmente o processo de descoberta.

O risco aqui é a \textbf{desatualização}: LLMs têm uma data de corte de treinamento e podem não conhecer bases de dados lançadas recentemente. Sempre verifique as informações obtidas diretamente nas fontes primárias.

\subsubsection{Apoio à limpeza e estruturação de dados}

LLMs --- e modelos de linguagem mais especializados --- podem ser usados para:

\begin{itemize}
  \item \textbf{Classificação de texto}: categorizar documentos, respostas abertas de surveys ou registros administrativos em categorias predefinidas. Por exemplo, classificar descrições de objetos de contratos públicos por setor de atividade, ou categorizar motivos de desligamento em dados de recursos humanos;
  \item \textbf{Extração de entidades}: identificar e extrair informações específicas de textos --- nomes de municípios, valores monetários, datas, nomes de programas --- em documentos não estruturados como atas e relatórios;
  \item \textbf{Padronização}: normalizar variações ortográficas de nomes de municípios, categorias e outras variáveis de texto que impedem o \textit{merge} entre bases.
\end{itemize}

Esses usos são mais robustos do que a geração de informações novas, porque o modelo está operando sobre textos fornecidos pelo pesquisador --- não gerando conteúdo a partir de seu treinamento.

\subsubsection{Revisão de literatura e síntese de evidências}

Ferramentas especializadas em revisão de literatura --- como \textit{Elicit} (\url{elicit.org}), \textit{Consensus} (\url{consensus.app}) e \textit{Semantic Scholar} --- usam modelos de IA para ajudar pesquisadores a:

\begin{itemize}
  \item Identificar estudos relevantes em uma área temática a partir de uma pergunta de pesquisa em linguagem natural;
  \item Extrair automaticamente informações estruturadas de artigos (tamanho da amostra, métodos, resultados principais);
  \item Mapear o estado da evidência sobre uma intervenção específica.
\end{itemize}

Essas ferramentas são úteis como ponto de partida para revisões sistemáticas, mas não substituem a leitura cuidadosa dos artigos identificados nem a avaliação crítica da qualidade metodológica dos estudos.

\subsection{Usos de maior risco: geração de dados e informações}

\subsubsection{Geração de dados sintéticos}

LLMs podem ser usados para gerar dados sintéticos --- dados artificiais com características estatísticas similares às de dados reais --- para fins de teste de código, treinamento de modelos de \textit{machine learning} ou preenchimento de lacunas em bases incompletas. No entanto, o uso de dados sintéticos em avaliações de políticas públicas exige cautela extrema:

\begin{itemize}
  \item Dados sintéticos gerados por LLMs podem reproduzir vieses presentes nos dados de treinamento do modelo;
  \item A semelhança estatística com os dados reais não garante que os dados sintéticos capturem as relações causais relevantes para a análise;
  \item Em análises para tomada de decisão de políticas públicas, o uso de dados sintéticos deve ser sempre claramente declarado e justificado.
\end{itemize}

\subsubsection{Respostas factuais sobre dados e políticas}

LLMs cometem erros factuais --- fenômeno conhecido como \textbf{alucinação} --- com frequência não trivial, especialmente em domínios especializados. Valores de indicadores, datas de implementação de políticas, nomes de programas e detalhes metodológicos de pesquisas são tipos de informação que LLMs podem apresentar com aparente confiança mas incorretamente.

\begin{boxatencao}{Atenção --- IA para organizar, não para certificar}
O uso de LLMs na pesquisa em políticas públicas deve seguir um princípio simples: \textbf{use IA para organizar, nunca para certificar}. Um LLM pode ajudá-lo a encontrar uma fonte, a entender a estrutura de uma base de dados ou a classificar um conjunto de textos --- mas qualquer informação factual gerada por um LLM deve ser verificada na fonte primária antes de ser usada em uma análise. Isso vale especialmente para valores numéricos, datas, nomes de programas e detalhes metodológicos. A responsabilidade pela qualidade dos dados usados em uma avaliação é sempre do pesquisador, não da ferramenta.
\end{boxatencao}

\subsection{Processamento de linguagem natural em fontes textuais}

Além dos LLMs de propósito geral, há uma classe de aplicações de Processamento de Linguagem Natural (PLN) especificamente desenvolvidas para análise de textos de políticas públicas. Essas aplicações podem transformar fontes textuais não estruturadas em dados analisáveis:

\begin{itemize}
  \item \textbf{Análise de sentimento em avaliações de programas}: classificar automaticamente respostas abertas de pesquisas de satisfação como positivas, negativas ou neutras;
  \item \textbf{Análise temática de documentos}: identificar os temas dominantes em grandes coleções de documentos --- relatórios de auditoria, atas de conselhos, processos de licenciamento;
  \item \textbf{Reconhecimento de entidades em processos administrativos}: extrair informações estruturadas (partes, valores, datas, objetos) de processos jurídicos e administrativos;
  \item \textbf{Análise de legislação}: identificar mudanças em políticas a partir de coleções de textos legais ao longo do tempo.
\end{itemize}

Para essas aplicações mais sofisticadas, modelos de linguagem especializados em português --- como o BERTimbau, treinado especificamente em textos em português brasileiro --- costumam superar modelos gerais em termos de precisão.

% ────────────────────────────────────────────────────────────
\section{Combinando APIs, \textit{scraping} e IA: um fluxo de trabalho integrado}
% ────────────────────────────────────────────────────────────

As três abordagens discutidas neste capítulo não são alternativas entre si --- são complementares e, usadas em conjunto, podem produzir fontes de dados que nenhuma delas criaria sozinha. Um fluxo de trabalho típico para coleta de dados de políticas públicas com apoio tecnológico pode ter a seguinte estrutura:

\begin{enumerate}
  \item \textbf{Descoberta}: usar LLMs e buscas avançadas para mapear quais fontes existem, quais têm API, quais requerem \textit{scraping} e quais exigem acesso manual;

  \item \textbf{Coleta via API}: para fontes com APIs abertas (IBGE, Banco Central, Portal da Transparência), escrever código de coleta automatizada com tratamento de erros e documentação da data de coleta;

  \item \textbf{Coleta via \textit{scraping}}: para fontes sem API, desenvolver scripts de raspagem com verificação periódica de consistência;

  \item \textbf{Extração de texto}: para documentos em PDF e HTML, usar ferramentas de extração com validação manual de uma amostra;

  \item \textbf{Estruturação com IA}: para textos não estruturados, usar modelos de PLN ou LLMs para classificação, extração de entidades e padronização --- sempre com validação manual de uma amostra dos resultados;

  \item \textbf{Integração e limpeza}: combinar as fontes coletadas, resolver inconsistências e documentar o processo de forma reproduzível;

  \item \textbf{Validação}: comparar indicadores derivados das fontes coletadas com indicadores independentes para verificar a coerência dos dados.
\end{enumerate}

% ────────────────────────────────────────────────────────────
\section{Considerações éticas e de reprodutibilidade}
% ────────────────────────────────────────────────────────────

O uso de tecnologias digitais na coleta e estruturação de dados para políticas públicas levanta questões éticas e de reprodutibilidade que merecem atenção explícita:

\subsection{Transparência sobre o processo de coleta}

Análises que usam dados coletados via API, \textit{scraping} ou processados por IA devem documentar claramente:

\begin{itemize}
  \item Quais fontes foram acessadas, por qual método e em que data;
  \item Quais transformações foram aplicadas (especialmente quando IA foi usada para classificação ou extração);
  \item Como a qualidade dos dados foi verificada.
\end{itemize}

\subsection{Reprodutibilidade}

O código de coleta e processamento deve ser disponibilizado junto com os resultados da análise, para que outros pesquisadores possam verificar e replicar o processo. Repositórios como GitHub são o padrão para compartilhamento de código; o Open Science Framework (OSF) permite registrar o plano de análise antes da coleta dos dados.

\subsection{Privacidade e LGPD}

A Lei Geral de Proteção de Dados (LGPD) impõe restrições ao tratamento de dados pessoais, mesmo quando esses dados são coletados de fontes públicas. Dados raspados de portais governamentais que contenham informações identificáveis de pessoas físicas --- como nomes em contratos ou beneficiários de programas --- devem ser tratados com os cuidados exigidos pela LGPD.

% ────────────────────────────────────────────────────────────
\section{Síntese do capítulo}
% ────────────────────────────────────────────────────────────

Este capítulo apresentou as principais abordagens tecnológicas para coleta e estruturação de dados para avaliação de políticas públicas. Os pontos centrais são:

\begin{itemize}
  \item \textbf{APIs governamentais} democratizaram o acesso a dados públicos no Brasil: o IBGE, o Banco Central, o Portal da Transparência e a \textit{Base dos Dados} oferecem acesso programático a dezenas de bases relevantes para avaliação;

  \item \textbf{\textit{Web scraping}} é a ferramenta adequada para fontes que disponibilizam dados em formato web sem API estruturada --- diários oficiais, portais de licitação, sistemas de transparência municipal. Exige cuidados com robustez do código, qualidade da extração e questões éticas e legais;

  \item \textbf{Inteligência artificial} tem usos bem estabelecidos em classificação de texto, extração de entidades e síntese de literatura, mas apresenta riscos sérios quando usada para gerar informações factuais ou dados novos. O princípio fundamental é usar IA para organizar, nunca para certificar;

  \item \textbf{Validação e documentação} são obrigatórias em qualquer processo de coleta tecnológica: sempre valide uma amostra dos dados coletados e documente o processo de forma reproduzível;

  \item Essas ferramentas não substituem as fontes tradicionais discutidas nos capítulos anteriores --- elas as complementam, preenchem lacunas e criam possibilidades analíticas novas quando usadas com rigor metodológico.
\end{itemize}

\medskip

Todas essas ferramentas e fontes --- tradicionais e emergentes, nacionais e internacionais --- convergem para um momento concreto: o de escolher um dado específico, obtê-lo, combiná-lo com outras bases e usá-lo para responder uma pergunta avaliativa real. O capítulo final reúne os fios condutores do guia em um roteiro prático para esse processo --- da pergunta à fonte, do acesso à documentação, da linkagem à conformidade com a LGPD.
% ============================================================
%  cap13_como_escolher_acessar.tex  [REVISADO]
%  Capítulo 13: Como Escolher, Acessar e Documentar
%               Fontes de Dados
% ============================================================

% ════════════════════════════════════════════════════════════
\chapter{Como Escolher, Acessar e Documentar Fontes de Dados}
% ════════════════════════════════════════════════════════════

\begin{boxdestaque}{Neste capítulo}
Este capítulo final reúne os fios condutores do guia em um conjunto de orientações práticas para o trabalho cotidiano com fontes de dados. Discutimos como selecionar a fonte certa para cada pergunta avaliativa, como navegar os diferentes mecanismos de acesso --- desde dados abertos até pedidos via Lei de Acesso à Informação e convênios institucionais ---, como documentar o uso de fontes para garantir reprodutibilidade, como combinar bases por meio de integração e quais são as obrigações e cuidados decorrentes da Lei Geral de Proteção de Dados (LGPD). O capítulo inclui também orientações específicas para o uso de bases obtidas por LAI ou convênio e uma síntese da distinção entre bases amostrais, censitárias e administrativas.
\end{boxdestaque}

% ────────────────────────────────────────────────────────────
\section{Da pergunta à fonte: um processo em cinco etapas}
% ────────────────────────────────────────────────────────────

Ao longo deste guia, apresentamos dezenas de fontes de dados organizadas por tema. O risco de um catálogo tão extenso é que o avaliador se perca na abundância de opções e perca de vista o que realmente importa: a pergunta avaliativa. Este capítulo propõe uma sequência de cinco etapas para ir da pergunta à fonte de forma sistemática:

\begin{enumerate}
  \item \textbf{Definir a pergunta com precisão} --- o quê, sobre quem, em que território e em que período.
  \item \textbf{Identificar a unidade e o recorte} da análise (aluno, escola, município, família etc.) e os filtros geográfico e temporal correspondentes.
  \item \textbf{Mapear as fontes candidatas}, percorrendo os tipos apresentados no guia (registros administrativos, pesquisas amostrais, censos, sistemas setoriais, dados fiscais, dados territoriais).
  \item \textbf{Aplicar o \textit{checklist} de adequação} de cada fonte candidata (os oito critérios do Capítulo~1).
  \item \textbf{Triangular quando necessário}, combinando fontes complementares para resultado, controles, denominadores e dados de implementação.
\end{enumerate}

As próximas subseções detalham cada etapa, com exemplos.

\subsection{Etapa 1 --- Defina a pergunta com precisão}

Antes de buscar qualquer dado, a pergunta avaliativa precisa ser formulada com precisão suficiente para orientar a busca. Uma pergunta vaga como ``qual é o impacto do programa X?'' não é operacionalizável. Uma pergunta precisa como ``o programa X reduziu a taxa de abandono escolar entre alunos do 6º ao 9º ano do ensino fundamental em municípios do Nordeste entre 2018 e 2022?'' define imediatamente:

\begin{itemize}
  \item \textbf{O indicador de resultado}: taxa de abandono escolar;
  \item \textbf{A população}: alunos do 6º ao 9º ano;
  \item \textbf{O recorte geográfico}: municípios do Nordeste;
  \item \textbf{O período}: 2018 a 2022.
\end{itemize}

Cada um desses elementos é um critério de seleção de fonte.

\subsection{Etapa 2 --- Identifique a unidade e o recorte}

Com a pergunta delimitada, defina a unidade de análise necessária (aluno, escola, município, família) e o recorte geográfico e temporal. Esses dois elementos são os filtros mais eficientes para eliminar fontes inadequadas:

\begin{itemize}
  \item Se a unidade é o aluno individual, pesquisas amostrais sem identificadores individuais são insuficientes; é preciso microdados do Censo Escolar ou equivalente;
  \item Se o recorte geográfico é municipal, fontes representativas apenas para UF ou Brasil não servem;
  \item Se o período é 2020, o Censo Demográfico 2010 pode estar muito defasado para a linha de base.
\end{itemize}

\subsection{Etapa 3 --- Mapeie as fontes candidatas}

Com os critérios de unidade, recorte e período definidos, percorra mentalmente (ou neste guia) as categorias de fontes potencialmente relevantes:

\begin{enumerate}
  \item Há algum registro administrativo que cobre exatamente essa população?
  \item Há alguma pesquisa amostral com representatividade para o recorte necessário?
  \item Há dados de avaliação de desempenho relevantes?
  \item Há dados fiscais ou de infraestrutura que precisam ser controlados?
  \item É necessário dado georreferenciado?
\end{enumerate}

\subsection{Etapa 4 --- Aplique o checklist de adequação}

Para cada fonte candidata, aplique o checklist apresentado no Capítulo~1:

\begin{enumerate}
  \item A cobertura geográfica é compatível com o recorte?
  \item A periodicidade permite capturar o período relevante (antes e depois)?
  \item A unidade de análise é a mesma da pergunta?
  \item É possível desagregar pelos grupos populacionais de interesse?
  \item A amostra (se for pesquisa amostral) é representativa para os domínios necessários?
  \item Os microdados estão acessíveis, ou há restrições a contornar?
  \item A documentação permite verificar o que cada variável mede de fato?
  \item Há quebras metodológicas relevantes na série histórica?
\end{enumerate}

\subsection{Etapa 5 --- Triangule quando necessário}

Raramente uma única fonte responde completamente a uma pergunta avaliativa. A etapa final é identificar quais dimensões precisam de triangulação:

\begin{itemize}
  \item \textbf{Resultado principal}: qual fonte tem o indicador de resultado mais próximo do que a pergunta avaliativa define?
  \item \textbf{Variáveis de controle}: quais fontes fornecem as características dos grupos de tratamento e controle que precisam ser controladas?
  \item \textbf{Denominadores}: quais fontes fornecem os denominadores populacionais para calcular taxas e proporções?
  \item \textbf{Dados de implementação}: quais fontes permitem verificar se a política foi implementada conforme planejado?
\end{itemize}

% ────────────────────────────────────────────────────────────
\section{Bases amostrais, censitárias e administrativas:
         uma distinção fundamental}
% ────────────────────────────────────────────────────────────

Antes de detalhar os mecanismos de acesso, vale fixar uma distinção conceitual que percorre todo o guia e que, quando ignorada, produz usos ingênuos das fontes disponíveis.

\textbf{Bases amostrais} (como PNAD Contínua, PNS e POF) são excelentes para inferência populacional quando o desenho amostral é respeitado. Elas não devem ser usadas para estimativas em domínios para os quais não foram planejadas --- em particular, a PNAD Contínua \textbf{não deve ser tratada como base de inferência municipal}. Recortes muito finos exigem cautela, eventualmente agregação temporal de múltiplos trimestres, e sempre transparência quanto à precisão estatística das estimativas. O uso correto exige não apenas a aplicação dos pesos amostrais, mas também a declaração do desenho amostral complexo --- estratos e unidades primárias de amostragem (PSUs) --- tanto para estimativas pontuais quanto para a inferência. Ignorar essa estrutura produz erros padrão sistematicamente subestimados.

\textbf{Bases censitárias} (como Censo Demográfico, Censo Escolar e RAIS) cobrem, em princípio, todo o universo de interesse, o que elimina o erro de amostragem e permite desagregações finas. A limitação é a periodicidade: censos são decenais ou anuais, e o dado pode estar desatualizado para fenômenos de rápida mudança.

\textbf{Bases administrativas} (como SIM, SINASC, CAGED, CadÚnico e registros policiais estaduais) podem ter granularidade temporal e espacial extraordinária, chegando ao nível de evento individual em tempo quase real. Mas refletem as regras institucionais de registro --- cobrem apenas o que passa pelo sistema formal, dependem da qualidade do preenchimento por operadores heterogêneos e estão sujeitas a mudanças de critério ao longo do tempo. Não representam necessariamente o universo completo do fenômeno de interesse.

\begin{boxdestaque}{Três perguntas para não confundir o tipo de base}
\begin{enumerate}
  \item \textbf{Para quem a estimativa é válida?} Bases amostrais têm domínios de representatividade definidos; bases censitárias e administrativas têm coberturas definidas pelo universo de registro.
  \item \textbf{O que fica de fora?} Toda base tem lacunas --- seja por desenho amostral, por subregistro ou por exclusão do sistema formal.
  \item \textbf{O que gera o dado?} Em bases amostrais, é um questionário aplicado a uma amostra; em bases censitárias, é um levantamento universal; em bases administrativas, é um processo institucional de rotina.
\end{enumerate}
\end{boxdestaque}

% ────────────────────────────────────────────────────────────
\section{Mecanismos de acesso a dados}
% ────────────────────────────────────────────────────────────

Uma vez identificadas as fontes adequadas, o próximo passo é obtê-las. Os dados públicos brasileiros estão disponíveis por quatro mecanismos principais, com diferentes implicações de prazo, custo e procedimento.

\subsection{Dados abertos: download direto e APIs}

A maior parte das fontes apresentadas neste guia tem algum nível de dados abertos --- microdados em CSV, APIs REST ou tabulações agregadas em portais de consulta. Para essas fontes, o acesso é imediato, gratuito e sem necessidade de cadastro ou autorização.

\begin{table}[H]
\centering
\caption{Principais repositórios de dados abertos para avaliação de políticas públicas}
\label{tab:repositorios}
\begin{tabularx}{\textwidth}{@{} >{\raggedright\arraybackslash}p{4.5cm} X @{}}
\toprule
\textbf{Repositório} & \textbf{O que contém} \\
\midrule
\url{basedosdados.org}       & Dezenas de bases tratadas (PNAD-C, Censo Escolar, RAIS, SIM, FINBRA, etc.) em BigQuery \\
\url{dados.gov.br}           & Catálogo de dados abertos federais (qualidade variável) \\
\url{sidra.ibge.gov.br}      & Tabelas de todas as pesquisas do IBGE \\
\url{datasus.saude.gov.br}   & Sistemas de saúde (SIM, SINASC, SIH, SIA, SINAN, CNES) \\
\url{gov.br/inep/microdados} & Microdados de todas as pesquisas do Inep \\
\url{siconfi.tesouro.gov.br} & Finanças municipais e estaduais (FINBRA/Siconfi) \\
\href{https://portaldatransparencia.gov.br}{\nolinkurl{portalda-transparencia.gov.br}} & Execução orçamentária federal, programas sociais, contratos \\
\url{ipeadata.gov.br}        & Séries históricas macroeconômicas e sociais \\
\url{cecad.cidadania.gov.br} & Dados agregados do CadÚnico e Bolsa Família por município \\
\url{dadosabertos.tse.jus.br} & Dados eleitorais e de financiamento de campanha \\
\bottomrule
\end{tabularx}
\fonte{FGV CLEAR, elaboração própria.}
\end{table}

\subsection{Acesso mediante cadastro}

Algumas fontes requerem cadastro gratuito antes de liberar o acesso aos microdados. O processo é geralmente simples, mas implica aceitar termos de uso que estabelecem restrições sobre como os dados podem ser utilizados e divulgados.

Exemplos de fontes com acesso mediante cadastro:
\begin{itemize}
  \item \textbf{RAIS anonimizada} (MTE): cadastro no portal do Ministério do Trabalho;
  \item \textbf{API do Portal da Transparência} (CGU): cadastro para obtenção de \textit{token};
  \item \textbf{Microdados do PISA para o Brasil} (Inep): cadastro no portal do Inep.
\end{itemize}

\subsection{Lei de Acesso à Informação (LAI)}

A Lei de Acesso à Informação (Lei n.º 12.527/2011) garante a qualquer cidadão o direito de solicitar informações a órgãos públicos. Para o avaliador, a LAI é um mecanismo fundamental para obter dados que não estão disponíveis publicamente mas existem nos sistemas administrativos dos órgãos governamentais.

\subsubsection{Antes de formular um pedido: pesquise o que já existe}

Antes de formular um novo pedido LAI, vale sempre verificar se o dado de interesse já foi disponibilizado em resposta a pedidos anteriores. Essa pesquisa prévia pode poupar tempo considerável e revelar precedentes úteis sobre como determinado órgão trata solicitações similares.

No âmbito federal, a ferramenta oficial \textbf{Busca de Pedidos e Respostas} (\url{www.gov.br/acessoainformacao/pt-br/assuntos/busca-de-pedidos-e-respostas}) permite consultar o histórico de pedidos e respostas de todos os órgãos do Poder Executivo federal. A plataforma \textbf{BuscaLAI} (\url{buscalai.cgu.gov.br}), mantida pela CGU, oferece uma interface de busca mais avançada para o mesmo acervo, facilitando a localização de respostas por tema, órgão e período. Essas consultas ajudam a identificar formatos de dados já fornecidos, entender os limites do que cada órgão costuma divulgar e calibrar a redação do pedido com base em experiências anteriores.

\subsubsection{Como fazer um pedido LAI eficaz}

\begin{itemize}
  \item \textbf{Seja específico}: pedidos vagos têm maior chance de serem negados ou respondidos de forma insatisfatória. Especifique exatamente quais variáveis, para qual período, em qual nível de desagregação e em qual formato;

  \item \textbf{Solicite formato adequado}: especifique que os dados devem ser entregues em formato aberto e legível por máquina (CSV, Excel), não em PDF impresso ou digitalizado;

  \item \textbf{Use o Fala.BR}: o canal oficial para pedidos LAI ao governo federal é atualmente o \textbf{Fala.BR --- Módulo Acesso à Informação} (\url{falabr.cgu.gov.br}), que substituiu e integrou o antigo e-SIC. Estados e municípios têm canais próprios, que devem ser consultados nos respectivos portais de transparência;

  \item \textbf{Conheça os prazos}: o órgão tem 20 dias corridos para responder, prorrogáveis por mais 10 dias mediante justificativa. Em caso de negativa ou resposta insatisfatória, o solicitante pode interpor recurso;

  \item \textbf{Documente tudo}: guarde os protocolos dos pedidos, as respostas recebidas e os recursos interpostos.
\end{itemize}

\subsubsection{Fluxo recursal no âmbito federal}

No caso do Poder Executivo Federal, o fluxo recursal administrativo tem três instâncias, e é importante conhecê-lo antes de desistir diante de uma negativa:

\begin{enumerate}
  \item \textbf{Recurso à autoridade hierárquica superior} do órgão que negou o acesso (1ª instância recursal);
  \item \textbf{Recurso à CGU} (Controladoria-Geral da União), que analisa a conformidade da negativa com a LAI (2ª instância recursal);
  \item \textbf{Recurso à CMRI} (Comissão Mista de Reavaliação de Informações), instância recursal final no âmbito administrativo federal, competente para revisar decisões da CGU em casos de negativa de acesso (\url{www.gov.br/casacivil/pt-br/assuntos/colegiados/comissao-mista-de-reavaliacao-de-informacoes-cmri}).
\end{enumerate}

Decisões da CGU e da CMRI sobre a aplicação da LAI estão disponíveis em \url{www.gov.br/acessoainformacao/pt-br/decisoes-da-cgu-e-da-cmri-sobre-a-aplicacao-da-lai} e constituem precedentes relevantes para antecipar o resultado de pedidos futuros sobre temas similares.

\begin{boxexemplo}{Exemplo prático --- Usando a LAI para obter dados de um programa municipal}
Uma pesquisadora queria avaliar o impacto de um programa municipal de distribuição de cestas básicas durante a pandemia. O município não disponibilizava a lista de beneficiários publicamente. Por meio de pedido LAI, solicitou: ``Lista de beneficiários do Programa X no período de abril de 2020 a dezembro de 2021, contendo: NIS ou CPF (anonimizado), bairro de residência, número de cestas recebidas por mês e composição familiar (número de membros). Formato: planilha CSV.'' O município respondeu em 18 dias com os dados solicitados, que foram então cruzados com dados do CadÚnico (obtidos via convênio com o MDS) para construir o grupo de controle da avaliação.

\textit{Fonte: FGV CLEAR, elaboração própria.}
\end{boxexemplo}

\subsection{Convênios e termos de cooperação técnica}

Algumas fontes com alto valor analítico têm acesso restrito que requer a formalização de convênio ou termo de cooperação técnica entre a instituição do pesquisador e o órgão detentor dos dados. Os principais exemplos são:

\begin{itemize}
  \item \textbf{RAIS identificada (com CPF)}: convênio com o Ministério do Trabalho e Emprego;
  \item \textbf{Microdados do CadÚnico e Bolsa Família}: convênio com o MDS;
  \item \textbf{SCR (Sistema de Informações de Crédito)}: convênio com o Banco Central;
  \item \textbf{SUIBE (dados do BPC)}: convênio com o INSS.
\end{itemize}

Convênios têm prazos de negociação que variam de dois a seis meses em média, exigem aprovação institucional, termos de sigilo assinados por todos os membros da equipe e, em alguns casos, análise por comitê de ética.

\begin{boxdica}{Dica --- Como acelerar a formalização de convênios}
\begin{itemize}
  \item Verifique se sua instituição já tem convênio vigente com o órgão de interesse --- a renovação é mais rápida do que a formalização inicial;
  \item Inicie o processo de convênio antes de finalizar o desenho da avaliação, para evitar que restrições de acesso forcem mudanças metodológicas tardias;
  \item Mantenha contato com a área técnica do órgão (não apenas a área jurídica) para esclarecer dúvidas sobre o escopo dos dados que podem ser compartilhados.
\end{itemize}
\end{boxdica}

% ────────────────────────────────────────────────────────────
\section{Boas práticas para bases obtidas por LAI ou convênio}
% ────────────────────────────────────────────────────────────

Bases administrativas obtidas por LAI, convênio ou cooperação técnica têm características distintas das fontes que chegam com documentação padronizada e estrutura conhecida. Quando você recebe um arquivo de uma prefeitura, de uma secretaria estadual ou de um ministério via pedido LAI, é comum que não venha acompanhado de dicionário de variáveis, notas metodológicas ou registro histórico de mudanças de sistema. Nesse contexto, a documentação precisa ser construída ativamente pelo pesquisador.

Para cada base obtida por esses mecanismos, registre sistematicamente:

\begin{itemize}
  \item \textbf{Universo coberto}: quem está incluído e quem está excluído? A base cobre apenas beneficiários ativos, ou também históricos? Cobre todos os municípios do estado, ou apenas os que aderiram ao sistema?

  \item \textbf{Unidade de observação}: cada linha representa um indivíduo, uma transação, um evento, uma família? Pode haver múltiplas linhas por pessoa (painel), ou cada pessoa aparece uma única vez?

  \item \textbf{Periodicidade e data de referência}: os dados são um corte transversal, uma série mensal, um registro de eventos? Qual é a data de referência do arquivo recebido?

  \item \textbf{Regras institucionais de geração do dado}: como o dado é produzido na prática --- por quem, em qual sistema, com qual grau de verificação? Há incentivos institucionais que possam distorcer o preenchimento?

  \item \textbf{Mudanças de sistema ao longo do tempo}: houve troca de software, reformulação de formulários ou mudança de nomenclatura de variáveis ao longo do período analisado? Essas mudanças introduzem quebras de série que precisam ser identificadas e tratadas;

  \item \textbf{Critérios de anonimização}: o órgão aplicou alguma anonimização antes de entregar os dados? Quais variáveis foram suprimidas ou mascaradas? Isso afeta as possibilidades de integração;

  \item \textbf{Restrições de uso impostas pelo termo}: quais análises são permitidas pelos termos do convênio ou da LAI? Há restrições sobre publicação, compartilhamento ou formas de divulgação dos resultados?
\end{itemize}

Essa documentação é especialmente importante quando a base será usada em análises longitudinais ou combinada com outras fontes. Sem ela, é difícil detectar anomalias nos dados e praticamente impossível para outros pesquisadores reproduzir ou verificar a análise.

% ────────────────────────────────────────────────────────────
\section{Integração de bases de dados}
% ────────────────────────────────────────────────────────────

A integração entre bases de dados --- a combinação de informações de fontes distintas usando um identificador comum --- é uma das técnicas mais poderosas da avaliação moderna de políticas públicas. Ela permite construir perfis multidimensionais de indivíduos, famílias ou estabelecimentos que nenhuma fonte isolada poderia fornecer.

\subsection{Identificadores disponíveis no Brasil}

O Brasil tem um conjunto relativamente rico de identificadores administrativos que podem ser usados para integração:

\begin{itemize}
  \item \textbf{CPF}: identificador universal de pessoas físicas. Presente na RAIS, no CadÚnico, no Censo Escolar e em muitos outros sistemas. É o identificador mais poderoso para integração, mas seu acesso é geralmente restrito;

  \item \textbf{NIS}: identificador do CadÚnico, presente nos dados de beneficiários do Bolsa Família, BPC e outros programas sociais;

  \item \textbf{CNPJ}: identificador de estabelecimentos e empresas. Presente na RAIS (base de estabelecimentos), CAGED, CNES e dados de contratos públicos;

  \item \textbf{Código INEP de escola}: identificador único de cada escola no sistema do Inep, presente no Censo Escolar, no SAEB e no IDEB;

  \item \textbf{Código IBGE de município}: código de 7 dígitos presente em praticamente todas as bases deste guia, base do identificador geográfico para integração municipal;

  \item \textbf{CNES}: código único de cada estabelecimento de saúde, presente no SIH, SIA, SINAN e CNES;

  \item \textbf{Identificadores geográficos}: CEP, código de setor censitário, coordenadas geográficas --- fundamentais para integração espacial entre bases que não compartilham identificadores administrativos comuns. A integração com malhas territoriais do IBGE (setores censitários, bairros, distritos) frequentemente é o único caminho para combinar fontes como cadastros imobiliários municipais com dados do Censo Demográfico ou com registros de saúde e educação.
\end{itemize}

\subsection{Boas práticas de integração}

\begin{itemize}
  \item \textbf{Documente a taxa de integração}: calcule e reporte a proporção de registros cruzados com sucesso. Uma taxa baixa pode indicar viés de seleção;

  \item \textbf{Verifique a consistência}: compare a distribuição de variáveis-chave na base cruzada com as distribuições nas bases originais;

  \item \textbf{Trate erros de medição nos identificadores}: CPFs inválidos, CNPJs com dígitos trocados e códigos de município desatualizados são causas comuns de falha na integração;

  \item \textbf{Integração probabilística}: quando identificadores exatos não estão disponíveis, técnicas probabilísticas combinam registros com base na similaridade de múltiplas variáveis (nome, data de nascimento, endereço), introduzindo incerteza adicional que deve ser reportada.
\end{itemize}

\begin{boxatencao}{Atenção --- Integração com CPF requer autorização específica}
A integração de dados pessoais de diferentes fontes cria um novo conjunto de dados com maior poder de identificação dos indivíduos --- o que exige, nos termos da LGPD, base legal específica. Integrações com CPF devem ser realizadas apenas sob os termos dos convênios firmados, que geralmente restringem o escopo das análises permitidas.
\end{boxatencao}

% ────────────────────────────────────────────────────────────
\section{Documentação para reprodutibilidade}
% ────────────────────────────────────────────────────────────

Uma avaliação que não pode ser reproduzida tem valor científico e de credibilidade substancialmente reduzido. A reprodutibilidade não é apenas uma exigência acadêmica; é uma condição para que os resultados sejam escrutinados, verificados e utilizados com confiança por gestores e formuladores de políticas.

\subsection{O que documentar}

Para cada fonte de dados utilizada, a documentação mínima deve incluir:

\begin{itemize}
  \item Identificação da fonte: nome completo, órgão produtor, versão ou edição, URL de acesso;
  \item Data de acesso ou download;
  \item Variáveis utilizadas, com nomes originais e transformações aplicadas;
  \item Filtros aplicados (período, território, população);
  \item Procedimentos de limpeza: valores ausentes, inconsistências, outliers;
  \item Integrações realizadas: bases combinadas, identificador usado, taxa de integração.
\end{itemize}

\subsection{Como documentar: ferramentas e práticas}

\begin{itemize}
  \item \textbf{Código comentado}: todo código de coleta, limpeza e análise deve conter comentários explicando as decisões tomadas em cada etapa;
  \item \textbf{\textit{README} do projeto}: arquivo descrevendo a estrutura do projeto, as fontes utilizadas e as instruções para reproduzir a análise;
  \item \textbf{Controle de versão}: use Git/GitHub para registrar todas as alterações no código e nos dados;
  \item \textbf{Dicionário de variáveis}: para bases construídas ou transformadas pelo pesquisador, mantenha um dicionário com a definição e as fontes de cada variável;
  \item \textbf{Pré-registro}: para avaliações prospectivas, o pré-registro no Open Science Framework (OSF) ou no AEA RCT Registry aumenta substancialmente a credibilidade dos resultados.
\end{itemize}

% ────────────────────────────────────────────────────────────
\section{LGPD e proteção de dados em avaliações de políticas públicas}
% ────────────────────────────────────────────────────────────

A Lei Geral de Proteção de Dados (Lei n.º 13.709/2018) estabelece regras para o tratamento de dados pessoais no Brasil, com impactos diretos sobre a forma como avaliadores coletam, armazenam e usam dados que contêm informações sobre pessoas identificáveis.

\subsection{Conceitos fundamentais}

\begin{itemize}
  \item \textbf{Dado pessoal}: qualquer informação relacionada a uma pessoa natural identificada ou identificável;
  \item \textbf{Dado sensível}: categoria especial que inclui origem racial ou étnica, convicção religiosa, opinião política, dados de saúde ou vida sexual, dados genéticos ou biométricos;
  \item \textbf{Anonimização}: processo pelo qual um dado perde a possibilidade de identificação. Dados anonimizados de forma irreversível estão fora do escopo da LGPD;
  \item \textbf{Pseudonimização}: substituição de identificadores diretos por códigos reversíveis. Dados pseudonimizados ainda são dados pessoais sob a LGPD.
\end{itemize}

\subsection{Bases legais para tratamento de dados em pesquisa}

Para avaliações de políticas públicas, as bases legais mais relevantes são:

\begin{itemize}
  \item \textbf{Execução de políticas públicas} (art. 7º, III): aplicável ao tratamento pelo poder público para execução de suas políticas;
  \item \textbf{Pesquisa} (art. 7º, IV): aplicável ao tratamento por órgão de pesquisa, garantida, quando possível, a anonimização;
  \item \textbf{Interesse legítimo} (art. 7º, IX): aplicável quando necessário para atender interesses legítimos do controlador, exceto quando prevalecem interesses fundamentais do titular.
\end{itemize}

\subsection{Práticas recomendadas para conformidade com a LGPD}

\begin{itemize}
  \item \textbf{Minimize os dados coletados}: colete apenas os dados estritamente necessários para a pergunta avaliativa;
  \item \textbf{Anonimize o mais cedo possível}: aplique anonimização ou pseudonimização assim que possível no fluxo de trabalho;
  \item \textbf{Controle o acesso}: dados pessoais devem ser acessíveis apenas aos membros da equipe que precisam deles;
  \item \textbf{Estabeleça prazo de retenção}: defina por quanto tempo os dados serão mantidos e garanta o descarte seguro após esse prazo;
  \item \textbf{Documente o tratamento}: mantenha um registro das atividades de tratamento, incluindo finalidade, base legal, tipos de dados e prazos.
\end{itemize}

% ────────────────────────────────────────────────────────────
\section{Checklist final: boas práticas no uso de fontes de dados}
% ────────────────────────────────────────────────────────────

\begin{boxdestaque}{Checklist de boas práticas no uso de fontes de dados para avaliação}

\textbf{Seleção da fonte}
\begin{itemize}
  \item[$\square$] A fonte foi selecionada a partir da pergunta avaliativa, não da familiaridade do analista com ela
  \item[$\square$] A cobertura geográfica, periodicidade e unidade de análise são compatíveis com a pergunta
  \item[$\square$] A documentação metodológica foi consultada para verificar o que cada variável mede de fato
  \item[$\square$] Possíveis quebras metodológicas na série histórica foram identificadas e consideradas
  \item[$\square$] O tipo de base (amostral, censitária, administrativa) foi considerado para o uso pretendido
\end{itemize}

\textbf{Acesso e coleta}
\begin{itemize}
  \item[$\square$] O mecanismo de acesso (download, API, LAI, convênio) foi identificado e os prazos planejados
  \item[$\square$] Para pedidos LAI, pedidos e respostas anteriores sobre o tema foram consultados
  \item[$\square$] A data de coleta foi registrada
  \item[$\square$] Uma cópia dos dados brutos foi salva em local seguro, separada das versões processadas
\end{itemize}

\textbf{Bases obtidas por LAI ou convênio}
\begin{itemize}
  \item[$\square$] O universo coberto foi documentado (quem está incluído e quem está excluído)
  \item[$\square$] As regras institucionais de geração do dado foram investigadas
  \item[$\square$] Mudanças de sistema ao longo do período analisado foram identificadas
  \item[$\square$] Os critérios de anonimização aplicados pelo órgão foram registrados
  \item[$\square$] As restrições de uso impostas pelo termo foram verificadas
\end{itemize}

\textbf{Limpeza e preparação}
\begin{itemize}
  \item[$\square$] Os valores ausentes foram identificados e tratados de forma documentada
  \item[$\square$] Inconsistências foram verificadas (valores implausíveis, distribuições inesperadas)
  \item[$\square$] Se foi realizada integração, a taxa de integração foi calculada e reportada
\end{itemize}

\textbf{Análise}
\begin{itemize}
  \item[$\square$] Para pesquisas amostrais, os pesos amostrais \textbf{e} o desenho de amostragem complexo (estratos e PSUs) foram corretamente especificados
  \item[$\square$] A representatividade da fonte para o recorte utilizado foi verificada (em especial, inferências municipais com PNAD Contínua foram evitadas)
  \item[$\square$] As limitações da fonte foram consideradas na interpretação dos resultados
  \item[$\square$] Os resultados foram triangulados com fontes independentes quando possível
\end{itemize}

\textbf{Documentação e reprodutibilidade}
\begin{itemize}
  \item[$\square$] Toda a cadeia de tratamento dos dados está documentada em código comentado
  \item[$\square$] As fontes utilizadas estão identificadas com URL, data de acesso e versão
  \item[$\square$] O código está versionado e disponível para verificação
\end{itemize}

\textbf{Proteção de dados}
\begin{itemize}
  \item[$\square$] Os dados pessoais foram tratados com a base legal adequada (LGPD)
  \item[$\square$] A anonimização ou pseudonimização foi aplicada o mais cedo possível
  \item[$\square$] O acesso aos dados pessoais foi restrito aos membros da equipe que precisam deles
\end{itemize}
\end{boxdestaque}

% ────────────────────────────────────────────────────────────
\section{Síntese do capítulo e do guia}
% ────────────────────────────────────────────────────────────

Este capítulo reuniu as orientações práticas para o trabalho com fontes de dados em avaliações de políticas públicas. Os pontos centrais são:

\begin{itemize}
  \item A seleção de fontes deve \textbf{partir da pergunta}, não da disponibilidade ou da familiaridade com determinadas bases. As cinco etapas propostas estruturam esse processo de forma sistemática;

  \item A distinção entre \textbf{bases amostrais, censitárias e administrativas} é fundamental: cada tipo tem lógica de cobertura, limitações e usos adequados distintos. Em particular, pesquisas amostrais exigem respeito ao desenho amostral complexo e têm domínios de representatividade que não devem ser ultrapassados;

  \item Os mecanismos de acesso --- download aberto, cadastro, LAI e convênio --- têm implicações diferentes de prazo, custo e procedimento. Para pedidos LAI, a pesquisa prévia de pedidos anteriores e o conhecimento do fluxo recursal completo (incluindo CGU e CMRI) são ferramentas importantes;

  \item Bases obtidas por LAI ou convênio exigem \textbf{documentação ativa} pelo pesquisador: universo coberto, regras de geração do dado, mudanças de sistema e restrições de uso precisam ser registrados sistematicamente;

  \item A \textbf{integração de bases}, especialmente com identificadores geográficos quando os administrativos não estão disponíveis, abre possibilidades analíticas que nenhuma fonte isolada oferece;

  \item A \textbf{documentação para reprodutibilidade} e a conformidade com a \textbf{LGPD} são condições para que os resultados da avaliação sejam verificáveis e utilizáveis com responsabilidade.
\end{itemize}

\medskip

Este capítulo encerra o corpo principal do guia. Os Anexos A e B, a seguir, oferecem ferramentas de referência rápida: uma tabela-catálogo consolidada de todas as fontes apresentadas e um glossário dos principais termos técnicos utilizados ao longo do texto.

\medskip

\begin{boxdestaque}{Uma última observação}
Dados não falam por si mesmos. A qualidade de uma avaliação depende menos do número de fontes consultadas do que da clareza da pergunta, do rigor com que cada fonte foi selecionada e da honestidade com que suas limitações foram declaradas. Este guia mapeia o território de dados disponíveis para o Brasil --- o que cada fonte cobre, como acessá-la, onde ela falha. O julgamento analítico de quem avalia é a parte que nenhum guia pode fornecer.
\end{boxdestaque}


\backmatter
\appendix
% ============================================================
%  anexo_a_tabela_catalogo.tex  [REVISADO]
%  Anexo A: Tabela-Catálogo Consolidada de Fontes de Dados
% ============================================================

% ════════════════════════════════════════════════════════════
\chapter*{Anexo A — Tabela-Catálogo de Fontes de Dados}
\addcontentsline{toc}{chapter}{Anexo A — Tabela-Catálogo de Fontes de Dados}
\markboth{Anexo A}{Anexo A}
% ════════════════════════════════════════════════════════════

Este anexo reúne em formato de consulta rápida todas as fontes de dados apresentadas ao longo do guia. As fontes estão organizadas por tema, com as informações essenciais para identificação, acesso e uso inicial. Para descrições completas, histórico, variáveis disponíveis e limitações de cada fonte, consulte o capítulo correspondente.

\medskip

\textbf{Como usar esta tabela}

\begin{itemize}
  \item A coluna \textbf{Unidade} indica a unidade de análise dos microdados (quando disponíveis);
  \item A coluna \textbf{Acesso} indica o mecanismo de acesso: \textbf{A} = aberto (download direto ou API); \textbf{C} = cadastro gratuito; \textbf{LAI} = solicitação via Lei de Acesso à Informação; \textbf{Conv} = convênio institucional;
  \item Os links foram verificados na data de elaboração deste guia e podem estar sujeitos a alterações.
\end{itemize}

% ────────────────────────────────────────────────────────────
\section*{A.1 Pesquisas Amostrais e Censos Demográficos}
\addcontentsline{toc}{section}{A.1 Pesquisas Amostrais e Censos Demográficos}
% ────────────────────────────────────────────────────────────

\begin{longtable}{@{} p{3cm} p{2.5cm} p{2cm} p{2cm} p{1.2cm} p{3.3cm} @{}}
\toprule
\textbf{Fonte} & \textbf{Órgão} & \textbf{Cobertura} & \textbf{Período} & \textbf{Acesso} & \textbf{URL} \\
\midrule
\endfirsthead
\multicolumn{6}{l}{\small\textit{(continuação)}} \\
\toprule
\textbf{Fonte} & \textbf{Órgão} & \textbf{Cobertura} & \textbf{Período} & \textbf{Acesso} & \textbf{URL} \\
\midrule
\endhead
\bottomrule
\endfoot

Censo Demográfico
  & IBGE
  & Brasil, UF, município, setor censitário
  & Decenal (1940--2022)
  & A
  & \url{ibge.gov.br/censo2022} \\

PNAD Contínua
  & IBGE
  & Brasil, UF, RM, capital
  & Trimestral (2012--)
  & A
  & \url{ibge.gov.br/pnadc} \\

PNAD (anterior)
  & IBGE
  & Brasil, UF, RM
  & Anual (1967--2015)
  & A
  & \url{ibge.gov.br/pnad} \\

Pesquisa Mensal de Emprego (PME)
  & IBGE
  & 6 RMs
  & Mensal (1980--2016)
  & A
  & \url{ibge.gov.br/pme} \\

Pesquisa de Orçamentos Familiares (POF)
  & IBGE
  & Brasil, UF
  & Irregular (1987--2018)
  & A
  & \url{ibge.gov.br/pof} \\

Censo Agropecuário
  & IBGE
  & Brasil, UF, município
  & Irregular (1920--2017)
  & A
  & \url{ibge.gov.br/censoagro} \\

\end{longtable}

% ────────────────────────────────────────────────────────────
\section*{A.2 Mercado de Trabalho e Renda}
\addcontentsline{toc}{section}{A.2 Mercado de Trabalho e Renda}
% ────────────────────────────────────────────────────────────

\begin{longtable}{@{} p{3cm} p{2.5cm} p{2cm} p{2cm} p{1.2cm} p{3.3cm} @{}}
\toprule
\textbf{Fonte} & \textbf{Órgão} & \textbf{Cobertura} & \textbf{Período} & \textbf{Acesso} & \textbf{URL} \\
\midrule
\endfirsthead
\multicolumn{6}{l}{\small\textit{(continuação)}} \\
\toprule
\textbf{Fonte} & \textbf{Órgão} & \textbf{Cobertura} & \textbf{Período} & \textbf{Acesso} & \textbf{URL} \\
\midrule
\endhead
\bottomrule
\endfoot

CAGED
  & MTE
  & Brasil, UF, município
  & Mensal (1992--2019)
  & A
  & \url{gov.br/trabalho} \\

Novo CAGED
  & MTE
  & Brasil, UF, município
  & Mensal (2020--)
  & A
  & \url{gov.br/trabalho} \\

RAIS (anonimizada)
  & MTE
  & Brasil, UF, município
  & Anual (1975--)
  & C
  & \url{gov.br/trabalho/rais} \\

RAIS (identificada)
  & MTE
  & Brasil, UF, município
  & Anual (1975--)
  & Conv
  & \url{gov.br/trabalho/rais} \\

eSocial
  & MTE/RFB/INSS
  & Brasil
  & Contínuo (2018--)
  & Conv
  & \url{gov.br/esocial} \\

Cadastro Único (CadÚnico) --- agregado
  & MDS
  & Brasil, UF, município
  & Mensal (2001--)
  & A
  & \url{cecad.cidadania.gov.br} \\

Cadastro Único (CadÚnico) --- microdados
  & MDS
  & Brasil, UF, município
  & Mensal (2001--)
  & Conv
  & \url{gov.br/mds} \\

\end{longtable}

% ────────────────────────────────────────────────────────────
\section*{A.3 Saúde}
\addcontentsline{toc}{section}{A.3 Saúde}
% ────────────────────────────────────────────────────────────

\begin{longtable}{@{} p{3cm} p{2.5cm} p{2cm} p{2cm} p{1.2cm} p{3.3cm} @{}}
\toprule
\textbf{Fonte} & \textbf{Órgão} & \textbf{Cobertura} & \textbf{Período} & \textbf{Acesso} & \textbf{URL} \\
\midrule
\endfirsthead
\multicolumn{6}{l}{\small\textit{(continuação)}} \\
\toprule
\textbf{Fonte} & \textbf{Órgão} & \textbf{Cobertura} & \textbf{Período} & \textbf{Acesso} & \textbf{URL} \\
\midrule
\endhead
\bottomrule
\endfoot

SIM --- Mortalidade
  & MS/DATASUS
  & Brasil, UF, município
  & Anual (1979--)
  & A
  & \url{datasus.saude.gov.br} \\

SINASC --- Nascidos Vivos
  & MS/DATASUS
  & Brasil, UF, município
  & Anual (1990--)
  & A
  & \url{datasus.saude.gov.br} \\

SIH --- Internações Hospitalares
  & MS/DATASUS
  & Brasil, UF, município
  & Mensal (1992--)
  & A
  & \url{datasus.saude.gov.br} \\

SIA --- Ambulatorial
  & MS/DATASUS
  & Brasil, UF, município
  & Mensal (1994--)
  & A
  & \url{datasus.saude.gov.br} \\

SINAN --- Agravos de Notificação
  & MS/DATASUS
  & Brasil, UF, município
  & Cont. (1993--)
  & A
  & \url{datasus.saude.gov.br} \\

CNES --- Estabelecimentos de Saúde
  & MS/DATASUS
  & Brasil, UF, município
  & Mensal (2005--)
  & A
  & \url{cnes.datasus.gov.br} \\

SIOPS --- Gasto em Saúde
  & MS
  & Brasil, UF, município
  & Anual (2000--)
  & A
  & \url{siops.datasus.gov.br} \\

SISVAN --- Vigilância Nutricional
  & MS/DATASUS
  & Brasil, UF, município
  & Cont. (2008--)
  & A
  & \url{sisvan.saude.gov.br} \\

Pesquisa Nacional de Saúde (PNS)
  & IBGE/MS
  & Brasil, UF, capital
  & Quinquenal (2013, 2019)
  & A
  & \url{ibge.gov.br/pns} \\

AMS --- Assist. Médico-Sanitária
  & IBGE
  & Brasil, UF, município
  & Irregular (1999--2009)
  & A
  & \url{ibge.gov.br/ams} \\

\end{longtable}

% ────────────────────────────────────────────────────────────
\section*{A.4 Educação}
\addcontentsline{toc}{section}{A.4 Educação}
% ────────────────────────────────────────────────────────────

\begin{longtable}{@{} p{3cm} p{2.5cm} p{2cm} p{2cm} p{1.2cm} p{3.3cm} @{}}
\toprule
\textbf{Fonte} & \textbf{Órgão} & \textbf{Cobertura} & \textbf{Período} & \textbf{Acesso} & \textbf{URL} \\
\midrule
\endfirsthead
\multicolumn{6}{l}{\small\textit{(continuação)}} \\
\toprule
\textbf{Fonte} & \textbf{Órgão} & \textbf{Cobertura} & \textbf{Período} & \textbf{Acesso} & \textbf{URL} \\
\midrule
\endhead
\bottomrule
\endfoot

Censo Escolar
  & Inep/MEC
  & Brasil, UF, município, escola
  & Anual (1995--) (microdados)
  & A
  & \url{gov.br/inep/censo-escolar} \\

IDEB
  & Inep/MEC
  & Brasil, UF, município, escola
  & Bienal (2005--)
  & A
  & \url{gov.br/inep/ideb} \\

SAEB / Prova Brasil
  & Inep/MEC
  & Brasil, UF, escola (Prova Brasil)
  & Bienal (1990--)
  & A
  & \url{gov.br/inep/saeb} \\

ENEM
  & Inep/MEC
  & Brasil (participantes)
  & Anual (1998--)
  & A
  & \url{gov.br/inep/enem} \\

Censo da Educação Superior
  & Inep/MEC
  & Brasil, UF, município, IES
  & Anual (2001--)
  & A
  & \url{gov.br/inep/censo-superior} \\

SIOPE --- Gasto em Educação
  & FNDE/MEC
  & Brasil, UF, município
  & Anual (1999--)
  & A
  & \url{fnde.gov.br/siope} \\

Plataforma Nilo Peçanha (PNP)
  & SETEC/MEC
  & Rede Federal de Ed. Profissional
  & Anual (2018--)
  & A
  & \url{plataformanilopecanha.mec.gov.br} \\

PISA
  & OCDE/Inep
  & Brasil (amostral nacional)
  & Trienal (2000--)
  & A
  & \url{gov.br/inep/pisa} \\

\end{longtable}

% ────────────────────────────────────────────────────────────
\section*{A.5 Assistência Social e Transferência de Renda}
\addcontentsline{toc}{section}{A.5 Assistência Social e Transferência de Renda}
% ────────────────────────────────────────────────────────────

\begin{longtable}{@{} p{3cm} p{2.5cm} p{2cm} p{2cm} p{1.2cm} p{3.3cm} @{}}
\toprule
\textbf{Fonte} & \textbf{Órgão} & \textbf{Cobertura} & \textbf{Período} & \textbf{Acesso} & \textbf{URL} \\
\midrule
\endfirsthead
\multicolumn{6}{l}{\small\textit{(continuação)}} \\
\toprule
\textbf{Fonte} & \textbf{Órgão} & \textbf{Cobertura} & \textbf{Período} & \textbf{Acesso} & \textbf{URL} \\
\midrule
\endhead
\bottomrule
\endfoot

Bolsa Família --- pagamentos
  & MDS / Portal Transparência
  & Brasil, UF, município
  & Mensal (2004--)
  & A
  & \url{portaldatransparencia.gov.br} \\

CECAD --- CadÚnico agregado
  & MDS
  & Brasil, UF, município
  & Mensal (2001--)
  & A
  & \url{cecad.cidadania.gov.br} \\

VISDATA --- Painéis MDS
  & MDS/SAGI
  & Brasil, UF, município
  & Variável
  & A
  & \url{aplicacoes.mds.gov.br/sagi/vis/data3} \\

BPC --- Beneficiários
  & MDS/INSS
  & Brasil, UF, município
  & Mensal (1996--)
  & A
  & \url{portaldatransparencia.gov.br} \\

Censo SUAS
  & MDS
  & Brasil, UF, município
  & Anual (2007--)
  & A
  & \url{gov.br/mds/censo-suas} \\

RMA --- Registro Mensal de Atendimentos
  & MDS
  & Brasil, UF, município
  & Mensal (2012--)
  & A
  & \url{gov.br/mds/rede-suas} \\

\end{longtable}

% ────────────────────────────────────────────────────────────
\section*{A.6 Segurança Pública}
\addcontentsline{toc}{section}{A.6 Segurança Pública}
% ────────────────────────────────────────────────────────────

\begin{longtable}{@{} p{3cm} p{2.5cm} p{2cm} p{2cm} p{1.2cm} p{3.3cm} @{}}
\toprule
\textbf{Fonte} & \textbf{Órgão} & \textbf{Cobertura} & \textbf{Período} & \textbf{Acesso} & \textbf{URL} \\
\midrule
\endfirsthead
\multicolumn{6}{l}{\small\textit{(continuação)}} \\
\toprule
\textbf{Fonte} & \textbf{Órgão} & \textbf{Cobertura} & \textbf{Período} & \textbf{Acesso} & \textbf{URL} \\
\midrule
\endhead
\bottomrule
\endfoot

Atlas da Violência
  & Ipea/FBSP
  & Brasil, UF, município
  & Anual (1980--)
  & A
  & \url{ipea.gov.br/atlasviolencia} \\

Anuário Brasileiro de Seg. Pública
  & FBSP
  & Brasil, UF
  & Anual (2007--)
  & A
  & \url{forumseguranca.org.br/anuario} \\

Ocorrências SSP-SP
  & SSP/SP
  & Município/delegacia (SP)
  & Mensal (2001--)
  & A
  & \url{ssp.sp.gov.br} \\

Ocorrências ISP-RJ
  & ISP/RJ
  & AISP (RJ)
  & Mensal (2003--)
  & A
  & \url{ispvisualizacao.rj.gov.br} \\

SISDEPEN / Infopen
  & DEPEN/MJSP
  & Brasil, UF, estabelecimento
  & Semestral (2004--)
  & A
  & \url{gov.br/senappen/sisdepen} \\

Vitimização PNAD 2009
  & IBGE
  & Brasil, UF
  & Pontual (2009)
  & A
  & \url{ibge.gov.br/pnad} \\

VIVA / SINAN Violências
  & MS/DATASUS
  & Brasil, UF, município
  & Cont. (2009--)
  & A
  & \url{datasus.saude.gov.br} \\

\end{longtable}

% ────────────────────────────────────────────────────────────
\section*{A.7 Infraestrutura, Território e Meio Ambiente}
\addcontentsline{toc}{section}{A.7 Infraestrutura, Território e Meio Ambiente}
% ────────────────────────────────────────────────────────────

\begin{longtable}{@{} p{3cm} p{2.5cm} p{2cm} p{2cm} p{1.2cm} p{3.3cm} @{}}
\toprule
\textbf{Fonte} & \textbf{Órgão} & \textbf{Cobertura} & \textbf{Período} & \textbf{Acesso} & \textbf{URL} \\
\midrule
\endfirsthead
\multicolumn{6}{l}{\small\textit{(continuação)}} \\
\toprule
\textbf{Fonte} & \textbf{Órgão} & \textbf{Cobertura} & \textbf{Período} & \textbf{Acesso} & \textbf{URL} \\
\midrule
\endhead
\bottomrule
\endfoot

SNIS --- Saneamento
  & MCidades
  & Brasil, UF, município
  & Anual (1995--)
  & A
  & \url{snis.gov.br} \\

Munic --- Informações Municipais
  & IBGE
  & Brasil, município
  & Irregular (1999--)
  & A
  & \url{ibge.gov.br/munic} \\

Estadic --- Informações Estaduais
  & IBGE
  & Brasil, UF
  & Irregular (2012--)
  & A
  & \url{ibge.gov.br/estadic} \\

ANEEL --- Energia Elétrica
  & ANEEL
  & Brasil, UF, município
  & Variável
  & A
  & \url{dadosabertos.aneel.gov.br} \\

ANATEL --- Telecomunicações
  & ANATEL
  & Brasil, UF, município
  & Trimestral
  & A
  & \url{dados.anatel.gov.br} \\

ANS --- Saúde Suplementar
  & ANS
  & Brasil, UF, município
  & Trimestral
  & A
  & \url{dadosabertos.ans.gov.br} \\

Malhas Territoriais IBGE
  & IBGE
  & Brasil (todos os níveis)
  & Atualização anual
  & A
  & \url{ibge.gov.br/malhas} \\

Setores Censitários
  & IBGE
  & Brasil
  & Decenal (Censo 2022)
  & A
  & \url{ibge.gov.br/censo2022} \\

CNEFE --- Cadastro de Endereços
  & IBGE
  & Brasil
  & Decenal (Censo 2022)
  & A
  & \url{ibge.gov.br/censo2022} \\

ITBI --- Transações Imobiliárias
  & Prefeituras municipais
  & Municipal
  & Variável (contínuo)
  & A/LAI
  & Portal de dados abertos ou LAI municipal \\

Cadastro Imobiliário / IPTU
  & Prefeituras municipais
  & Municipal
  & Anual
  & A/LAI
  & Portal de dados abertos ou LAI municipal \\

Planta Genérica de Valores (PGV)
  & Prefeituras municipais
  & Municipal
  & Variável (por lei municipal)
  & A/LAI
  & Portal de dados abertos ou LAI municipal \\

Licenciamento Urbanístico / Habite-se
  & Prefeituras municipais
  & Municipal
  & Variável (contínuo)
  & A/LAI
  & Portal de dados abertos ou LAI municipal \\

Zoneamento e Uso do Solo
  & Prefeituras municipais
  & Municipal
  & Variável (por lei municipal)
  & A/LAI
  & Portal de dados abertos ou LAI municipal \\

PRODES --- Desmatamento Amazônia
  & INPE
  & Amazônia Legal
  & Anual (1988--)
  & A
  & \url{data.inpe.br/biomasbr/prodes} \\

DETER --- Alertas de Desmatamento
  & INPE
  & Amazônia Legal e Cerrado
  & Contínuo
  & A
  & \url{data.inpe.br/biomasbr/deter} \\

MapBiomas --- Uso do Solo
  & MapBiomas
  & Brasil
  & Anual (1985--)
  & A
  & \url{mapbiomas.org} \\

CAR --- Cadastro Ambiental Rural
  & SFB/MMA
  & Brasil
  & Cont. (2013--)
  & A
  & \url{car.gov.br} \\

\end{longtable}

% ────────────────────────────────────────────────────────────
\section*{A.8 Finanças Públicas}
\addcontentsline{toc}{section}{A.8 Finanças Públicas}
% ────────────────────────────────────────────────────────────

\begin{longtable}{@{} p{3cm} p{2.5cm} p{2cm} p{2cm} p{1.2cm} p{3.3cm} @{}}
\toprule
\textbf{Fonte} & \textbf{Órgão} & \textbf{Cobertura} & \textbf{Período} & \textbf{Acesso} & \textbf{URL} \\
\midrule
\endfirsthead
\multicolumn{6}{l}{\small\textit{(continuação)}} \\
\toprule
\textbf{Fonte} & \textbf{Órgão} & \textbf{Cobertura} & \textbf{Período} & \textbf{Acesso} & \textbf{URL} \\
\midrule
\endhead
\bottomrule
\endfoot

FINBRA / Siconfi
  & STN
  & Brasil, UF, município
  & Anual (1986--)
  & A
  & \url{siconfi.tesouro.gov.br} \\

RREO --- Execução Bimestral
  & STN / entes federativos
  & Brasil, UF, município
  & Bimestral (2000--)
  & A
  & \url{siconfi.tesouro.gov.br} \\

RGF --- Gestão Fiscal (LRF)
  & STN / entes federativos
  & Brasil, UF, município
  & Quadrimestral (2000--)
  & A
  & \url{siconfi.tesouro.gov.br} \\

Portal da Transparência Federal
  & CGU
  & Governo Federal
  & Diário (2004--)
  & A
  & \url{portaldatransparencia.gov.br} \\

SIOPE --- Gasto em Educação
  & FNDE/MEC
  & Brasil, UF, município
  & Anual (1999--)
  & A
  & \url{fnde.gov.br/siope} \\

SIOPS --- Gasto em Saúde
  & MS
  & Brasil, UF, município
  & Anual (2000--)
  & A
  & \url{siops.datasus.gov.br} \\

TSE --- Dados Eleitorais
  & TSE
  & Brasil, UF, município
  & Por eleição (2002--)
  & A
  & \url{dadosabertos.tse.jus.br} \\

BCB/SGS --- Séries Temporais
  & BCB
  & Brasil (macro)
  & Variável
  & A
  & \url{api.bcb.gov.br} \\

BCB/Estban --- Crédito Municipal
  & BCB
  & Brasil, UF, município
  & Mensal
  & A
  & \url{bcb.gov.br/estban} \\

Ipeadata
  & Ipea
  & Brasil, UF, município
  & Variável
  & A
  & \url{ipeadata.gov.br} \\

\end{longtable}

% ────────────────────────────────────────────────────────────
\section*{A.9 Fontes Internacionais}
\addcontentsline{toc}{section}{A.9 Fontes Internacionais}
% ────────────────────────────────────────────────────────────

\begin{longtable}{@{} p{3cm} p{2.5cm} p{2cm} p{2cm} p{1.2cm} p{3.3cm} @{}}
\toprule
\textbf{Fonte} & \textbf{Órgão} & \textbf{Cobertura} & \textbf{Período} & \textbf{Acesso} & \textbf{URL} \\
\midrule
\endfirsthead
\multicolumn{6}{l}{\small\textit{(continuação)}} \\
\toprule
\textbf{Fonte} & \textbf{Órgão} & \textbf{Cobertura} & \textbf{Período} & \textbf{Acesso} & \textbf{URL} \\
\midrule
\endhead
\bottomrule
\endfoot

World Development Indicators (WDI)
  & Banco Mundial
  & 200+ países
  & Anual (1960--)
  & A
  & \url{databank.worldbank.org/wdi} \\

Poverty \& Inequality Platform (PIP)
  & Banco Mundial
  & 100+ países
  & Por pesquisa
  & A
  & \url{pip.worldbank.org} \\

OECD Stats
  & OCDE
  & 38+ países
  & Anual
  & A
  & \url{stats.oecd.org} \\

PISA
  & OCDE/Inep
  & 80+ países
  & Trienal (2000--)
  & A
  & \url{oecd.org/pisa} \\

ILOSTAT
  & OIT
  & 200+ países
  & Anual
  & A
  & \url{ilostat.ilo.org} \\

Global Health Observatory (GHO)
  & OMS
  & 194 países
  & Anual
  & A
  & \url{who.int/data/gho} \\

Human Development Reports (HDR)
  & PNUD
  & 190+ países
  & Anual
  & A
  & \url{hdr.undp.org/data-center} \\

IDHM --- IDH Municipal Brasil
  & PNUD Brasil
  & Brasil, município
  & Decenal (Censos)
  & A
  & \url{undp.org/pt/brazil} \\

CEPALSTAT / BADEHOG
  & CEPAL
  & 33 países AL\&C
  & Anual / por pesquisa
  & A
  & \url{statistics.cepal.org} \\

FAOSTAT
  & FAO
  & 245+ países
  & Anual (1961--)
  & A
  & \url{fao.org/faostat} \\

V-Dem
  & Univ. Gotemburgo
  & 200+ países
  & Anual (1789--)
  & A
  & \url{v-dem.net} \\

\end{longtable}

% ────────────────────────────────────────────────────────────
\section*{A.10 Acesso à Informação: canais e recursos}
\addcontentsline{toc}{section}{A.10 Acesso à Informação: canais e recursos}
% ────────────────────────────────────────────────────────────

Esta seção reúne os principais canais e ferramentas para acesso a dados via Lei de Acesso à Informação (LAI) e para pesquisa de pedidos anteriores.

\begin{longtable}{@{} p{3.5cm} p{2.5cm} p{5cm} p{1.2cm} p{2.8cm} @{}}
\toprule
\textbf{Canal / Ferramenta} & \textbf{Mantenedor} & \textbf{Função} & \textbf{Acesso} & \textbf{URL} \\
\midrule
\endfirsthead
\multicolumn{5}{l}{\small\textit{(continuação)}} \\
\toprule
\textbf{Canal / Ferramenta} & \textbf{Mantenedor} & \textbf{Função} & \textbf{Acesso} & \textbf{URL} \\
\midrule
\endhead
\bottomrule
\endfoot

Fala.BR --- Módulo Acesso à Informação
  & CGU
  & Canal oficial de pedidos LAI ao Poder Executivo Federal (substitui o e-SIC)
  & A
  & \url{falabr.cgu.gov.br} \\

Busca de Pedidos e Respostas
  & CGU
  & Consulta ao histórico de pedidos e respostas LAI do Poder Executivo Federal
  & A
  & \url{gov.br/acessoainformacao/pt-br/assuntos/busca-de-pedidos-e-respostas} \\

BuscaLAI
  & CGU
  & Busca avançada de pedidos e respostas LAI por tema, órgão e período
  & A
  & \url{buscalai.cgu.gov.br} \\

Recursos LAI --- fluxo recursal
  & CGU
  & Informações sobre os recursos administrativos previstos na LAI
  & A
  & \url{gov.br/acessoainformacao/pt-br/assuntos/recursos} \\

CMRI --- Comissão Mista de Reavaliação de Informações
  & Casa Civil
  & Instância recursal final no âmbito administrativo federal (após CGU)
  & A
  & \url{gov.br/casacivil/pt-br/assuntos/colegiados/comissao-mista-de-reavaliacao-de-informacoes-cmri} \\

Decisões CGU e CMRI sobre a LAI
  & CGU
  & Precedentes de decisões sobre aplicação da LAI por órgãos federais
  & A
  & \url{gov.br/acessoainformacao/pt-br/decisoes-da-cgu-e-da-cmri-sobre-a-aplicacao-da-lai} \\

\end{longtable}

% ────────────────────────────────────────────────────────────
\section*{A.11 Plataformas e Repositórios Agregadores}
\addcontentsline{toc}{section}{A.11 Plataformas e Repositórios Agregadores}
% ────────────────────────────────────────────────────────────

\begin{longtable}{@{} p{3cm} p{2.5cm} p{4cm} p{1.2cm} p{3.3cm} @{}}
\toprule
\textbf{Plataforma} & \textbf{Mantenedor} & \textbf{O que agrega} & \textbf{Acesso} & \textbf{URL} \\
\midrule
\endfirsthead
\multicolumn{5}{l}{\small\textit{(continuação)}} \\
\toprule
\textbf{Plataforma} & \textbf{Mantenedor} & \textbf{O que agrega} & \textbf{Acesso} & \textbf{URL} \\
\midrule
\endhead
\bottomrule
\endfoot

Base dos Dados
  & Base dos Dados
  & PNAD-C, Censo, RAIS, SIM, FINBRA e dezenas de outras em BigQuery
  & A
  & \url{basedosdados.org} \\

Ipeadata
  & Ipea
  & Séries macroeconômicas e sociais históricas
  & A
  & \url{ipeadata.gov.br} \\

SIDRA
  & IBGE
  & Todas as pesquisas IBGE em tabelas
  & A
  & \url{sidra.ibge.gov.br} \\

TabNet
  & DATASUS/MS
  & Todos os sistemas do DATASUS
  & A
  & \url{tabnet.datasus.gov.br} \\

dados.gov.br
  & CGU
  & Catálogo de dados abertos federais
  & A
  & \url{dados.gov.br} \\

QEdu
  & Meritt/Fundação Lemann
  & Dados educacionais (Censo Escolar, SAEB, IDEB) com visualização
  & A
  & \url{qedu.org.br} \\

Atlas da Violência
  & Ipea/FBSP
  & Indicadores de violência por município (base SIM)
  & A
  & \url{ipea.gov.br/atlasviolencia} \\

\end{longtable}

% ============================================================
%  anexo_b_glossario.tex
%  Anexo B: Glossário de Termos Técnicos
% ============================================================

% ════════════════════════════════════════════════════════════
\chapter*{Anexo B — Glossário de Termos Técnicos}
\addcontentsline{toc}{chapter}{Anexo B — Glossário de Termos Técnicos}
\markboth{Anexo B}{Anexo B}
% ════════════════════════════════════════════════════════════

Este glossário reúne os principais termos técnicos utilizados ao longo do guia, com definições concisas orientadas ao uso prático em avaliação de políticas públicas. Os termos estão organizados em ordem alfabética.

\medskip

\newcommand{\termo}[2]{%
  \noindent\textbf{#1}\par
  \noindent #2
  \medskip
}

% ── A ────────────────────────────────────────────────────────

\termo{Alucinação (IA)}{
Fenômeno pelo qual modelos de inteligência artificial geram informações factualmente incorretas com aparente confiança. Especialmente problemático em contextos que exigem precisão factual, como valores numéricos, datas e referências metodológicas. Ver Capítulo~12.
}

\termo{Amostra complexa}{
Desenho amostral que combina estratificação, conglomeração e probabilidades desiguais de seleção. Pesquisas domiciliares como a PNAD Contínua usam amostras complexas, o que exige procedimentos estatísticos específicos (declaração do desenho amostral, uso de pesos) para produzir estimativas válidas.
}

\termo{Amostragem probabilística}{
Método de seleção de amostras em que cada unidade da população tem uma probabilidade conhecida e não nula de ser selecionada. Permite a inferência estatística para a população a partir da amostra.
}

\termo{Anonimização}{
Processo pelo qual dados pessoais perdem a possibilidade de identificação direta ou indireta de uma pessoa natural. Dados anonimizados de forma irreversível estão fora do escopo da LGPD. Distinto de pseudonimização.
}

\termo{API (\textit{Application Programming Interface})}{
Interface que permite a comunicação padronizada entre sistemas de computador. APIs governamentais permitem acesso programático a dados públicos sem necessidade de download manual. Ver Capítulo~12.
}

\termo{Avaliação de impacto}{
Tipo de avaliação que busca estimar o efeito causal de uma política ou programa sobre um indicador de resultado, isolando esse efeito de outros fatores que poderiam explicar mudanças observadas. Requer a construção de um contrafactual.
}

% ── B ────────────────────────────────────────────────────────

\termo{Base de dados longitudinal}{
Base de dados em que as mesmas unidades (indivíduos, domicílios, escolas) são observadas em múltiplos pontos no tempo. Permite análises de trajetória e controle de características não observadas fixas no tempo.
}

\termo{Benchmarking}{
Processo de comparação do desempenho de uma unidade (país, estado, município, escola) com outras unidades de referência, para identificar brechas de desempenho e aprender com boas práticas.
}

% ── C ────────────────────────────────────────────────────────

\termo{CID-10 (Classificação Internacional de Doenças, 10ª Revisão)}{
Sistema de classificação de doenças e causas de morte da Organização Mundial da Saúde, adotado pelo Brasil a partir de 1996. Usado no SIM e no SINAN para codificação padronizada de causas de óbito e agravos notificados.
}

\termo{Cifra oculta (\textit{dark figure of crime})}{
Diferença entre o crime que efetivamente ocorre e o crime que é registrado nos sistemas oficiais. Estimativas sugerem que apenas uma fração dos crimes é reportada à polícia no Brasil. Ver Capítulo~8.
}

\termo{CNAE (Classificação Nacional de Atividades Econômicas)}{
Sistema de classificação de atividades econômicas adotado pelo IBGE e utilizado em registros administrativos como CAGED e RAIS para identificar o setor de atividade de estabelecimentos e trabalhadores.
}

\termo{Contrafactual}{
Situação hipotética que representa o que teria acontecido com os beneficiários de uma política caso ela não tivesse existido. Central para a avaliação de impacto causal.
}

\termo{Cobertura (de uma fonte de dados)}{
Proporção da população ou fenômeno de interesse que é efetivamente registrada por uma fonte. Fontes com baixa cobertura produzem viés de seleção se as unidades não cobertas têm características sistemáticas diferentes das cobertas.
}

\termo{Convênio de dados}{
Acordo formal entre uma instituição de pesquisa e um órgão governamental que regula o acesso a dados restritos. Geralmente requer proposta de pesquisa, termo de sigilo e aprovação institucional. Ver Capítulo~13.
}

% ── D ────────────────────────────────────────────────────────

\termo{DALY (Disability-Adjusted Life Year)}{
Medida de carga de doença que combina anos de vida perdidos por morte prematura e anos vividos com incapacidade. Usada pela OMS e pelo Banco Mundial para comparações de saúde entre países.
}

\termo{Dados administrativos}{
Dados produzidos como subproduto da operação rotineira de serviços e programas públicos, não para fins de pesquisa. Exemplos: CAGED, RAIS, Censo Escolar, SIM. Cobrem apenas o que passa pelo sistema formal.
}

\termo{Dados primários}{
Dados coletados diretamente pelo pesquisador ou avaliador para os fins específicos de uma pesquisa. Oferecem maior personalização, mas têm custo e tempo de coleta maiores.
}

\termo{Dados secundários}{
Dados coletados por outros agentes, para outros fins, e disponibilizados para uso. Incluem pesquisas do IBGE, registros administrativos e bases internacionais.
}

\termo{Dados sintéticos}{
Dados artificialmente gerados com características estatísticas similares às de dados reais. Usados para testes de código e treinamento de modelos, mas devem ser declarados explicitamente em análises de políticas públicas.
}

\termo{Deflacionamento}{
Processo de converter valores monetários nominais (em preços correntes) em valores reais (em preços constantes de um ano-base), usando um índice de preços como o IPCA. Necessário para comparações de gasto ao longo do tempo.
}

\termo{Desagregação}{
Capacidade de uma fonte de apresentar dados separados por subgrupos (sexo, raça/cor, faixa etária, território). Fundamental para análises de equidade e heterogeneidade de impacto.
}

\termo{Diferença-em-diferenças (DiD)}{
Método de avaliação de impacto que compara a mudança no indicador de resultado entre o grupo tratado e o grupo de controle, antes e depois da implementação da política. Ver o \textit{Guia de Avaliação de Impacto} da Série Avaliação na Prática.
}

\termo{Domínio (de uma pesquisa amostral)}{
Subgrupo da população para o qual a amostra foi planejada para produzir estimativas com precisão adequada. Estimativas para domínios não planejados têm erros padrão maiores e podem ser não confiáveis.
}

% ── E ────────────────────────────────────────────────────────

\termo{Efeito \textit{flypaper}}{
Fenômeno pelo qual transferências intergovernamentais (como o FPM) têm um efeito multiplicador sobre o gasto público local maior do que um aumento equivalente na renda privada dos moradores --- como se o dinheiro ``grudasse'' onde cai.
}

\termo{Erro de amostragem}{
Diferença entre a estimativa obtida a partir de uma amostra e o verdadeiro valor populacional. Pode ser quantificado por intervalos de confiança quando o desenho amostral é probabilístico.
}

\termo{Erro de medição}{
Diferença entre o valor registrado em uma fonte de dados e o verdadeiro valor do fenômeno de interesse. Pode ser aleatório (ruído) ou sistemático (viés).
}

\termo{Experimento natural}{
Situação em que variações exógenas nas condições de tratamento ocorrem de forma não intencional, mas que o pesquisador pode explorar como se fosse um experimento controlado para estimar efeitos causais.
}

% ── G ────────────────────────────────────────────────────────

\termo{Georreferenciamento}{
Processo de associar informações a coordenadas geográficas (latitude e longitude). Permite a produção de mapas e a realização de análises espaciais.
}

\termo{Gini (índice de)}{
Medida de desigualdade que varia de 0 (igualdade perfeita) a 1 (desigualdade máxima). Amplamente utilizado para medir desigualdade de renda. No Brasil, é calculado principalmente a partir dos microdados da PNAD Contínua e da POF.
}

% ── I ────────────────────────────────────────────────────────

\termo{Identificador administrativo}{
Código único que identifica uma unidade (pessoa, estabelecimento, município) em um ou mais sistemas administrativos. Exemplos: CPF (pessoa física), CNPJ (pessoa jurídica), código IBGE (município), CNES (estabelecimento de saúde).
}

\termo{Instrumento (variável instrumental)}{
Variável usada em estimação por variáveis instrumentais (IV) que afeta o indicador de resultado apenas por meio da variável de tratamento e é exógena ao processo estudado. Ver o \textit{Guia de Avaliação de Impacto} da Série.
}

\termo{Integração de bases (\textit{record linkage})}{
Processo de combinar informações de duas ou mais bases de dados usando um identificador comum (como CPF, CNPJ ou código IBGE). Também referida como \textit{linkagem} ou cruzamento de bases. Permite construir perfis multidimensionais de unidades de análise. Ver Capítulo~13.
}

% ── L ────────────────────────────────────────────────────────

\termo{LAI (Lei de Acesso à Informação)}{
Lei n.º 12.527/2011, que garante a qualquer cidadão o direito de solicitar informações a órgãos e entidades públicas. Principal mecanismo para acesso a dados governamentais não disponibilizados publicamente. Ver Capítulo~13.
}

\termo{LGPD (Lei Geral de Proteção de Dados)}{
Lei n.º 13.709/2018, que estabelece regras para o tratamento de dados pessoais no Brasil. Aplica-se a avaliações de políticas públicas que envolvem dados identificáveis de pessoas. Ver Capítulo~13.
}

\termo{Linha de base}{
Conjunto de dados que caracteriza o estado de um indicador antes da implementação de uma política. Fundamental para comparações antes-depois e para a construção de grupos de comparação em avaliações de impacto.
}

\termo{LLM (\textit{Large Language Model})}{
Modelo de inteligência artificial treinado em grandes volumes de texto, capaz de gerar, resumir e estruturar linguagem natural. Exemplos: GPT-4, Claude. Ver Capítulo~12.
}

% ── M ────────────────────────────────────────────────────────

\termo{Microdados}{
Arquivos de dados em que cada linha corresponde a uma unidade de análise individual (pessoa, domicílio, escola, estabelecimento), antes de qualquer agregação. Permitem análises customizadas que não são possíveis com tabelas pré-agregadas.
}

% ── N ────────────────────────────────────────────────────────

\termo{Não-resposta diferencial}{
Situação em que certos grupos de unidades respondem menos a uma pesquisa do que outros de forma sistemática, o que pode introduzir viés nas estimativas mesmo em amostras probabilísticas.
}

\termo{NIS (Número de Identificação Social)}{
Identificador único de cada pessoa cadastrada no CadÚnico. Usado para integração entre o CadÚnico e os registros de programas sociais como o Bolsa Família e o BPC.
}

% ── P ────────────────────────────────────────────────────────

\termo{Painel de dados}{
Base de dados que combina as dimensões de corte transversal (múltiplas unidades) e série temporal (múltiplos períodos), observando as mesmas unidades ao longo do tempo.
}

\termo{Pareamento por escore de propensão (\textit{propensity score matching})}{
Método de avaliação de impacto que constrói um grupo de controle a partir de unidades não tratadas com probabilidade de tratamento similar à das unidades tratadas, com base em características observáveis. Ver o \textit{Guia de Matching} da Série.
}

\termo{Periodicidade}{
Frequência com que uma fonte de dados é coletada e/ou publicada. Pode ser contínua, mensal, trimestral, anual, quinquenal, decenal ou irregular.
}

\termo{Peso amostral}{
Fator de expansão que indica quantas unidades da população cada unidade amostrada representa. Deve ser usado em todas as análises com pesquisas amostrais complexas para produzir estimativas representativas da população.
}

\termo{PLN (Processamento de Linguagem Natural)}{
Área da inteligência artificial voltada para a análise e geração de linguagem humana por computadores. Aplicações incluem classificação de texto, extração de entidades e análise de sentimento. Ver Capítulo~12.
}

\termo{Poder estatístico}{
Probabilidade de que um teste estatístico detecte um efeito real quando ele existe. Depende do tamanho da amostra, da magnitude do efeito e do nível de significância adotado. Ver o \textit{Guia de Poder Estatístico} da Série.
}

\termo{PPC (Paridade do Poder de Compra)}{
Taxa de câmbio ajustada pelo nível de preços de cada país, que permite comparações de renda e bem-estar entre países de forma mais precisa do que taxas de câmbio de mercado.
}

\termo{Pré-registro}{
Prática de registrar publicamente, antes da coleta ou análise dos dados, o desenho de pesquisa, as hipóteses e os métodos de análise planejados. Aumenta a credibilidade dos resultados ao prevenir o ajuste retroativo das hipóteses (\textit{HARKing}).
}

\termo{Pseudonimização}{
Substituição de identificadores diretos (como CPF) por códigos que não permitem identificação direta, mas que podem ser revertidos com acesso a uma tabela de correspondência. Dados pseudonimizados ainda são dados pessoais sob a LGPD.
}

% ── Q ────────────────────────────────────────────────────────

\termo{Quebra metodológica}{
Alteração no método de coleta, definição de variáveis ou universo de cobertura de uma fonte de dados que compromete a comparabilidade da série histórica antes e depois da mudança. Exemplos: transição da PNAD para PNAD Contínua; Novo CAGED em 2020.
}

% ── R ────────────────────────────────────────────────────────

\termo{Regressão descontínua (RDD)}{
Método de avaliação de impacto que explora descontinuidades em regras de elegibilidade (como limiares de renda ou de idade) para estimar efeitos causais, comparando unidades imediatamente acima e abaixo do limiar. Ver o \textit{Guia de Avaliação de Impacto} da Série.
}

\termo{Representatividade}{
Propriedade de uma amostra de refletir as características da população de interesse. Uma amostra representativa permite inferências válidas sobre a população a partir dos dados amostrados.
}

\termo{Reprodutibilidade}{
Capacidade de um estudo ser replicado por outro pesquisador com os mesmos dados e métodos, obtendo os mesmos resultados. Depende da documentação do processo de coleta, limpeza e análise. Ver Capítulo~13.
}

\termo{Registro administrativo}{
Dado produzido como subproduto da operação rotineira de serviços públicos, não para fins de pesquisa. Cobre apenas o que passa pelo sistema formal; pode ter qualidade de preenchimento variável.
}

% ── S ────────────────────────────────────────────────────────

\termo{\textit{Scraping} (web scraping)}{
Técnica de extração automatizada de informações de páginas web. Usada quando dados de interesse estão disponíveis em formato HTML mas não há API ou download estruturado. Ver Capítulo~12.
}

\termo{Setor censitário}{
Menor unidade territorial do sistema estatístico brasileiro, usada como unidade de coleta do Censo Demográfico. Agrupa aproximadamente 200 a 300 domicílios em áreas urbanas.
}

\termo{SIG (Sistema de Informações Geográficas)}{
Sistema para captura, armazenamento, análise e visualização de dados georreferenciados. Softwares como QGIS e bibliotecas como \texttt{sf} (R) e \texttt{geopandas} (Python) são amplamente usados em análises espaciais de políticas públicas.
}

\termo{Subregistro}{
Fenômeno pelo qual eventos reais (nascimentos, óbitos, crimes, doenças) não são registrados nos sistemas administrativos competentes. Pode introduzir viés nas análises se o subregistro for diferencial entre grupos ou regiões.
}

\termo{Subnotificação}{
Caso específico de subregistro em que agravos de saúde de notificação compulsória não são comunicados às autoridades sanitárias. Prevalente no SINAN, especialmente para doenças com maior estigma ou em períodos de alta demanda dos serviços.
}

% ── T ────────────────────────────────────────────────────────

\termo{Taxa de integração (ou taxa de \textit{linkage})}{
Proporção de registros de uma base que foram vinculados com sucesso a registros de outra base em um processo de integração. Taxas baixas podem introduzir viés se os registros não cruzados têm características sistemáticas diferentes.
}

\termo{Triangulação}{
Uso de múltiplas fontes de dados ou métodos para verificar a consistência de um resultado. Quando diferentes fontes convergem para o mesmo resultado, a credibilidade da evidência aumenta.
}

\termo{TRI (Teoria de Resposta ao Item)}{
Modelo psicométrico usado para calcular proficiências em avaliações educacionais (como SAEB e ENEM) a partir das respostas a itens individuais. Permite comparações de desempenho entre diferentes edições de uma avaliação, mesmo com diferentes conjuntos de itens.
}

% ── U ────────────────────────────────────────────────────────

\termo{Unidade de análise}{
Entidade sobre a qual os dados são organizados e as análises são conduzidas --- indivíduo, domicílio, escola, município, estabelecimento. A unidade de análise da fonte deve ser compatível com a pergunta avaliativa.
}

% ── V ────────────────────────────────────────────────────────

\termo{Validade externa}{
Grau em que os resultados de uma avaliação podem ser generalizados para além da população, do período e do contexto específicos em que foram obtidos.
}

\termo{Validade interna}{
Grau em que uma avaliação estima corretamente o efeito causal da política sobre os beneficiários estudados, sem confundimento por outros fatores.
}

\termo{Variável de controle}{
Variável incluída em um modelo estatístico para controlar por diferenças entre grupos que poderiam confundir a estimativa do efeito de interesse.
}

\termo{Viés de seleção}{
Situação em que os grupos comparados em uma avaliação diferem sistematicamente em características que afetam o indicador de resultado, além da exposição à política avaliada. Principal ameaça à validade interna das avaliações.
}

\termo{Viés de cobertura}{
Tipo de viés de seleção em que determinados grupos populacionais são sistematicamente sub-representados ou ausentes em uma fonte de dados. Exemplo: trabalhadores informais estão ausentes da RAIS e do CAGED.
}

% Bibliografia recomendada: lista todas as entradas do .bib mesmo
% sem citação no texto (este é um guia de fontes, não um paper).
\renewcommand{\bibname}{Bibliografia recomendada}
\nocite{*}
\bibliography{referencias}

\end{document}
